\documentclass[11pt]{article}
\usepackage[margin=1in]{geometry}
\usepackage[T1]{fontenc}
\usepackage[utf8]{inputenc}
\usepackage[french]{babel}
\usepackage{amsmath,amssymb,amsthm}
\usepackage[colorlinks=true,linkcolor=blue,urlcolor=blue]{hyperref}
\newtheorem{problem}{Problème}
\newenvironment{refs}{\par\small\noindent\emph{Références :}\begin{itemize}\setlength\itemsep{0pt}}{\end{itemize}\normalsize}
\title{Hors de la Caverne \\ \Large La salle des lycéens}
\author{}
\date{}
\begin{document}
\maketitle
\noindent\emph{Des problèmes, pas de réponses.} Chaque problème ci-dessous est posé
sans solution. Les références indiquent où trouver les connaissances nécessaires.
\medskip


\section*{Logique, théorie des ensembles et fondements}

\begin{problem}[{Chevaliers et valets --- Lycée}]
\textbf{FND-01.} Sur une certaine île, tout habitant est un \textit{chevalier},
qui dit toujours la vérité, ou un \textit{valet}, qui ment toujours.
\begin{enumerate}
\item[(a)] Vous rencontrez A et B. A déclare :
\og{} Au moins l'un de nous deux est un valet. \fg{} Déterminez, avec démonstration, ce que sont A et B.
\item[(b)] Vous rencontrez C et D.
C déclare : \og{} Nous sommes tous les deux des valets. \fg{} Déterminez ce que sont C et D.
\item[(c)] Vous arrivez à une bifurcation ; une
branche mène au village, l'autre au marécage. Un habitant isolé, de type inconnu, se tient
à la bifurcation. Trouvez une unique question fermée (oui/non) que vous puissiez poser de sorte que
la réponse vous indique la bonne route, et démontrez que votre question fonctionne, que
l'habitant soit chevalier ou valet.
\end{enumerate}

\end{problem}

\noindent Raymond Smullyan a rendu ces énigmes célèbres --- et forgé les noms de chevalier
et de valet --- dans son recueil de 1978 \textit{What Is the Name of This Book?} C'est de la logique
propositionnelle déguisée : chaque énigme revient à décider quelles attributions de valeurs de
vérité satisfont une petite formule, et (c) est un premier avant-goût de la conception d'une
formule dont la valeur de vérité est invariante sous le comportement d'un adversaire. Les
résoudre par une analyse de cas systématique, plutôt que par un éclair d'intuition, est le
premier pas vers la démonstration.

\begin{refs}
\item Contexte : Chevaliers et valets --- Wikipédia (en anglais) (\url{https://en.wikipedia.org/wiki/Knights_and_Knaves})
\item Lecture : Raymond Smullyan, \textit{What Is the Name of This Book?}
\end{refs}

\bigskip

\begin{problem}[{Seize fonctions de vérité --- Lycée}]
\textbf{FND-02.} \begin{enumerate}
\item[(a)] Écrivez les tables de vérité de $ \neg P $,
$ P \wedge Q $, $ P \vee Q $, $ P \to Q $ et $ P \leftrightarrow Q $, et vérifiez par
table de vérité que $ ((P \to Q) \wedge (Q \to R)) \to (P \to R) $ est une tautologie (vraie
sous toute attribution de valeurs).
\item[(b)] Montrez que $ P \to Q $ n'est \textit{pas} logiquement équivalente à sa
réciproque $ Q \to P $, mais qu'elle l'\textit{est} à sa contraposée
$ \neg Q \to \neg P $.
\item[(c)] Comptez combien il existe de fonctions de vérité distinctes de deux variables,
et démontrez que chacune d'elles peut s'exprimer à l'aide des seuls connecteurs $ \neg $
et $ \wedge $.
\end{enumerate}

\end{problem}

\noindent Les tables de vérité ont pris leur forme reconnaissable chez Emil Post et
Ludwig Wittgenstein vers 1920, avec des racines chez Frege, Peirce et Boole. La partie (c) est un
premier théorème de \textit{complétude fonctionnelle} --- le constat qu'un tout petit jeu de
connecteurs engendre toute la logique propositionnelle. Henry Sheffer a montré en 1913 qu'un seul
connecteur (la barre de Sheffer, NON-ET) suffit, fait câblé en dur dans chaque puce
d'ordinateur.

\begin{refs}
\item Contexte : Table de vérité --- Wikipédia (\url{https://fr.wikipedia.org/wiki/Table_de_v%C3%A9rit%C3%A9})
\item Contexte : Complétude fonctionnelle --- Wikipédia (en anglais) (\url{https://en.wikipedia.org/wiki/Functional_completeness})
\item Lecture : Daniel Velleman, \textit{How to Prove It: A Structured Approach}, ch. 1
\end{refs}

\bigskip

\begin{problem}[{Le jeu de la négation des quantificateurs --- Lycée}]
\textbf{FND-03.} Toutes les variables parcourent les entiers relatifs $\mathbb{Z}$.
\begin{enumerate}
\item[(a)] Écrivez la
négation de chacun des énoncés ci-dessous sous une forme où aucun signe de négation ne précède
un quantificateur, puis déterminez --- avec démonstration --- lesquels des six énoncés (les trois
originaux, les trois négations) sont vrais : (i) $ \forall x\, \exists y\, (y \cdot y = x) $ ;
(ii) $ \exists y\, \forall x\, (x + y = x) $ ; (iii) tout nombre pair supérieur à 2 est
la somme de deux nombres premiers.
\item[(b)] Exhibez une propriété $ P(x,y) $ des couples d'entiers pour
laquelle $ \forall x\, \exists y\, P(x,y) $ est vraie mais $ \exists y\, \forall x\, P(x,y) $
est fausse, et démontrez que, pour toute propriété $ P $, le second énoncé implique le premier.
\end{enumerate}

\end{problem}

\noindent Apprendre que \og{} pour tout x il existe un y \fg{} et \og{} il existe un y pour tout x \fg{}
ne disent pas la même chose est le pas le plus important du calcul vers l'analyse : les
définitions en $\varepsilon$--$\delta$ de Weierstrass vivent tout entières dans les alternances de quantificateurs.
L'énoncé (iii) est la conjecture de Goldbach (1742), toujours non démontrée --- et c'est
exactement pourquoi on vous demande seulement de la nier correctement, non de la trancher ;
voyez la page des Problèmes ouverts. Depuis Hintikka, les logiciens voient les quantificateurs
alternés comme un jeu entre un vérificateur et un falsificateur.

\begin{refs}
\item Contexte : Quantification (logique) --- Wikipédia (\url{https://fr.wikipedia.org/wiki/Quantification_(logique)})
\item Contexte : Calcul des prédicats (logique du premier ordre) --- Wikipédia (\url{https://fr.wikipedia.org/wiki/Calcul_des_pr%C3%A9dicats})
\item Lecture : Daniel Velleman, \textit{How to Prove It: A Structured Approach}, ch. 2
\end{refs}

\bigskip

\begin{problem}[{La récurrence, et une démonstration qui démontre trop --- Lycée}]
\textbf{FND-04.} \begin{enumerate}
\item[(a)] Démontrez par récurrence que
$ 1 + 2 + \cdots + n = \frac{n(n+1)}{2} $ et que
$ 1^3 + 2^3 + \cdots + n^3 = (1 + 2 + \cdots + n)^2 $.
\item[(b)] Démontrez par récurrence forte que
tout entier $ n \ge 2 $ est un produit de nombres premiers.
\item[(c)] Voici une célèbre \og{} démonstration \fg{} que tous
les chevaux ont la même couleur. \textit{Cas de base :} tout ensemble de 1 cheval n'a qu'une couleur.
\textit{Pas de récurrence :} étant donné $ n+1 $ chevaux, retirez-en un ; les $ n $ restants
partagent une couleur par hypothèse. Remettez-le, puis retirez-en un autre ; de nouveau les $ n $
restants partagent une couleur. Les deux ensembles se chevauchant, les $ n+1 $ chevaux doivent
tous partager une même couleur. Identifiez le point exact où cet argument échoue, et expliquez
pourquoi.
\end{enumerate}

\end{problem}

\noindent La récurrence a été utilisée implicitement pendant des siècles, énoncée
clairement par Pascal (1654), nommée par De Morgan (1838), et érigée en axiome de l'arithmétique
par Peano (1889). Le sophisme des chevaux, popularisé par George Pólya comme mise en garde,
enseigne la leçon que tout mathématicien en exercice finit par apprendre à ses dépens : une
récurrence ne vaut que ce que vaut le cas où le pas de récurrence s'enclenche pour la première
fois.

\begin{refs}
\item Contexte : Raisonnement par récurrence --- Wikipédia (\url{https://fr.wikipedia.org/wiki/Raisonnement_par_r%C3%A9currence})
\item Contexte : Axiomes de Peano --- Wikipédia (\url{https://fr.wikipedia.org/wiki/Axiomes_de_Peano})
\item Lecture : Daniel Velleman, \textit{How to Prove It: A Structured Approach}, ch. 6
\end{refs}

\bigskip

\begin{problem}[{Le barbier, et l'ensemble de tous les ensembles --- Lycée}]
\textbf{FND-05.} \begin{enumerate}
\item[(a)] Dans un certain village, le barbier rase exactement les
villageois qui ne se rasent pas eux-mêmes. Démontrez qu'un tel barbier ne peut pas exister.
\item[(b)] La théorie naïve des ensembles
autorise la formation de l'ensemble $ \{x : P(x)\} $ pour toute propriété $ P $. Formez
$ R = \{x : x \notin x\} $ et déduisez-en une contradiction.
\item[(c)] Expliquez précisément quelle étape
de formation de (b) les axiomes de Zermelo--Fraenkel interdisent, énoncez le schéma d'axiomes de
séparation (compréhension) qui la remplace, et servez-vous de la séparation ainsi que de votre
argument de (b) pour démontrer qu'il n'existe pas d'ensemble de tous les ensembles.
\end{enumerate}

\end{problem}

\noindent Bertrand Russell a trouvé ce paradoxe en 1901 et l'a communiqué à Frege en
1902, au moment précis où le second volume de l'\oe{}uvre d'une vie de Frege, les
\textit{Grundgesetze}, était sous presse ; l'appendice où Frege reconnaît l'effondrement est l'un
des documents les plus poignants de la mathématique. Zermelo avait remarqué le paradoxe de son
côté. La réparation --- les axiomes de Zermelo de 1908, étendus par Fraenkel --- est la théorie des
ensembles (ZFC) dans laquelle on fait des mathématiques aujourd'hui.

\begin{refs}
\item Contexte : Paradoxe de Russell --- Wikipédia (\url{https://fr.wikipedia.org/wiki/Paradoxe_de_Russell})
\item Contexte : Paradoxe du barbier --- Wikipédia (\url{https://fr.wikipedia.org/wiki/Paradoxe_du_barbier})
\item Contexte : Théorie des ensembles de Zermelo--Fraenkel --- Wikipédia (\url{https://fr.wikipedia.org/wiki/Th%C3%A9orie_des_ensembles_de_Zermelo-Fraenkel})
\item Lecture : Paul Halmos, \textit{Naive Set Theory}
\end{refs}

\bigskip

\begin{problem}[{Une phrase qu'aucun insulaire ne peut prononcer --- Lycée, short}]
\textbf{FND-24.} Sur l'île des chevaliers et des valets de FND-01, démontrez
qu'aucun habitant --- chevalier ou valet --- ne peut jamais dire \og{} Je suis un valet \fg{}.
\end{problem}

\noindent C'est le paradoxe du menteur désamorcé : la phrase \og{} Je mens \fg{}, attribuée dans
l'Antiquité à Épiménide le Crétois, ne peut être affirmée de façon cohérente par quiconque est
astreint à la vérité ou au mensonge. Smullyan a fait de cette observation même le moteur de ses
recueils d'énigmes --- quand un personnage dit quelque chose qu'aucun insulaire ne pourrait dire,
ce sont les hypothèses elles-mêmes qui doivent céder.

\begin{refs}
\item Contexte : Paradoxe du menteur --- Wikipédia (\url{https://fr.wikipedia.org/wiki/Paradoxe_du_menteur})
\item Lecture : Raymond Smullyan, \textit{What Is the Name of This Book?}
\end{refs}

\bigskip

\begin{problem}[{Compter les parties --- Lycée, short}]
\textbf{FND-25.} Démontrez qu'un ensemble à $ n $ éléments possède
exactement $ 2^n $ parties, en comptant l'ensemble vide et l'ensemble tout entier.
\end{problem}

\noindent La collection de toutes les parties de $ A $ est l'\textit{ensemble des parties}
$ \mathcal{P}(A) $, et ce dénombrement est la raison de la notation $ 2^A $ : une partie,
c'est la même chose qu'une réponse par oui ou non pour chaque élément. Le théorème de Cantor
(FND-07), c'est ce qu'il advient de ce fait innocent quand $ n $ devient infini --- l'ensemble
des parties dépasse l'ensemble, à jamais.

\begin{refs}
\item Contexte : Ensemble des parties d'un ensemble --- Wikipédia (\url{https://fr.wikipedia.org/wiki/Ensemble_des_parties_d%27un_ensemble})
\item Lecture : Daniel Velleman, \textit{How to Prove It: A Structured Approach}, ch. 6
\end{refs}

\bigskip

\begin{problem}[{L'hôtel de Hilbert --- Lycée}]
\textbf{FND-36.} Un hôtel a des chambres numérotées $ 1, 2, 3, \ldots $ --- une par entier
strictement positif --- et toutes les chambres sont occupées. Dans chaque partie ci-dessous, donnez
une règle explicite de réattribution des clients aux chambres et \textit{démontrez} que votre règle
est une bijection, de sorte que personne ne soit expulsé et que personne ne partage sa chambre.
\begin{enumerate}
\item[(a)] Un nouveau client arrive. Logez tout le monde.
\item[(b)] Un car transportant une infinité de nouveaux clients, numérotés $ 1, 2, 3, \ldots $,
arrive. Logez tout le monde.
\item[(c)] Une infinité de cars arrivent, chacun transportant une infinité de clients, ceux-ci étant
repérés par des couples (numéro du car, numéro du siège). Logez tout le monde.
\item[(d)] On dit qu'un ensemble $ A $ est \textit{infini au sens de Dedekind} s'il existe une bijection
entre $ A $ et une partie propre de $ A $. Démontrez que $ \mathbb{N} $ est infini au sens
de Dedekind, qu'aucun ensemble fini ne l'est, et que tout ensemble infini au sens de Dedekind
contient une partie en bijection avec $ \mathbb{N} $.
\end{enumerate}

\end{problem}

\noindent Hilbert s'est servi de cet hôtel dans une conférence de 1924--25, \og{} Über das
Unendliche \fg{}, pour faire voir que \og{} l'hôtel est plein \fg{} et \og{} on ne peut plus loger personne \fg{}
cessent d'être la même phrase dès que les chambres sont en nombre infini ; George Gamow l'a porté
devant le grand public dans \textit{One Two Three\dots{} Infinity} (1947). La mathématique est plus
ancienne et plus tranchante que la plaisanterie : Dedekind avait déjà pris la propriété de
l'hôtel comme \textit{définition} de l'infini dans \textit{Was sind und was sollen die Zahlen?}
(1888) --- un ensemble est infini quand on peut le mettre en correspondance avec une partie de
lui-même. La partie (d) cache une subtilité qu'il vaut la peine de connaître : démontrer que tout
ensemble infini est infini au sens de Dedekind exige une forme faible de l'axiome du choix
(FND-09), et dans ZF seul les deux notions se séparent.

\begin{refs}
\item Contexte : Hôtel de Hilbert --- Wikipédia (\url{https://fr.wikipedia.org/wiki/H%C3%B4tel_de_Hilbert})
\item Contexte : Ensemble infini au sens de Dedekind --- Wikipédia (en anglais) (\url{https://en.wikipedia.org/wiki/Dedekind-infinite_set})
\item Lecture : Paul Halmos, \textit{Naive Set Theory}
\end{refs}

\bigskip

\begin{problem}[{La récurrence et le principe du plus petit élément --- Lycée, short}]
\textbf{FND-37.} Le \textit{principe du plus petit élément} dit que toute partie
non vide de l'ensemble des entiers naturels contient un plus petit élément. Démontrez qu'il est
équivalent au principe de récurrence --- déduisez chacun de l'autre.
\end{problem}

\noindent Tout argument du type \og{} prenons un contre-exemple minimal \fg{} et toute récurrence
sont la même démonstration sous des habits différents, et c'est le théorème qui le dit ; la
méthode de descente infinie de Fermat (NT-29, sur la page de théorie des nombres) est le principe
du plus petit élément employé comme une arme. Peano a pris la récurrence pour axiome en 1889, si
bien que chez lui le principe du plus petit élément est un théorème --- mais ce choix est une
convention, non un fait sur les entiers. Un avertissement pour plus tard : restreints à des
formules de complexité logique bornée, les deux schémas cessent d'être interchangeables, et
démêler exactement en quoi ils diffèrent est un vrai sujet, l'étude des fragments de
l'arithmétique.

\begin{refs}
\item Contexte : Principe du bon ordre --- Wikipédia (en anglais) (\url{https://en.wikipedia.org/wiki/Well-ordering_principle})
\item Contexte : Raisonnement par récurrence --- Wikipédia (\url{https://fr.wikipedia.org/wiki/Raisonnement_par_r%C3%A9currence})
\end{refs}

\bigskip


\section*{Algèbre abstraite}

\begin{problem}[{La preuve par neuf et l'horloge des puissances --- Lycée}]
\textbf{ALG-01.} Deux entiers $a$ et $b$ sont dits \textit{congrus
modulo $n$}, ce que l'on écrit $a \equiv b \pmod{n}$, lorsque $n$ divise $a-b$.
\begin{enumerate}
\item[(i)] Démontrez que la
congruence modulo $n$ respecte l'addition et la multiplication : si $a \equiv a' \pmod{n}$ et
$b \equiv b' \pmod{n}$, alors $a+b \equiv a'+b' \pmod{n}$ et $ab \equiv a'b' \pmod{n}$.
\item[(ii)] À l'aide de (i), démontrez qu'un entier strictement positif est divisible par 9 si et seulement si la somme de ses
chiffres décimaux l'est.
\item[(iii)] Déterminez, en le démontrant, le dernier chiffre décimal de
$7^{2026}$.
\end{enumerate}

\end{problem}

\noindent L'arithmétique modulaire a été rendue systématique par Gauss dans les premières pages de ses
\textit{Disquisitiones Arithmeticae} (1801), écrites alors qu'il avait 24 ans ; la notation $\equiv$ est de lui. Le
critère par la somme des chiffres --- la \og{} preuve par neuf \fg{} --- est bien plus ancien : les comptables du Moyen Âge
s'en servaient pour vérifier leurs livres. Derrière la question (iii) se cache le premier exemple de groupe de ce
recueil : les résidus modulo 10 qui possèdent un inverse multiplicatif, tournant en rond sous la multiplication.

\begin{refs}
\item Contexte : Arithmétique modulaire --- Wikipédia (\url{https://fr.wikipedia.org/wiki/Arithm%C3%A9tique_modulaire})
\item Contexte : Disquisitiones arithmeticae --- Wikipédia (\url{https://fr.wikipedia.org/wiki/Disquisitiones_arithmeticae})
\item Lecture : J. Fraleigh, \textit{A First Course in Abstract Algebra}, sections 0 et 4
\end{refs}

\bigskip

\begin{problem}[{Les vingt-quatre rotations du cube --- Lycée}]
\textbf{ALG-02.} \begin{enumerate}
\item[(i)] Démontrez qu'un cube possède exactement 24 symétries de rotation
--- les rotations de l'espace qui envoient le cube sur lui-même --- et justifiez que votre liste est complète.
\item[(ii)] Un cube possède quatre grandes diagonales (les segments joignant deux sommets opposés). Montrez que toute
symétrie de rotation permute ces quatre diagonales, et que deux rotations distinctes donnent deux permutations
distinctes.
\item[(iii)] Concluez que les rotations du cube réalisent \textit{toutes} les
$4! = 24$ réorganisations possibles des quatre diagonales.
\end{enumerate}

\end{problem}

\noindent Dénombrer des symétries, c'est l'endroit où la théorie des groupes se laisse voir à l'\oe{}il nu.
Les solides réguliers ont été étudiés par les Grecs --- les \textit{Éléments} d'Euclide s'achèvent sur eux --- mais
l'idée que leurs symétries forment un système clos d'opérations composables appartient au XIXe siècle,
et la question (iii) dit que le groupe des rotations du cube \og{} est \fg{} le groupe symétrique $S_4$ déguisé.
Le programme d'Erlangen, énoncé par Klein en 1872, proposait de définir la géométrie elle-même par ses groupes de symétrie.

\begin{refs}
\item Contexte : Symétrie octaédrique --- Wikipédia (en anglais) (\url{https://en.wikipedia.org/wiki/Octahedral_symmetry})
\item Contexte : Groupe de symétrie --- Wikipédia (\url{https://fr.wikipedia.org/wiki/Groupe_de_sym%C3%A9trie})
\item Lecture : M. Artin, \textit{Algebra}, ch. 6 (la symétrie)
\end{refs}

\bigskip

\begin{problem}[{Quand $n^4 + 4$ est-il premier ? --- Lycée}]
\textbf{ALG-03.} Déterminez, en le démontrant, tous les entiers strictement positifs $n$ pour lesquels
$n^4 + 4$ est un nombre premier. Plus généralement, montrez que $n^4 + 4m^4$ est composé pour tous les
entiers $n, m \ge 2$.
\end{problem}

\noindent La clé est une factorisation dont la plupart des gens jurent qu'elle ne peut pas exister --- une somme
qui se factorise. Elle porte le nom de Sophie Germain (1776--1831), qui apprit seule les mathématiques
dans la bibliothèque de son père pendant la Révolution française et correspondit avec Gauss sous un
pseudonyme masculin, jusqu'à ce qu'il découvre, ravi, qui était réellement \og{} Monsieur Le Blanc \fg{}. L'identité
qui gouverne ce problème traverse toutes les mathématiques olympiques et intervient dans ses travaux sur le
dernier théorème de Fermat.

\begin{refs}
\item Contexte : Identité de Sophie Germain --- Wikipédia (\url{https://fr.wikipedia.org/wiki/Identit%C3%A9_de_Sophie_Germain})
\item Contexte : Sophie Germain --- Wikipédia (\url{https://fr.wikipedia.org/wiki/Sophie_Germain})
\item Lecture : A. Engel, \textit{Problem-Solving Strategies}, ch. 6 (théorie des nombres)
\end{refs}

\bigskip

\begin{problem}[{Le théorème de la racine rationnelle, et pourquoi $\sqrt[3]{2}$ est irrationnel --- Lycée}]
\textbf{ALG-04.} \begin{enumerate}
\item[(i)] Démontrez le théorème de la racine rationnelle : si un polynôme
$a_n x^n + \cdots + a_1 x + a_0$ à coefficients entiers admet une racine rationnelle $p/q$ écrite sous
forme irréductible, alors $p$ divise $a_0$ et $q$ divise $a_n$.
\item[(ii)] Déduisez-en que
$\sqrt[3]{2}$ est irrationnel.
\item[(iii)] Montrez que $x^3 - x - 1$ n'admet aucune racine rationnelle, et
expliquez pourquoi, pour un polynôme de degré trois, cela démontre déjà qu'il ne peut pas se factoriser en
polynômes de degré strictement inférieur à coefficients rationnels.
\end{enumerate}

\end{problem}

\noindent C'est l'humble ancêtre de tout ce qui figure dans la bande Master, plus bas : l'irrationalité
de $\sqrt[3]{2}$ est la première obstruction, et la plus facile, du problème vieux de deux mille ans de la
duplication du cube (ALG-18). Le théorème lui-même revient pour l'essentiel à Descartes et était devenu
routinier du temps d'Euler ; ce qui a changé ensuite, c'est que \og{} n'a pas de racine rationnelle \fg{} est devenu la
notion bien plus profonde d'\textit{irréductibilité}, le concept-atome de la théorie des corps.

\begin{refs}
\item Contexte : Théorème de la racine rationnelle --- Wikipédia (\url{https://fr.wikipedia.org/wiki/Racine_%C3%A9vidente})
\item Contexte : Polynôme irréductible --- Wikipédia (\url{https://fr.wikipedia.org/wiki/Polyn%C3%B4me_irr%C3%A9ductible})
\item Lecture : J. Fraleigh, \textit{A First Course in Abstract Algebra}, section 23
\end{refs}

\bigskip

\begin{problem}[{Les symétries d'un rectangle --- Lycée, short}]
\textbf{ALG-26.} Énumérez les quatre symétries d'un rectangle non carré, dressez leur table de composition, et démontrez que tout élément autre que l'identité est son propre inverse.
\end{problem}

\noindent C'est le groupe de Klein, le plus petit groupe qui ne soit pas cyclique, et on peut le découvrir avec un morceau de carton avant même que le mot \og{} groupe \fg{} ait été défini. Construire la table à la main, puis en remarquer le motif, c'est exactement ainsi que les axiomes abstraits ont été dégagés historiquement --- les exemples sont venus d'abord.

\begin{refs}
\item Contexte : Groupe de Klein --- Wikipédia (\url{https://fr.wikipedia.org/wiki/Groupe_de_Klein})
\end{refs}

\bigskip

\begin{problem}[{Une horloge qui perd des nombres --- Lycée, short}]
\textbf{ALG-27.} Dans l'arithmétique modulo $ 12 $, déterminez exactement quels résidus possèdent un inverse multiplicatif, et démontrez que votre critère est correct pour un module $ n $ quelconque.
\end{problem}

\noindent L'arithmétique modulaire ressemble à l'arithmétique ordinaire jusqu'au moment où la division échoue, et déterminer précisément quand elle échoue --- exactement lorsque le résidu partage un facteur avec le module --- est le premier véritable théorème que la plupart des gens démontrent sur les anneaux. C'est aussi la raison pour laquelle les modules premiers se comportent tellement mieux, et donc celle pour laquelle la cryptographie les préfère.

\begin{refs}
\item Contexte : Inverse modulaire --- Wikipédia (\url{https://fr.wikipedia.org/wiki/Inverse_modulaire})
\end{refs}

\bigskip

\begin{problem}[{La parité, et le 14-15 qui ne sort jamais --- Lycée}]
\textbf{ALG-34.} Une \textit{transposition} est une permutation qui échange deux éléments et fixe
tous les autres ; toute permutation d'un ensemble fini est un produit de transpositions. Le \textit{taquin}
est un plateau $4 \times 4$ portant les pièces $1, 2, \dots, 15$ et une case vide ; un coup fait glisser
dans la case vide une pièce qui lui est adjacente horizontalement ou verticalement.
\begin{enumerate}
\item[(a)] Démontrez que si une permutation $\sigma$ s'écrit de deux façons comme produit de
transpositions, les deux nombres de facteurs ont la même parité. La \textit{signature}
$\operatorname{sgn}(\sigma) = \pm 1$ est donc bien définie ; démontrez que
$\operatorname{sgn}(\sigma\tau) = \operatorname{sgn}(\sigma)\operatorname{sgn}(\tau)$.
\item[(b)] Déduisez-en que pour $n \ge 2$ les permutations paires forment un sous-groupe de $S_n$ contenant
exactement la moitié de ses éléments.
\item[(c)] Lisez une position du taquin comme une permutation $\sigma$ des seize cases (la case
vide comptant pour une seizième pièce), la case vide occupant la ligne $i$ et la colonne $j$.
Démontrez que
\[ \operatorname{sgn}(\sigma) \cdot (-1)^{\,i+j} \]
est inchangé par tout coup légal.
\item[(d)] Concluez que la position obtenue à partir de la position résolue en échangeant les pièces 14 et 15 ne
peut jamais être atteinte, et démontrez qu'au plus la moitié des $16!$ dispositions sont accessibles à
partir de l'une quelconque d'entre elles.
\end{enumerate}

\end{problem}

\noindent Le casse-tête fut imaginé vers 1874 par Noyes Palmer Chapman, receveur des postes à
Canastota, dans l'État de New York, et fit fureur en 1880 ; Sam Loyd, qui n'avait rien à voir avec son
invention, offrit plus tard $\$1000$ pour une solution de la position 14-15 et revendiqua le casse-tête
comme sien à partir de 1891 --- revendication démolie par l'histoire qu'en ont écrite Slocum et Sonneveld en
2006. Le prix ne risquait rien, car dès 1879 William Woolsey Johnson et William E. Story avaient publié
l'invariant de la question (c) dans l'\textit{American Journal of Mathematics} : l'un des tout premiers
arguments où un invariant algébrique démontre qu'une chose est impossible, au lieu de simplement échouer à
la faire. La signature d'une permutation, née des travaux du XVIIIe siècle sur les déterminants,
est le même homomorphisme qui produit le groupe alterné d'ALG-11 et, pour finir, l'insolubilité de
l'équation du cinquième degré.

\begin{refs}
\item Contexte : Taquin --- Wikipédia (\url{https://fr.wikipedia.org/wiki/Taquin})
\item Contexte : Signature d'une permutation --- Wikipédia (\url{https://fr.wikipedia.org/wiki/Signature_d%27une_permutation})
\item Article : Wm. W. Johnson \& W. E. Story, \og{} Notes on the ``15'' Puzzle \fg{} (1879), American Journal of Mathematics 2, 397--404
\item Article : A. F. Archer, \og{} A modern treatment of the 15 puzzle \fg{} (1999), American Mathematical Monthly 106, 793--799
\end{refs}

\bigskip

\begin{problem}[{Toute table de groupe est un carré latin --- Lycée, short}]
\textbf{ALG-35.} Démontrez que dans la table de multiplication d'un groupe fini,
chaque élément du groupe apparaît exactement une fois dans chaque ligne et exactement une fois dans chaque
colonne. Montrez ensuite, en en exhibant un, qu'un tableau carré ayant cette propriété n'est pas
nécessairement la table de multiplication d'un groupe.
\end{problem}

\noindent La table est une invention de Cayley : dans son article de 1854 sur les groupes abstraits,
il en dresse une et observe exactement ce motif, qui n'est rien d'autre que la règle de simplification rendue
visible. Les carrés remplis de sorte que chaque symbole figure une fois par ligne et une fois par colonne
avaient déjà été étudiés par Euler en 1782 sous le nom de \textit{carrés latins}, et l'écart entre les deux
notions est instructif : un carré latin encode la simplification, mais l'associativité ne laisse aucune trace
visible dans la grille --- c'est pourquoi la réciproque est fausse, et pourquoi les objets qui ne satisfont que
la condition sur les lignes et les colonnes (les quasi-groupes) forment un sujet à part.

\begin{refs}
\item Contexte : Table de Cayley --- Wikipédia (\url{https://fr.wikipedia.org/wiki/Table_de_Cayley})
\item Contexte : Carré latin --- Wikipédia (\url{https://fr.wikipedia.org/wiki/Carr%C3%A9_latin})
\end{refs}

\bigskip


\section*{Algèbre linéaire}

\begin{problem}[{Deux droites, une réponse --- le plus souvent --- Lycée}]
\textbf{LIN-01.} Considérons le système de deux équations
$ax + by = e$, $cx + dy = f$, où $a, b, c, d, e, f$ sont des nombres réels donnés.
\begin{enumerate}
\item[(i)] Démontrer que le système possède exactement une solution $(x, y)$ si et seulement si
$ad - bc \ne 0$.
\item[(ii)] Lorsque $ad - bc = 0$, déterminer --- avec démonstration --- précisément quand le système
n'a aucune solution et quand il en a une infinité.
\item[(iii)] Interpréter géométriquement chaque cas comme un
énoncé portant sur deux droites du plan.
\end{enumerate}

\end{problem}

\noindent Résoudre des systèmes d'équations est plus ancien que l'algèbre elle-même : le
chapitre 8 du classique chinois \textit{Les Neuf Chapitres sur l'art mathématique} (compilé au
Ier siècle de notre ère) résout de tels systèmes en manipulant le tableau des
coefficients --- le pivot de Gauss, dix-huit siècles avant Gauss, qui le réinventa en 1809 pour
suivre l'astéroïde Cérès. La quantité $ad - bc$ qui fait sa première apparition en (i) est le
déterminant, dont cette page suit la longue carrière.

\begin{refs}
\item Contexte : Système d'équations linéaires --- Wikipédia (\url{https://fr.wikipedia.org/wiki/Syst%C3%A8me_d%27%C3%A9quations_lin%C3%A9aires})
\item Contexte : Les Neuf Chapitres sur l'art mathématique --- Wikipédia (\url{https://fr.wikipedia.org/wiki/Les_Neuf_Chapitres_sur_l%27art_math%C3%A9matique})
\item Lecture : G. Strang, \textit{Introduction to Linear Algebra}, ch. 1--2
\end{refs}

\bigskip

\begin{problem}[{Les rotations vues comme des matrices --- Lycée}]
\textbf{LIN-02.} \begin{enumerate}
\item[(i)] Montrer que la rotation du plan autour de l'origine d'un angle
$\theta$ envoie le point $(x, y)$ sur
$(x \cos\theta - y \sin\theta, \; x \sin\theta + y \cos\theta)$.
\item[(ii)] Démontrer qu'effectuer
la rotation d'angle $\alpha$ puis la rotation d'angle $\beta$ donne la rotation d'angle
$\alpha + \beta$, et en déduire les formules d'addition du sinus et du cosinus.
\item[(iii)] Une réflexion
du plan par rapport à une droite passant par l'origine déplace elle aussi les points linéairement.
Trouver sa formule en fonction de l'angle de la droite, et démontrer que la composée de deux
réflexions est une rotation.
\end{enumerate}

\end{problem}

\noindent Voici la porte par laquelle la géométrie entre dans l'algèbre : les transformations
deviennent des tableaux de nombres, et la composition devient une multiplication. La notation
matricielle et la règle de multiplication ont été fixées par Cayley dans son \textit{Memoir on the Theory
of Matrices} de 1858, précisément pour rendre automatique une comptabilité comme celle de (ii).
La surprise agréable que les identités trigonométriques soient des théorèmes \textit{sur les matrices}
est, pour bien des étudiants, le premier indice qu'une notation peut penser à votre place.

\begin{refs}
\item Contexte : Matrice de rotation --- Wikipédia (\url{https://fr.wikipedia.org/wiki/Matrice_de_rotation})
\item Contexte : Matrice d'une application linéaire --- Wikipédia (\url{https://fr.wikipedia.org/wiki/Matrice_d%27une_application_lin%C3%A9aire})
\item Lecture : G. Strang, \textit{Introduction to Linear Algebra}, ch. 2
\end{refs}

\bigskip

\begin{problem}[{Le déterminant est une aire --- Lycée}]
\textbf{LIN-03.} Soient $u = (a, b)$ et $v = (c, d)$ deux vecteurs du
plan.
\begin{enumerate}
\item[(i)] Démontrer que l'aire du parallélogramme de côtés $u$ et $v$ vaut
$|ad - bc|$.
\item[(ii)] Montrer que le signe de $ad - bc$ enregistre l'orientation : il est positif
exactement lorsque le plus court tour de $u$ vers $v$ se fait dans le sens direct.
\item[(iii)] Utiliser (i) et (ii)
pour démontrer la \og{} formule du lacet \fg{} : l'aire d'un polygone de sommets
$(x_1, y_1), \dots, (x_n, y_n)$, énumérés dans le sens direct, vaut
$\frac{1}{2} \sum_{i=1}^{n} (x_i y_{i+1} - x_{i+1} y_i)$, les indices étant pris modulo $n$.
\end{enumerate}

\end{problem}

\noindent Les déterminants précèdent les matrices : Seki Takakazu au Japon et Leibniz en
Europe les rencontrèrent tous deux en 1683, comme conditions de compatibilité d'un système
d'équations. Que la même expression mesure une aire --- et, en dimension supérieure, un volume --- est
l'âme géométrique de la théorie, développée par Lagrange et Cauchy. La formule du lacet, outil de
base des arpenteurs, est habituellement attribuée à Meister (1769) et à Gauss ; c'est le déterminant
faisant honnêtement du travail de plein air.

\begin{refs}
\item Contexte : Déterminant (mathématiques) --- Wikipédia (\url{https://fr.wikipedia.org/wiki/D%C3%A9terminant_%28math%C3%A9matiques%29})
\item Contexte : Formule du lacet --- Wikipédia (en anglais) (\url{https://en.wikipedia.org/wiki/Shoelace_formula})
\item Lecture : G. Strang, \textit{Introduction to Linear Algebra}, ch. 5
\end{refs}

\bigskip

\begin{problem}[{Trace et déterminant cachés dans un polynôme --- Lycée, short}]
\textbf{LIN-22.} Pour la matrice
$A = \begin{pmatrix} a & b \\ c & d \end{pmatrix}$, démontrer que
\[ \det(A - \lambda I) = \lambda^2 - (a + d)\lambda + (ad - bc). \]
En déduire que si $A$ possède deux valeurs propres réelles, leur somme est la trace $a + d$ et
leur produit le déterminant $ad - bc$.
\end{problem}

\noindent Les deux nombres que tout le monde apprend à calculer pour une matrice
$2 \times 2$ se révèlent être les seuls coefficients de son polynôme caractéristique --- ce sont donc
exactement les données qu'une valeur propre peut voir. Cayley et Sylvester, qui baptisèrent la trace
et le déterminant dans les années 1850, cherchaient précisément de telles quantités : des expressions
des coefficients qui ne changent pas lorsque la matrice est réécrite dans une autre base. En
dimension $n$, le même polynôme a $n$ coefficients, et ce sont les fonctions symétriques
élémentaires des valeurs propres.

\begin{refs}
\item Contexte : Polynôme caractéristique --- Wikipédia (\url{https://fr.wikipedia.org/wiki/Polyn%C3%B4me_caract%C3%A9ristique})
\item Contexte : Trace (algèbre) --- Wikipédia (\url{https://fr.wikipedia.org/wiki/Trace_%28alg%C3%A8bre%29})
\end{refs}

\bigskip

\begin{problem}[{Des matrices dont le carré est nul --- Lycée, short}]
\textbf{LIN-23.} Déterminer, avec démonstration, toutes les matrices réelles
$2 \times 2$ $N$ vérifiant $N^2 = 0$. Montrer qu'une telle matrice $N$ non nulle écrase le
plan tout entier sur une droite passant par l'origine, puis écrase cette droite sur l'origine.
\end{problem}

\noindent Un nombre non nul ne peut avoir un carré nul, mais une matrice non nulle le peut ---
premier signe que les matrices sont une arithmétique véritablement neuve plutôt que des nombres
déguisés. De telles matrices sont dites \textit{nilpotentes} ; le plus petit exemple,
$\begin{pmatrix} 0 & 1 \\ 0 & 0 \end{pmatrix}$, est le cisaillement qui pousse le plan de
côté, et c'est l'unique brique élémentaire dont est assemblée la réduction de Jordan de LIN-14.

\begin{refs}
\item Contexte : Matrice nilpotente --- Wikipédia (\url{https://fr.wikipedia.org/wiki/Matrice_nilpotente})
\item Lecture : G. Strang, \textit{Introduction to Linear Algebra}, ch. 2
\end{refs}

\bigskip

\begin{problem}[{Combien y a-t-il de carrés magiques, vraiment ? --- Lycée}]
\textbf{LIN-24.} Appelons \textit{carré magique} un tableau $3 \times 3$ de nombres réels dont les
trois sommes par ligne, les trois sommes par colonne et les deux sommes diagonales sont toutes égales
à un même nombre $S$, la somme magique. Les coefficients peuvent être des réels quelconques,
positifs ou non.
\begin{enumerate}
\item[(i)] Démontrer que l'ensemble $M$ de tous les carrés magiques $3 \times 3$ est un espace
vectoriel : une somme de carrés magiques est magique, et tout multiple scalaire aussi.
\item[(ii)] Démontrer que le coefficient central d'un carré magique vaut toujours $S/3$.
\item[(iii)] Déterminer $\dim M$, avec démonstration, et écrire explicitement une base. (De façon
équivalente : combien de coefficients peut-on choisir librement avant que les autres soient
imposés ?)
\item[(iv)] Démontrer que tout carré magique dont les neuf coefficients sont exactement les entiers
$1, 2, \dots, 9$ s'obtient à partir du carré Lo Shu
$\begin{pmatrix} 4 & 9 & 2 \\ 3 & 5 & 7 \\ 8 & 1 & 6 \end{pmatrix}$
par rotation ou réflexion de la grille --- il y en a donc exactement huit.
\end{enumerate}

\end{problem}

\noindent Les carrés magiques sont anciens --- le motif Lo Shu apparaît dans des textes chinois
vieux de plus de deux mille ans, et Dürer en plaça un $4 \times 4$ dans \textit{Melencolia I} (1514),
datant la gravure dans sa dernière ligne. Pendant presque toute cette histoire, on les a traités
comme des curiosités. La lecture par l'algèbre linéaire change entièrement la question : les huit
conditions magiques sont des équations linéaires homogènes, donc les carrés magiques forment le
noyau d'une application linéaire, et demander \og{} combien y en a-t-il \fg{} revient à demander une
dimension. La partie (iv) est le dénombrement classique, et il est frappant que la réponse soit huit
plutôt que huit mille.

\begin{refs}
\item Contexte : Carré magique (mathématiques) --- Wikipédia (\url{https://fr.wikipedia.org/wiki/Carr%C3%A9_magique_%28math%C3%A9matiques%29})
\item Contexte : Noyau (algèbre) --- Wikipédia (\url{https://fr.wikipedia.org/wiki/Noyau_%28alg%C3%A8bre%29})
\item Lecture : G. Strang, \textit{Introduction to Linear Algebra}, ch. 3
\end{refs}

\bigskip

\begin{problem}[{Le produit vectoriel, et le volume qu'il mesure --- Lycée}]
\textbf{LIN-34.} Pour des vecteurs $u = (u_1, u_2, u_3)$ et $v = (v_1, v_2, v_3)$ de l'espace, on pose
\[ u \times v = (u_2 v_3 - u_3 v_2,\; u_3 v_1 - u_1 v_3,\; u_1 v_2 - u_2 v_1). \]
\begin{enumerate}
\item[(a)] Démontrer que $u \times v$ est orthogonal à la fois à $u$ et à $v$, et démontrer l'identité
de Lagrange
\[ |u \times v|^2 = |u|^2 |v|^2 - (u \cdot v)^2 . \]
En déduire que $|u \times v| = |u|\,|v| \sin\theta$, l'aire du parallélogramme engendré par $u$
et $v$, et que $u \times v = 0$ exactement lorsque $u$ et $v$ sont colinéaires.
\item[(b)] Démontrer que le produit mixte $u \cdot (v \times w)$ est égal au déterminant de la matrice
$3 \times 3$ de lignes $u, v, w$, que sa valeur absolue est le volume du parallélépipède qu'ils
engendrent, et qu'il change de signe lorsqu'on échange deux des trois vecteurs.
\item[(c)] Démontrer que le produit vectoriel n'est pas associatif, en trouvant $u, v, w$ tels que
$(u \times v) \times w \ne u \times (v \times w)$. Démontrer ensuite le développement
$u \times (v \times w) = v\,(u \cdot w) - w\,(u \cdot v)$ et l'identité de Jacobi
\[ u \times (v \times w) + v \times (w \times u) + w \times (u \times v) = 0 . \]
\item[(d)] En déduire que trois vecteurs appartiennent à un même plan passant par l'origine si et seulement
si leur produit mixte est nul, et donner une explication purement géométrique de la raison pour
laquelle le déterminant de la partie (b) devait être un volume.
\end{enumerate}

\end{problem}

\noindent Le produit vectoriel est un fragment d'une structure plus vaste. En 1843, Hamilton
inventa les quaternions et, avec eux, les mots \textit{vecteur} et \textit{scalaire} ; multiplier deux
quaternions \og{} purs \fg{} produit à la fois le produit scalaire et le produit vectoriel, l'un comme partie
scalaire et l'autre comme partie vectorielle. Grassmann disposait dès 1844 d'un produit plus général,
et fut ignoré. Ce sont Gibbs, dans des notes de cours imprimées à compte d'auteur en 1881, et
Heaviside indépendamment, qui scindèrent le produit de Hamilton en les deux opérations familières ---
provoquant une querelle publique d'une décennie avec les quaternionistes --- et c'est E. B. Wilson,
étudiant de Gibbs, qui transforma ces notes en le manuel \textit{Vector Analysis} de 1901, lequel fixa
la notation encore en usage aujourd'hui. Une curiosité qui mérite d'être connue : un produit
bilinéaire de deux vecteurs obéissant aux règles de (a) n'existe qu'en dimensions trois et sept,
celui de dimension sept provenant des octonions.

\begin{refs}
\item Contexte : Produit vectoriel --- Wikipédia (\url{https://fr.wikipedia.org/wiki/Produit_vectoriel})
\item Contexte : Produit mixte --- Wikipédia (\url{https://fr.wikipedia.org/wiki/Produit_mixte})
\item Lecture : E. B. Wilson, \textit{Vector Analysis} (1901), d'après les cours de J. W. Gibbs, ch. 2
\item Lecture : M. J. Crowe, \textit{A History of Vector Analysis}
\end{refs}

\bigskip

\begin{problem}[{Des lignes dont la somme vaut un --- Lycée, short}]
\textbf{LIN-35.} Soit $P$ une matrice $n \times n$ à coefficients positifs
ou nuls dont chaque ligne a pour somme $1$. Démontrer que le vecteur dont toutes les coordonnées
valent un est vecteur propre de $P$ pour la valeur propre $1$, et démontrer que toute valeur
propre réelle $\lambda$ de $P$ vérifie $|\lambda| \le 1$.
\end{problem}

\noindent Une telle matrice est la table de transition d'une marche aléatoire : le
coefficient $(i,j)$ est la probabilité de sauter de l'état $i$ à l'état $j$. Andreï Markov
introduisit ces chaînes en 1906 pour montrer que la loi des grands nombres survit à la dépendance
entre les épreuves, et il éprouva sa théorie en comptant l'alternance des voyelles et des consonnes
dans \textit{Eugène Onéguine} de Pouchkine. La valeur propre $1$ que vous trouvez ici est la raison
pour laquelle une longue marche aléatoire finit par se stabiliser sur une distribution fixe, et la
version renforcée de cet énoncé --- Perron--Frobenius --- est ce qui classe les pages web en LIN-15.

\begin{refs}
\item Contexte : Matrice stochastique --- Wikipédia (\url{https://fr.wikipedia.org/wiki/Matrice_stochastique})
\item Contexte : Chaîne de Markov --- Wikipédia (\url{https://fr.wikipedia.org/wiki/Cha%C3%AEne_de_Markov})
\end{refs}

\bigskip


\section*{Théorie des nombres}

\begin{problem}[{Irrationalité de $\surd$2 --- Lycée}]
\textbf{NT-01.} Démontrer que $\sqrt{2}$ est irrationnel : il n'existe pas
d'entiers strictement positifs $p, q$ tels que $p^2 = 2q^2$. Déterminer ensuite, avec
démonstration, exactement quels entiers strictement positifs $n$ ont $\sqrt{n}$ irrationnel.
\end{problem}

\noindent Voici la découverte qui brisa le credo pythagoricien selon lequel tout est nombre
entier et rapport : la diagonale d'un carré est incommensurable avec son côté. La légende ---
rapportée des siècles plus tard et dont il ne faut pas croire le détail --- veut que le découvreur,
parfois nommé Hippase de Métaponte (Ve siècle av. J.-C.), ait été noyé en mer pour l'avoir
révélée. Théodore de Cyrène démontra plus tard l'irrationalité de $\sqrt{3}, \sqrt{5}, \ldots$
jusqu'à $\sqrt{17}$, et les \textit{Éléments} d'Euclide conservent l'argument classique. C'est
l'archétype du raisonnement par l'absurde et le premier théorème de ce que nous appelons aujourd'hui
la théorie des nombres.

\begin{refs}
\item Contexte : Racine carrée de deux --- Wikipédia (\url{https://fr.wikipedia.org/wiki/Racine_carr%C3%A9e_de_deux})
\item Contexte : Nombre irrationnel --- Wikipédia (\url{https://fr.wikipedia.org/wiki/Nombre_irrationnel})
\item Lecture : G. H. Hardy et E. M. Wright, \textit{An Introduction to the Theory of Numbers}, ch. 4
\end{refs}

\bigskip

\begin{problem}[{Euclide : une infinité de nombres premiers --- Lycée}]
\textbf{NT-02.} Démontrer qu'il existe une infinité de nombres premiers.
Examiner ensuite les nombres $N_k = p_1 p_2 \cdots p_k + 1$, où
$p_1 < p_2 < \cdots < p_k$ sont les $k$ premiers nombres premiers : $N_k$ doit-il
lui-même être premier ? Trancher la question par une démonstration ou par le plus petit
contre-exemple.
\end{problem}

\noindent Proposition 20 du livre IX des \textit{Éléments} d'Euclide (v. 300 av. J.-C.) :
\og{} les nombres premiers sont plus nombreux que toute multitude assignée de nombres premiers \fg{}.
L'argument d'Euclide lui-même est direct et n'est pas --- contrairement à ce qu'on retient souvent ---
un raisonnement par l'absurde, et on le cite couramment comme le modèle de ce que doit être une
démonstration. La seconde partie prémunit contre le contresens le plus répandu sur cet argument. On
connaît aujourd'hui des dizaines d'autres démonstrations --- par les nombres de Fermat (Goldbach), par
la topologie (Furstenberg), par la divergence de $\sum 1/p$ (Euler) --- et les comparer est une
leçon sur le nombre de portes que peut avoir une seule vérité.

\begin{refs}
\item Contexte : Théorème d'Euclide sur les nombres premiers --- Wikipédia (\url{https://fr.wikipedia.org/wiki/Th%C3%A9or%C3%A8me_d%27Euclide_sur_les_nombres_premiers})
\item Lecture : M. Aigner et G. M. Ziegler, \textit{Proofs from THE BOOK}, ch. 1 (six démonstrations)
\item Lecture : G. H. Hardy et E. M. Wright, \textit{An Introduction to the Theory of Numbers}, ch. 1--2
\end{refs}

\bigskip

\begin{problem}[{L'algorithme d'Euclide et le théorème de Bachet-Bézout --- Lycée}]
\textbf{NT-03.} Démontrer que l'algorithme d'Euclide --- qui remplace
répétitivement le couple $(a, b)$ par $(b,\ a \bmod b)$ --- termine et calcule $\gcd(a, b)$.
Démontrer l'identité de Bézout : $\gcd(a, b)$ est le plus petit entier strictement positif qui
s'exprime sous la forme $ax + by$ avec $x, y$ entiers. En déduire le lemme d'Euclide --- si un
nombre premier $p$ divise $ab$ alors $p$ divise $a$ ou $p$ divise $b$ --- et, de là,
l'unicité de la décomposition en facteurs premiers.
\end{problem}

\noindent L'algorithme figure au livre VII des \textit{Éléments} et on l'appelle souvent le
plus ancien algorithme non trivial encore d'usage quotidien. L'identité porte le nom d'Étienne
Bézout, qui démontra l'analogue polynomial au XVIIIe siècle ; pour les entiers, elle fut
publiée par Bachet de Méziriac en 1624. La chaîne de déductions demandée ici --- algorithme, identité,
lemme, factorisation unique --- est la colonne vertébrale logique de la théorie élémentaire des
nombres, bien que le théorème fondamental de l'arithmétique n'ait été énoncé et démontré avec toute
la rigueur voulue que par Gauss dans les \textit{Disquisitiones Arithmeticae} (1801).

\begin{refs}
\item Contexte : Algorithme d'Euclide --- Wikipédia (\url{https://fr.wikipedia.org/wiki/Algorithme_d%27Euclide})
\item Contexte : Théorème de Bachet-Bézout --- Wikipédia (\url{https://fr.wikipedia.org/wiki/Th%C3%A9or%C3%A8me_de_Bachet-B%C3%A9zout})
\item Contexte : Théorème fondamental de l'arithmétique --- Wikipédia (\url{https://fr.wikipedia.org/wiki/Th%C3%A9or%C3%A8me_fondamental_de_l%27arithm%C3%A9tique})
\item Lecture : J. H. Silverman, \textit{A Friendly Introduction to Number Theory}, ch. 5--7
\end{refs}

\bigskip

\begin{problem}[{Congruences et preuve par neuf --- Lycée}]
\textbf{NT-04.} Pour des entiers $a, b$ et $m \ge 1$, on écrit
$a \equiv b \pmod{m}$ lorsque $m$ divise $a - b$. Démontrer que la congruence respecte
l'addition et la multiplication. En déduire les critères classiques sur les chiffres : $9$ divise
$n$ exactement lorsque $9$ divise la somme des chiffres de $n$, et $11$ divise $n$
exactement lorsque $11$ divise la somme alternée des chiffres. Expliquer précisément quelles
erreurs de calcul la vérification médiévale de la \og{} preuve par neuf \fg{} détecte, et en exhiber une
qu'elle laisse passer.
\end{problem}

\noindent L'arithmétique modulo m se pratiquait bien avant d'être nommée : la preuve par
neuf apparaît dans l'arithmétique indienne et islamique il y a plus de mille ans, comme contrôle de
comptable sur un calcul à la main. La notation moderne et la théorie systématique ouvrent les
\textit{Disquisitiones Arithmeticae} de Gauss (1801), écrites alors qu'il n'avait qu'une vingtaine
d'années ; son signe de congruence $\equiv$ transforma des astuces éparses sur les chiffres en un calcul.
Presque tout ce qui figure sur cette page parle cette langue.

\begin{refs}
\item Contexte : Arithmétique modulaire --- Wikipédia (\url{https://fr.wikipedia.org/wiki/Arithm%C3%A9tique_modulaire})
\item Contexte : Preuve par neuf --- Wikipédia (\url{https://fr.wikipedia.org/wiki/Preuve_par_neuf})
\item Lecture : H. Davenport, \textit{The Higher Arithmetic}, ch. II
\end{refs}

\bigskip

\begin{problem}[{Les premiers $\equiv$ 3 (mod 4), et pourquoi 1 (mod 4) est plus difficile --- Lycée}]
\textbf{NT-05.} Démontrer que tout entier $n \equiv 3 \pmod 4$ possède un
facteur premier $\equiv 3 \pmod 4$. En déduire une adaptation de l'argument d'Euclide --- considérer
$4p_1 p_2 \cdots p_k - 1$ --- et démontrer qu'il existe une infinité de nombres premiers
$\equiv 3 \pmod 4$. Essayer ensuite de mener le même argument pour les premiers
$\equiv 1 \pmod 4$, et identifier précisément l'étape où il s'effondre.
\end{problem}

\noindent Premier avant-goût des nombres premiers en progression arithmétique. L'astuce
d'Euclide se généralise à quelques progressions puis cale, et ce blocage est instructif : un produit
de nombres $\equiv 1 \pmod 4$ est encore $\equiv 1 \pmod 4$, donc rien ne force l'apparition d'un
facteur premier de la bonne espèce. La classe résiduelle 1 (mod 4) exige véritablement une idée
nouvelle --- les résidus quadratiques --- et le problème NT-12 achève cette histoire. Le théorème
complet, selon lequel toute progression $a, a+m, a+2m, \ldots$ avec $\gcd(a, m) = 1$ contient
une infinité de nombres premiers, fut démontré par Dirichlet en 1837, qui inventa pour cela la
théorie analytique des nombres.

\begin{refs}
\item Contexte : Théorème de la progression arithmétique --- Wikipédia (\url{https://fr.wikipedia.org/wiki/Th%C3%A9or%C3%A8me_de_la_progression_arithm%C3%A9tique})
\item Contexte : Théorème d'Euclide sur les nombres premiers --- Wikipédia (\url{https://fr.wikipedia.org/wiki/Th%C3%A9or%C3%A8me_d%27Euclide_sur_les_nombres_premiers})
\item Lecture : H. Davenport, \textit{The Higher Arithmetic}, ch. I--II
\end{refs}

\bigskip

\begin{problem}[{Le petit théorème de Fermat --- Lycée}]
\textbf{NT-06.} Soit $p$ un nombre premier. Démontrer que
$a^p \equiv a \pmod p$ pour tout entier $a$, et donc que
$a^{p-1} \equiv 1 \pmod p$ lorsque $p$ ne divise pas $a$. Donner deux démonstrations
différentes : l'une par récurrence sur $a$ à l'aide de la formule du binôme, l'autre en montrant
que la multiplication par $a$ permute les résidus non nuls modulo $p$. En guise d'application
d'échauffement, trouver le reste de $3^{1000}$ dans la division par $101$.
\end{problem}

\noindent Fermat énonça le théorème dans une lettre à Frénicle de Bessy datée du
18 octobre 1640, en ajoutant qu'il en enverrait la démonstration \og{} si je ne craignais d'être trop
long \fg{}. Euler publia la première démonstration en 1736 (Leibniz en avait une plus tôt, non publiée).
Le théorème est le moteur de presque tous les tests de primalité modernes --- les problèmes NT-19,
NT-21 et NT-23 ci-dessous en sont les descendants --- et, via sa généralisation par Euler (NT-11), il
fonde le cryptosystème RSA. Une petite congruence à très longue portée.

\begin{refs}
\item Contexte : Petit théorème de Fermat --- Wikipédia (\url{https://fr.wikipedia.org/wiki/Petit_th%C3%A9or%C3%A8me_de_Fermat})
\item Lecture : J. H. Silverman, \textit{A Friendly Introduction to Number Theory}, ch. 9
\item Cours : MIT OCW 18.781 Theory of Numbers (en anglais) (\url{https://ocw.mit.edu/courses/18-781-theory-of-numbers-spring-2012/})
\end{refs}

\bigskip

\begin{problem}[{Les triplets pythagoriciens --- Lycée}]
\textbf{NT-07.} Un triplet pythagoricien $(a, b, c)$ d'entiers strictement
positifs vérifie $a^2 + b^2 = c^2$, et il est primitif si $\gcd(a, b, c) = 1$.
Démontrer la paramétrisation d'Euclide : les triplets primitifs sont exactement
$(m^2 - n^2,\ 2mn,\ m^2 + n^2)$ avec $m > n \ge 1$ premiers entre eux et de parités
opposées, à l'échange de $a$ et $b$ près. Démontrer ensuite que dans tout triplet pythagoricien,
$3$ divise l'un des côtés de l'angle droit, $4$ divise l'un d'eux, et $5$ divise l'un des
côtés.
\end{problem}

\noindent La tablette babylonienne Plimpton 322 (v. 1800 av. J.-C.) énumère des nombres que
la plupart des spécialistes lisent comme des triplets pythagoriciens, compilés plus d'un millénaire
avant Pythagore. La paramétrisation figure au livre X des \textit{Éléments} d'Euclide. C'est la
première équation diophantienne résolue de façon non triviale, et le modèle de toutes celles qui
suivirent : Fermat griffonna son grand théorème dans la marge, à côté exactement de ce problème,
dans son exemplaire de Diophante, et le problème des nombres congruents (NT-29) commence par les
triangles rationnels construits à partir de ces triplets.

\begin{refs}
\item Contexte : Triplet pythagoricien --- Wikipédia (\url{https://fr.wikipedia.org/wiki/Triplet_pythagoricien})
\item Contexte : Plimpton 322 --- Wikipédia (\url{https://fr.wikipedia.org/wiki/Plimpton_322})
\item Lecture : J. H. Silverman, \textit{A Friendly Introduction to Number Theory}, ch. 2--3
\end{refs}

\bigskip

\begin{problem}[{La fraction continue de $\surd$2 --- Lycée}]
\textbf{NT-08.} Montrer que
\[ \sqrt{2} \;=\; 1 + \cfrac{1}{2 + \cfrac{1}{2 + \cfrac{1}{2 + \cdots}}}, \]
la fraction continue infinie dont tous les quotients partiels valent $2$ après le premier.
Calculer les premières réduites
$1,\ \tfrac{3}{2},\ \tfrac{7}{5},\ \tfrac{17}{12},\ \tfrac{41}{29}, \ldots$ et démontrer
que toute réduite $p/q$ vérifie $p^2 - 2q^2 = \pm 1$ --- les réduites fournissent donc une
infinité de solutions approchées oscillant autour de $\sqrt{2}$.
\end{problem}

\noindent Les Grecs connaissaient ces fractions sous le nom de \og{} nombres latéraux et
diagonaux \fg{} (Théon de Smyrne, IIe siècle de notre ère) : la récurrence qui les engendre
était leur façon d'approcher par des rapports la diagonale incommensurable de NT-01. La théorie
systématique des fractions continues vint avec Euler et Lagrange au XVIIIe siècle, et
l'identité $p^2 - 2q^2 = \pm 1$ est la porte d'entrée de l'équation de Pell (NT-14). Les fractions
continues reviennent deux fois plus bas : pour le nombre e (NT-26) et comme moteur de la
transcendance (NT-25).

\begin{refs}
\item Contexte : Fraction continue --- Wikipédia (\url{https://fr.wikipedia.org/wiki/Fraction_continue})
\item Contexte : Suite de Pell --- Wikipédia (\url{https://fr.wikipedia.org/wiki/Suite_de_Pell})
\item Lecture : H. Davenport, \textit{The Higher Arithmetic}, ch. IV
\end{refs}

\bigskip

\begin{problem}[{Le dernier chiffre d'une tour --- Lycée, short}]
\textbf{NT-36.} Déterminer le dernier chiffre de $ 7^{100} $, et démontrer votre réponse sans calculer le nombre.
\end{problem}

\noindent Le dernier chiffre d'une puissance est son résidu modulo $ 10 $, et les puissances de $ 7 $ sont cycliques de période quatre --- toute la question se ramène donc à $ 100 \bmod 4 $. C'est le plus petit avant-goût possible de l'ordre d'un élément dans un groupe, une idée qui deviendra plus tard le théorème de Lagrange et le théorème d'Euler.

\begin{refs}
\item Contexte : Arithmétique modulaire --- Wikipédia (\url{https://fr.wikipedia.org/wiki/Arithm%C3%A9tique_modulaire})
\end{refs}

\bigskip

\begin{problem}[{Pourquoi la preuve par neuf fonctionne --- Lycée, short}]
\textbf{NT-37.} Démontrer qu'un entier strictement positif est divisible par $ 9 $ si et seulement si la somme de ses chiffres décimaux est divisible par $ 9 $. Énoncer et démontrer ensuite la règle analogue pour $ 11 $.
\end{problem}

\noindent On enseigne cette règle à tous les écoliers et on n'en montre presque à aucun la raison ; la démonstration tient en trois lignes dès qu'on écrit le nombre sous la forme $ \sum d_k 10^k $ et qu'on remarque que $ 10 \equiv 1 \pmod 9 $. Le cas $ 11 $ oblige à traiter $ 10 \equiv -1 $, ce qui explique l'alternance des signes des chiffres --- une petite mais authentique surprise.

\begin{refs}
\item Contexte : Critère de divisibilité --- Wikipédia (\url{https://fr.wikipedia.org/wiki/Crit%C3%A8re_de_divisibilit%C3%A9})
\end{refs}

\bigskip

\begin{problem}[{Une infinité de premiers $\equiv$ 5 (mod 6) --- Lycée, short}]
\textbf{NT-44.} Démontrer que tout entier $ n \equiv 5 \pmod 6 $ possède un
facteur premier $ \equiv 5 \pmod 6 $. En déduire, à la manière d'Euclide, qu'il existe une infinité
de nombres premiers congrus à $ 5 $ modulo $ 6 $.
\end{problem}

\noindent C'est le pendant de NT-05, et il fonctionne pour la même raison : la classe
résiduelle $ 5 \pmod 6 $ n'est stable par rien d'utile, si bien qu'un nombre qui s'y trouve est
contraint d'exposer un premier de sa propre espèce. L'astuce d'Euclide atteint ainsi une poignée de
progressions puis s'arrête ; l'énoncé général --- toute progression $ a, a+m, a+2m, \ldots $ avec
$ \gcd(a,m)=1 $ contient une infinité de nombres premiers --- est le théorème de Dirichlet de 1837,
et NT-50 plus bas montre ce qu'il en coûte de le démontrer.

\begin{refs}
\item Contexte : Théorème de la progression arithmétique --- Wikipédia (\url{https://fr.wikipedia.org/wiki/Th%C3%A9or%C3%A8me_de_la_progression_arithm%C3%A9tique})
\item Contexte : Théorème d'Euclide sur les nombres premiers --- Wikipédia (\url{https://fr.wikipedia.org/wiki/Th%C3%A9or%C3%A8me_d%27Euclide_sur_les_nombres_premiers})
\end{refs}

\bigskip

\begin{problem}[{La conjecture de Catalan, et ce qui l'a suivie --- Lycée}]
\textbf{NT-45.} Appelons $ m $ une \textit{puissance parfaite} si $ m = a^k $ pour des entiers
$ a \ge 2 $ et $ k \ge 2 $. En 1844, Eugène Catalan conjectura que $ 8 $ et $ 9 $
sont les deux seules puissances parfaites consécutives.
\begin{enumerate}
\item[(a)] Trouver, avec démonstration, toutes les solutions en entiers positifs ou nuls de
\[ \left| 2^{a} - 3^{b} \right| = 1 . \]
(Travailler modulo de petits nombres ; $ 8 $ et $ 9 $ constituent l'une des solutions, et il faut
montrer qu'il n'y en a pas d'autres.)
\item[(b)] Démontrer que pour chaque $ k \ge 2 $ fixé, deux puissances $ k $-ièmes distinctes d'entiers
strictement positifs ne diffèrent jamais de $ 1 $ ; en particulier $ x^2 - y^2 = 1 $ n'a pas de
solution avec $ x > y \ge 1 $.
\item[(c)] Énoncer précisément le théorème de Mih\u{a}ilescu, puis énoncer la conjecture de Pillai --- selon
laquelle tout entier strictement positif n'apparaît qu'un nombre fini de fois comme différence de
deux puissances parfaites. Dire exactement lequel des deux est un théorème et lequel est ouvert.
\item[(d)] Rechercher l'énoncé du théorème de Tijdeman de 1976 sur l'équation de Catalan et expliquer, avec
vos propres mots, précisément ce qu'il a établi et précisément ce qu'il a laissé à faire à
Mih\u{a}ilescu.
\end{enumerate}

\end{problem}

\noindent Catalan posa la question dans une note au journal de Crelle en 1844 et n'y revint
jamais. Le premier progrès véritable est bien plus ancien : Lévi ben Gerson (Gersonide), écrivant en
Provence en 1343, démontra exactement la partie (a) --- que $ 8 $ et $ 9 $ sont les seules
puissances consécutives de $ 2 $ et de $ 3 $ --- à la demande, dit-on, d'un musicien qui voulait
savoir pourquoi les rapports $ 9/8 $ et $ 3/2 $ se comportent comme ils le font. Lebesgue régla
le cas $ x^p - y^2 = 1 $ en 1850 ; Tijdeman démontra en 1976 qu'il n'y a en tout qu'un nombre fini
de solutions, mais avec une borne beaucoup trop grande pour être vérifiée ; et en avril 2002 Preda
Mih\u{a}ilescu combla l'écart par un argument dans les corps cyclotomiques, publié au journal de Crelle
en 2004 --- 158 ans après que la question fut posée. La généralisation, la conjecture de Pillai de
1931, est toujours ouverte, et découlerait de la conjecture abc (NT-35) : on croit que les écarts
entre puissances parfaites consécutives croissent, et personne ne sait le démontrer.

\begin{refs}
\item Contexte : Conjecture de Catalan --- Wikipédia (\url{https://fr.wikipedia.org/wiki/Conjecture_de_Catalan})
\item Contexte : Puissance parfaite --- Wikipédia (\url{https://fr.wikipedia.org/wiki/Puissance_parfaite})
\item Article : R. Tijdeman, \og{} On the equation of Catalan \fg{}, \textit{Acta Arithmetica} 29 (1976), 197--209
\item Article : P. Mih\u{a}ilescu, \og{} Primary cyclotomic units and a proof of Catalan's conjecture \fg{}, \textit{Journal für die reine und angewandte Mathematik} (2004)
\end{refs}

\bigskip


\section*{Analyse réelle et théorie de la mesure}

\begin{problem}[{Bornes supérieures --- Lycée}]
\textbf{RA-01.} Pour un ensemble $ S $ de nombres réels, un \textit{majorant}
est un nombre $ b $ tel que $ x \le b $ pour tout $ x \in S $ ; la \textit{borne supérieure}
$ \sup S $ est le plus petit des majorants, et la \textit{borne inférieure} $ \inf S $ le plus
grand des minorants. Déterminez, avec démonstration, la borne supérieure et la borne inférieure de
chacun des ensembles suivants, et décidez dans chaque cas si elles appartiennent à l'ensemble :
\begin{enumerate}
\item[(a)] $ \{ 1/n : n = 1, 2, 3, \dots \} $ ;
\item[(b)] $ \{ n/(n+1) : n = 1, 2, 3, \dots \} $ ;
\item[(c)] $ \{ x \in \mathbb{Q} : x > 0 \text{ et } x^2 < 2 \} $.
Pour (c), expliquez pourquoi la borne supérieure existe dans $ \mathbb{R} $ alors que l'ensemble
n'a pas de borne supérieure à l'intérieur de $ \mathbb{Q} $.
\end{enumerate}

\end{problem}

\noindent L'\textit{axiome de la borne supérieure} --- toute partie non vide et majorée de $\mathbb{R}$
admet un plus petit majorant --- est l'unique propriété qui sépare $\mathbb{R}$ de $\mathbb{Q}$, et toute l'analyse repose
sur elle. Richard Dedekind l'a isolée en 1858 en préparant un cours de calcul différentiel à
Zurich, mécontent du flou géométrique entretenu autour de la droite numérique, et il a publié en
1872 sa construction de $\mathbb{R}$ par \og{} coupures \fg{}. La question (c) est exactement le trou de $\mathbb{Q}$ que
$ \sqrt{2} $ vient combler.

\begin{refs}
\item Contexte : Borne supérieure et borne inférieure --- Wikipédia (\url{https://fr.wikipedia.org/wiki/Borne_sup%C3%A9rieure_et_borne_inf%C3%A9rieure})
\item Contexte : Complétude des nombres réels --- Wikipédia (en anglais) (\url{https://en.wikipedia.org/wiki/Completeness_of_the_real_numbers})
\item Lecture : Abbott, \textit{Understanding Analysis}, ch. 1
\end{refs}

\bigskip

\begin{problem}[{Pourquoi 0,999\dots{} = 1 --- Lycée}]
\textbf{RA-02.} Le symbole $ 0{,}999\ldots $ désigne, par définition, la
limite de la suite $ 0{,}9,\ 0{,}99,\ 0{,}999,\ \dots $ --- autrement dit la somme de la série
$ \frac{9}{10} + \frac{9}{100} + \frac{9}{1000} + \cdots $. Démontrez que
$ 0{,}999\ldots = 1 $. Donnez deux démonstrations : l'une en sommant la série géométrique, l'autre
directement à partir de la définition d'une limite (montrez que pour tout $ \varepsilon > 0 $,
les termes de la suite finissent par se trouver à une distance inférieure à $ \varepsilon $ de 1).
Expliquez ensuite précisément ce qui cloche dans l'objection \og{} $ 0{,}999\ldots $ est plus petit que
1 parce qu'il n'atteint jamais 1 \fg{}.
\end{problem}

\noindent C'est l'égalité la plus discutée d'internet, et la discussion porte toujours sur
les définitions : un développement décimal illimité \textit{est} une limite, et une fois cela admis
la démonstration est courte. La question est réellement instructive --- elle force la prise de
conscience, rendue précise pour la première fois par Cauchy dans son \textit{Cours d'analyse} (1821),
que les processus infinis ne prennent sens qu'à travers la notion de limite. C'est aussi une
première rencontre avec le fait qu'un nombre réel peut avoir deux écritures décimales.

\begin{refs}
\item Contexte : 0,999\dots{} --- Wikipédia (\url{https://fr.wikipedia.org/wiki/0,999%E2%80%A6})
\item Lecture : Abbott, \textit{Understanding Analysis}, ch. 2
\end{refs}

\bigskip

\begin{problem}[{La série harmonique diverge --- Lycée}]
\textbf{RA-03.} Démontrez que la série harmonique
$ 1 + \frac{1}{2} + \frac{1}{3} + \frac{1}{4} + \cdots $ diverge : les sommes partielles
$ H_n = 1 + \frac{1}{2} + \cdots + \frac{1}{n} $ dépassent tout majorant. Montrez ensuite
quantitativement que $ H_{2^k} \ge 1 + k/2 $, et servez-vous-en pour estimer combien de termes
sont nécessaires avant que les sommes partielles ne dépassent 10 pour la première fois.
\end{problem}

\noindent Nicole Oresme a trouvé l'argument classique de regroupement vers 1350 --- une
démonstration médiévale que l'on ne peut toujours pas améliorer en élégance ---, après quoi elle fut
oubliée puis redémontrée par Pietro Mengoli (1650) et Jakob Bernoulli (1689). Le résultat est un
avertissement permanent : les termes d'une série peuvent tendre vers zéro pendant que la somme
croît sans limite. La question compagne, la valeur exacte de $ \sum 1/n^2 $, est le problème de
Bâle, traité dans L'art de la démonstration (\url{proofs.html}).

\begin{refs}
\item Contexte : Série harmonique --- Wikipédia (\url{https://fr.wikipedia.org/wiki/S%C3%A9rie_harmonique})
\item Lecture : Abbott, \textit{Understanding Analysis}, ch. 2
\item Cours : MIT OCW 18.100A Real Analysis (\url{https://ocw.mit.edu/courses/18-100a-real-analysis-fall-2020/})
\end{refs}

\bigskip

\begin{problem}[{Les deux limites qu'une suite possède toujours --- Lycée, short}]
\textbf{RA-24.} Pour la suite $ a_n = (-1)^n \left( 1 + \frac{1}{n}
\right) $, déterminez $ \limsup_{n \to \infty} a_n $ et $ \liminf_{n \to \infty} a_n $, avec
démonstration. Démontrez ensuite qu'une suite bornée de nombres réels converge si et seulement si sa
limite supérieure et sa limite inférieure coïncident.
\end{problem}

\noindent Une suite bornée n'a aucune raison de converger, mais elle possède toujours une
plus grande et une plus petite \og{} valeur finale \fg{} --- et, contrairement à la limite, ces deux-là
existent toujours. Cauchy travaillait déjà avec la plus grande valeur d'adhérence d'une suite ; en
faire une définition plutôt qu'une façon de parler est ce qui permet à l'analyse d'énoncer des
théorèmes (la règle de Cauchy, le lemme de Fatou) valables sans aucune hypothèse de convergence.

\begin{refs}
\item Contexte : Limite supérieure et limite inférieure --- Wikipédia (\url{https://fr.wikipedia.org/wiki/Limite_sup%C3%A9rieure_et_limite_inf%C3%A9rieure})
\item Lecture : Abbott, \textit{Understanding Analysis}, ch. 2
\end{refs}

\bigskip

\begin{problem}[{D'où vient e --- Lycée, short}]
\textbf{RA-25.} Démontrez que la suite
$ a_n = \left( 1 + \frac{1}{n} \right)^{n} $ est strictement croissante et majorée par
$ 3 $, de sorte qu'elle converge ; sa limite est par définition le nombre $ e $. Déduisez-en
que $ 2 < e < 3 $.
\end{problem}

\noindent Jacob Bernoulli a rencontré cette suite en 1683 en se demandant ce qu'il advient
des intérêts composés lorsqu'on les compose de plus en plus souvent ; les intérêts croissent, mais
pas sans limite, et le plafond est $ e $. Euler a donné à la constante sa lettre vers 1731 et l'a
calculée avec dix-huit décimales. La démonstration est une pure application du principe selon lequel
une suite croissante et majorée converge --- le premier endroit où un nombre réel est \textit{créé} par
une limite au lieu d'être écrit.

\begin{refs}
\item Contexte : e (nombre) --- Wikipédia (\url{https://fr.wikipedia.org/wiki/E_(nombre)})
\item Lecture : Rudin, \textit{Principles of Mathematical Analysis}, ch. 3
\end{refs}

\bigskip

\begin{problem}[{Segments emboîtés, et comment Cantor a dénombré la droite --- Lycée}]
\textbf{RA-26.} Ne supposez de $ \mathbb{R} $ rien d'autre que d'être un corps totalement
ordonné dans lequel toute partie non vide majorée admet un plus petit majorant (RA-01).
\begin{enumerate}
\item[(a)] Démontrez le théorème des segments emboîtés : si $ I_1 \supseteq I_2 \supseteq I_3 \supseteq \cdots $
sont des intervalles fermés bornés $ I_n = [a_n, b_n] $, alors $ \bigcap_{n=1}^{\infty} I_n $ n'est
pas vide. Montrez par un exemple que la conclusion peut tomber en défaut si les intervalles sont
ouverts, ou s'ils ne sont pas bornés.
\item[(b)] Déduisez-en le théorème de Bolzano--Weierstrass : toute suite bornée de nombres réels admet une
sous-suite convergente.
\item[(c)] Démontrez que $ \mathbb{R} $ n'est pas dénombrable, en utilisant (a) plutôt que les
développements décimaux : étant donné une suite quelconque $ x_1, x_2, x_3, \dots $ de nombres
réels, construisez des intervalles fermés emboîtés $ I_1 \supseteq I_2 \supseteq \cdots $ avec
$ x_n \notin I_n $, et concluez qu'aucune suite ne peut énumérer tous les nombres réels.
\item[(d)] Démontrez que les nombres algébriques --- les racines des polynômes non nuls à coefficients
entiers --- forment un ensemble dénombrable. Combinez avec (c) pour en déduire que les nombres
transcendants existent. Remarquez bien que cet argument n'en nomme aucun.
\end{enumerate}

\end{problem}

\noindent L'article de Georg Cantor de 1874, \og{} Sur une propriété de la collection de tous
les nombres algébriques réels \fg{}, contenait exactement l'argument des questions (c) et (d) --- le
célèbre argument diagonal n'est venu que dix-sept ans plus tard, en 1891. Trente ans auparavant,
Liouville avait construit explicitement un nombre transcendant, à la main, au prix d'un effort
énorme ; Cantor a montré en deux pages que presque tout nombre est transcendant sans en exhiber un
seul. Kronecker trouvait ce raisonnement répugnant et tenta d'en bloquer la publication. La question
(a) est l'humble outil qui fait tout fonctionner, et c'est aussi la forme la plus limpide de la
complétude : c'est elle qui distingue $ \mathbb{R} $ de $ \mathbb{Q} $.

\begin{refs}
\item Contexte : Théorème des fermés emboîtés --- Wikipédia (\url{https://fr.wikipedia.org/wiki/Th%C3%A9or%C3%A8me_des_ferm%C3%A9s_embo%C3%AEt%C3%A9s})
\item Contexte : Argument de la diagonale de Cantor --- Wikipédia (\url{https://fr.wikipedia.org/wiki/Argument_de_la_diagonale_de_Cantor})
\item Contexte : Nombre transcendant --- Wikipédia (\url{https://fr.wikipedia.org/wiki/Nombre_transcendant})
\item Lecture : Abbott, \textit{Understanding Analysis}, ch. 1
\end{refs}

\bigskip

\begin{problem}[{La propriété d'Archimède, et comment $\mathbb{Q}$ remplit $\mathbb{R}$ --- Lycée}]
\textbf{RA-36.} Ne supposez de $ \mathbb{R} $ rien d'autre que d'être un corps totalement
ordonné dans lequel toute partie non vide majorée admet un plus petit majorant.
\begin{enumerate}
\item[(a)] Démontrez la propriété d'Archimède : $ \mathbb{N} $ n'est pas majoré dans $ \mathbb{R} $,
donc pour tout réel $ x $ il existe un entier positif $ n > x $. (Supposez $ \mathbb{N} $
majoré et examinez $ \sup \mathbb{N} - 1 $.) Déduisez-en que pour tout
$ \varepsilon > 0 $ il existe un $ n $ tel que $ 1/n < \varepsilon $.
\item[(b)] Démontrez que les rationnels sont denses : dès que $ a < b $ il existe un rationnel $ q $
tel que $ a < q < b $. (Choisissez $ n $ tel que $ 1/n < b-a $ et prenez le plus petit
entier $ m $ tel que $ m/n > a $ ; expliquez pourquoi un tel plus petit entier existe.)
\item[(c)] Démontrez que $ \sqrt{2} $ est irrationnel, et déduisez-en que les irrationnels sont denses
eux aussi.
\item[(d)] Démontrez que tout réel $ x > 0 $ possède exactement une racine $ n $-ième positive, en
montrant que $ y = \sup\{ t > 0 : t^n < x \} $ vérifie $ y^n = x $ --- éliminez tour à tour
$ y^n < x $ et $ y^n > x $. Signalez chaque endroit où (a) est utilisé.
\item[(e)] Montrez que (a) ne découle pas des seuls axiomes de corps totalement ordonné : ordonnez le corps
des fractions rationnelles $ \mathbb{R}(t) $ en déclarant $ p/q > 0 $ lorsque les coefficients
dominants de $ p $ et de $ q $ sont de même signe, vérifiez que cela en fait un corps totalement
ordonné, et constatez que $ t $ y est plus grand que tout entier.
\end{enumerate}

\end{problem}

\noindent La propriété de (a) est plus ancienne que son nom : Archimède l'énonce comme un
postulat dans \textit{De la sphère et du cylindre} et en attribue l'idée à Eudoxe, qui s'en servit pour
rendre rigoureuse la méthode d'exhaustion au quatrième siècle avant notre ère. La question (e) est la
raison pour laquelle il s'agit d'un véritable théorème sur $ \mathbb{R} $ et non d'une trivialité ---
il existe de parfaits corps totalement ordonnés contenant des éléments infiniment grands, et c'est la
complétude qui les exclut ici. Tout ce qui, en analyse, commence par \og{} choisissons $ n $ assez
grand \fg{} est une application de (a), et (d) est l'endroit où le nombre $ \sqrt{2} $ que RA-01 avait
localisé comme borne supérieure devient enfin un nombre.

\begin{refs}
\item Contexte : Archimédien (propriété d'Archimède) --- Wikipédia (\url{https://fr.wikipedia.org/wiki/Archim%C3%A9dien})
\item Contexte : Partie dense --- Wikipédia (\url{https://fr.wikipedia.org/wiki/Partie_dense})
\item Contexte : Racine d'un nombre --- Wikipédia (\url{https://fr.wikipedia.org/wiki/Racine_d'un_nombre})
\item Lecture : Rudin, \textit{Principles of Mathematical Analysis}, ch. 1
\item Lecture : Abbott, \textit{Understanding Analysis}, ch. 1
\end{refs}

\bigskip

\begin{problem}[{Une récurrence qui monte vers une limite --- Lycée, short}]
\textbf{RA-37.} Soient $ a_1 = \sqrt{2} $ et $ a_{n+1} = \sqrt{2 + a_n} $
pour $ n \ge 1 $. Démontrez que la suite est strictement croissante et majorée par $ 2 $, donc
convergente, et déterminez sa limite --- en justifiant soigneusement pourquoi la limite doit vérifier
l'équation que vous écrivez, et pourquoi cette équation n'a qu'une seule racine admissible.
\end{problem}

\noindent La démonstration, c'est le théorème de convergence monotone et rien d'autre, mais
c'est dans la dernière clause que se trouve la réflexion : passer à la limite à l'intérieur de
$ \sqrt{2+\cdot} $ utilise la continuité, et écarter la seconde racine de l'équation du second
degré obtenue utilise l'encadrement que vous avez démontré. Ces radicaux imbriqués sont aussi le
mécanisme de la formule de Viète pour $\pi$ (1593), le premier produit infini jamais écrit, où les mêmes
expressions apparaissent comme cosinus d'angles divisés en deux encore et encore.

\begin{refs}
\item Contexte : Théorème de convergence monotone --- Wikipédia (\url{https://fr.wikipedia.org/wiki/Th%C3%A9or%C3%A8me_de_convergence_monotone})
\item Contexte : Radical imbriqué --- Wikipédia (\url{https://fr.wikipedia.org/wiki/Radical_imbriqu%C3%A9})
\item Lecture : Abbott, \textit{Understanding Analysis}, ch. 2
\end{refs}

\bigskip


\section*{Analyse complexe}

\begin{problem}[{De Moivre et les racines de l'unité --- Lycée}]
\textbf{CA-01.} \begin{enumerate}
\item[(a)] Démontrez la formule de Moivre par récurrence : pour tout
entier $ n \ge 1 $,
$ (\cos\theta + i\sin\theta)^n = \cos n\theta + i\sin n\theta $.
\item[(b)] Servez-vous-en pour trouver toutes les solutions complexes de $ z^n = 1 $, montrez que ce sont
$ n $ points régulièrement espacés sur le cercle unité (les sommets d'un $ n $-gone régulier), et
démontrez que pour $ n \ge 2 $ leur somme vaut 0.
\item[(c)] Extrayez de la trigonométrie réelle à partir de l'arithmétique complexe : établissez
la formule donnant $ \cos 3\theta $ comme polynôme en $ \cos\theta $, et calculez
$ \cos 36^\circ $ exactement à l'aide des racines cinquièmes de l'unité.
\end{enumerate}

\end{problem}

\noindent Abraham de Moivre a utilisé la formule dès 1707 dans ses travaux sur les
équations, et Roger Cotes (1714) est passé à un cheveu de la forme exponentielle qu'Euler a
finalement écrite $ e^{i\theta} = \cos\theta + i\sin\theta $. Les racines de l'unité sont le point
de rencontre de l'algèbre et de la géométrie : la construction du 17-gone régulier par Gauss en
1796, à dix-neuf ans --- la découverte qui lui a fait choisir les mathématiques ---, c'est au fond la
question (c) poussée jusqu'à $ n = 17 $.

\begin{refs}
\item Contexte : Formule de Moivre --- Wikipédia (\url{https://fr.wikipedia.org/wiki/Formule_de_Moivre})
\item Contexte : Racine de l'unité --- Wikipédia (\url{https://fr.wikipedia.org/wiki/Racine_de_l'unit%C3%A9})
\item Lecture : Needham, \textit{Visual Complex Analysis}, ch. 1
\end{refs}

\bigskip

\begin{problem}[{Multiplier, c'est faire tourner --- Lycée}]
\textbf{CA-02.} Représentez $ z = a + bi $ par le point $ (a, b) $ du
plan, de module $ |z| = \sqrt{a^2 + b^2} $ et d'argument l'angle avec le demi-axe réel
positif.
\begin{enumerate}
\item[(a)] Démontrez que $ |zw| = |z|\,|w| $ et que les arguments s'ajoutent par multiplication,
de sorte que multiplier par $ z $ fait tourner le plan de $ \arg z $ et le dilate de $ |z| $.
\item[(b)] Montrez que l'identité $ |zw| = |z|\,|w| $, écrite en coordonnées, est exactement
l'identité classique des deux carrés :
$ (a^2 + b^2)(c^2 + d^2) = (ac - bd)^2 + (ad + bc)^2 $.
\item[(c)] Démontrez l'inégalité triangulaire $ |z + w| \le |z| + |w| $ et déterminez exactement quand
il y a égalité.
\end{enumerate}

\end{problem}

\noindent Les nombres complexes sont entrés en mathématiques par la bande : Cardan et
Bombelli (1545--1572) en avaient besoin pour résoudre les équations du troisième degré à trois
racines réelles. Ils ne sont devenus respectables qu'avec le foyer géométrique que leur ont donné
Wessel (1797), Argand (1806) et Gauss. La question (b) montre que la géométrie se cachait à la vue
de tous depuis des siècles --- l'identité des deux carrés était connue de Diophante et de Brahmagupta
mille ans avant que quiconque ne dessine un nombre complexe.

\begin{refs}
\item Contexte : Plan complexe --- Wikipédia (\url{https://fr.wikipedia.org/wiki/Plan_complexe})
\item Contexte : Identité de Brahmagupta--Fibonacci --- Wikipédia (\url{https://fr.wikipedia.org/wiki/Identit%C3%A9_de_Brahmagupta})
\item Lecture : Needham, \textit{Visual Complex Analysis}, ch. 1
\end{refs}

\bigskip

\begin{problem}[{La règle du parallélogramme --- Lycée, short}]
\textbf{CA-20.} Démontrez que pour tous nombres complexes $ z $ et $ w $,
\[ |z + w|^2 + |z - w|^2 = 2|z|^2 + 2|w|^2, \]
et expliquez ce que cette identité dit des diagonales et des côtés d'un parallélogramme.
\end{problem}

\noindent La démonstration en une ligne --- développer $ |a|^2 = a\bar{a} $ et regarder les
termes croisés s'annuler --- est une première démonstration que la conjugaison fait le travail que
faisaient autrefois les coordonnées. L'identité est aussi l'empreinte algébrique exacte d'une norme
issue d'un produit scalaire : une norme sur un espace vectoriel réel ou complexe vérifie la règle du
parallélogramme si et seulement si elle est induite par un produit scalaire, ce qui est le théorème
de Jordan--von Neumann de 1935.

\begin{refs}
\item Contexte : Règle du parallélogramme --- Wikipédia (\url{https://fr.wikipedia.org/wiki/R%C3%A8gle_du_parall%C3%A9logramme})
\item Contexte : Nombre complexe --- Wikipédia (\url{https://fr.wikipedia.org/wiki/Nombre_complexe})
\end{refs}

\bigskip

\begin{problem}[{Racines carrées sans trigonométrie --- Lycée, short}]
\textbf{CA-21.} Trouvez les deux racines carrées complexes de $ 3 + 4i $ en
résolvant $ (x + iy)^2 = 3 + 4i $ avec $ x, y $ réels, sans utiliser la trigonométrie. Démontrez
ensuite que tout nombre complexe non nul possède exactement deux racines carrées, et donnez-en une
formule en fonction de $ a $, $ b $ et $ \sqrt{a^2 + b^2} $ lorsque le nombre est
$ a + bi $.
\end{problem}

\noindent Comparer parties réelles et parties imaginaires transforme le problème en deux
équations réelles, et le module fournit la troisième relation qui les rend solubles --- petite
illustration de la façon dont l'algèbre complexe tient sa propre comptabilité. Que tout nombre
possède exactement $ n $ racines $ n $-ièmes est le premier indice du théorème fondamental de
l'algèbre (CA-07), et la raison pour laquelle les mathématiciens ont cessé de s'excuser pour
$ \sqrt{-1} $.

\begin{refs}
\item Contexte : Racine carrée --- Wikipédia (\url{https://fr.wikipedia.org/wiki/Racine_carr%C3%A9e})
\item Contexte : Nombre complexe --- Wikipédia (\url{https://fr.wikipedia.org/wiki/Nombre_complexe})
\end{refs}

\bigskip

\begin{problem}[{Les nombres complexes comme moteur de géométrie --- Lycée}]
\textbf{CA-22.} Identifiez le plan à $ \mathbb{C} $, et notez $ \omega = e^{2\pi i/3} $
une racine cubique primitive de l'unité.
\begin{enumerate}
\item[(a)] Démontrez que la rotation du plan d'angle $ \theta $ autour du point $ p $ est
l'application $ z \mapsto p + e^{i\theta}(z - p) $, et que la composée de deux rotations est une
rotation ou une translation --- déterminez exactement quand chaque cas se produit.
\item[(b)] Démontrez que le triangle de sommets distincts $ a, b, c $ est équilatéral si et seulement si
\[ a^2 + b^2 + c^2 = ab + bc + ca, \]
et montrez que cela équivaut à $ a + \omega b + \omega^2 c = 0 $ ou
$ a + \omega^2 b + \omega c = 0 $, selon l'orientation du triangle.
\item[(c)] Démontrez le théorème de Napoléon : construisez vers l'extérieur des triangles équilatéraux sur
les trois côtés d'un triangle quelconque ; alors les centres de ces trois triangles équilatéraux
sont eux-mêmes les sommets d'un triangle équilatéral.
\item[(d)] Démontrez que le triangle de Napoléon de (c) a le même centre de gravité que le triangle
initial.
\end{enumerate}

\end{problem}

\noindent Dès lors qu'une rotation n'est plus qu'une multiplication par $ e^{i\theta} $,
tout un genre de géométrie plane devient une algèbre qu'un élève patient peut mener sans dessin. Le
théorème de Napoléon est attribué --- sans aucune bonne raison --- à l'empereur ; il apparaît pour la
première fois par écrit en 1825, quatre ans après sa mort, et reste la meilleure publicité pour la
méthode des nombres complexes. La question (b) mérite d'être retenue : une équation sans contenu
géométrique apparent se révèle dire \og{} équilatéral \fg{}, et c'est ce qui donne à la méthode son air de
machine.

\begin{refs}
\item Contexte : Théorème de Napoléon --- Wikipédia (\url{https://fr.wikipedia.org/wiki/Th%C3%A9or%C3%A8me_de_Napol%C3%A9on})
\item Contexte : Racine de l'unité --- Wikipédia (\url{https://fr.wikipedia.org/wiki/Racine_de_l'unit%C3%A9})
\item Lecture : Needham, \textit{Visual Complex Analysis}, ch. 1
\item Lecture : Hahn, \textit{Complex Numbers and Geometry} (MAA, 1994)
\end{refs}

\bigskip

\begin{problem}[{Sommer des cosinus avec une série géométrique --- Lycée, short}]
\textbf{CA-32.} Pour $ \theta $ non multiple de $ 2\pi $, sommez la série
géométrique $ \sum_{k=0}^{n-1} e^{ik\theta} $ et prenez parties réelle et imaginaire pour obtenir
des formules closes pour $ \sum_{k=0}^{n-1} \cos k\theta $ et $ \sum_{k=0}^{n-1} \sin k\theta $.
Déduisez-en l'identité du noyau de Dirichlet
\[ 1 + 2\sum_{k=1}^{n} \cos k\theta \;=\; \frac{\sin\!\left(\left(n + \tfrac{1}{2}\right)\theta\right)}{\sin(\theta/2)} . \]
\end{problem}

\noindent Des sommes trigonométriques d'apparence redoutable ne sont que des séries
géométriques déguisées dès que l'on lit $ \cos k\theta $ comme la partie réelle de
$ (e^{i\theta})^k $ ; la formule d'Euler transforme une page de chasse aux angles en une ligne
d'algèbre. Les identités étaient connues de Lagrange au dix-huitième siècle, et le noyau de la
seconde formule est celui que Dirichlet a placé au centre de sa démonstration de 1829 de la
convergence de la série de Fourier d'une fonction convenable --- il réapparaît, avec le même visage,
dans l'analyse fonctionnelle de ce site.

\begin{refs}
\item Contexte : Formule d'Euler --- Wikipédia (\url{https://fr.wikipedia.org/wiki/Formule_d'Euler})
\item Contexte : Série géométrique --- Wikipédia (\url{https://fr.wikipedia.org/wiki/S%C3%A9rie_g%C3%A9om%C3%A9trique}), Noyau de Dirichlet --- Wikipédia (\url{https://fr.wikipedia.org/wiki/Noyau_de_Dirichlet})
\item Lecture : Needham, \textit{Visual Complex Analysis}, ch. 1
\end{refs}

\bigskip

\begin{problem}[{Le théorème de Ptolémée en une identité --- Lycée}]
\textbf{CA-33.} Soient $ a, b, c, d $ quatre nombres complexes quelconques, vus comme les
points $ A, B, C, D $ du plan.
\begin{enumerate}
\item[(a)] Vérifiez en développant les deux membres l'identité
\[ (a - b)(c - d) + (a - d)(b - c) \;=\; (a - c)(b - d) . \]
\item[(b)] Déduisez-en l'inégalité de Ptolémée : pour quatre points quelconques du plan,
$ AC \cdot BD \le AB \cdot CD + AD \cdot BC $, où $ XY $ désigne la distance entre
$ X $ et $ Y $.
\item[(c)] Placez maintenant les quatre points sur le cercle unité, $ a = e^{i\alpha} $, $ b = e^{i\beta} $,
$ c = e^{i\gamma} $, $ d = e^{i\delta} $ avec $ 0 \le \alpha < \beta < \gamma < \delta < 2\pi $.
Démontrez la formule de la corde $ |e^{i\theta} - e^{i\varphi}| = 2\left|\sin\frac{\theta - \varphi}{2}\right| $,
et montrez que (b) devient alors une \textit{égalité} --- le théorème de Ptolémée pour un quadrilatère
inscriptible --- en vérifiant l'identité obtenue entre produits de sinus à l'aide des formules de
transformation de produits en sommes.
\item[(d)] Appliquez (c) à quatre sommets d'un pentagone régulier pour démontrer que le rapport d'une
diagonale à un côté est le nombre d'or $ \varphi = \frac{1 + \sqrt{5}}{2} $. (CA-01(c) atteint la
même constante par une autre route, via les racines cinquièmes de l'unité.)
\end{enumerate}

\end{problem}

\noindent Claude Ptolémée a démontré le théorème dans l'\textit{Almageste} vers 150 de notre
ère, et ce n'était pas un ornement : l'égalité de (c) est précisément la règle d'addition et de
soustraction des cordes, la récurrence avec laquelle il a calculé la table des cordes qui a porté
l'astronomie grecque et islamique pendant quatorze siècles --- de la trigonométrie avant l'existence
des fonctions trigonométriques. La démonstration par les nombres complexes comprime tout cela dans
l'identité scolaire de la question (a), qui ne demande ni dessin ni discussion de cas, alors que la
démonstration synthétique classique exige une construction auxiliaire qu'il faut deviner. La
question (b) montre que l'inégalité vaut pour \textit{quatre points quelconques}, en toute
configuration ; le cercle est exactement le cas d'égalité.

\begin{refs}
\item Contexte : Théorème de Ptolémée --- Wikipédia (\url{https://fr.wikipedia.org/wiki/Th%C3%A9or%C3%A8me_de_Ptol%C3%A9m%C3%A9e})
\item Contexte : Inégalité de Ptolémée --- Wikipédia (\url{https://fr.wikipedia.org/wiki/In%C3%A9galit%C3%A9_de_Ptol%C3%A9m%C3%A9e}), Table des cordes de Ptolémée --- Wikipédia (en anglais) (\url{https://en.wikipedia.org/wiki/Ptolemy%27s_table_of_chords})
\item Lecture : Needham, \textit{Visual Complex Analysis}, ch. 1
\item Lecture : Hahn, \textit{Complex Numbers and Geometry} (MAA, 1994)
\end{refs}

\bigskip


\section*{Topologie}

\begin{problem}[{Les ponts de Königsberg --- Lycée}]
\textbf{TOP-01.} La ville de Königsberg comptait sept ponts reliant deux îles fluviales aux deux
rives de la Pregel : la plus grande île était reliée par deux ponts à chaque rive et par un pont à la
plus petite île, elle-même reliée par un pont à chaque rive. Les habitants se demandaient si l'on
pouvait faire une promenade traversant chaque pont exactement une fois.
\begin{enumerate}
\item[(a)] Démontrez qu'aucune telle promenade n'existe.
\item[(b)] Modélisez la situation par un graphe --- les masses de terre comme sommets, les ponts comme arêtes
--- et déterminez, avec démonstration, exactement quels graphes connexes admettent une promenade
empruntant chaque arête une fois : en fonction du nombre de sommets rencontrant un nombre impair
d'arêtes, quand une telle promenade existe-t-elle, et quand peut-elle en outre revenir à son point de
départ ?
\end{enumerate}

\end{problem}

\noindent Euler a réglé la question en 1736 dans un article qu'il décrivait comme relevant de
la \og{} géométrie de position \fg{} --- une géométrie sans aucune longueur ni aucun angle, rien que des
connexions. C'est la date de naissance conventionnelle de la théorie des graphes autant que de la
topologie. Euler a démontré la nécessité de sa condition ; la moitié \og{} suffisance \fg{} de la question (b)
n'a été rédigée soigneusement que par Carl Hierholzer en 1873. Le problème récompense ce que la
topologie récompense toujours : trouver la seule quantité qui ne peut pas changer.

\begin{refs}
\item Contexte : Problème des sept ponts de Königsberg --- Wikipédia (\url{https://fr.wikipedia.org/wiki/Probl%C3%A8me_des_sept_ponts_de_K%C3%B6nigsberg})
\item Contexte : Graphe eulérien (chemin eulérien) --- Wikipédia (\url{https://fr.wikipedia.org/wiki/Graphe_eul%C3%A9rien})
\item Lecture : M. A. Armstrong, \textit{Basic Topology}, ch. 1
\end{refs}

\bigskip

\begin{problem}[{La formule des polyèdres d'Euler et les cinq solides de Platon --- Lycée}]
\textbf{TOP-02.} Soit un polyèdre convexe à $ V $ sommets, $ E $ arêtes et $ F $
faces.
\begin{enumerate}
\item[(a)] Démontrez la formule d'Euler
\[ V - E + F = 2 . \]
(Une voie honnête : aplatissez le polyèdre en un réseau plan connexe et raisonnez par récurrence sur
le nombre d'arêtes.)
\item[(b)] Un solide de Platon est un polyèdre convexe dont les faces sont des $ p $-gones réguliers
congruents, avec le même nombre $ q $ de faces se rencontrant en chaque sommet. À l'aide de (a),
démontrez que $ (p, q) $ doit vérifier $ \frac{1}{p} + \frac{1}{q} > \frac{1}{2} $, et
déduisez-en qu'exactement cinq couples sont possibles.
\end{enumerate}

\end{problem}

\noindent Euler a annoncé la formule dans une lettre à Goldbach en 1750, en avouant qu'il ne
savait pas encore la démontrer ; Legendre en a donné une démonstration en 1794 et Cauchy une autre en
1813, et la longue querelle sur ce qui compte exactement comme un \og{} polyèdre \fg{} est devenue le sujet du
dialogue classique de Lakatos, \textit{Proofs and Refutations}. Les cinq solides eux-mêmes closent les
\textit{Éléments} d'Euclide (livre XIII). La quantité $ V - E + F $ est la caractéristique d'Euler --- le
tout premier invariant topologique découvert, et l'ancêtre de tout ce que l'on trouve en TOP-19 et
TOP-20.

\begin{refs}
\item Contexte : Caractéristique d'Euler --- Wikipédia (\url{https://fr.wikipedia.org/wiki/Caract%C3%A9ristique_d'Euler})
\item Contexte : Graphe planaire --- Wikipédia (\url{https://fr.wikipedia.org/wiki/Graphe_planaire})
\item Lecture : M. A. Armstrong, \textit{Basic Topology}, ch. 1
\end{refs}

\bigskip

\begin{problem}[{Expériences avec un ruban de Möbius, rendues rigoureuses --- Lycée}]
\textbf{TOP-03.} Fabriquez un ruban de Möbius : donnez un demi-tour à une bande de papier et collez
les extrémités. Modélisez-le mathématiquement par le carré $ [0,1] \times [0,1] $ dont les bords
gauche et droit sont recollés par $ (0, y) \sim (1, 1-y) $.
\begin{enumerate}
\item[(a)] Démontrez que le bord du ruban de Möbius est une seule courbe fermée (le cylindre ordinaire en a
deux).
\item[(b)] Prédisez ce que l'on obtient en coupant le long de la ligne médiane $ y = \tfrac12 $ ;
démontrez-le ensuite : montrez que la bande coupée est connexe, et déterminez le nombre de demi-tours
qu'elle porte.
\item[(c)] Faites de même pour la coupe le long de la ligne $ y = \tfrac13 $ : démontrez que le résultat a
exactement deux morceaux et décrivez, avec démonstration à partir du modèle de recollement, comment
ils sont joints.
\item[(d)] Montrez qu'une fourmi qui fait un tour complet le long de la ligne médiane revient inversée comme
dans un miroir --- précisez, à l'aide de la règle de recollement, ce que signifie \og{} le ruban n'a qu'une
seule face \fg{}.
\end{enumerate}

\end{problem}

\noindent August Möbius et Johann Listing ont découvert le ruban indépendamment en 1858
(Listing, qui a aussi forgé le mot \og{} topologie \fg{}, a publié le premier). C'est la plus simple des
surfaces non orientables, et les expériences aux ciseaux --- familières des spectacles de magie ---
deviennent des théorèmes dès que le modèle de papier est remplacé par le carré recollé, technique même
qui construit le tore et la bouteille de Klein en TOP-09. La non-orientabilité, rencontrée ici pour la
première fois, est l'un des deux invariants qui classifient toutes les surfaces (TOP-19).

\begin{refs}
\item Contexte : Ruban de Möbius --- Wikipédia (\url{https://fr.wikipedia.org/wiki/Ruban_de_M%C3%B6bius})
\item Lecture : M. A. Armstrong, \textit{Basic Topology}, ch. 1
\end{refs}

\bigskip

\begin{problem}[{Antipodes sur l'équateur --- Lycée}]
\textbf{TOP-04.} Supposez que la température varie continûment le long de
l'équateur terrestre. Démontrez qu'à tout instant donné il existe deux points diamétralement opposés
(antipodaux) sur l'équateur ayant exactement la même température. Formellement : si
$ f : [0, 2\pi] \to \mathbb{R} $ est continue avec $ f(0) = f(2\pi) $, montrez qu'il existe
$ t $ tel que $ f(t) = f(t + \pi) $ (les angles étant pris modulo $ 2\pi $). Vous pouvez
utiliser le théorème des valeurs intermédiaires : une fonction continue qui change de signe sur un
intervalle y admet un zéro.
\end{problem}

\noindent C'est le cas unidimensionnel du théorème de Borsuk--Ulam (TOP-16), et c'est
l'illustration la plus nette possible de la façon dont la topologie démontre l'existence : aucune
formule ne localise jamais les deux points, et pourtant un argument de deux lignes montre qu'ils sont
forcément là. Le théorème des valeurs intermédiaires lui-même a été démontré rigoureusement pour la
première fois par Bolzano en 1817, dans un article explicitement consacré au remplacement de l'évidence
géométrique par la démonstration --- le geste fondateur de l'analyse comme de la topologie
ensembliste.

\begin{refs}
\item Contexte : Théorème des valeurs intermédiaires --- Wikipédia (\url{https://fr.wikipedia.org/wiki/Th%C3%A9or%C3%A8me_des_valeurs_interm%C3%A9diaires})
\item Contexte : Théorème de Borsuk-Ulam --- Wikipédia (\url{https://fr.wikipedia.org/wiki/Th%C3%A9or%C3%A8me_de_Borsuk-Ulam})
\end{refs}

\bigskip

\begin{problem}[{Un point fixe sur l'intervalle --- Lycée, short}]
\textbf{TOP-25.} Soit $ f : [0,1] \to [0,1] $ continue. Démontrez qu'il existe
$ x \in [0,1] $ tel que $ f(x) = x $.
\end{problem}

\noindent C'est le théorème du point fixe de Brouwer en dimension un, et toute la
démonstration tient en une ligne dès que l'on regarde $ g(x) = f(x) - x $ aux deux extrémités.
Brouwer a démontré le cas général vers 1910 ; l'énoncé en dimension supérieure (TOP-14) demande de la
machinerie, mais le cas de l'intervalle ne demande que le théorème des valeurs intermédiaires de
Bolzano --- un bon endroit pour sentir comment une hypothèse purement qualitative force un objet à
exister.

\begin{refs}
\item Contexte : Théorème du point fixe de Brouwer --- Wikipédia (\url{https://fr.wikipedia.org/wiki/Th%C3%A9or%C3%A8me_du_point_fixe_de_Brouwer})
\item Contexte : Théorème des valeurs intermédiaires --- Wikipédia (\url{https://fr.wikipedia.org/wiki/Th%C3%A9or%C3%A8me_des_valeurs_interm%C3%A9diaires})
\end{refs}

\bigskip

\begin{problem}[{Dedans et dehors, comptés par parité --- Lycée, short}]
\textbf{TOP-26.} Soit $ P $ un polygone simple fermé du plan (un nombre fini de
segments, se rencontrant seulement en des extrémités communes, sans auto-intersection) et soit $ q $
un point hors de $ P $. Démontrez que le nombre de fois qu'une demi-droite issue de $ q $ traverse
$ P $ est pair pour toute telle demi-droite, ou impair pour toute telle demi-droite --- la parité ne
dépend pas de la demi-droite que vous lancez, pourvu qu'elle ne passe par aucun sommet de
$ P $.
\end{problem}

\noindent Cette parité est exactement ce que signifie \og{} dedans \fg{} pour un polygone, et la
démontrer est le noyau polygonal honnête du théorème de Jordan : Camille Jordan a énoncé le théorème
général en 1887, et Oswald Veblen en a donné en 1905 une démonstration largement tenue pour la première
complète. Le même argument, sous le nom de lancer de rayon, est ce qu'un moteur graphique exécute chaque
fois qu'il décide si votre curseur est à l'intérieur d'une forme.

\begin{refs}
\item Contexte : Théorème de Jordan --- Wikipédia (\url{https://fr.wikipedia.org/wiki/Th%C3%A9or%C3%A8me_de_Jordan})
\item Contexte : Point dans un polygone --- Wikipédia (en anglais) (\url{https://en.wikipedia.org/wiki/Point_in_polygon})
\end{refs}

\bigskip

\begin{problem}[{Eau, gaz, électricité --- Lycée}]
\textbf{TOP-27.} Trois maisons doivent chacune être reliées à trois services --- eau, gaz,
électricité --- par des conduites tracées dans le plan, sans qu'aucune ne se croise. On note
$ K_{3,3} $ ce graphe et $ K_5 $ le graphe à cinq sommets muni de ses dix arêtes. Vous pouvez
utiliser la formule d'Euler $ V - E + F = 2 $ pour un graphe connexe tracé dans le plan sans
croisement, $ F $ comptant la région extérieure.
\begin{enumerate}
\item[(a)] Démontrez qu'un graphe simple connexe planaire avec $ V \ge 3 $ vérifie
\[ E \le 3V - 6 , \]
en comptant les incidences arête--face, et déduisez-en que $ K_5 $ ne peut pas être tracé dans le
plan sans croisement.
\item[(b)] Démontrez que si un tel graphe ne contient de plus aucun triangle, alors $ E \le 2V - 4 $, et
déduisez-en que $ K_{3,3} $ ne peut pas non plus être tracé --- l'énigme n'a pas de solution.
\item[(c)] Tracez $ K_5 $ et $ K_{3,3} $ sur la surface d'un beignet (un tore) sans croisement, et dites
exactement quelle étape de votre argument de comptage y échoue.
\end{enumerate}

\end{problem}

\noindent Henry Dudeney a publié l'énigme \og{} eau, gaz et électricité \fg{} en 1917 et la disait déjà
ancienne ; des générations de solveurs ont cherché le tracé astucieux qui n'existe pas. La démonstration
qu'il n'existe pas est un petit chef-d'\oe{}uvre de la méthode topologique : rien sur les conduites, un
simple décompte de sommets, d'arêtes et de régions. Kuratowski a démontré en 1930 que $ K_5 $ et
$ K_{3,3} $ sont les \textit{seules} obstructions --- un graphe est planaire exactement quand il ne
contient de subdivision ni de l'un ni de l'autre --- et la question (c) est le premier indice que
\og{} planaire \fg{} n'a jamais concerné le graphe seul, mais la surface sur laquelle on le trace, dont la
caractéristique d'Euler (TOP-02) fait le vrai travail.

\begin{refs}
\item Contexte : Énigme des trois maisons --- Wikipédia (\url{https://fr.wikipedia.org/wiki/%C3%89nigme_des_trois_maisons})
\item Contexte : Graphe planaire --- Wikipédia (\url{https://fr.wikipedia.org/wiki/Graphe_planaire}), Théorème de Kuratowski --- Wikipédia (en anglais) (\url{https://en.wikipedia.org/wiki/Kuratowski%27s_theorem})
\item Lecture : R. Diestel, \textit{Graph Theory}, ch. 4
\end{refs}

\bigskip

\begin{problem}[{Deux alpinistes, une seule altitude --- Lycée}]
\textbf{TOP-37.} Deux profils de montagne sont donnés par des fonctions continues
$ f, g : [0,1] \to [0,1] $ avec $ f(0) = g(0) = 0 $, $ f(1) = g(1) = 1 $ et
$ 0 < f(x), g(x) < 1 $ pour $ 0 < x < 1 $. Un alpiniste parcourt le profil $ f $ et
l'autre le profil $ g $ ; une \textit{traversée conjointe} est un couple d'horaires continus
$ x_1, x_2 : [0,1] \to [0,1] $ avec $ x_1(0) = x_2(0) = 0 $, $ x_1(1) = x_2(1) = 1 $, et
\[ f\big(x_1(u)\big) = g\big(x_2(u)\big) \quad \text{pour tout } u \in [0,1] , \]
de sorte que les deux sont à la même altitude à chaque instant.
\begin{enumerate}
\item[(a)] Esquissez deux profils pour lesquels aucune traversée conjointe ne permet aux deux alpinistes de
n'aller que vers l'est --- à un moment donné, l'un d'eux doit redescendre.
\item[(b)] Supposez maintenant $ f $ et $ g $ affines par morceaux, avec un nombre fini de sommets et de
creux, et telles qu'aucune valeur de maximum ou de minimum local de $ f $ n'égale une valeur de
maximum ou de minimum local de $ g $. Posons
\[ S = \{ (x,y) \in [0,1] \times [0,1] : f(x) = g(y) \} . \]
Démontrez que $ S $ est une réunion finie de segments de droite qui s'assemblent en un nombre fini
d'arcs et de boucles fermées, et que les seules extrémités d'arcs sont les deux coins $ (0,0) $ et
$ (1,1) $.
\item[(c)] Déduisez-en que le morceau de $ S $ contenant $ (0,0) $ est un arc aboutissant à
$ (1,1) $, et lisez une traversée conjointe sur un paramétrage de cet arc --- les alpinistes peuvent
donc toujours le faire.
\item[(d)] Montrez que la continuité seule ne suffit pas : décrivez des profils pour lesquels $ f $ est
constante sur un intervalle tandis que $ g $ oscille une infinité de fois à cette même hauteur, et
expliquez pourquoi le second alpiniste devrait osciller indéfiniment.
\end{enumerate}

\end{problem}

\noindent Tatsuo Homma a résolu une version de ce problème en 1952 et James Whittaker l'a nommé
et posé sous sa forme actuelle en 1966 ; John Philip Huneke a démontré en 1969 qu'une traversée existe
dès qu'aucun des deux profils n'est constant sur un intervalle, et le contre-exemple esquissé en (d)
montre que cette hypothèse ne peut pas être simplement abandonnée. L'exposé de Goodman, Pach et Yap de
1989 a reçu le prix Lester R. Ford de la Mathematical Association of America et a relié l'énigme à la
planification de mouvement des robots et au problème du carré inscrit. La méthode est celle que la
topologie emploie partout : cessez de raisonner sur les alpinistes et raisonnez sur la forme de
l'ensemble de toutes les positions licites.

\begin{refs}
\item Contexte : Problème de l'escalade en montagne --- Wikipédia (en anglais) (\url{https://en.wikipedia.org/wiki/Mountain_climbing_problem})
\item Contexte : Connexité (mathématiques) --- Wikipédia (\url{https://fr.wikipedia.org/wiki/Connexit%C3%A9_(math%C3%A9matiques)})
\item Article : J. E. Goodman, J. Pach, C.-K. Yap, \og{} Mountain climbing, ladder moving, and the ring-width of a polygon \fg{} (1989), \textit{American Mathematical Monthly} 96, 494--510
\end{refs}

\bigskip

\begin{problem}[{Le polyèdre qui casse la formule d'Euler --- Lycée, short}]
\textbf{TOP-38.} Construisez un cadre de tableau polyédrique --- un bloc
rectangulaire percé de part en part d'un trou rectangulaire --- découpé de sorte que chaque face soit un
polygone plat sans trou, puis comptez $ V $, $ E $ et $ F $ et vérifiez que
\[ V - E + F = 0 , \quad \text{et non } 2 . \]
Identifiez exactement quelle étape de la démonstration usuelle de la formule d'Euler (TOP-02) échoue
pour ce solide, et dites ce que le nombre $ V - E + F $ mesure à la place.
\end{problem}

\noindent Simon L'Huilier a publié précisément de telles exceptions en 1812--13, montrant qu'un
solide percé de $ g $ tunnels vérifie $ V - E + F = 2 - 2g $ ; Hessel a ajouté d'autres \og{} monstres \fg{}
en 1832. Imre Lakatos a transformé cette longue querelle --- sur la question de savoir si de tels objets
sont des polyèdres, ou si la définition doit être réparée pour les exclure --- en
\textit{Proofs and Refutations}, le meilleur livre jamais écrit sur la manière dont les mathématiques se
développent réellement. Notez que si vous autorisez le dessus et le dessous à être des faces annulaires
d'un seul tenant, vous retrouvez $ 2 $, et c'est exactement le genre d'échappatoire dont on discutait
au dix-neuvième siècle.

\begin{refs}
\item Contexte : Polyèdre toroïdal --- Wikipédia (en anglais) (\url{https://en.wikipedia.org/wiki/Toroidal_polyhedron})
\item Contexte : Caractéristique d'Euler --- Wikipédia (\url{https://fr.wikipedia.org/wiki/Caract%C3%A9ristique_d'Euler})
\item Lecture : I. Lakatos, \textit{Proofs and Refutations}
\end{refs}

\bigskip


\section*{Géométrie}

\begin{problem}[{Le théorème de l'angle inscrit --- Lycée}]
\textbf{GEO-01.} Soient $ A $ et $ B $ deux points d'un cercle de centre $ O $,
et soit $ P $ un point quelconque du grand arc $ AB $.
\begin{enumerate}
\item[(a)] Démontrez que l'angle au centre vaut le double de l'angle inscrit :
\[ \angle AOB = 2\,\angle APB. \]
\item[(b)] Déduisez-en le théorème de l'angle inscrit dans un demi-cercle : un angle inscrit dans un demi-cercle est droit.
\item[(c)] Déduisez-en que les angles opposés d'un quadrilatère inscriptible ont pour somme $ 180^\circ $.
\end{enumerate}

\end{problem}

\noindent Le théorème de l'angle inscrit est la proposition 20 du livre III des
\textit{Éléments} d'Euclide, et la tradition attribue à Thalès de Milet (VIe siècle
av. J.-C.) le cas du demi-cercle --- peut-être le tout premier théorème de l'histoire attribué
à une personne nommée. C'est le moteur de la \og{} chasse aux angles \fg{} : presque tout problème
sérieux de géométrie du cercle en olympiade tourne grâce à lui. La démonstration exige une
disjonction de cas ; trouver la bonne décomposition, c'est tout l'enjeu.

\begin{refs}
\item Contexte : Angle inscrit --- Wikipédia (\url{https://fr.wikipedia.org/wiki/Angle_inscrit})
\item Contexte : Angle inscrit dans un demi-cercle --- Wikipédia (\url{https://fr.wikipedia.org/wiki/Angle_inscrit_dans_un_demi-cercle})
\item Lecture : H. S. M. Coxeter, \textit{Introduction to Geometry}, ch. 1
\end{refs}

\bigskip

\begin{problem}[{Le théorème de Ptolémée --- Lycée}]
\textbf{GEO-02.} Soit $ ABCD $ un quadrilatère inscrit dans un cercle.
\begin{enumerate}
\item[(a)] Démontrez le théorème de Ptolémée :
\[ AC \cdot BD = AB \cdot CD + AD \cdot BC. \]
\item[(b)] Déduisez-en le théorème de Pythagore comme cas particulier.
\item[(c)] Pour aller plus loin, démontrez la forme en inégalité : pour \textit{quatre points quelconques}
$ A, B, C, D $ du plan, $ AC \cdot BD \le AB \cdot CD + AD \cdot BC $, avec
égalité exactement lorsque $ ABCD $ est inscriptible dans cet ordre.
\end{enumerate}

\end{problem}

\noindent Claude Ptolémée a démontré cette identité dans l'\textit{Almageste}
(IIe siècle apr. J.-C.) : elle était le c\oe{}ur calculatoire de sa table des cordes ---
en pratique une table des sinus construite pour l'astronomie. Appliquée à des quadrilatères
bien choisis, elle fournit les formules d'addition des sinus, et c'est ainsi que les identités
trigonométriques ont d'abord été obtenues. La démonstration classique introduit un point
auxiliaire ; le trouver est un rite de passage.

\begin{refs}
\item Contexte : Théorème de Ptolémée --- Wikipédia (\url{https://fr.wikipedia.org/wiki/Th%C3%A9or%C3%A8me_de_Ptol%C3%A9m%C3%A9e})
\item Lecture : H. S. M. Coxeter, \textit{Introduction to Geometry}, ch. 1
\end{refs}

\bigskip

\begin{problem}[{La droite d'Euler --- Lycée}]
\textbf{GEO-03.} Dans un triangle non équilatéral quelconque, soient $ O $ le centre du
cercle circonscrit (intersection des médiatrices), $ G $ le centre de gravité (intersection
des médianes) et $ H $ l'orthocentre (intersection des hauteurs).
\begin{enumerate}
\item[(a)] Démontrez que chacun de ces trois points existe bel et bien --- c'est-à-dire que les triplets
de droites concernés sont concourants.
\item[(b)] Démontrez que $ O, G, H $ sont alignés et que $ OG : GH = 1 : 2 $.
\end{enumerate}

\end{problem}

\noindent Euler a démontré cet alignement en 1765, deux millénaires après Euclide :
rappel saisissant que la géométrie élémentaire du triangle n'était pas close. La démonstration
moderne, élégante, repose sur une homothétie bien choisie de centre $ G $ ; découvrir cette
application enseigne davantage sur les transformations qu'un chapitre de théorie. La droite
porte aujourd'hui le nom d'Euler et passe par plusieurs autres centres remarquables.

\begin{refs}
\item Contexte : Droite d'Euler --- Wikipédia (\url{https://fr.wikipedia.org/wiki/Droite_d%27Euler})
\item Lecture : H. S. M. Coxeter, \textit{Introduction to Geometry}, ch. 1
\end{refs}

\bigskip

\begin{problem}[{Le cercle des neuf points --- Lycée}]
\textbf{GEO-04.} Considérez un triangle quelconque, d'orthocentre $ H $ et de centre du
cercle circonscrit $ O $.
\begin{enumerate}
\item[(a)] Démontrez que les neuf points suivants sont sur un même cercle : les milieux des trois
côtés, les pieds des trois hauteurs, et les milieux des trois segments joignant les sommets
à $ H $.
\item[(b)] Démontrez de plus que le centre de ce cercle est le milieu de $ OH $ (sur la droite
d'Euler de GEO-03) et que son rayon vaut la moitié du rayon du cercle circonscrit.
\end{enumerate}

\end{problem}

\noindent Le théorème complet a été publié par Karl Feuerbach en 1822, après des
versions partielles dues à Brianchon et Poncelet ; Feuerbach a ensuite montré que ce cercle
est tangent au cercle inscrit et aux trois cercles exinscrits, l'un des plus beaux résultats
de la géométrie classique. Une démonstration propre compose l'idée d'homothétie de GEO-03 avec
le théorème de l'angle inscrit --- la récompense d'avoir traité honnêtement les deux problèmes
précédents.

\begin{refs}
\item Contexte : Cercle d'Euler (cercle des neuf points) --- Wikipédia (\url{https://fr.wikipedia.org/wiki/Cercle_d%27Euler})
\item Lecture : H. S. M. Coxeter, \textit{Introduction to Geometry}, ch. 1
\end{refs}

\bigskip

\begin{problem}[{Isopérimétrie, premiers pas --- Lycée}]
\textbf{GEO-05.} Le but à l'horizon est le théorème isopérimétrique : parmi toutes les
courbes fermées de longueur donnée $ L $, le cercle enferme la plus grande aire, à savoir
$ L^2/4\pi $. Montez-y par étapes.
\begin{enumerate}
\item[(a)] Parmi tous les rectangles de périmètre donné, démontrez que le carré a la plus grande
aire.
\item[(b)] Parmi tous les triangles de base donnée et de périmètre donné, démontrez que l'isocèle a
la plus grande aire, et déduisez-en que parmi tous les triangles de périmètre donné c'est
l'équilatéral qui maximise l'aire.
\item[(c)] Démontrez qu'une région d'aire maximale à périmètre fixé est nécessairement convexe.
\item[(d)] Expliquez précisément ce qui manque encore à (a)--(c) si l'on veut le théorème
isopérimétrique complet.
\end{enumerate}

\end{problem}

\noindent Le problème isopérimétrique est antique --- la tradition grecque le rattache
à la reine Didon, à qui l'on accorda autant de terre qu'elle pourrait en entourer avec une
peau de b\oe{}uf découpée en lanières. Zénodore (IIe siècle av. J.-C.) comparait les
polygones ; Jakob Steiner (1838) donna de superbes arguments de symétrisation, dont l'astuce
de réflexion qui résout (c) --- mais l'argument de Steiner supposait, on le sait bien, l'existence
d'un maximiseur, lacune qui ne fut comblée que par Weierstrass. La question (d) vous demande de
trouver cette lacune vous-même : c'est une première leçon sur l'importance des démonstrations
d'existence.

\begin{refs}
\item Contexte : Théorème isopérimétrique --- Wikipédia (\url{https://fr.wikipedia.org/wiki/Th%C3%A9or%C3%A8me_isop%C3%A9rim%C3%A9trique})
\item Lecture : H. S. M. Coxeter, \textit{Introduction to Geometry}
\end{refs}

\bigskip

\begin{problem}[{La puissance d'un point --- Lycée, short}]
\textbf{GEO-23.} Soit $ P $ un point extérieur à un cercle, et soit une droite passant par $ P $ qui coupe le cercle en $ A $ et $ B $. Démontrez que le produit $ PA \cdot PB $ est le même pour toute droite de ce type passant par $ P $.
\end{problem}

\noindent Cet invariant est la puissance du point, et il transforme une famille de configurations apparemment sans lien en un seul nombre attaché à $ P $ seul. C'est le moteur de l'axe radical et de la plupart des problèmes d'olympiade mettant en jeu deux cercles, et la démonstration se réduit à une application des triangles semblables.

\begin{refs}
\item Contexte : Puissance d'un point par rapport à un cercle --- Wikipédia (\url{https://fr.wikipedia.org/wiki/Puissance_d%27un_point_par_rapport_%C3%A0_un_cercle})
\end{refs}

\bigskip

\begin{problem}[{Compter les points entiers dans un polygone --- Lycée, short}]
\textbf{GEO-24.} Vérifiez le théorème de Pick --- un polygone simple dont les sommets sont sur le réseau des points entiers a pour aire $ I + B/2 - 1 $, où $ I $ est le nombre de points entiers intérieurs et $ B $ le nombre de ceux du bord --- sur cinq polygones de votre invention, dont un non convexe. Démontrez-le ensuite pour les rectangles à côtés parallèles aux axes.
\end{problem}

\noindent Le théorème de Pick, démontré en 1899, relie une quantité continue (l'aire) à un dénombrement purement discret : c'est le genre de pont qui donnera plus tard toute la géométrie des points d'un réseau. Le vérifier sur un exemple non convexe est important : l'intuition de la plupart des gens suppose discrètement la convexité, alors que le théorème n'en a pas besoin.

\begin{refs}
\item Contexte : Théorème de Pick --- Wikipédia (\url{https://fr.wikipedia.org/wiki/Th%C3%A9or%C3%A8me_de_Pick})
\end{refs}

\bigskip

\begin{problem}[{Le théorème des trisectrices de Morley --- Lycée}]
\textbf{GEO-31.} Dans un triangle $ ABC $, partagez chacun des trois angles intérieurs en trois
parts égales. Pour chaque côté, prenez les deux trisectrices les plus proches de lui --- une issue de
chaque extrémité de ce côté --- et soient $ P_A, P_B, P_C $ les points où ces couples se rencontrent,
$ P_A $ étant celui associé au côté $ BC $.
\begin{enumerate}
\item[(a)] Démontrez que chacun de ces couples de trisectrices se rencontre effectivement à l'intérieur du
triangle, de sorte que les trois points existent.
\item[(b)] Démontrez le théorème de Morley : le triangle $ P_A P_B P_C $ est équilatéral, quel que soit le
triangle de départ. (Une voie trigonométrique utilise la loi des sinus jointe à l'identité
$ \sin 3\theta = 4 \sin\theta \, \sin(60^\circ + \theta) \, \sin(60^\circ - \theta) $.)
\item[(c)] Démontrez que le côté du triangle de Morley a pour longueur
\[ 8R \, \sin\tfrac{A}{3} \, \sin\tfrac{B}{3} \, \sin\tfrac{C}{3} , \]
où $ R $ est le rayon du cercle circonscrit à $ ABC $ et $ A, B, C $ ses angles.
\item[(d)] Autorisez également les trisectrices des angles \textit{extérieurs}. Exhibez au moins un autre
triangle équilatéral obtenu de cette façon --- il y a en tout dix-huit triangles de Morley --- et dites
lesquels de vos arguments de (b) se transposent sans changement.
\end{enumerate}

\end{problem}

\noindent Frank Morley, un Anglais qui fit toute sa carrière à Haverford puis à Johns Hopkins,
remarqua ce fait vers 1899 en étudiant la géométrie des cardioïdes tangentes aux côtés d'un triangle ;
il en parla à des amis et le résultat circula de bouche à oreille bien avant sa publication, la
première démonstration imprimée étant due à M. T. Naraniengar en 1909. Ce retard n'a rien d'un
hasard : la trisection de l'angle est impossible à la règle et au compas, si bien que la configuration
ne peut être tracée par les moyens classiques et qu'aucun Grec n'aurait pu tomber dessus. Le théorème
attire encore de nouvelles démonstrations --- Alain Connes en a publié une, algébrique, en 1998, qui
vaut sur tout corps de caractéristique différente de deux et de trois.

\begin{refs}
\item Contexte : Théorème de Morley --- Wikipédia (\url{https://fr.wikipedia.org/wiki/Th%C3%A9or%C3%A8me_de_Morley})
\item Contexte : Trisection de l'angle --- Wikipédia (\url{https://fr.wikipedia.org/wiki/Trisection_de_l%27angle})
\item Lecture : H. S. M. Coxeter et S. L. Greitzer, \textit{Geometry Revisited}, ch. 1
\item Lecture : H. S. M. Coxeter, \textit{Introduction to Geometry}, ch. 1
\end{refs}

\bigskip

\begin{problem}[{Deux bissectrices de même longueur --- Lycée, short}]
\textbf{GEO-32.} Démontrez le théorème de Steiner--Lehmus : si les bissectrices
intérieures issues de deux sommets d'un triangle ont la même longueur, alors le triangle est
isocèle. (La réciproque est une conséquence immédiate de la symétrie ; toute la difficulté est dans
ce sens-ci.)
\end{problem}

\noindent C. L. Lehmus demanda à Sturm, dans une lettre de 1840, une démonstration purement
géométrique ; Sturm fit circuler la question, et Jakob Steiner fut parmi les premiers à répondre ---
depuis, le problème a suscité bien plus d'une centaine de démonstrations publiées. Presque toutes
raisonnent par l'absurde, en prouvant plutôt qu'un triangle non isocèle a des bissectrices de
longueurs différentes, et la possibilité d'une démonstration véritablement \og{} directe \fg{} est devenue
une petite controverse logique à part entière : John Conway soutenait qu'il n'en existe aucune,
Victor Pambuccian a démontré ensuite que des démonstrations directes doivent exister en logique
classique comme en logique intuitionniste, et l'on en a depuis écrit une vérifiée par
machine.

\begin{refs}
\item Contexte : Théorème de Steiner-Lehmus --- Wikipédia (\url{https://fr.wikipedia.org/wiki/Th%C3%A9or%C3%A8me_de_Steiner-Lehmus})
\item Lecture : H. S. M. Coxeter et S. L. Greitzer, \textit{Geometry Revisited}, ch. 1
\end{refs}

\bigskip


\section*{Combinatoire et théorie des graphes}

\begin{problem}[{Classiques du principe des tiroirs --- Lycée}]
\textbf{CMB-01.} Le principe des tiroirs dit : si l'on range plus de $ n $ objets dans
$ n $ tiroirs, un tiroir en contient au moins deux. Servez-vous-en pour démontrer :
\begin{enumerate}
\item[(a)] parmi $ n+1 $ entiers quelconques choisis dans $ \{1, 2, \dots, 2n\} $, deux sont
consécutifs, donc deux sont premiers entre eux ;
\item[(b)] parmi $ n+1 $ entiers quelconques choisis dans $ \{1, 2, \dots, 2n\} $, l'un divise un
autre ;
\item[(c)] parmi $ n+1 $ points quelconques d'un triangle équilatéral de côté 2, deux sont à distance
au plus $ 2/\sqrt{n} $\dots{} en fait, démontrez le cas propre : parmi 5 points quelconques d'un
triangle équilatéral de côté 2, deux sont à distance au plus 1 l'un de l'autre.
\end{enumerate}

\end{problem}

\noindent Dirichlet a utilisé ce principe dans les années 1830--1840 (il l'appelait
\textit{Schubfachprinzip}, \og{} principe des tiroirs \fg{}) pour démontrer des faits profonds sur
l'approximation des irrationnels par des fractions. Les questions (a) et (b) étaient des
favorites de Paul Erd\H{o}s, qui aimait tester les enfants prodiges avec (b) --- le jeune Lajos Pósa
l'aurait résolue pendant le dîner. Le principe a l'air d'une trivialité ; apprendre tout ce qu'on
peut lui faire porter est une première leçon de levier mathématique.

\begin{refs}
\item Contexte : Principe des tiroirs --- Wikipédia (\url{https://fr.wikipedia.org/wiki/Principe_des_tiroirs})
\item Lecture : M. Aigner et G. M. Ziegler, \textit{Proofs from THE BOOK}, ch. \og{} Pigeon-hole and double counting \fg{}
\item Lecture : M. Bóna, \textit{A Walk Through Combinatorics}, ch. 1
\end{refs}

\bigskip

\begin{problem}[{Dénombrer les sous-ensembles, avec démonstration --- Lycée}]
\textbf{CMB-02.} \begin{enumerate}
\item[(a)] Démontrez qu'un ensemble à $ n $ éléments a exactement $ 2^n $ sous-ensembles --- donnez
deux démonstrations véritablement différentes, l'une par construction d'une bijection avec les
mots binaires, l'autre par récurrence.
\item[(b)] Démontrez que $ \binom{n}{0} + \binom{n}{1} + \cdots + \binom{n}{n} = 2^n $ en expliquant
ce que dénombrent les deux membres.
\item[(c)] Démontrez que pour $ n \ge 1 $ exactement la moitié des sous-ensembles sont de cardinal
pair, autrement dit $ \binom{n}{0} - \binom{n}{1} + \binom{n}{2} - \cdots \pm \binom{n}{n} = 0 $,
en exhibant une bijection entre les sous-ensembles de cardinal pair et ceux de cardinal
impair.
\end{enumerate}

\end{problem}

\noindent Les coefficients binomiaux sont anciens --- la prosodie indienne de Pingala
(vers 200 av. J.-C.), le triangle \og{} de Pascal \fg{} connu en Chine (Yang Hui, XIIIe siècle)
et en Perse (al-Karaji, vers l'an mille) bien avant le traité de Pascal de 1654. Le vrai contenu
de ce problème est méthodologique : une identité combinatoire se démontre au mieux en faisant
compter la même chose aux deux membres, et un énoncé de parité par un appariement. Ces deux
gestes --- le double comptage et la bijection --- sont le pain quotidien du domaine.

\begin{refs}
\item Contexte : Coefficient binomial --- Wikipédia (\url{https://fr.wikipedia.org/wiki/Coefficient_binomial})
\item Lecture : M. Bóna, \textit{A Walk Through Combinatorics}, ch. 3--4
\item Cours : MIT OCW 6.042J Mathematics for Computer Science (\url{https://ocw.mit.edu/})
\end{refs}

\bigskip

\begin{problem}[{Le lemme des poignées de main --- Lycée}]
\textbf{CMB-03.} Un graphe est un ensemble de sommets dont certaines paires sont jointes par
des arêtes ; le degré d'un sommet est le nombre d'arêtes qui y aboutissent.
\begin{enumerate}
\item[(a)] Démontrez le lemme des poignées de main : la somme de tous les degrés vaut le double du
nombre d'arêtes, de sorte que dans toute soirée le nombre de personnes ayant serré un nombre
impair de mains est pair.
\item[(b)] Démontrez que dans toute soirée de $ n \ge 2 $ personnes, deux invités ont serré le même
nombre de mains au sein de la soirée.
\item[(c)] Montrez par un exemple que (b) peut tomber en défaut si l'on exige le \og{} même nombre \fg{} pour
\textit{trois} invités.
\end{enumerate}

\end{problem}

\noindent Le lemme des poignées de main est implicite dans l'article de 1736 d'Euler sur
Königsberg --- le premier théorème de la théorie des graphes, démontré avant que le mot \og{} graphe \fg{}
n'existe (le nom vient de Sylvester, en 1878). La question (b) est un exercice classique
d'application du principe des tiroirs aux graphes, avec une petite subtilité qui vous force à
regarder de près les degrés extrêmes 0 et $ n-1 $. Compter un ensemble de choses de deux
manières, comme en (a), ne cesse jamais d'être utile.

\begin{refs}
\item Contexte : Lemme des poignées de main --- Wikipédia (\url{https://fr.wikipedia.org/wiki/Lemme_des_poign%C3%A9es_de_main})
\item Lecture : D. West, \textit{Introduction to Graph Theory}, ch. 1
\end{refs}

\bigskip

\begin{problem}[{L'échiquier mutilé --- Lycée}]
\textbf{CMB-04.} Retirez deux cases d'angle diagonalement opposées d'un
échiquier ordinaire $ 8 \times 8 $. Démontrez que les 62 cases restantes ne peuvent pas être
pavées par 31 dominos couvrant chacun deux cases adjacentes. Démontrez ensuite le théorème plus
fin de Gomory : si vous retirez à la place une case noire quelconque et une case blanche
quelconque, les 62 cases restantes \textit{peuvent} toujours être pavées.
\end{problem}

\noindent Popularisé par Max Black (1946) puis par les chroniques de Martin Gardner,
l'échiquier mutilé est l'exemple canonique d'argument par invariant : une seule quantité bien
choisie, inchangée par chaque coup légal, tue d'un coup une infinité de pavages candidats. C'est
aussi une pierre de touche en intelligence artificielle et en théorie de la démonstration --- John
McCarthy l'a proposé comme défi aux démonstrateurs automatiques précisément parce que la
démonstration humaine tient en une phrase et que la recherche exhaustive est astronomique.

\begin{refs}
\item Contexte : Problème de l'échiquier mutilé --- Wikipédia (\url{https://fr.wikipedia.org/wiki/Probl%C3%A8me_de_l%27%C3%A9chiquier_mutil%C3%A9})
\item Lecture : M. Aigner et G. M. Ziegler, \textit{Proofs from THE BOOK}
\end{refs}

\bigskip

\begin{problem}[{Pavages par triominos --- Lycée}]
\textbf{CMB-05.} Un triomino en L est la forme faite de trois carrés unités
disposés en L. Démontrez que pour tout $ n \ge 1 $, un plateau $ 2^n \times 2^n $ privé d'une
case quelconque peut être pavé par des triominos en L. Décidez ensuite, avec démonstration : pour
quelles positions de la case retirée un plateau $ 5 \times 5 $ privé d'une case peut-il être
pavé par des triominos en L ?
\end{problem}

\noindent Le résultat sur les plateaux $ 2^n \times 2^n $ est dû à Solomon Golomb
(1954), qui a forgé le mot \og{} polyomino \fg{} et lancé une petite industrie de mathématiques du pavage
que Martin Gardner a portée jusqu'à des millions de lecteurs. La démonstration attendue est un
bijou de manuel sur la récurrence --- le pas inductif contient une idée véritablement créative, et
la trouver est tout l'enjeu. La partie $ 5 \times 5 $ montre qu'à côté des constructions
astucieuses il faut des arguments de coloration (CMB-04) pour démontrer une impossibilité.

\begin{refs}
\item Contexte : Triomino --- Wikipédia (\url{https://fr.wikipedia.org/wiki/Triomino})
\item Contexte : Polyomino --- Wikipédia (\url{https://fr.wikipedia.org/wiki/Polyomino})
\item Lecture : M. Bóna, \textit{A Walk Through Combinatorics}, ch. 2 (récurrence)
\end{refs}

\bigskip

\begin{problem}[{Chemins dans un réseau, étoiles et barres --- Lycée}]
\textbf{CMB-06.} \begin{enumerate}
\item[(a)] Démontrez que le nombre de plus courts chemins de $ (0,0) $ à $ (m,n) $ par pas unités
vers la droite ou vers le haut vaut $ \binom{m+n}{n} $, au moyen d'une bijection explicite avec
des mots.
\item[(b)] Démontrez le théorème des \og{} étoiles et barres \fg{} : le nombre de façons d'écrire $ n $ comme
somme ordonnée de $ k $ entiers positifs ou nuls vaut $ \binom{n+k-1}{k-1} $.
\item[(c)] Utilisez (a) pour démontrer l'identité $ \binom{m+n}{n} = \sum_i \binom{m}{i}\binom{n}{i} $\dots{}
non --- démontrez plutôt la relation de Pascal, plus simple,
$ \binom{m+n}{n} = \binom{m+n-1}{n} + \binom{m+n-1}{n-1} $, en partageant les chemins selon leur
dernier pas.
\end{enumerate}

\end{problem}

\noindent Le dénombrement de chemins est la porte la plus accueillante vers
l'énumération sérieuse : la même figure démontre la relation de Pascal, l'identité de Vandermonde
(CMB-07) et --- après une astuce de réflexion --- le théorème du scrutin et les nombres de Catalan
(CMB-09). La méthode des étoiles et des barres fut popularisée par le classique de probabilités
de William Feller (1950) ; c'est le calcul d'occupation qui sous-tend une bonne part de la
mécanique statistique.

\begin{refs}
\item Contexte : Chemin dans un réseau --- Wikipédia (\url{https://fr.wikipedia.org/wiki/Chemin_dans_un_r%C3%A9seau})
\item Contexte : Méthode des étoiles et des barres --- Wikipédia (\url{https://fr.wikipedia.org/wiki/M%C3%A9thode_des_%C3%A9toiles_et_des_barres})
\item Lecture : M. Bóna, \textit{A Walk Through Combinatorics}, ch. 3
\end{refs}

\bigskip

\begin{problem}[{Deux personnes ayant le même nombre d'amis --- Lycée, short}]
\textbf{CMB-27.} Démontrez que dans tout groupe de $ n \geq 2 $ personnes, au moins deux d'entre elles ont le même nombre d'amis au sein du groupe (l'amitié étant réciproque).
\end{problem}

\noindent Le principe des tiroirs marche presque immédiatement --- il y a $ n $ personnes et $ n $ nombres d'amis possibles, de $ 0 $ à $ n-1 $ --- et tout le problème consiste à remarquer que $ 0 $ et $ n-1 $ ne peuvent pas coexister. Cette observation, et non le principe des tiroirs lui-même, est le contenu mathématique.

\begin{refs}
\item Contexte : Principe des tiroirs --- Wikipédia (\url{https://fr.wikipedia.org/wiki/Principe_des_tiroirs})
\end{refs}

\bigskip

\begin{problem}[{Une poignée de main qui doit être paire --- Lycée, short}]
\textbf{CMB-28.} Démontrez que dans tout graphe fini le nombre de sommets de degré impair est pair. Servez-vous-en ensuite pour montrer qu'il est impossible qu'exactement neuf personnes à une soirée serrent chacune la main d'exactement trois autres.
\end{problem}

\noindent Le lemme des poignées de main est le premier théorème de la théorie des graphes et découle du décompte des extrémités d'arêtes de deux façons, technique --- le double comptage --- qui continuera de rapporter des années durant. L'application illustre l'habitude utile : une démonstration d'impossibilité obtenue en trouvant un invariant que la configuration souhaitée devrait violer.

\begin{refs}
\item Contexte : Lemme des poignées de main --- Wikipédia (\url{https://fr.wikipedia.org/wiki/Lemme_des_poign%C3%A9es_de_main})
\end{refs}

\bigskip

\begin{problem}[{Des poignées de main qui ne se croisent jamais --- Lycée, short}]
\textbf{CMB-35.} Placez $ 2n $ personnes en des points distincts d'un cercle et appariez-les, chacune serrant la main d'exactement une autre ; tracez chaque poignée de main comme une corde rectiligne. Démontrez que le nombre de tels appariements où deux cordes ne se croisent jamais est le nombre de Catalan $ C_n = \frac{1}{n+1}\binom{2n}{n} $.
\end{problem}

\noindent Les diagrammes de cordes non croisées sont un membre de plus de l'énorme famille dénombrée par les nombres de Catalan --- le livre que Richard Stanley leur a consacré recense 214 interprétations combinatoires, et celle-ci ne figure pas parmi les trois de CMB-09, si bien qu'il vous faut la relier aux autres par un argument que vous trouverez vous-même. Les deux voies fonctionnent : construire une bijection avec les parenthésages équilibrés, ou découvrir la récurrence en se demandant ce que la corde d'une personne donnée fait au reste du cercle. Les structures non croisées ne sont pas une curiosité : le treillis des partitions non croisées, dégagé par Germain Kreweras en 1972, s'est révélé des décennies plus tard être l'ossature combinatoire des probabilités libres.

\begin{refs}
\item Contexte : Nombre de Catalan --- Wikipédia (\url{https://fr.wikipedia.org/wiki/Nombre_de_Catalan})
\item Contexte : Partition non croisée --- Wikipédia (\url{https://fr.wikipedia.org/wiki/Partition_non_crois%C3%A9e})
\item Lecture : R. P. Stanley, \textit{Catalan Numbers} (2015)
\end{refs}

\bigskip

\begin{problem}[{Dénombrer les colorations : chemins, cycles, arbres --- Lycée}]
\textbf{CMB-36.} Pour un graphe fini $ G $ à $ n $ sommets et un entier strictement positif $ k $, soit $ P_G(k) $ le nombre de façons de donner à chaque sommet l'une de $ k $ couleurs de sorte que deux sommets adjacents reçoivent toujours des couleurs différentes.
\begin{enumerate}
\item[(a)] Démontrez que $ P_T(k) = k(k-1)^{n-1} $ pour tout arbre $ T $ à $ n $ sommets.
\item[(b)] Démontrez que pour le cycle $ C_n $ avec $ n \ge 3 $,
\[ P_{C_n}(k) = (k-1)^n + (-1)^n (k-1). \]
\item[(c)] Démontrez la règle de suppression--contraction : pour toute arête $ e $ de $ G $,
\[ P_G(k) = P_{G-e}(k) - P_{G/e}(k), \]
où $ G-e $ est $ G $ privé de $ e $ et $ G/e $ est $ G $ dont les deux extrémités de $ e $ ont été fusionnées en un seul sommet.
\item[(d)] Déduisez de (c), par récurrence sur le nombre d'arêtes, que $ P_G(k) $ coïncide avec un polynôme en $ k $ de degré $ n $, de coefficient dominant 1, dont les coefficients alternent en signe, et confrontez-y votre formule de (b).
\end{enumerate}

\end{problem}

\noindent George Birkhoff introduisit cette fonction de dénombrement en 1912 avec un espoir précis : si l'on parvenait à montrer que $ P_G(4) > 0 $ pour tout $ G $ planaire, le théorème des quatre couleurs (CMB-17) découlerait de l'algèbre plutôt que d'une analyse de cas. L'espoir a échoué, mais le polynôme chromatique a survécu et prospéré --- Hassler Whitney a dégagé le sens de ses coefficients en 1932, et en 1968 Ronald Read a conjecturé qu'ils croissent toujours puis décroissent sans creux. Cette conjecture d'apparence innocente a tenu quarante-quatre ans, jusqu'à ce que June Huh démontre en 2012 la log-concavité qui la sous-tend au moyen de la géométrie algébrique, travaux qui ont contribué à sa médaille Fields de 2022. Tout ce problème se traite au crayon ; l'objet qu'il produit a été trois fois une frontière de la recherche.

\begin{refs}
\item Contexte : Polynôme chromatique --- Wikipédia (\url{https://fr.wikipedia.org/wiki/Polyn%C3%B4me_chromatique})
\item Contexte : Formule de suppression--contraction --- Wikipédia (en anglais) (\url{https://en.wikipedia.org/wiki/Deletion%E2%80%93contraction_formula})
\item Lecture : D. West, \textit{Introduction to Graph Theory}, ch. 5
\end{refs}

\bigskip


\section*{Calcul différentiel et intégral}

\begin{problem}[{La limite, à partir des premiers principes --- Lycée}]
\textbf{CAL-01.} Calculez
$ \lim_{x\to 1} \frac{x^2-1}{x-1} $,
en expliquant pourquoi la réponse n'est pas \og{} $0/0$ \fg{} et pourquoi simplifier par $ x-1 $ est
légitime bien que la fonction ne soit pas définie en $ x=1 $. Énoncez ensuite la définition
$\varepsilon$--$\delta$ de $ \lim_{x\to a} f(x)=L $, et démontrez directement à partir d'elle :
\begin{enumerate}
\item[(i)] que la limite que vous
avez calculée est correcte ;
\item[(ii)] la loi de la somme
$ \lim (f+g) = \lim f + \lim g $ et la loi du produit pour les limites --- les deux lois que
tout calcul ultérieur de cette page utilise en silence.
\end{enumerate}

\end{problem}

\noindent Newton et Leibniz calculaient avec des quantités \og{} infiniment petites \fg{}
et furent moqués pour cela --- Berkeley (1734) appelait leurs accroissements évanouissants
\og{} les fantômes de quantités défuntes \fg{}. La réponse moderne a demandé 150 ans : le
\textit{Cours d'analyse} de Cauchy (1821) a placé la limite au fondement, et la formulation
$\varepsilon$--$\delta$ de Weierstrass (cours berlinois des années 1860) l'a rendue inattaquable. Démontrer les lois
des limites une fois pour toutes est le droit d'entrée à tout ce qui suit.

\begin{refs}
\item Contexte : Limite d'une fonction --- Wikipédia (\url{https://fr.wikipedia.org/wiki/Limite_%28math%C3%A9matiques_%C3%A9l%C3%A9mentaires%29})
\item Lecture : Michael Spivak, \textit{Calculus}, ch. 5 (les limites)
\item Cours : MIT OCW 18.01SC Single Variable Calculus (\url{https://ocw.mit.edu/courses/18-01sc-single-variable-calculus-fall-2010/})
\item Vidéo : 3Blue1Brown, Essence of Calculus (\url{https://www.3blue1brown.com/lessons/essence-of-calculus/})
\end{refs}

\bigskip

\begin{problem}[{La dérivée est une limite --- Lycée}]
\textbf{CAL-02.} Posez
\[ f'(a) = \lim_{h\to 0} \frac{f(a+h)-f(a)}{h}. \]
Travaillez uniquement à partir de cette définition --- aucune règle.
\begin{enumerate}
\item[(a)] Calculez les dérivées de $ f(x)=x^2 $, $ f(x)=\sqrt{x} $ et
$ f(x)=1/x $.
\item[(b)] Démontrez la règle du produit $ (fg)' = f'g + fg' $ directement à partir de la définition, en précisant
exactement quelles lois des limites de CAL-01 vous invoquez.
\item[(c)] Démontrez que la dérivabilité de $ f $ en un point y entraîne sa continuité ---
la partie (b) en a besoin.
\end{enumerate}

\end{problem}

\noindent Les premiers manuscrits privés de Leibniz (1675) ont brièvement supposé que
$ d(uv) = du \cdot dv $ avant qu'il ne se corrige --- la règle du produit est bel et bien un
théorème, non une trivialité de notation. Newton appelait les dérivées des \og{} fluxions \fg{} de
\og{} fluentes \fg{} ; la définition par le taux d'accroissement ci-dessus est celle de Cauchy. Faire
ces calculs honnêtement, une fois, est ce qui sépare le fait de savoir l'analyse du fait de s'en souvenir.

\begin{refs}
\item Contexte : Dérivée --- Wikipédia (\url{https://fr.wikipedia.org/wiki/D%C3%A9riv%C3%A9e})
\item Lecture : Michael Spivak, \textit{Calculus}, ch. 9--10
\item Cours : MIT OCW 18.01SC Single Variable Calculus (\url{https://ocw.mit.edu/courses/18-01sc-single-variable-calculus-fall-2010/})
\item Vidéo : 3Blue1Brown, Essence of Calculus (\url{https://www.3blue1brown.com/lessons/essence-of-calculus/})
\end{refs}

\bigskip

\begin{problem}[{Les tangentes avant l'analyse --- Lycée}]
\textbf{CAL-03.} Travaillez uniquement avec de l'algèbre --- ni limites, ni dérivées.
\begin{enumerate}
\item[(a)] Trouvez la tangente à la parabole $ y=x^2 $ au point $ (a, a^2) $ : exigez
que la droite $ y = mx + c $ rencontre la parabole en une racine double en $ x=a $, et résolvez
en $ m $ et $ c $.
\item[(b)] Faites de même pour $ y=x^3 $ en $ (a, a^3) $.
\item[(c)] Vérifiez que les deux réponses s'accordent avec la dérivée, puis expliquez précisément pourquoi
la méthode de la racine double fonctionne pour les polynômes mais ne peut pas définir la tangente à
$ y=\sin x $.
\end{enumerate}

\end{problem}

\noindent C'est ainsi que l'on trouvait les tangentes avant l'existence de l'analyse. La
méthode des normales de Descartes (\textit{La Géométrie}, 1637) et l'\textit{adégalité} de Fermat (diffusée
dans les années 1630) résolvaient des problèmes de tangente et d'extremum pour les courbes algébriques une génération avant que
Newton ne naisse ; leur limitation aux courbes algébriques est exactement ce que la notion de limite
a levé. La méthode de Fermat anticipe la dérivée d'assez près pour que Lagrange lui en ait attribué
l'invention.

\begin{refs}
\item Contexte : Adégalité --- Wikipédia (\url{https://fr.wikipedia.org/wiki/Ad%C3%A9galit%C3%A9})
\item Contexte : Tangente --- Wikipédia (\url{https://fr.wikipedia.org/wiki/Tangente_%28g%C3%A9om%C3%A9trie%29})
\item Lecture : Carl B. Boyer, \textit{The History of the Calculus and Its Conceptual Development}
\end{refs}

\bigskip

\begin{problem}[{Le plus grand champ qu'une clôture puisse enfermer --- Lycée}]
\textbf{CAL-04.} Vous disposez de $ L $ mètres de clôture.
\begin{enumerate}
\item[(i)] Démontrez que
parmi tous les rectangles de périmètre $ L $, le carré enferme la plus grande aire --- deux fois :
une fois avec l'analyse, et une fois sans analyse du tout, à l'aide de l'inégalité arithmético-géométrique
$ \sqrt{xy} \le \frac{x+y}{2} $ (démontrez-la également).
\item[(ii)] Cette fois la clôture ferme un
champ rectangulaire contre le long mur droit d'une grange, de sorte que seuls trois côtés sont à clôturer :
trouvez les proportions optimales, avec démonstration.
\item[(iii)] Expliquez pourquoi trouver un point où
$ f'=0 $ n'est \textit{pas} à soi seul une démonstration d'optimalité, et complétez la justification
dans les deux cas (comportement au bord, dérivée seconde, ou concavité).
\end{enumerate}

\end{problem}

\noindent L'optimisation est la plus ancienne publicité pour l'analyse. L'intuition
isopérimétrique --- plus c'est symétrique, plus cela contient --- remonte à l'Antiquité (le problème de Didon dans
l'\textit{Énéide} de Virgile ; Zénodore, vers 180 av. J.-C.), et la \textit{Nova stereometria
doliorum vinariorum} de Kepler (1615) optimisait les proportions des tonneaux avant l'existence des dérivées.
La partie (iii) est l'âme du problème : un point critique est un suspect, pas un verdict.

\begin{refs}
\item Contexte : Théorème isopérimétrique --- Wikipédia (\url{https://fr.wikipedia.org/wiki/Th%C3%A9or%C3%A8me_isop%C3%A9rim%C3%A9trique})
\item Lecture : James Stewart, \textit{Calculus}, ch. 4 (applications de la dérivation)
\item Lecture : Michael Spivak, \textit{Calculus}, ch. 11
\end{refs}

\bigskip

\begin{problem}[{La boîte de conserve la plus économique --- Lycée}]
\textbf{CAL-05.} Une boîte cylindrique fermée doit contenir un volume fixé
$ V $. Démontrez que l'aire totale de sa surface (le côté plus les deux couvercles) est minimale exactement lorsque
la hauteur égale le diamètre, $ h = 2r $. Justifiez que votre point critique est un véritable
minimum global sur $ (0,\infty) $ --- examinez ce qu'il advient de l'aire lorsque $ r\to 0^+ $
puis lorsque $ r\to\infty $. Mesurez ensuite une vraie boîte de votre cuisine et déterminez à quelle distance
de l'optimum elle se trouve ; suggérez, en une phrase, quels coûts autres que la matière pourraient expliquer
l'écart.
\end{problem}

\noindent Le problème de la boîte est l'optimisation sous contrainte classique à une variable :
une contrainte (le volume) élimine une variable et laisse une honnête minimisation à une
variable. Les vraies boîtes sont nettement plus hautes que l'optimum mathématique --- une leçon précoce :
un modèle démontre des théorèmes sur lui-même, et l'étape de modélisation est une responsabilité
distincte.

\begin{refs}
\item Contexte : Optimisation (mathématiques) --- Wikipédia (\url{https://fr.wikipedia.org/wiki/Optimisation_%28math%C3%A9matiques%29})
\item Lecture : James Stewart, \textit{Calculus}, ch. 4.7 (problèmes d'optimisation)
\item Cours : MIT OCW 18.01SC Single Variable Calculus (\url{https://ocw.mit.edu/courses/18-01sc-single-variable-calculus-fall-2010/})
\end{refs}

\bigskip

\begin{problem}[{Les taux liés, honnêtement --- Lycée}]
\textbf{CAL-06.} Une échelle de 5~m est appuyée contre un mur ; son pied
glisse en s'éloignant à 1~m/s.
\begin{enumerate}
\item[(i)] Trouvez la vitesse du sommet lorsque le pied est à 3~m
du mur, en précisant exactement l'étape de dérivation des fonctions composées que votre calcul utilise et le théorème qui
l'autorise.
\item[(ii)] Montrez que lorsque le pied s'approche de 5~m du mur, votre formule
prédit que la vitesse du sommet tend vers l'infini ; expliquez ce que cela démontre au sujet du modèle
(le sommet d'une vraie échelle reste-t-il contre le mur ?).
\item[(iii)] De l'eau s'écoule d'un
cône renversé de demi-angle 30$^\circ$ : reliez $ dV/dt $ au taux de variation de la hauteur d'eau,
en dérivant --- et non en citant --- la relation de la formule du volume $ V = \tfrac{1}{3}\pi r^2 h $
que vous utilisez.
\end{enumerate}

\end{problem}

\noindent Les problèmes de taux liés sont un classique depuis les manuels britanniques du
dix-neuvième siècle issus de la tradition des \og{} mathématiques mixtes \fg{}. Le paradoxe de l'échelle
qui tombe, en (ii), est un échec véritablement instructif : la mathématique est correcte, c'est la
prémisse (le sommet collé au mur) qui cède --- et l'analyse elle-même détecte le moment
où elle doit céder.

\begin{refs}
\item Contexte : Taux liés --- Wikipédia (en anglais) (\url{https://en.wikipedia.org/wiki/Related_rates})
\item Contexte : Dérivation des fonctions composées --- Wikipédia (\url{https://fr.wikipedia.org/wiki/Th%C3%A9or%C3%A8me_de_d%C3%A9rivation_des_fonctions_compos%C3%A9es})
\item Lecture : James Stewart, \textit{Calculus}, ch. 3.9 (taux liés)
\end{refs}

\bigskip

\begin{problem}[{Le théorème du compteur de vitesse --- Lycée}]
\textbf{CAL-07.} Une voiture passe un premier péage à midi et un second,
90~km plus loin, à 13~h. Démontrez qu'à un certain instant son compteur indiquait
exactement 90~km/h. Pour le faire correctement : énoncez le théorème de Rolle ; déduisez-en le
théorème des accroissements finis --- pour $ f $ convenable il existe $ c\in(a,b) $ tel que
$ f'(c) = \frac{f(b)-f(a)}{b-a} $ --- et énoncez les hypothèses de modélisation sur la fonction
position dont votre démonstration a besoin. Utilisez ensuite ce théorème pour démontrer qu'une fonction dont la dérivée
s'annule sur un intervalle y est constante, et expliquez en un paragraphe pourquoi toute la
théorie des primitives (le \og{} $ +\,C $ \fg{}) repose sur ce fait d'apparence anodine.
\end{problem}

\noindent Michel Rolle a énoncé son théorème pour les polynômes en 1691 --- tout en étant
l'un des critiques les plus acérés de l'analyse ; il la traitait de recueil de sophismes ingénieux.
Cauchy a fait du théorème des accroissements finis le moteur de l'analyse rigoureuse : c'est le pont qui mène de
l'information locale (la dérivée en chaque point) aux conclusions globales (ce que la fonction
a fait sur tout l'intervalle). Presque tous les théorèmes de cette bande traversent ce pont.

\begin{refs}
\item Contexte : Théorème des accroissements finis --- Wikipédia (\url{https://fr.wikipedia.org/wiki/Th%C3%A9or%C3%A8me_des_accroissements_finis})
\item Contexte : Théorème de Rolle --- Wikipédia (\url{https://fr.wikipedia.org/wiki/Th%C3%A9or%C3%A8me_de_Rolle})
\item Lecture : Michael Spivak, \textit{Calculus}, ch. 11
\item Cours : MIT OCW 18.01SC Single Variable Calculus (\url{https://ocw.mit.edu/courses/18-01sc-single-variable-calculus-fall-2010/})
\end{refs}

\bigskip

\begin{problem}[{L'aire à partir des premiers principes --- Lycée}]
\textbf{CAL-08.} En utilisant la définition de l'intégrale comme limite de
sommes de Riemann, calculez $ \int_0^1 x^2\,dx $ directement : découpez $ [0,1] $ en $ n $
morceaux égaux, écrivez les sommes supérieures et inférieures, évaluez-les à l'aide de l'identité
\[ 1^2 + 2^2 + \cdots + n^2 = \frac{n(n+1)(2n+1)}{6} \]
(démontrez-la par récurrence), et montrez que les deux sommes convergent vers la même limite. Calculez ensuite
$ \int_0^b x^k\,dx $ pour tout entier $ k \ge 1 $ par l'astuce de Fermat : découpez
$ [0,b] $ non pas régulièrement mais aux points géométriques $ b, br, br^2, \dots $ avec
$ r<1 $, sommez les rectangles obtenus comme une série géométrique, et faites tendre $ r\to 1 $.
\end{problem}

\noindent Archimède a obtenu l'aire sous une parabole vers 250 av. J.-C. par la
méthode d'exhaustion --- la même réponse, deux millénaires avant que les limites ne soient définies. L'astuce
du découpage géométrique de Fermat (années 1630) a quadraturé toutes les courbes $ y=x^k $ d'un coup, encore
sans aucun théorème fondamental : l'intégration est plus ancienne que la dérivation, et la pratiquer
une fois à la dure est ce qui vous fait sentir ce que le théorème fondamental (CAL-09) épargne réellement.

\begin{refs}
\item Contexte : Somme de Riemann --- Wikipédia (\url{https://fr.wikipedia.org/wiki/Somme_de_Riemann})
\item Contexte : La Quadrature de la parabole --- Wikipédia (\url{https://fr.wikipedia.org/wiki/La_Quadrature_de_la_parabole_%28Archim%C3%A8de%29})
\item Lecture : Tom M. Apostol, \textit{Calculus}, vol. 1, ch. 1
\end{refs}

\bigskip

\begin{problem}[{Pourquoi le théorème fondamental est un théorème --- Lycée}]
\textbf{CAL-09.} Énoncez précisément les deux parties du théorème fondamental
de l'analyse, avec leurs hypothèses.
\begin{enumerate}
\item[(i)] Démontrez la partie 1 : si $ f $ est continue sur
$ [a,b] $ et $ F(x) = \int_a^x f(t)\,dt $, alors $ F'(x) = f(x) $ --- servez-vous de la continuité
de $ f $ et encadrez le taux d'accroissement entre le minimum et le maximum de $ f $
sur un intervalle qui rétrécit.
\item[(ii)] Déduisez-en la partie 2, la règle d'évaluation
$ \int_a^b f = F(b) - F(a) $ pour toute primitive $ F $, en utilisant le théorème
de constance de CAL-07.
\item[(iii)] Voyez maintenant pourquoi les hypothèses gagnent leur salaire : posez
$ F(x) = x^2 \sin(1/x^2) $ pour $ x\ne 0 $ et $ F(0)=0 $. Montrez que $ F $ est
dérivable partout, et pourtant que $ F' $ n'est pas bornée sur $ (0,1] $, de sorte que
$ \int_0^1 F' $ n'existe pas comme intégrale de Riemann --- la partie 2 peut être en défaut pour une dérivée
qui n'est pas intégrable.
\end{enumerate}

\end{problem}

\noindent Que le problème de la tangente et le problème de l'aire soient inverses l'un de l'autre
est une \textit{découverte}, entrevue géométriquement par Barrow et exploitée systématiquement par
Newton et Leibniz ; la première démonstration rigoureuse pour des intégrandes continues est celle de Cauchy
(1823). La partie (iii) montre que le théorème est véritablement conditionnel --- et la quête de la
\og{} bonne \fg{} intégrale qui le répare mène à Lebesgue (voir CAL-21 et
Analyse réelle (\url{pure-real-analysis.html})).

\begin{refs}
\item Contexte : Théorème fondamental de l'analyse --- Wikipédia (\url{https://fr.wikipedia.org/wiki/Th%C3%A9or%C3%A8me_fondamental_de_l%27analyse})
\item Lecture : Michael Spivak, \textit{Calculus}, ch. 14
\item Cours : MIT OCW 18.01SC Single Variable Calculus (\url{https://ocw.mit.edu/courses/18-01sc-single-variable-calculus-fall-2010/})
\item Vidéo : 3Blue1Brown, Essence of Calculus (\url{https://www.3blue1brown.com/lessons/essence-of-calculus/})
\end{refs}

\bigskip

\begin{problem}[{La longueur d'arc, avec sa mise en place démontrée --- Lycée}]
\textbf{CAL-10.} Définissez la longueur d'une courbe comme la borne supérieure des
longueurs des lignes polygonales inscrites.
\begin{enumerate}
\item[(i)] Pour $ f $ continûment dérivable, démontrez
la formule
\[ L = \int_a^b \sqrt{1 + f'(x)^2}\,dx \]
à partir de cette définition --- appliquez le théorème des accroissements finis sur chaque sous-intervalle d'une subdivision et
reconnaissez la somme de Riemann obtenue.
\item[(ii)] Servez-vous-en pour calculer la longueur exacte de
$ y = \tfrac{2}{3}x^{3/2} $ de $ x=0 $ à $ x=3 $.
\item[(iii)] Défi : calculez la
longueur d'arc de la parabole $ y = x^2/2 $ de $ x=0 $ à $ x=1 $ (l'intégrande
$ \sqrt{1+x^2} $ vous donnera du fil à retordre).
\end{enumerate}

\end{problem}

\noindent Descartes pensait que le rapport entre lignes courbes et lignes droites ne pourrait
jamais être connu exactement. Il avait tort : en 1657 William Neile a rectifié la parabole
semi-cubique $ y^2 = x^3 $ --- essentiellement la courbe de (ii), la première courbe algébrique dont la longueur
fut trouvée exactement --- et Wren a rectifié la cycloïde peu après. L'étape du théorème des accroissements finis en (i) est
le c\oe{}ur du problème : la formule est habituellement affirmée avec un geste de la main vers des
\og{} triangles rectangles infinitésimaux \fg{}, et vous transformez ici ce geste en démonstration.

\begin{refs}
\item Contexte : Longueur d'un arc --- Wikipédia (\url{https://fr.wikipedia.org/wiki/Longueur_d%27un_arc})
\item Contexte : Parabole semi-cubique --- Wikipédia (\url{https://fr.wikipedia.org/wiki/Parabole_semi-cubique})
\item Lecture : James Stewart, \textit{Calculus}, ch. 8.1
\end{refs}

\bigskip

\begin{problem}[{La série harmonique diverge --- Lycée}]
\textbf{CAL-11.} Démontrez que
$ 1 + \frac12 + \frac13 + \frac14 + \cdots $ diverge, deux fois :
\begin{enumerate}
\item[(i)] par le
groupement des termes d'Oresme en blocs dépassant chacun $ \frac12 $ ;
\item[(ii)] en comparant la somme
partielle $ H_n = \sum_{k=1}^n \frac1k $ avec $ \int_1^n \frac{dx}{x} $. Démontrez ensuite que
$ \sum \frac{1}{n^2} $ \textit{converge} par comparaison avec la série télescopique
$ \sum \frac{1}{n(n-1)} $. Enfin, à partir de vos encadrements de (ii), démontrez que
$ H_n - \ln n $ est décroissante et minorée, donc convergente. Sa limite est un nombre
$ \gamma \approx 0{,}5772 $ --- retenez-le : il revient en CAL-22 avec un problème ouvert
attaché.
\end{enumerate}

\end{problem}

\noindent Nicole Oresme a démontré la divergence vers 1350 ; la démonstration fut oubliée
puis redécouverte par Mengoli et les Bernoulli trois siècles plus tard. La série diverge
si lentement (les $ 10^{43} $ premiers termes n'atteignent pas 100) qu'aucune expérience n'aurait jamais
pu en révéler la vérité --- un pur triomphe de la démonstration sur l'évidence. La valeur de
$ \sum 1/n^2 $ est le problème de Bâle, dont l'histoire vit sur
L'art de la démonstration (\url{proofs.html}).

\begin{refs}
\item Contexte : Série harmonique --- Wikipédia (\url{https://fr.wikipedia.org/wiki/S%C3%A9rie_harmonique})
\item Lecture : Michael Spivak, \textit{Calculus}, chapitres sur les séries infinies
\item Cours : MIT OCW 18.01SC Single Variable Calculus (\url{https://ocw.mit.edu/courses/18-01sc-single-variable-calculus-fall-2010/})
\end{refs}

\bigskip

\begin{problem}[{Les critères de convergence sont des théorèmes --- Lycée}]
\textbf{CAL-12.} Énoncez et démontrez :
\begin{enumerate}
\item[(i)] le critère de comparaison ;
\item[(ii)] la
règle de d'Alembert (par comparaison avec une série géométrique) ;
\item[(iii)] le critère des séries alternées,
avec la majoration de l'erreur $ |S - S_n| \le a_{n+1} $ ;
\item[(iv)] le test de condensation
de Cauchy : pour $ a_n \ge 0 $ décroissante, $ \sum a_n $ converge si et seulement si
$ \sum 2^k a_{2^k} $ converge. Décidez ensuite, avec démonstration, de la convergence de
\[ \sum_{n\ge 1} \frac{n!}{n^n}, \qquad \sum_{n\ge 1} \frac{(-1)^n}{\sqrt{n}},
\qquad \sum_{n\ge 2} \frac{1}{n \ln n}. \]
\end{enumerate}

\end{problem}

\noindent Les critères ont été assemblés pièce à pièce --- la règle du rapport par d'Alembert,
le critère des séries alternées par Leibniz dans sa correspondance, la condensation par Cauchy, dont le
\textit{Cours d'analyse} (1821) a traité la convergence systématiquement pour la première fois, après un siècle
de manipulations confiantes de séries divergentes qui avaient produit merveilles et absurdités.
Chaque critère est une petite machine ; ici vous ouvrez le boîtier et voyez l'argument de comparaison
à l'intérieur de chacun.

\begin{refs}
\item Contexte : Règle de d'Alembert --- Wikipédia (\url{https://fr.wikipedia.org/wiki/R%C3%A8gle_de_d%27Alembert})
\item Contexte : Critère de convergence des séries alternées --- Wikipédia (\url{https://fr.wikipedia.org/wiki/Crit%C3%A8re_de_convergence_des_s%C3%A9ries_altern%C3%A9es})
\item Contexte : Test de condensation de Cauchy --- Wikipédia (\url{https://fr.wikipedia.org/wiki/Test_de_condensation_de_Cauchy})
\item Lecture : Tom M. Apostol, \textit{Calculus}, vol. 1
\end{refs}

\bigskip

\begin{problem}[{Approcher e, avec un certificat --- Lycée}]
\textbf{CAL-13.} Énoncez le théorème de Taylor avec le reste sous forme de
Lagrange :
\[ f(x) = \sum_{k=0}^{n} \frac{f^{(k)}(0)}{k!}x^k + \frac{f^{(n+1)}(c)}{(n+1)!}x^{n+1}
\quad\text{pour un certain } c \text{ entre } 0 \text{ et } x. \]
(i) Démontrez d'abord, par comparaison avec une série géométrique, que $ e < 3 $. (ii) À l'aide
du polynôme de Taylor de $ e^x $ en $ x=1 $, trouvez le plus petit $ n $ pour lequel vous
pouvez \textit{démontrer} que l'erreur d'approximation est inférieure à $ 10^{-3} $, et calculez $ e $ avec
cette précision à la main --- l'enjeu n'est pas les décimales mais la majoration certifiée de l'erreur.
(iii) Défi : poussez le même développement plus loin et démontrez que $ e $ est irrationnel
(supposez $ e = p/q $ et multipliez par $ q! $).
\end{problem}

\noindent Brook Taylor a publié sa série en 1715 ; Lagrange a fourni le
reste qui transforme une série formelle en garantie numérique. Euler a calculé $ e $ avec
de nombreuses décimales exactement par la méthode de (ii). La démonstration d'irrationalité de (iii) est celle de Fourier
(1815) et c'est l'un des plus beaux arguments en deux lignes des mathématiques une fois trouvé ---
le trouver est le problème.

\begin{refs}
\item Contexte : Théorème de Taylor --- Wikipédia (\url{https://fr.wikipedia.org/wiki/Th%C3%A9or%C3%A8me_de_Taylor})
\item Contexte : e (constante mathématique) --- Wikipédia (\url{https://fr.wikipedia.org/wiki/E_%28nombre%29})
\item Lecture : Michael Spivak, \textit{Calculus}, ch. 20 (approximation par des fonctions polynomiales)
\item Vidéo : 3Blue1Brown, les séries de Taylor (\url{https://www.3blue1brown.com/lessons/taylor-series-geometric-view/})
\end{refs}

\bigskip

\begin{problem}[{L'Hôpital, justifié --- Lycée}]
\textbf{CAL-14.} \begin{enumerate}
\item[(i)] Énoncez et démontrez le théorème des accroissements
finis généralisé de Cauchy : pour $ f, g $ continues sur $ [a,b] $ et dérivables sur $ (a,b) $,
il existe $ c $ tel que $ (f(b)-f(a))\,g'(c) = (g(b)-g(a))\,f'(c) $.
\item[(ii)] Servez-vous-en pour démontrer
la règle de L'Hôpital dans le cas $ 0/0 $ lorsque $ x\to a $, avec des
hypothèses soignées.
\item[(iii)] Calculez $ \lim_{x\to 0} \frac{1-\cos x}{x^2} $ avec la règle, puis
de nouveau sans elle.
\item[(iv)] Montrez que la règle est à sens unique : exhibez $ f, g $ tels que
$ \lim_{x\to\infty} f(x)/g(x) $ existe alors que $ \lim f'(x)/g'(x) $ n'existe pas (essayez
$ f(x) = x + \sin x $, $ g(x) = x $).
\item[(v)] Expliquez pourquoi appliquer la règle à
$ \lim_{x\to 0} \frac{\sin x}{x} $ est circulaire si vous dérivez $ \sin $ de la
manière habituelle.
\end{enumerate}

\end{problem}

\noindent La règle est apparue dans le tout premier manuel d'analyse imprimé ---
l'\textit{Analyse des infiniment petits} de L'Hôpital (1696) --- mais elle avait été découverte par
Jean Bernoulli, que le marquis s'était attaché par contrat pour l'exclusivité de ses
découvertes mathématiques. Bernoulli s'en est plaint amèrement après la mort de L'Hôpital ;
la correspondance a survécu et tranche l'affaire. Le contenu mathématique du
problème est le point (i) : le théorème des accroissements finis de Cauchy est le moteur honnête sous le raccourci populaire.

\begin{refs}
\item Contexte : Règle de L'Hôpital --- Wikipédia (\url{https://fr.wikipedia.org/wiki/R%C3%A8gle_de_L%27H%C3%B4pital})
\item Lecture : Michael Spivak, \textit{Calculus}, ch. 11 et exercices
\item Cours : MIT OCW 18.01SC Single Variable Calculus (\url{https://ocw.mit.edu/courses/18-01sc-single-variable-calculus-fall-2010/})
\end{refs}

\bigskip

\begin{problem}[{L'aire sous la courbe en cloche --- un défi --- Lycée}]
\textbf{CAL-15.} Démontrez que
\[ \int_{-\infty}^{\infty} e^{-x^2}\,dx = \sqrt{\pi}. \]
Démontrez d'abord que l'intégrale impropre converge (comparez $ e^{-x^2} $ avec
$ e^{-|x|} $ pour $ |x|\ge 1 $). Évaluez-la ensuite. Avertissement loyal, qui relève du contexte et
non de l'indice : $ e^{-x^2} $ n'a pas de primitive élémentaire (c'est un théorème --- voir
CAL-23), de sorte qu'aucune chasse à la primitive ne peut aboutir ; il faut quelque chose de
structurellement différent.
\end{problem}

\noindent Cette intégrale est l'aire totale sous la courbe normale, la raison pour laquelle le
facteur $\surd$$\pi$ hante les probabilités. De Moivre a rencontré la courbe en 1733 en approchant des
comptages de pile ou face ; Laplace a évalué l'intégrale le premier (années 1770), et Gauss a attaché son nom à la
loi en 1809. Lord Kelvin aurait dit à une classe qu'un mathématicien est quelqu'un
pour qui cette identité est aussi évidente que deux fois deux font quatre. Elle n'est pas évidente. Elle est
en revanche trouvable --- et inoubliable une fois trouvée. Sa signification probabiliste est reprise dans
Probabilités (\url{applied-probability.html}) (PRB-14).

\begin{refs}
\item Contexte : Intégrale de Gauss --- Wikipédia (\url{https://fr.wikipedia.org/wiki/Int%C3%A9grale_de_Gauss})
\item Contexte : Loi normale --- Wikipédia (\url{https://fr.wikipedia.org/wiki/Loi_normale})
\item Lecture : Michael Spivak, \textit{Calculus} (intégrales impropres)
\end{refs}

\bigskip

\begin{problem}[{Une limite sans rien à simplifier --- Lycée, short}]
\textbf{CAL-25.} Démontrez, directement à partir de la définition de la dérivée, que la dérivée de $ f(x) = \sqrt{x} $ en un point $ a > 0 $ vaut $ 1/(2\sqrt{a}) $. Précisez exactement où vous utilisez que $ a \neq 0 $.
\end{problem}

\noindent Le taux d'accroissement d'une racine carrée ne se simplifie pas par factorisation ; il faut multiplier par la quantité conjuguée, une astuce qui revient constamment en analyse. La seconde phrase est le véritable contenu --- une formule de dérivée qui échoue silencieusement en une extrémité est le genre de chose qui devient plus tard une erreur sérieuse.

\begin{refs}
\item Contexte : Taux d'accroissement --- Wikipédia (en anglais) (\url{https://en.wikipedia.org/wiki/Difference_quotient})
\end{refs}

\bigskip

\begin{problem}[{Une série qui se télescope --- Lycée, short}]
\textbf{CAL-26.} Démontrez que
\[ \sum_{n=1}^{\infty} \frac{1}{n(n+1)} = 1 \]
en trouvant une formule pour la somme partielle $ \sum_{n=1}^{N} \frac{1}{n(n+1)} $ puis en passant à la limite. Justifiez le passage à la limite.
\end{problem}

\noindent Le télescopage est la seule technique de sommation qui donne une réponse exacte sans aucune machinerie, et voici son cas le plus net. Comparez-le à la série harmonique, qui lui ressemble et diverge : remarquer pourquoi l'une s'effondre et l'autre non, c'est commencer à comprendre la convergence.

\begin{refs}
\item Contexte : Somme télescopique --- Wikipédia (\url{https://fr.wikipedia.org/wiki/Somme_t%C3%A9lescopique})
\end{refs}

\bigskip

\begin{problem}[{La boîte de conserve la plus économe --- Lycée, short}]
\textbf{CAL-27.} Une boîte cylindrique doit contenir un volume fixé $ V $. Déterminez, en démontrant que votre réponse est bien un minimum et non un simple point critique, le rayon et la hauteur qui minimisent l'aire totale de sa surface.
\end{problem}

\noindent C'est le problème d'optimisation classique, et presque tous ceux qui le résolvent s'arrêtent après avoir annulé la dérivée. L'obligation de démonstration est justement le c\oe{}ur du sujet : une dérivée qui s'annule repère des candidats, et il faut quelque chose de plus --- la dérivée seconde, une étude de signe, ou un argument de compacité --- avant que le mot \og{} minimum \fg{} soit honnête.

\begin{refs}
\item Contexte : Optimisation (mathématiques) --- Wikipédia (\url{https://fr.wikipedia.org/wiki/Optimisation_%28math%C3%A9matiques%29})
\item Lecture : Michael Spivak, \textit{Calculus}, ch. 11
\end{refs}

\bigskip

\begin{problem}[{Le changement de variable, démontré --- Lycée}]
\textbf{CAL-34.} La règle de substitution dit que si $ g $ est continûment dérivable sur
$ [a,b] $ et si $ f $ est continue sur un intervalle contenant les valeurs de $ g $, alors
\[ \int_a^b f\big(g(x)\big)\,g'(x)\,dx \;=\; \int_{g(a)}^{g(b)} f(u)\,du. \]
\begin{enumerate}
\item[(a)] Démontrez-la à partir de la dérivation des fonctions composées et du théorème fondamental de l'analyse, et dites exactement où
chaque hypothèse sert. Notez que $ g $ n'est \textit{pas} supposée injective, et expliquez
pourquoi le théorème y survit.
\item[(b)] Servez-vous-en pour évaluer $ \int_0^{\pi} \sin^3 x \,dx $ et $ \int_0^1 x\sqrt{1-x^2}\,dx $,
et vérifiez la seconde réponse indépendamment en interprétant l'intégrale comme une aire.
\item[(c)] Voici maintenant le sens qui, lui, exige l'injectivité : énoncez et démontrez la formule de changement
de variable sous la forme $ \int_c^d f(u)\,du = \int_{g^{-1}(c)}^{g^{-1}(d)} f(g(x))g'(x)\,dx $,
et exhibez un $ g $ non monotone pour lequel cette lecture inverse est en défaut.
\item[(d)] Deux pièges. Premièrement, dans $ \int_{-1}^{1} \frac{dx}{1+x^2} $ substituez $ x = 1/u $ ; vous
obtiendrez l'équation $ I = -I $, donc $ I = 0 $, pour l'intégrale d'une fonction strictement
positive. Deuxièmement, appliquez la substitution en tangente de l'arc moitié $ t = \tan(x/2) $ à
$ \int_0^{2\pi} \frac{dx}{2+\cos x} $ et regardez-la renvoyer $ 0 $ elle aussi. Dans chaque cas, trouvez
l'hypothèse de (a) qui a été violée, et réparez le calcul.
\end{enumerate}

\end{problem}

\noindent La notation $ dx $ de Leibniz a été conçue pour que le changement de variable ait l'air
d'une simplification de symboles, et c'est la principale raison pour laquelle sa notation l'a emporté sur celle de Newton. Mais la
simplification est un théorème, et la partie (d) montre ce qui arrive quand on lui fait confiance comme à une comptabilité :
les deux sophismes sont de vieux favoris des salles de classe, et tous deux viennent d'une substitution qui n'est pas
dérivable --- voire pas même définie --- quelque part à l'intérieur de l'intervalle. La substitution en arc moitié
de (d) est habituellement appelée substitution de Weierstrass, une attribution que personne n'a
jamais su documenter ; elle apparaît chez Euler.

\begin{refs}
\item Contexte : Intégration par changement de variable --- Wikipédia (\url{https://fr.wikipedia.org/wiki/Int%C3%A9gration_par_changement_de_variable})
\item Contexte : Substitution en tangente de l'arc moitié --- Wikipédia (en anglais) (\url{https://en.wikipedia.org/wiki/Tangent_half-angle_substitution})
\item Lecture : Tom M. Apostol, \textit{Calculus}, vol. 1 (le théorème de substitution)
\item Cours : MIT OCW 18.01SC Single Variable Calculus (\url{https://ocw.mit.edu/courses/18-01sc-single-variable-calculus-fall-2010/})
\end{refs}

\bigskip

\begin{problem}[{Aires en polaires et arche de la cycloïde --- Lycée}]
\textbf{CAL-35.} Une courbe peut être donnée sous forme polaire $ r = r(\theta) $, ou paramétriquement par
$ \big(x(t), y(t)\big) $.
\begin{enumerate}
\item[(a)] Démontrez la formule de l'aire en polaires
\[ A \;=\; \frac{1}{2}\int_{\alpha}^{\beta} r(\theta)^2 \, d\theta \]
pour $ r \ge 0 $ continue : approchez par des secteurs circulaires plutôt que par des rectangles, écrivez
les sommes supérieures et inférieures, et montrez qu'elles convergent vers la même limite.
\item[(b)] Calculez, avec démonstration, l'aire enfermée par la cardioïde $ r = 1+\cos\theta $ et l'aire
d'un pétale de la rosace $ r = \cos 3\theta $ ; expliquez ensuite pourquoi cette rosace a trois pétales
tandis que $ r = \cos 2\theta $ en a quatre.
\item[(c)] Calculez l'aire balayée par la spirale d'Archimède $ r = \theta $ pour
$ 0 \le \theta \le 2\pi $, et comparez-la à l'aire du cercle $ r = 2\pi $ --- la
comparaison qu'Archimède a démontrée dans \textit{Des spirales} sans aucune intégrale.
\item[(d)] Démontrez la formule de longueur d'arc en paramétriques
$ L = \int_{t_0}^{t_1} \sqrt{x'(t)^2 + y'(t)^2}\,dt $ pour $ x, y $ continûment
dérivables, et servez-vous-en, avec la formule de l'aire
$ A = \int y \, dx $, pour trouver la longueur d'une arche de la cycloïde
$ x = r(t-\sin t) $, $ y = r(1-\cos t) $, $ 0 \le t \le 2\pi $, et l'aire entre cette
arche et le sol.
\item[(e)] L'arche possède un point de rebroussement : montrez que $ x'(0) = y'(0) = 0 $, de sorte que la courbe n'y a pas
de direction tangente, et expliquez pourquoi la formule de longueur d'arc de (d) reste néanmoins valable.
\end{enumerate}

\end{problem}

\noindent Jacob Bernoulli a publié le premier usage systématique des coordonnées polaires dans les
\textit{Acta Eruditorum} en 1691 ; les rosaces ont été nommées et étudiées par Guido Grandi dans les
années 1720. La cycloïde fut l'obsession du dix-septième siècle --- \og{} l'Hélène des
géomètres \fg{}, ainsi nommée pour les querelles qu'elle a provoquées. Roberval a trouvé l'aire sous une arche en
1634, Christopher Wren a rectifié l'arche exactement en 1658, Huygens a découvert en 1659 qu'elle est
la tautochrone et a bâti des horloges à pendule sur ce fait, et le défi lancé par Jean Bernoulli en 1696
l'a révélée brachistochrone, la courbe de descente la plus rapide. Tout cela a été fait
avec la machinerie que vous démontrez ici, ou avec moins.

\begin{refs}
\item Contexte : Coordonnées polaires --- Wikipédia (\url{https://fr.wikipedia.org/wiki/Coordonn%C3%A9es_polaires})
\item Contexte : Cycloïde --- Wikipédia (\url{https://fr.wikipedia.org/wiki/Cyclo%C3%AFde})
\item Contexte : Courbe tautochrone --- Wikipédia (\url{https://fr.wikipedia.org/wiki/Courbe_tautochrone})
\item Lecture : James Stewart, \textit{Calculus}, ch. 10 (équations paramétriques et coordonnées polaires)
\item Lecture : E. Hairer \& G. Wanner, \textit{Analysis by Its History} (Springer)
\end{refs}

\bigskip

\begin{problem}[{Les fonctions qui mesurent une chaîne suspendue --- Lycée, short}]
\textbf{CAL-36.} Posez
$ \cosh x = \frac{e^x + e^{-x}}{2} $ et $ \sinh x = \frac{e^x - e^{-x}}{2} $. Démontrez que
$ \cosh^2 x - \sinh^2 x = 1 $, que $ \frac{d}{dx}\cosh x = \sinh x $, et que la longueur
d'arc de la courbe $ y = \cosh x $ de $ x = 0 $ à $ x = a $ vaut exactement
$ \sinh a $.
\end{problem}

\noindent Vincenzo Riccati a introduit ces combinaisons dans les années 1760 et Lambert
les a fait entrer dans l'usage courant ; l'identité de la première tâche est la raison pour laquelle on les dit hyperboliques
--- le point $ (\cosh t, \sinh t) $ parcourt l'hyperbole $ x^2 - y^2 = 1 $ exactement comme
$ (\cos t, \sin t) $ parcourt le cercle. La dernière tâche est le petit miracle derrière CAL-37 :
l'unique courbe dont l'intégrale de longueur d'arc s'effondre est celle selon laquelle une chaîne pend réellement.

\begin{refs}
\item Contexte : Fonction hyperbolique --- Wikipédia (\url{https://fr.wikipedia.org/wiki/Fonction_hyperbolique})
\item Contexte : Chaînette --- Wikipédia (\url{https://fr.wikipedia.org/wiki/Cha%C3%AEnette})
\end{refs}

\bigskip


\section*{Probabilités}

\begin{problem}[{Les dés de Galilée --- Lycée}]
\textbf{PRB-01.} On lance trois dés équilibrés. Les sommes 9 et 10 s'écrivent
chacune comme somme de trois valeurs de dés d'exactement six façons
($ 9 = 1{+}2{+}6 = 1{+}3{+}5 = \dots $, et de même pour 10) --- ce qui laisse croire qu'elles
devraient être équiprobables. Les joueurs savaient d'expérience que le 10 sort plus souvent.
Donnez raison aux joueurs : construisez l'univers correct des $ 6^3 = 216 $ issues ordonnées
équiprobables, montrez que
$ P(\text{somme}=10) = \frac{27}{216} > \frac{25}{216} = P(\text{somme}=9) $, et expliquez
exactement où l'argument des \og{} six façons pour chacune \fg{} se trompe. Déduisez-en ensuite la loi
complète de la somme de \textit{deux} dés et démontrez que 7 en est la valeur la plus probable.
\end{problem}

\noindent Le grand-duc de Toscane soumit l'énigme du 9 contre le 10 à Galilée, qui la
trancha dans une courte note (\textit{Sopra le scoperte dei dadi}, vers 1620) ; Cardan était
parvenu à des idées voisines dans son \textit{Liber de ludo aleae} un demi-siècle plus tôt, resté
inédit jusqu'en 1663. La leçon --- compter les issues \textit{équiprobables}, et non les descriptions
d'issues --- est la première idée sérieuse de la théorie des probabilités, et les joueurs de dés
avaient relevé empiriquement un écart de 1/108 avant que les mathématiques ne sachent
l'expliquer.

\begin{refs}
\item Contexte : Histoire des probabilités --- Wikipédia (\url{https://fr.wikipedia.org/wiki/Histoire_des_probabilit%C3%A9s})
\item Lecture : William Feller, \textit{An Introduction to Probability Theory and Its Applications}, vol. 1, ch. I
\item Cours : MIT OCW 18.600 Probability and Random Variables (\url{https://ocw.mit.edu/courses/18-600-probability-and-random-variables-fall-2019/})
\end{refs}

\bigskip

\begin{problem}[{Les deux paris du Chevalier --- Lycée}]
\textbf{PRB-02.} Le chevalier de Méré gagnait de l'argent en pariant qu'il
obtiendrait au moins un six en quatre lancers d'un dé, puis en perdit en pariant qu'il obtiendrait
au moins un double-six en vingt-quatre lancers de deux dés --- alors qu'il tenait les deux paris
pour équivalents (\og{} 4 est à 6 ce que 24 est à 36 \fg{}).
\begin{enumerate}
\item[(i)] Démontrez que le premier pari est favorable :
$ 1 - (5/6)^4 > 1/2 $.
\item[(ii)] Démontrez que le second ne l'est pas :
$ 1 - (35/36)^{24} < 1/2 $, et déterminez, avec démonstration, le plus petit nombre de lancers
qui rend favorable le pari sur le double-six.
\item[(iii)] Mettez le doigt sur l'étape fausse du
raisonnement de proportionnalité.
\item[(iv)] Le problème des partis : deux joueurs misent la même somme sur un jeu de pile ou face
équilibré où le premier à remporter un nombre fixé de manches emporte tout ; ils doivent
s'arrêter alors qu'il manque encore 2 victoires au joueur A et 3 au joueur B. Déterminez, avec
démonstration, le partage équitable des mises.
\end{enumerate}

\end{problem}

\noindent Voilà les questions que de Méré apporta à Pascal en 1654, et la partie (iv) est
\textit{le} problème des partis --- partager équitablement les mises exige de raisonner sur des
parties qui n'ont jamais été jouées, et les lettres de Pascal et Fermat qui le résolvent sont les
actes fondateurs de la théorie des probabilités. Fermat dénombrait ; Pascal récursait ; tous deux
obtinrent la même réponse, et chacune des deux méthodes est devenue un pilier de la
discipline.

\begin{refs}
\item Contexte : Problème des partis --- Wikipédia (\url{https://fr.wikipedia.org/wiki/Probl%C3%A8me_des_partis})
\item Lecture : Joseph K. Blitzstein \& Jessica Hwang, \textit{Introduction to Probability}, ch. 1 (en libre accès (\url{https://probabilitybook.net}))
\item Lecture : William Feller, \textit{An Introduction to Probability Theory and Its Applications}, vol. 1
\end{refs}

\bigskip

\begin{problem}[{Le paradoxe des anniversaires --- Lycée}]
\textbf{PRB-03.} On suppose les anniversaires indépendants et uniformes sur
365 jours.
\begin{enumerate}
\item[(i)] Établissez la probabilité que $ n $ personnes aient toutes des anniversaires distincts :
$ p_n = \prod_{k=0}^{n-1} \big(1 - \frac{k}{365}\big) $, et démontrez que $ n = 23 $ est la
plus petite taille de groupe pour laquelle un anniversaire commun est plus probable
qu'improbable.
\item[(ii)] Expliquez maintenant le phénomène, et pas seulement le nombre : avec $ N $ anniversaires
possibles, utilisez l'inégalité $ 1 - x \le e^{-x} $ (démontrez-la) pour montrer qu'une
collision devient probable dès que $ n \approx 1{,}18\sqrt{N} $, puis utilisez une inégalité
correspondante dans l'autre sens pour montrer que l'ordre $ \sqrt{N} $ est bien le seuil --- la
réponse varie comme la racine carrée de la longueur de l'année, ce qui explique que 23 paraisse
trop petit.
\end{enumerate}

\end{problem}

\noindent Le problème fut posé par Richard von Mises (1939) et popularisé par le manuel
de Feller ; on dit que Harold Davenport le connaissait plus tôt. La loi en $\surd$N de la partie (ii)
en est le contenu véritable : le même calcul régit les collisions dans les tables de hachage et
l'\og{} attaque des anniversaires \fg{} sur les fonctions de hachage cryptographiques, où $ N $ est
astronomiquement grand et où 1,18$\surd$N reste le budget dont un attaquant a besoin.

\begin{refs}
\item Contexte : Paradoxe des anniversaires --- Wikipédia (\url{https://fr.wikipedia.org/wiki/Paradoxe_des_anniversaires})
\item Lecture : Joseph K. Blitzstein \& Jessica Hwang, \textit{Introduction to Probability}, ch. 1 (en libre accès (\url{https://probabilitybook.net}))
\item Cours : MIT OCW 18.600 Probability and Random Variables (\url{https://ocw.mit.edu/courses/18-600-probability-and-random-variables-fall-2019/})
\end{refs}

\bigskip

\begin{problem}[{Monty Hall, démontré --- et non plaidé --- Lycée}]
\textbf{PRB-04.} Une voiture est cachée uniformément au hasard derrière
l'une de trois portes ; les deux autres cachent des chèvres. Vous choisissez la porte 1. Le
présentateur --- qui sait où est la voiture, ouvre toujours une porte que vous n'avez pas choisie,
révèle toujours une chèvre, et choisit uniformément au hasard quand ses deux options cachent des
chèvres --- ouvre la porte 3.
\begin{enumerate}
\item[(i)] Calculez, à partir de la
définition de la probabilité conditionnelle, $ P(\text{voiture derrière la porte 2} \mid
\text{le présentateur ouvre la porte 3}) $, et démontrez que changer de porte fait gagner avec
probabilité $ 2/3 $.
\item[(ii)] Changez maintenant une règle : le présentateur ignore où est la voiture, ouvre uniformément
au hasard l'une des deux portes non choisies, et \textit{se trouve} révéler une chèvre. Démontrez
que la probabilité de gagner en changeant vaut désormais $ 1/2 $.
\item[(iii)] Dites en une phrase ce que ce couple de
calculs démontre au sujet de \og{} la \fg{} réponse à une énigme probabiliste dont le protocole n'est pas
énoncé.
\end{enumerate}

\end{problem}

\noindent Steve Selvin posa le problème dans \textit{The American Statistician} (1975) ; il
explosa en 1990 quand Marilyn vos Savant publia la réponse \og{} changez de porte \fg{} dans
\textit{Parade} et reçut des milliers de lettres --- beaucoup de docteurs en mathématiques ---
soutenant qu'elle avait tort. Paul Erd\H{o}s serait resté incrédule jusqu'à ce qu'on lui montre une
simulation. C'est la démonstration canonique que la plausibilité verbale ne vaut rien en
probabilités : seuls le modèle et le calcul tranchent.

\begin{refs}
\item Contexte : Problème de Monty Hall --- Wikipédia (\url{https://fr.wikipedia.org/wiki/Probl%C3%A8me_de_Monty_Hall})
\item Lecture : Joseph K. Blitzstein \& Jessica Hwang, \textit{Introduction to Probability}, ch. 2 (en libre accès (\url{https://probabilitybook.net}))
\item Cours : Harvard Stat 110 sur edX (Blitzstein) (\url{https://www.edx.org/bio/joseph-blitzstein})
\end{refs}

\bigskip

\begin{problem}[{L'erreur du parieur et la longueur des séries --- Lycée}]
\textbf{PRB-05.} On lance indéfiniment une pièce équilibrée, les lancers étant indépendants.
\begin{enumerate}
\item[(i)] Démontrez, à partir de la définition de l'indépendance, qu'après n'importe quelle suite
d'issues observée --- y compris dix \og{} face \fg{} d'affilée --- la probabilité que le lancer suivant donne
\og{} face \fg{} vaut encore $ 1/2 $.
\item[(ii)] Ceux qui fabriquent de fausses suites \og{} aléatoires \fg{} évitent les longues séries, et c'est
ainsi qu'on les démasque : démontrez qu'en 200 lancers, la probabilité d'observer une série d'au
moins 6 \og{} face \fg{} consécutifs dépasse $ 2/5 $. (Découpez 198 lancers en 33 blocs disjoints de
six et majorez la probabilité qu'aucun bloc ne soit entièrement \og{} face \fg{}.)
\item[(iii)] Dites
précisément ce que la loi des grands nombres ne promet \textit{pas} : montrez que l'espérance de
(nombre de \og{} face \fg{} $ - \; n/2 $) vaut $ 0 $ pour tout $ n $, mais que son écart type vaut
$ \sqrt{n}/2 $ --- la fréquence converge, l'écart absolu non ; les déviations sont diluées, jamais
\og{} corrigées \fg{}.
\end{enumerate}

\end{problem}

\noindent Le 18 août 1913, la roulette de Monte-Carlo sortit le noir 26 fois de suite, et
des joueurs se ruinèrent en misant sur le rouge \og{} parce qu'il était dû \fg{}. Les psychologues Tversky
et Kahneman ont mesuré plus tard la même erreur chez des scientifiques (\og{} Belief in the law of
small numbers \fg{}, 1971). La majoration sur la longueur des séries de la partie (ii) est un
véritable outil médico-légal, et la partie (iii) est l'énoncé honnête du théorème que l'erreur
déforme.

\begin{refs}
\item Contexte : Erreur du parieur --- Wikipédia (\url{https://fr.wikipedia.org/wiki/Erreur_du_parieur})
\item Lecture : William Feller, \textit{An Introduction to Probability Theory and Its Applications}, vol. 1, ch. XIII
\item Cours : MIT OCW 18.600 Probability and Random Variables (\url{https://ocw.mit.edu/courses/18-600-probability-and-random-variables-fall-2019/})
\end{refs}

\bigskip

\begin{problem}[{Bayes au cabinet médical --- Lycée}]
\textbf{PRB-06.} Une maladie touche 1 personne sur 1000. Un test la
détecte avec probabilité $ 0{,}99 $ lorsqu'elle est présente (sensibilité) et donne une fausse
alerte avec probabilité $ 0{,}05 $ lorsqu'elle est absente. Une personne tirée au hasard dans la
population a un test positif.
\begin{enumerate}
\item[(i)] Démontrez le théorème de Bayes
$ P(A\mid B) = \frac{P(B\mid A)\,P(A)}{P(B)} $ à partir de la définition de la probabilité
conditionnelle, y compris la formule des probabilités totales dont vous avez besoin au
dénominateur.
\item[(ii)] Calculez $ P(\text{maladie} \mid \text{test positif}) $ exactement, et expliquez avec des
mots --- en utilisant les effectifs dans une population imaginaire de 100 000 personnes ---
pourquoi la réponse est si loin en dessous de $ 0{,}99 $.
\item[(iii)] Recalculez en supposant que la personne a été testée à cause de symptômes qui portent la
probabilité a priori à $ 1/10 $, et énoncez en une phrase la morale sur les lois a priori.
\end{enumerate}

\end{problem}

\noindent L'essai de Thomas Bayes fut publié à titre posthume en 1763 ; Laplace retrouva
la règle indépendamment et l'utilisa dans toutes les sciences. Dans une célèbre enquête de 1978 à
la faculté de médecine de Harvard, la plupart des médecins répondaient à une version de (ii) par
$\approx$ 95 \% --- l'oubli de la fréquence de base, chez ceux-là mêmes qui prescrivent les tests. La
reformulation en \og{} fréquences naturelles \fg{} de Gigerenzer, celle de la partie (ii), est aujourd'hui
enseignée en médecine parce qu'elle rend la bonne réponse aussi évidente que la mauvaise le
paraissait.

\begin{refs}
\item Contexte : Théorème de Bayes --- Wikipédia (\url{https://fr.wikipedia.org/wiki/Th%C3%A9or%C3%A8me_de_Bayes})
\item Contexte : Oubli de la fréquence de base --- Wikipédia (\url{https://fr.wikipedia.org/wiki/Oubli_de_la_fr%C3%A9quence_de_base})
\item Lecture : Joseph K. Blitzstein \& Jessica Hwang, \textit{Introduction to Probability}, ch. 2 (en libre accès (\url{https://probabilitybook.net}))
\item Cours : Harvard Stat 110 sur edX (Blitzstein) (\url{https://www.edx.org/bio/joseph-blitzstein})
\end{refs}

\bigskip

\begin{problem}[{Deux enfants, deux protocoles --- Lycée, short}]
\textbf{PRB-23.} Une famille a deux enfants, chacun étant indépendamment
garçon ou fille avec probabilité $ 1/2 $. Démontrez que si tout ce qu'on vous dit est qu'au moins
l'un des deux est un garçon, la probabilité conditionnelle que les deux soient des garçons vaut
$ 1/3 $, tandis que si vous rencontrez l'un des deux enfants --- choisi uniformément au hasard --- et
que cet enfant est un garçon, la probabilité que les deux soient des garçons vaut $ 1/2 $.
\end{problem}

\noindent Martin Gardner posa l'énigme dans sa chronique de \textit{Scientific American} en
1959, donna la réponse $ 1/3 $, puis concéda que la question telle qu'elle est formulée ne
détermine pas la réponse : ce qui compte n'est pas le fait que vous apprenez, mais la procédure
qui l'a produit. Les deux calculs présentés ici sont la plus petite illustration complète de la
morale de PRB-04 --- en probabilités, une énigme sans protocole énoncé n'est pas encore un
problème.

\begin{refs}
\item Contexte : Paradoxe des deux enfants --- Wikipédia (\url{https://fr.wikipedia.org/wiki/Paradoxe_des_deux_enfants})
\item Lecture : Joseph K. Blitzstein \& Jessica Hwang, \textit{Introduction to Probability}, ch. 2 (en libre accès (\url{https://probabilitybook.net}))
\end{refs}

\bigskip

\begin{problem}[{Des dés qui se battent en rond --- Lycée, short}]
\textbf{PRB-24.} Quatre dés équilibrés à six faces portent les valeurs
$ A = (4,4,4,4,0,0) $, $ B = (3,3,3,3,3,3) $, $ C = (6,6,2,2,2,2) $,
$ D = (5,5,5,1,1,1) $ ; deux joueurs lancent chacun un dé et le plus grand nombre l'emporte.
Démontrez que
\[ P(A \text{ bat } B) = P(B \text{ bat } C) = P(C \text{ bat } D)
 = P(D \text{ bat } A) = \tfrac{2}{3}, \]
de sorte que \og{} a plus de chances de gagner contre \fg{} n'est pas une relation transitive.
\end{problem}

\noindent Ce sont les dés de Bradley Efron, rendus célèbres par Martin Gardner dans
\textit{Scientific American} en 1970. Ils comptent parce que presque tout le monde suppose en
silence que \og{} meilleur que \fg{} doit être transitif --- l'hypothèse qui sous-tend le classement des
joueurs, des produits et des options de vote --- et voici un contre-exemple en deux lignes ; l'astuce
est qu'une probabilité compare deux lois, et que de telles comparaisons ne s'enchaînent pas
nécessairement.

\begin{refs}
\item Contexte : Dés non transitifs --- Wikipédia (\url{https://fr.wikipedia.org/wiki/D%C3%A9s_non_transitifs})
\item Lecture : Martin Gardner, \textit{The Colossal Book of Mathematics} (chapitre sur les paradoxes non transitifs)
\end{refs}

\bigskip

\begin{problem}[{Le problème des chapeaux --- Lycée}]
\textbf{PRB-25.} À une soirée, $ n $ invités déposent leur chapeau au vestiaire, et un
préposé distrait les leur rend dans un ordre uniformément aléatoire --- c'est-à-dire selon une
permutation uniformément aléatoire de $ \{1, 2, \dots, n\} $. On dit qu'un invité est une
\textit{rencontre} s'il reçoit son propre chapeau.
\begin{enumerate}
\item[(i)] Par inclusion-exclusion, démontrez que la probabilité qu'il n'y ait aucune rencontre vaut
\[ p_n = \sum_{k=0}^{n} \frac{(-1)^k}{k!} , \]
de sorte que le nombre de permutations sans point fixe (un \textit{dérangement}) vaut
$ D_n = n!\, p_n $.
\item[(ii)] Déduisez-en que $ p_n \to e^{-1} $ et que $ |p_n - e^{-1}| < \frac{1}{(n+1)!} $ --- la
réponse est donc essentiellement la même pour 10 invités et pour 10 millions.
\item[(iii)] Démontrez la récurrence $ D_n = (n-1)\big( D_{n-1} + D_{n-2} \big) $ pour $ n \ge 3 $ par
un argument de dénombrement, et non par un calcul algébrique sur la formule.
\item[(iv)] Démontrez que la probabilité d'avoir exactement $ m $ rencontres vaut
$ \frac{1}{m!} \sum_{k=0}^{n-m} \frac{(-1)^k}{k!} $, et montrez que lorsque $ n \to \infty $
cette quantité tend vers $ \frac{e^{-1}}{m!} $ --- la loi de Poisson de PRB-28 de paramètre 1, en
accord avec le nombre moyen de rencontres calculé en PRB-07.
\end{enumerate}

\end{problem}

\noindent Pierre Rémond de Montmort publia ceci sous le nom de \textit{problème des
rencontres} dans son \textit{Essay d'analyse sur les jeux de hazard} (1708), où il modélisait un
jeu de cartes appelé le treize ; Nicolas Bernoulli puis Euler y revinrent, et la somme alternée
de (i) est l'une des premières apparitions imprimées du principe d'inclusion-exclusion. La chute
de (ii) fait la célébrité du problème : une probabilité qui refuse obstinément de dépendre de la
taille de la soirée, fixée à $ 1/e $ dès $ n = 6 $ environ et jusqu'aux limites de la
mesure.

\begin{refs}
\item Contexte : Problème des rencontres --- Wikipédia (\url{https://fr.wikipedia.org/wiki/Probl%C3%A8me_des_rencontres})
\item Contexte : Principe d'inclusion-exclusion --- Wikipédia (\url{https://fr.wikipedia.org/wiki/Principe_d'inclusion-exclusion})
\item Lecture : William Feller, \textit{An Introduction to Probability Theory and Its Applications}, vol. 1, ch. IV
\item Cours : MIT OCW 18.600 Probability and Random Variables (\url{https://ocw.mit.edu/courses/18-600-probability-and-random-variables-fall-2019/})
\end{refs}

\bigskip

\begin{problem}[{Ce que trois axiomes imposent déjà --- Lycée, short}]
\textbf{PRB-35.} Les axiomes de Kolmogorov disent seulement ceci : tout
événement $ A $ reçoit un nombre $ P(A) \ge 0 $, l'événement certain vérifie
$ P(\Omega) = 1 $, et $ P(A_1 \cup A_2 \cup \cdots) = \sum_{n} P(A_n) $ dès que les événements
$ A_n $ sont deux à deux disjoints. En n'utilisant rien d'autre, démontrez la règle du
complémentaire $ P(A^c) = 1 - P(A) $, la croissance ($ A \subseteq B $ entraîne
$ P(A) \le P(B) $), la formule à deux événements
$ P(A \cup B) = P(A) + P(B) - P(A \cap B) $ et l'inégalité de Boole
$ P\big( \bigcup_n A_n \big) \le \sum_n P(A_n) $ ; démontrez ensuite qu'aucune probabilité
n'attribue le \textit{même} nombre à chacun des événements $ \{1\}, \{2\}, \{3\}, \dots $ sur
l'univers $ \Omega = \{1,2,3,\dots\} $, de sorte que \og{} tirer un entier naturel au hasard \fg{} ne
décrit rien du tout.
\end{problem}

\noindent Andreï Kolmogorov publia ces axiomes dans \textit{Grundbegriffe der
Wahrscheinlichkeitsrechnung} (1933), réglant pour les probabilités la part du sixième problème
de Hilbert (1900) qui réclamait un traitement axiomatique de la discipline ; en une décennie, tout
résultat de ce domaine devint un théorème de théorie de la mesure. La dernière partie du problème
montre que le troisième axiome --- l'additivité sur une infinité \textit{dénombrable} d'événements
disjoints --- est un engagement réel et non une commodité comptable : c'est exactement lui qui tue
la loi uniforme sur les entiers, et c'est pourquoi Bruno de Finetti (PRB-18) plaidait pour une
théorie seulement finiment additive. Le même axiome impose une restriction de plus, invisible à ce
niveau : aucune probabilité dénombrablement additive et invariante par translation ne peut être
définie sur \textit{toutes} les parties du cercle, comme Giuseppe Vitali l'a montré en 1905, ce qui
explique que les axiomes parlent d'une famille d'événements choisie plutôt que de tous les
ensembles sans exception.

\begin{refs}
\item Contexte : Axiomes des probabilités --- Wikipédia (\url{https://fr.wikipedia.org/wiki/Axiomes_des_probabilit%C3%A9s})
\item Contexte : Ensemble de Vitali --- Wikipédia (\url{https://fr.wikipedia.org/wiki/Ensemble_de_Vitali})
\item Lecture : A. N. Kolmogorov, \textit{Foundations of the Theory of Probability} (1933 ; traduction anglaise, Chelsea, 1956)
\end{refs}

\bigskip

\begin{problem}[{Le problème du scrutin --- Lycée}]
\textbf{PRB-36.} Lors d'une élection, le candidat A recueille $ p $ voix et le candidat B en
recueille $ q $, avec $ p > q $. Les $ p + q $ bulletins sont dépouillés un à un dans un
ordre uniformément aléatoire, les $ \binom{p+q}{p} $ ordres étant équiprobables. On représente un
dépouillement par un chemin qui monte à chaque voix pour A et descend à chaque voix pour B.
\begin{enumerate}
\item[(a)] Démontrez le théorème du scrutin : la probabilité que A soit strictement en tête de B à chaque
étape du dépouillement vaut
\[ \frac{p-q}{p+q} . \]
(Réfléchissez le segment initial de chaque mauvais chemin par rapport à l'axe horizontal jusqu'au
premier instant d'égalité entre les deux candidats, et montrez que cela apparie bijectivement les
mauvais chemins commençant par une voix pour A avec ceux commençant par une voix pour B.)
\item[(b)] Démontrez-le une seconde fois par le lemme du cycle : écrivez les $ p+q $ bulletins autour
d'un cercle et montrez que, parmi les $ p+q $ rotations d'une disposition fixée, exactement
$ p - q $ donnent un dépouillement où A est en tête du début à la fin.
\item[(c)] Démontrez la forme faible --- le nombre d'ordres dans lesquels A n'est jamais \textit{derrière}
vaut $ \frac{p+1-q}{p+1}\binom{p+q}{q} $ --- et déduisez-en que pour $ p = q = n $ le nombre de
tels ordres est le nombre de Catalan $ \frac{1}{n+1}\binom{2n}{n} $.
\item[(d)] Remarquez ce dont la réponse ne dépend pas. Calculez la probabilité que le vainqueur soit en
tête du premier au dernier bulletin dans un électorat d'un million de personnes se partageant
600 000 contre 400 000, et comparez avec un village de dix électeurs se partageant 6 contre
4 ; montrez ensuite qu'une marge d'une seule voix ($ p = q+1 $) rend une avance de bout en bout
essentiellement impossible, avec probabilité $ 1/(p+q) $.
\end{enumerate}

\end{problem}

\noindent Joseph Bertrand posa la question dans les \textit{Comptes rendus} en 1887 et la
résolut par une récurrence, en remarquant que la simplicité de la réponse appelait une
démonstration directe ; W. A. Whitworth avait en fait publié le résultat dès 1878. Quelques
semaines après Bertrand, Désiré André fournit un argument direct, et tous les manuels depuis
l'appellent \og{} le principe de réflexion \fg{} --- mais Marc Renault, lisant l'original français en 2008, a
montré qu'André n'avait rien réfléchi : sa bijection permute les voix, et la démonstration par
réflexion n'entre dans les archives qu'en 1923. Le lemme du cycle de la partie (b) est dû à
Dvoretzky et Motzkin (1947), et le théorème lui-même est le germe de la théorie des fluctuations
de PRB-38 et des apparitions des nombres de Catalan dans toute la combinatoire.

\begin{refs}
\item Contexte : Problème du scrutin --- Wikipédia (\url{https://fr.wikipedia.org/wiki/Probl%C3%A8me_du_scrutin})
\item Contexte : Nombre de Catalan --- Wikipédia (\url{https://fr.wikipedia.org/wiki/Nombre_de_Catalan})
\item Article : M. Renault, \og{} Lost (and found) in translation: André's actual method and its application to the generalized ballot problem \fg{}, Amer. Math. Monthly 115 (2008), 358--363
\item Lecture : William Feller, \textit{An Introduction to Probability Theory and Its Applications}, vol. 1, ch. III
\item Voir aussi : PUZ-51 (\url{puzzles.html#PUZ-51}) --- la page Casse-tête, où le même dénombrement est posé comme une énigme.
\end{refs}

\bigskip


\section*{Statistique}

\begin{problem}[{Moyenne contre médiane --- Lycée}]
\textbf{STA-01.} Pour des réels $ x_1, \dots, x_n $, on pose
\[ f(c) = \sum_{i=1}^{n} (x_i - c)^2 , \qquad g(c) = \sum_{i=1}^{n} |x_i - c| . \]
\begin{enumerate}
\item[(a)] Démontrez que $ f $ atteint son minimum en un unique point, la moyenne des nombres.
\item[(b)] Démontrez que $ g $ atteint son minimum en toute médiane (tout point ayant au moins la
moitié des données de chaque côté), et déterminez exactement quand le minimiseur est unique.
\item[(c)] Construisez une liste de dix \og{} revenus des ménages \fg{} dont la moyenne dépasse dix fois la
médiane, et argumentez laquelle des deux est le \og{} revenu typique \fg{} honnête pour votre liste ---
avec une raison mathématique, pas un slogan.
\end{enumerate}

\end{problem}

\noindent La tension entre les deux moyennes est ancienne : la moyenne appartient à la
tradition des moindres carrés de Legendre et Gauss, tandis que la médiane fut défendue par Laplace
(qui démontra qu'elle minimise l'erreur absolue) puis par Edgeworth. Le choix tranche encore
aujourd'hui des débats publics --- revenus, prix de l'immobilier, délais d'attente à l'hôpital ---
parce que la moyenne court après les valeurs extrêmes et la médiane non. La partie (b) est une
véritable démonstration sur une fonction affine par morceaux plutôt que dérivable, ce qui est
exactement pourquoi elle mérite d'être rédigée.

\begin{refs}
\item Contexte : Médiane (statistiques) --- Wikipédia (\url{https://fr.wikipedia.org/wiki/M%C3%A9diane_(statistiques)})
\item Lecture : D. Freedman, R. Pisani \& R. Purves, \textit{Statistics}, ch. 4
\item Cours : MIT OCW 18.650 Statistics for Applications (\url{https://ocw.mit.edu/courses/18-650-statistics-for-applications-fall-2016/})
\end{refs}

\bigskip

\begin{problem}[{Construisez votre propre paradoxe de Simpson --- Lycée}]
\textbf{STA-02.} Deux hôpitaux, A et B, traitent chacun exactement 1000 patients, chaque
patient étant classé comme cas bénin ou cas grave.
\begin{enumerate}
\item[(a)] Construisez des données entières telles que : l'hôpital A ait un taux de survie strictement
supérieur à celui de B chez les cas bénins, A ait un taux de survie strictement supérieur à celui
de B chez les cas graves, et pourtant B ait un taux de survie global strictement supérieur.
\item[(b)] Démontrez le résultat d'impossibilité suivant : si les deux hôpitaux traitent le \textit{même
nombre} de cas bénins (et donc le même nombre de cas graves), aucune inversion de ce genre ne
peut se produire.
\end{enumerate}

\end{problem}

\noindent L'inversion porte le nom d'E. H. Simpson (1951), bien que Karl Pearson et Udny
Yule l'aient remarquée un demi-siècle plus tôt. Le cas réel le plus célèbre est celui des données
d'admission en troisième cycle à Berkeley en 1973, analysées par Bickel, Hammel et O'Connell :
l'université semblait globalement défavoriser les femmes, alors que département par département ce
n'était pas le cas --- les femmes avaient postulé aux départements les plus sélectifs. Le paradoxe
est un théorème sur les moyennes pondérées, et votre démonstration d'impossibilité montrera
exactement d'où vient la liberté d'inverser.

\begin{refs}
\item Contexte : Paradoxe de Simpson --- Wikipédia (\url{https://fr.wikipedia.org/wiki/Paradoxe_de_Simpson})
\item Lecture : J. Pearl \& D. Mackenzie, \textit{The Book of Why}, ch. 6
\item Article : P. J. Bickel, E. A. Hammel \& J. W. O'Connell, \og{} Sex bias in graduate admissions: Data from Berkeley \fg{}, \textit{Science} 187 (1975), 398--404
\end{refs}

\bigskip

\begin{problem}[{Corrélation n'est pas causalité --- sous forme de théorème --- Lycée}]
\textbf{STA-03.} La corrélation empirique $ r $ de données appariées mesure l'association
linéaire et vérifie toujours $ -1 \le r \le 1 $ (vous pouvez l'admettre).
\begin{enumerate}
\item[(a)] Soient $ Z, U, V $ des quantités aléatoires indépendantes, chacune d'espérance 0 et de
variance 1, et posons $ X = Z + U $ et $ Y = Z + V $. Montrez que la corrélation
(théorique) de $ X $ et $ Y $ vaut $ 1/2 $, alors que ni $ X $ ni $ Y $ n'influence
l'autre --- ils partagent seulement la cause commune $ Z $.
\item[(b)] Construisez une variable aléatoire $ X $ prenant un nombre fini de valeurs, telle que
$ X $ et $ Y = X^2 $ aient une corrélation exactement nulle bien que $ Y $ soit entièrement
déterminée par $ X $. Démontrez les deux affirmations.
\item[(c)] Concluez, en un paragraphe soigneux, sur ce qu'une corrélation non nulle établit et n'établit
pas, et sur ce qu'une corrélation nulle établit et n'établit pas.
\end{enumerate}

\end{problem}

\noindent \og{} Corrélation n'implique pas causalité \fg{} est d'ordinaire laissé à l'état de
proverbe ; il devient ici deux petits théorèmes. La partie (a) est le \textit{facteur de confusion} :
les ventes de glaces et les noyades sont corrélées parce que l'été entraîne les deux. La partie (b)
montre le piège réciproque : corrélation nulle n'est pas indépendance. Le coefficient de
corrélation lui-même est dû à Francis Galton et Karl Pearson dans les années 1880--1890 ; la
mathématique moderne des conditions sous lesquelles une corrélation \textit{peut} soutenir une
affirmation causale est l'objet de l'inférence causale, l'un des domaines les plus vivants de la
statistique d'aujourd'hui.

\begin{refs}
\item Contexte : Corrélation n'implique pas causalité --- Wikipédia (\url{https://fr.wikipedia.org/wiki/Corr%C3%A9lation_n'implique_pas_causalit%C3%A9})
\item Contexte : Coefficient de corrélation de Pearson --- Wikipédia (en anglais) (\url{https://en.wikipedia.org/wiki/Pearson_correlation_coefficient})
\item Lecture : D. Freedman, R. Pisani \& R. Purves, \textit{Statistics}, ch. 9--12
\end{refs}

\bigskip

\begin{problem}[{Pourquoi les fils des pères grands sont plus petits --- Lycée, short}]
\textbf{STA-19.} Soient $ X $ et $ Y $ les tailles standardisées du père
et du fils (chacune d'espérance 0 et d'écart type 1), de corrélation $ r $, et supposons la
moyenne conditionnelle linéaire : $ \mathbb{E}[Y \mid X = x] = a + bx $. Démontrez que
$ a = 0 $ et $ b = r $, de sorte que $ |\mathbb{E}[Y \mid X = x]| <|x| $ dès que
$ |r| <1 $ et $ x \ne 0 $, puis expliquez --- par un argument de symétrie, pas par un slogan ---
pourquoi cela ne signifie \textit{pas} que la dispersion des tailles se réduit de génération en
génération.
\end{problem}

\noindent Francis Galton découvrit l'effet dans les années 1880 et l'appela \og{} régression
vers la médiocrité \fg{} ; toute la discipline moderne de l'analyse de régression tire son nom de cette
seule observation biologique. Le point de la seconde moitié du problème est que le même argument
mené à l'envers prédit que les fils grands avaient des pères plus petits, ce qu'aucune théorie de
l'hérédité ne peut expliquer --- l'effet est une propriété de la corrélation imparfaite, pas de la
biologie. C'est aussi le piège classique de l'évaluation des interventions : les écoles les moins
performantes, les patients les plus malades et les pilotes les plus malchanceux s'améliorent tout
seuls, et quelque chose d'autre en récolte toujours le crédit.

\begin{refs}
\item Contexte : Régression vers la moyenne --- Wikipédia (\url{https://fr.wikipedia.org/wiki/R%C3%A9gression_vers_la_moyenne})
\item Article : F. Galton, \og{} Regression towards mediocrity in hereditary stature \fg{}, \textit{Journal of the Anthropological Institute} 15 (1886), 246--263
\item Lecture : D. Freedman, R. Pisani \& R. Purves, \textit{Statistics}, ch. 10
\end{refs}

\bigskip

\begin{problem}[{Une garantie qui n'a besoin d'aucune courbe en cloche --- Lycée, short}]
\textbf{STA-20.} Soit $ x_1, \dots, x_n $ une liste quelconque de nombres,
de moyenne $ \bar{x} $ et d'écart type
$ s = \sqrt{\frac{1}{n}\sum_i (x_i - \bar{x})^2} $, et soit $ k >1 $. Démontrez que la
proportion de la liste située à une distance au moins $ ks $ de $ \bar{x} $ est au plus
$ 1/k^2 $, et montrez que la borne ne peut pas être améliorée en exhibant une liste de huit
nombres pour laquelle cette proportion vaut exactement $ 1/4 $ lorsque $ k = 2 $.
\end{problem}

\noindent C'est l'inégalité de Bienaymé-Tchebychev dépouillée de toute probabilité : la
démonstration tient en une ligne d'arithmétique --- chaque terme de la somme des carrés est positif,
alors jetez les petits. La fameuse règle \og{} 68-95-99,7 \fg{} est bien plus fine, mais elle ne vaut que
pour la courbe normale ; la borne de Tchebychev vaut pour des données de revenus, des magnitudes de
séismes et n'importe quelle liste que vous rencontrerez, ce qui explique que Tchebychev s'en soit
servi en 1867 pour démontrer la loi des grands nombres sans rien supposer sur la forme de la
loi.

\begin{refs}
\item Contexte : Inégalité de Bienaymé-Tchebychev --- Wikipédia (\url{https://fr.wikipedia.org/wiki/In%C3%A9galit%C3%A9_de_Bienaym%C3%A9-Tchebychev})
\item Contexte : Règle 68-95-99,7 --- Wikipédia (\url{https://fr.wikipedia.org/wiki/R%C3%A8gle_68-95-99,7})
\item Lecture : D. Freedman, R. Pisani \& R. Purves, \textit{Statistics}, ch. 4--5
\end{refs}

\bigskip

\begin{problem}[{Compter des poissons que l'on ne peut pas compter --- Lycée}]
\textbf{STA-21.} Un lac contient un nombre inconnu $ N $ de poissons. Vous en capturez
$ K $, vous les marquez et vous les relâchez ; plus tard vous capturez un second échantillon de
$ n $ poissons, dont $ k $ se révèlent marqués. On suppose que tout sous-ensemble de $ n $
poissons a la même probabilité d'être le second échantillon.
\begin{enumerate}
\item[(a)] Écrivez la probabilité d'observer exactement $ k $ poissons marqués (une probabilité
hypergéométrique), en fonction de l'inconnue $ N $.
\item[(b)] En examinant le rapport $ L(N)/L(N-1) $ de vraisemblances consécutives, démontrez que cette
probabilité est maximale en $ \hat{N} = \lfloor Kn/k \rfloor $ lorsque $ k \ge 1 $ ---
l'estimateur du maximum de vraisemblance est exactement la règle de proportionnalité naïve
$ k/n \approx K/N $.
\item[(c)] Montrez que $ \hat{N} $ est inutilisable tel quel : $ P(k = 0) >0 $ pour tout
$ N \ge K + n $, et sur cet événement l'estimation est infinie. Démontrez que l'estimateur de
Chapman
\[ \tilde{N} = \frac{(K+1)(n+1)}{k+1} - 1 \]
est exactement sans biais dès que $ K + n \ge N $. (Une identité de type Vandermonde pour les
coefficients binomiaux fait le travail.)
\item[(d)] Calculez les deux estimations pour $ K = 100 $, $ n = 100 $, $ k = 20 $. Énumérez ensuite
les hypothèses de modélisation --- population fermée, aucune marque perdue, tous les poissons
également capturables --- et pour chacune dites dans quel sens son échec ferait dévier
l'estimation.
\end{enumerate}

\end{problem}

\noindent Laplace utilisa la même arithmétique en 1802 pour estimer la population de la
France à partir des registres de naissances, et la méthode fut réinventée pour les animaux par
C. G. J. Petersen pour les poissons (1896) et F. C. Lincoln pour les oiseaux d'eau bagués (1930) ;
la correction de biais de (c) est celle de D. G. Chapman, en 1951. La technique sert aujourd'hui
bien au-delà de l'écologie : estimer le sous-dénombrement d'un recensement, le nombre de défauts
logiciels non découverts et --- en croisant plusieurs listes incomplètes --- le nombre de victimes d'un
conflit armé. C'est l'exemple le plus net, en statistique élémentaire, de la façon d'apprendre la
taille d'une chose que l'on ne pourra jamais voir en entier.

\begin{refs}
\item Contexte : Capture-marquage-recapture --- Wikipédia (\url{https://fr.wikipedia.org/wiki/Capture-marquage-recapture})
\item Contexte : Loi hypergéométrique --- Wikipédia (\url{https://fr.wikipedia.org/wiki/Loi_hyperg%C3%A9om%C3%A9trique})
\item Lecture : G. A. F. Seber, \textit{The Estimation of Animal Abundance and Related Parameters}, 2e éd.
\end{refs}

\bigskip

\begin{problem}[{La dame qui déguste du thé --- Lycée, short}]
\textbf{STA-31.} On prépare huit tasses de thé, quatre où le lait a été versé
en premier et quatre où c'est le thé ; on informe la dégustatrice qu'il y en a quatre de chaque
sorte et elle doit désigner les quatre tasses au lait versé en premier, de sorte que sous
l'hypothèse nulle où elle devine au hasard, les $ \binom{8}{4} = 70 $ sélections possibles sont
équiprobables. Démontrez que la probabilité de désigner les quatre correctement par hasard vaut
$ 1/70 $, calculez la probabilité d'en trouver au moins trois sur quatre, et déduisez-en qu'au
seuil de 5 \% rien de moins qu'un sans-faute ne peut compter comme preuve --- puis dites ce qui change
si l'expérience est menée avec six tasses, trois de chaque sorte.
\end{problem}

\noindent Ronald Fisher ouvre \textit{The Design of Experiments} (1935) par cette
expérience, et elle a réellement eu lieu : à Rothamsted, dans les années 1920, l'algologue Muriel
Bristol affirmait pouvoir sentir au goût si le lait avait été versé dans la tasse en premier, et
Fisher conçut le test sur-le-champ. Ce chapitre est le premier endroit où la randomisation,
l'hypothèse nulle et la valeur p exacte apparaissent ensemble, et le dénombrement que vous venez de
faire est le \textit{test exact de Fisher} pour un tableau $ 2 \times 2 $ --- pas d'approximation
normale, pas de théorie asymptotique, rien que de la combinatoire. Remarquez ce que révèle la
dernière partie : le plus petit effet détectable est fixé par le plan d'expérience avant qu'une
seule tasse ne soit versée, et c'est pourquoi Fisher insistait pour que l'on consulte le
statisticien avant l'expérience plutôt qu'après.

\begin{refs}
\item Contexte : La dame dégustant du thé --- Wikipédia (\url{https://fr.wikipedia.org/wiki/The_lady_tasting_tea})
\item Contexte : Test exact de Fisher --- Wikipédia (\url{https://fr.wikipedia.org/wiki/Test_exact_de_Fisher})
\item Lecture : R. A. Fisher, \textit{The Design of Experiments} (1935), ch. 2
\item Lecture : D. Salsburg, \textit{The Lady Tasting Tea: How Statistics Revolutionized Science in the Twentieth Century} (2001)
\end{refs}

\bigskip

\begin{problem}[{Deux millions de mauvaises réponses --- Lycée}]
\textbf{STA-32.} Une population comporte une fraction $ p $ d'électeurs favorables au candidat
R et $ 1 - p $ favorables au candidat L. Un magazine envoie des bulletins par la poste à une
immense liste de noms. Indépendamment de tout le reste, un électeur favorable à R renvoie le
bulletin avec probabilité $ r_1 $, et un électeur favorable à L le renvoie avec probabilité
$ r_0 $ ; le magazine publie la fraction de partisans de R parmi les bulletins qui lui
reviennent.
\begin{enumerate}
\item[(a)] Démontrez que lorsque le nombre de bulletins envoyés croît, la fraction publiée se stabilise
non pas en $ p $ mais en
\[ p^{*} = \frac{p\, r_1}{p\, r_1 + (1-p)\, r_0} , \]
et que $ p^{*} = p $ si et seulement si $ r_1 = r_0 $. Remarquez que le nombre de bulletins
n'apparaît nulle part dans $ p^{*} $ : l'erreur ne diminue pas quand le sondage grossit.
\item[(b)] Injectez les chiffres historiques. En 1936, le \textit{Literary Digest} envoya environ dix
millions de bulletins, en reçut environ 2,38 millions en retour, et prédit que Landon obtiendrait
57 \% des voix ; Roosevelt en obtint 60,8 \%. À l'aide de (a), déterminez le rapport
$ r_1/r_0 $ des taux de réponse qui expliquerait à lui seul l'écart, et commentez la modestie de
ce différentiel.
\item[(c)] Montrez que l'erreur d'échantillonnage n'était pas le problème. Pour un échantillon réellement
aléatoire de taille $ n $ tiré de la population, démontrez que l'écart type de la fraction
observée vaut $ \sqrt{p(1-p)/n} \le 1/(2\sqrt{n}) $, et évaluez-le pour
$ n = 2\,380\,000 $. Comparez le résultat à l'erreur de 20 points du magazine, puis trouvez la
taille d'échantillon à laquelle la seule erreur d'échantillonnage aléatoire atteint un point de
pourcentage.
\item[(d)] Le \textit{Digest} a commis deux erreurs distinctes : il a bâti son fichier d'adresses à partir
d'annuaires téléphoniques, de listes de clubs et d'immatriculations d'automobiles, et il a accepté
la réponse de quiconque voulait bien répondre. Expliquez, à l'aide de (a), pourquoi chacune est
fatale à elle seule, et démontrez qu'envoyer vingt millions de bulletins au lieu de dix n'aurait
remédié ni à l'une ni à l'autre.
\end{enumerate}

\end{problem}

\noindent Le \textit{Literary Digest} avait correctement annoncé toutes les élections
présidentielles depuis 1920 et était si confiant dans le volume brut qu'il publia sa prévision de
1936 comme une quasi-certitude ; George Gallup, travaillant sur un échantillon d'environ cinquante
mille personnes choisies pour représenter l'électorat, annonça le résultat correct, et le magazine
disparut en moins de deux ans. La reconstitution de Peverill Squire en 1988 a identifié la
non-réponse comme le plus grave des deux échecs : les partisans de Roosevelt renvoyaient tout
simplement moins de bulletins. La leçon n'est pas seulement historique. Xiao-Li Meng a montré en
2018 que l'erreur d'une moyenne d'échantillon se factorise en un terme de qualité des données, un
terme de quantité de données et la variance de la population, de sorte qu'une corrélation minuscule
entre le fait de répondre et la préférence rend un échantillon énorme aussi mauvais qu'un petit
échantillon aléatoire --- son \og{} paradoxe des mégadonnées \fg{}, appliqué à l'élection américaine de 2016,
est STA-32(a) réécrit pour le XXIe siècle.

\begin{refs}
\item Contexte : The Literary Digest --- Wikipédia (en anglais) (\url{https://en.wikipedia.org/wiki/The_Literary_Digest})
\item Contexte : Échantillon biaisé --- Wikipédia (\url{https://fr.wikipedia.org/wiki/%C3%89chantillon_biais%C3%A9})
\item Article : P. Squire, \og{} Why the 1936 Literary Digest poll failed \fg{}, \textit{Public Opinion Quarterly} 52 (1988), 125--133
\item Article : X.-L. Meng, \og{} Statistical paradises and paradoxes in big data (I) \fg{}, \textit{Annals of Applied Statistics} 12 (2018), 685--726
\end{refs}

\bigskip


\section*{Équations différentielles et systèmes dynamiques}

\begin{problem}[{La croissance exponentielle à partir des principes premiers --- Lycée}]
\textbf{ODE-01.} (a) Démontrer que les \textit{seules} fonctions dérivables
vérifiant $ y' = k y $ sur la droite réelle sont $ y(x) = C e^{kx} $ avec $ C $ constante.
(On pourra utiliser les règles élémentaires de dérivation et le fait qu'une fonction de dérivée
nulle sur un intervalle y est constante.)
(b) En déduire la loi du temps de doublement : une solution positive avec $ k >0 $ double en
un même temps $ (\ln 2)/k $, quel que soit l'instant où l'on déclenche le chronomètre.
(c) Datation au radiocarbone : le carbone 14 se désintègre avec une demi-vie de 5730 ans. Un
parchemin conserve 70 \% de son carbone 14 initial. Déterminer son âge, en le justifiant
entièrement à partir de (a).
\end{problem}

\noindent C'est la première équation différentielle jamais comprise --- elle régit les
intérêts composés (Jacques Bernoulli, années 1690), le refroidissement, la désintégration
radioactive (Rutherford et Soddy, 1902) et la croissance des populations (Malthus, 1798). La
démonstration d'unicité en (a) est la partie importante : c'est le cas le plus simple du fait
qu'une équation différentielle assortie d'une condition initiale détermine complètement l'avenir,
le thème de toute cette page. Willard Libby a reçu le prix Nobel de chimie 1960 pour avoir fait
de (c) un instrument d'archéologie.

\begin{refs}
\item Contexte : Croissance exponentielle --- Wikipédia (\url{https://fr.wikipedia.org/wiki/Croissance_exponentielle})
\item Lecture : W. E. Boyce \& R. C. DiPrima, \textit{Elementary Differential Equations}, chap. 1--2
\item Cours : MIT OCW 18.03SC Differential Equations (\url{https://ocw.mit.edu/courses/18-03sc-differential-equations-fall-2011/})
\end{refs}

\bigskip

\begin{problem}[{La loi de refroidissement de Newton, à partir de données --- Lycée}]
\textbf{ODE-02.} La loi de refroidissement de Newton affirme que la température $ T $ d'un
corps placé dans un milieu à température constante $ A $ obéit à $ T' = -k (T - A) $ pour une
certaine constante $ k >0 $.
\begin{enumerate}
\item[(a)] Résoudre l'équation (utiliser ODE-01 après la substitution $ u = T - A $).
\item[(b)] Une tasse de café à 90$^\circ$C se trouve dans une pièce à 20$^\circ$C. Elle mesure 70$^\circ$C au
bout de 5 minutes et environ 55,7$^\circ$C au bout de 10 minutes. Montrer que ces données sont
compatibles avec le modèle, déterminer $ k $, puis prédire la température au bout de 20 minutes
ainsi que l'instant où le café atteint 40$^\circ$C.
\item[(c)] Démontrer que, selon le modèle, le café n'atteint jamais réellement la température de la
pièce en temps fini --- et expliquer pourquoi le modèle est néanmoins utile.
\end{enumerate}

\end{problem}

\noindent Newton a publié la loi de refroidissement anonymement en 1701, dans un article
sur une échelle de chaleur allant jusqu'à la fusion du fer ; confronter un modèle d'équation
différentielle à des données, comme en (b), est exactement ce qu'il faisait. Cette loi est le
premier exemple canonique d'ajustement d'une équation différentielle à des mesures --- les
médecins légistes s'en servent encore pour estimer l'heure du décès --- et la partie (c) apprend
la bonne façon de lire une approche exponentielle de l'équilibre.

\begin{refs}
\item Contexte : Loi de refroidissement de Newton --- Wikipédia (\url{https://fr.wikipedia.org/wiki/Loi_de_refroidissement_de_Newton})
\item Lecture : W. E. Boyce \& R. C. DiPrima, \textit{Elementary Differential Equations}, chap. 1
\item Cours : MIT OCW 18.03SC Differential Equations (\url{https://ocw.mit.edu/courses/18-03sc-differential-equations-fall-2011/})
\end{refs}

\bigskip

\begin{problem}[{La courbe logistique --- Lycée}]
\textbf{ODE-03.} Une population limitée par les ressources est modélisée par l'équation
logistique $ P' = r P \left( 1 - \frac{P}{K} \right) $ avec des constantes
$ r, K >0 $.
\begin{enumerate}
\item[(a)] Sans résoudre, montrer que toute solution avec $ 0 <P(0) <K $ est croissante, et que
toute solution avec $ P(0) >K $ est décroissante.
\item[(b)] Résoudre l'équation par séparation des variables et décomposition en éléments simples, puis
vérifier que $ P(t) \to K $ quand $ t \to \infty $ pour toute valeur initiale
positive.
\item[(c)] Démontrer qu'une solution partant en dessous de $ K/2 $ a sa croissance la plus forte
exactement lorsque $ P = K/2 $, et que la courbe solution est symétrique par rapport à ce
point. Esquisser la courbe en S.
\end{enumerate}

\end{problem}

\noindent Le mathématicien belge Pierre-François Verhulst a proposé l'équation logistique
en 1838 pour corriger Malthus : la croissance ne peut pas rester exponentielle. La même courbe en
S modélise aujourd'hui les épidémies, l'adoption des technologies et l'autocatalyse chimique ;
sa cousine à temps discret reviendra en ODE-16, chargée de chaos. Verhulst s'en est servi pour
prédire un plafond à la population belge --- un rappel que les paramètres du modèle, contrairement
à ses mathématiques, sont des conjectures.

\begin{refs}
\item Contexte : Fonction logistique (Verhulst) --- Wikipédia (\url{https://fr.wikipedia.org/wiki/Fonction_logistique_(Verhulst)})
\item Lecture : S. H. Strogatz, \textit{Nonlinear Dynamics and Chaos}, chap. 2
\item Cours : MIT OCW 18.03SC Differential Equations (\url{https://ocw.mit.edu/courses/18-03sc-differential-equations-fall-2011/})
\end{refs}

\bigskip

\begin{problem}[{La vitesse limite --- Lycée, short}]
\textbf{ODE-22.} Un corps de masse $ m $ tombe sans vitesse initiale sous
l'effet de la pesanteur, contre une résistance de l'air proportionnelle à la vitesse, de sorte
que sa vitesse vérifie
\[ m v' = mg - k v, \qquad v(0) = 0, \]
avec des constantes $ m, g, k >0 $. Résoudre en $ v(t) $ et démontrer que $ v $ croît
strictement vers la vitesse limite $ mg/k $ sans jamais l'atteindre.
\end{problem}

\noindent Le frottement linéaire est le modèle honnête d'une chute lente --- une
gouttelette de brume, une bille d'acier dans l'huile --- et Stokes en a établi la loi en 1851. La
moitié qualitative de ce problème est la plus précieuse : la limite $ mg/k $ se lit sur
l'équation avant toute intégration, en demandant où $ v' = 0 $, ce qui est le premier geste de
toute la théorie qualitative vers laquelle cette page tend.

\begin{refs}
\item Contexte : Vitesse terminale --- Wikipédia (\url{https://fr.wikipedia.org/wiki/Vitesse_terminale})
\item Lecture : W. E. Boyce \& R. C. DiPrima, \textit{Elementary Differential Equations}, chap. 1
\end{refs}

\bigskip

\begin{problem}[{Le réservoir qui se vide en temps fini --- Lycée, short}]
\textbf{ODE-23.} De l'eau s'écoule par un petit trou percé au fond d'un
réservoir cylindrique ; la loi de Torricelli impose à la hauteur $ h(t) $ d'obéir à
\[ h' = -c\sqrt{h}, \qquad h(0) = h_0 >0, \]
avec une constante $ c >0 $ déterminée par le trou. Déterminer, avec démonstration, l'instant
exact où le réservoir est vide, et démontrer que la hauteur est véritablement nulle à partir de
cet instant, et non pas seulement tendant vers zéro.
\end{problem}

\noindent Evangelista Torricelli --- élève de Galilée et inventeur du baromètre --- a publié
la loi en 1644 : l'eau quitte le trou à la vitesse qu'elle atteindrait en tombant librement depuis
la surface. Le temps de vidange fini est tout l'intérêt. Les solutions de $ y' = y $ mettent
l'éternité à atteindre quoi que ce soit, mais $ \sqrt{h} $ n'est pas lipschitzienne en $ 0 $,
et c'est ce défaut qui permet à une solution d'arriver à l'origine en temps fini --- le même défaut
qui détruit l'unicité en ODE-08.

\begin{refs}
\item Contexte : Formule de Torricelli --- Wikipédia (\url{https://fr.wikipedia.org/wiki/Formule_de_Torricelli})
\item Cours : MIT OCW 18.03SC Differential Equations (\url{https://ocw.mit.edu/courses/18-03sc-differential-equations-fall-2011/})
\end{refs}

\bigskip

\begin{problem}[{Champs de tangentes et méthode d'Euler --- Lycée}]
\textbf{ODE-24.} On considère le problème de Cauchy
\[ y' = y - t, \qquad y(0) = \tfrac{1}{2}. \]
\begin{enumerate}
\item[(a)] Tracer le champ de tangentes sur le rectangle $ 0 \le t \le 2 $, $ -1 \le y \le 3 $.
Vérifier que $ y = t + 1 $ est solution, et utiliser le signe de $ y - t $ au-dessus et
en dessous de cette droite pour prédire l'allure de toutes les autres solutions.
\item[(b)] Vérifier que la solution du problème de Cauchy est
$ y(t) = t + 1 - \tfrac{1}{2}e^{t} $, et confirmer qu'elle correspond à votre prédiction.
\item[(c)] Appliquer la méthode d'Euler $ y_{n+1} = y_n + h\,(y_n - t_n) $ avec les pas
$ h = 0{,}5,\ 0{,}25,\ 0{,}125 $ pour estimer $ y(1) $. Dresser le tableau des trois erreurs
et observer que chaque division par deux de $ h $ divise à peu près l'erreur par deux.
\item[(d)] Expliquer cette observation. Montrer par un développement de Taylor qu'un pas d'Euler à
partir d'une valeur exacte commet une erreur au plus $ \tfrac{1}{2}h^2 \max|y''| $, et
argumenter qu'en cumulant $ 1/h $ pas de ce type on obtient une erreur totale d'ordre $ h $ ---
la méthode d'Euler est donc du \textit{premier ordre}.
\end{enumerate}

\end{problem}

\noindent Leonhard Euler a publié cette méthode dans ses \textit{Institutionum calculi
integralis} (1768) comme moyen de calculer avec des équations que personne ne savait résoudre.
Elle est l'ancêtre de tous les solveurs numériques utilisés aujourd'hui, et aussi un instrument
théorique : Cauchy dans les années 1820, et Peano en 1886, ont démontré que les solutions
\textit{existent} en montrant que les approximations polygonales d'Euler convergent quand
$ h \to 0 $ (comparer à la voie différente d'ODE-07). Le champ de tangentes de (a) est l'autre
moitié de la leçon --- une équation qu'on ne sait pas résoudre peut tout de même se lire.

\begin{refs}
\item Contexte : Méthode d'Euler --- Wikipédia (\url{https://fr.wikipedia.org/wiki/M%C3%A9thode_d%27Euler})
\item Contexte : Champ de tangentes --- Wikipédia (en anglais) (\url{https://en.wikipedia.org/wiki/Slope_field})
\item Lecture : W. E. Boyce \& R. C. DiPrima, \textit{Elementary Differential Equations}, chap. 1--2
\item Cours : MIT OCW 18.03SC Differential Equations (\url{https://ocw.mit.edu/courses/18-03sc-differential-equations-fall-2011/})
\end{refs}

\bigskip

\begin{problem}[{La chaîne suspendue --- Lycée}]
\textbf{ODE-34.} Une chaîne souple homogène de poids $ w $ par unité de longueur pend au
repos entre deux supports. Soit $ y(x) $ sa forme ; on place le point le plus bas en
$ x = 0 $, et l'on note $ H $ la traction horizontale en ce point.
\begin{enumerate}
\item[(a)] Écrire l'équilibre des forces sur le morceau de chaîne compris entre le point le plus bas
et un point courant. Montrer que la composante horizontale de la tension vaut $ H $ en tout
point, que la composante verticale vaut $ w\,s $ où $ s $ est l'abscisse curviligne mesurée
depuis le point le plus bas, et donc que $ y'(x) = w\,s(x)/H $. Dériver, en utilisant
$ \frac{ds}{dx} = \sqrt{1 + (y')^2} $, pour obtenir
\[ y'' = \frac{w}{H}\,\sqrt{1 + (y')^2}\,. \]
\item[(b)] Résoudre l'équation : poser $ p = y' $, séparer les variables, et montrer que la chaîne
telle que $ y'(0) = 0 $ et $ y(0) = H/w $ est
\[ y(x) = \frac{H}{w}\,\cosh\!\left( \frac{w x}{H} \right). \]
Confirmer par dérivation directe que cette fonction satisfait l'équation.
\item[(c)] Donner tort à Galilée : montrer qu'aucune parabole $ y = ax^2 + bx + c $ avec $ a \ne 0 $
ne peut satisfaire l'équation, en comparant les deux membres. Puis développer
$ \cosh u = 1 + \frac{u^2}{2} + \frac{u^4}{24} + \cdots $ et expliquer pourquoi une chaîne peu
tendue paraît néanmoins parabolique.
\item[(d)] Changer le chargement : suspendre un tablier au câble, de sorte que la charge soit $ w_0 $
par unité de longueur \textit{horizontale}, le poids propre du câble étant négligeable. Refaire
l'équilibre des forces, établir la nouvelle équation, et démontrer que ce câble est exactement
une parabole. Laquelle de ces deux courbes suit le câble porteur d'un pont suspendu, et laquelle
suit un collier ?
\end{enumerate}

\end{problem}

\noindent Galilée affirmait dans les \textit{Discours concernant deux sciences nouvelles}
(1638) qu'une chaîne suspendue est une parabole ; Joachim Jungius a trouvé l'affirmation fausse,
et en 1690 Jacques Bernoulli a publiquement mis les mathématiciens d'Europe au défi de produire
la vraie courbe. Huygens, Leibniz et Jean Bernoulli ont chacun publié une solution dans les
\textit{Acta Eruditorum} en 1691 --- Huygens a forgé le nom de \textit{catenaria}, du latin
\textit{catena}, chaîne --- et la réponse, qui exige le cosinus hyperbolique, fut une victoire
précoce du calcul nouveau sur un problème que les Anciens ne pouvaient que deviner. Robert Hooke
avait vu la conséquence pour l'ingénierie dès 1675 et l'avait publiée sous forme d'anagramme
latine : comme pend un câble souple, ainsi, renversés, se tiennent les claveaux d'une arche. La
partie (d) explique pourquoi les deux formes sont sans cesse confondues --- c'est la charge, non la
courbe, qui diffère.

\begin{refs}
\item Contexte : Chaînette --- Wikipédia (\url{https://fr.wikipedia.org/wiki/Cha%C3%AEnette})
\item Contexte : Fonction hyperbolique --- Wikipédia (\url{https://fr.wikipedia.org/wiki/Fonction_hyperbolique})
\item Lecture : G. F. Simmons, \textit{Differential Equations with Applications and Historical Notes}, chap. 1
\item Cours : MIT OCW 18.03SC Differential Equations (\url{https://ocw.mit.edu/courses/18-03sc-differential-equations-fall-2011/})
\end{refs}

\bigskip

\begin{problem}[{La vitesse de libération --- Lycée, short}]
\textbf{ODE-35.} Un projectile est tiré verticalement depuis la surface
d'une planète sans atmosphère, de masse $ M $ et de rayon $ R $, à la vitesse
$ v_0 > 0 $, puis se déplace sous la seule action de la gravitation :
\[ \frac{d^2 r}{dt^2} = -\frac{GM}{r^2}, \qquad r(0) = R, \qquad r'(0) = v_0 . \]
En posant $ v = dr/dt $ et en utilisant $ \frac{dv}{dt} = v\,\frac{dv}{dr} $, déterminer avec
démonstration la plus petite valeur de $ v_0 $ pour laquelle $ r(t) $ croît sans borne, et
démontrer que pour toute valeur plus petite de $ v_0 $ le projectile atteint une hauteur
maximale puis retombe.
\end{problem}

\noindent Newton a dessiné l'image dans son \textit{Traité du système du monde} : un canon
sur une très haute montagne, tirant de plus en plus fort, jusqu'à ce que le boulet ne retombe
plus. La substitution demandée ici est toute l'astuce --- elle transforme une équation du second
ordre en une équation du premier ordre en la variable $ r $, et ce qu'elle produit est la
conservation de l'énergie, la même intégrale première qui donne la période du pendule en ODE-12.
Le nombre qu'elle fournit vaut environ 11,2 km/s pour la Terre et 2,4 km/s pour la Lune, ce qui
explique en grande partie pourquoi la Lune n'a pas d'atmosphère à perdre.

\begin{refs}
\item Contexte : Vitesse de libération --- Wikipédia (\url{https://fr.wikipedia.org/wiki/Vitesse_de_lib%C3%A9ration})
\item Contexte : Canon de Newton --- Wikipédia (\url{https://fr.wikipedia.org/wiki/Canon_de_Newton})
\item Lecture : W. E. Boyce \& R. C. DiPrima, \textit{Elementary Differential Equations}, chap. 1--2
\end{refs}

\bigskip


\section*{Analyse numérique}

\begin{problem}[{Pourquoi 0,1 + 0,2 $\ne$ 0,3 sur un ordinateur --- Lycée}]
\textbf{NUM-01.} Les ordinateurs stockent les nombres réels en virgule
flottante binaire : un \og{} double \fg{} IEEE 754 est un nombre
$\pm(1{,}b_1 b_2 \ldots b_{52})_2 \times 2^{e}$ --- un signe, 52 chiffres binaires après
le 1 initial et un exposant entier --- et tout résultat arithmétique est arrondi au
nombre de cette forme le plus proche.
\begin{enumerate}
\item[(a)] Démontrez que $1/10$ n'a pas de développement binaire fini
(regardez le dénominateur) et montrez que son développement infini vaut
$(0{,}0\overline{0011})_2$.
\item[(b)] Montrez que le double le plus proche de $1/10$ est exactement
$3602879701896397/2^{55}$, qui dépasse $1/10$ d'environ $5{,}5 \times 10^{-18}$.
\item[(c)] Trouvez les doubles les plus proches de $0{,}2$ et de $0{,}3$, effectuez l'addition
machine des valeurs stockées de $0{,}1$ et $0{,}2$ (additionnez exactement, puis
arrondissez) et montrez que le résultat diffère du $0{,}3$ stocké d'exactement une
unité au dernier rang, $2^{-54}$ --- de sorte que l'ordinateur déclare honnêtement
\texttt{0.1 + 0.2 == 0.3} faux.
\item[(d)] Expliquez pourquoi les programmeurs comparent donc
$|a - b| < \text{tolérance}$ plutôt que $a = b$.
\end{enumerate}

\end{problem}

\noindent Rien n'est cassé : la machine additionne parfaitement les mauvais
nombres. La norme IEEE 754 (1985, menée par William Kahan, qui en a reçu le prix
Turing) a rendu ce comportement d'arrondi précis et universel, ce qui est justement
ce qui rend possible un raisonnement comme celui de (b)--(c). Tous les langages sur
tous les processeurs modernes présentent cette fameuse surprise ; la comprendre
exactement, au lieu de hausser les épaules, est le premier pas de la pensée
numérique.

\begin{refs}
\item Contexte : Virgule flottante --- Wikipédia (\url{https://fr.wikipedia.org/wiki/Virgule_flottante})
\item Contexte : Format de virgule flottante en double précision --- Wikipédia (en anglais) (\url{https://en.wikipedia.org/wiki/Double-precision_floating-point_format})
\item Article : D. Goldberg, \og{} What Every Computer Scientist Should Know About Floating-Point Arithmetic \fg{} (1991), ACM Computing Surveys 23
\item Lecture : N. J. Higham, \textit{Accuracy and Stability of Numerical Algorithms}, ch. 1--2
\end{refs}

\bigskip

\begin{problem}[{Annulation catastrophique dans la formule du second degré --- Lycée}]
\textbf{NUM-02.} Considérez $x^2 - 10^8 x + 1 = 0$.
\begin{enumerate}
\item[(a)] Avec une
arithmétique qui conserve 10 chiffres décimaux significatifs à chaque étape, évaluez
la formule des manuels $x = \big(-b - \sqrt{b^2 - 4ac}\big)/(2a)$ pour la plus petite
racine et montrez que la réponse calculée vaut $0$ --- une erreur relative de 100 \% ---
alors que la vraie petite racine vaut environ $10^{-8}$.
\item[(b)] Diagnostiquez le mal : montrez que
$\sqrt{b^2 - 4ac}$ coïncide avec $|b|$ sur ses premiers chiffres, de sorte que la
soustraction annule presque tous les chiffres corrects et promeut l'erreur d'arrondi
au premier rang. C'est l'\textit{annulation catastrophique}.
\item[(c)] Le remède : démontrez les identités algébriques
qui fondent l'algorithme stable --- calculez d'abord la grande racine (sans annulation),
puis obtenez la petite par la relation de Viète $x_1 x_2 = c/a$. Vérifiez qu'il
retrouve la petite racine presque avec toute la précision, dans la même arithmétique
à 10 chiffres.
\item[(d)] Montrez que
$\sqrt{x+1} - \sqrt{x}$ souffre du même mal pour $x$ grand, et corrigez-le en
rationalisant.
\end{enumerate}

\end{problem}

\noindent L'annulation est la manière classique dont une formule
mathématiquement exacte devient un algorithme numériquement ruineux : la formule et
l'algorithme sont deux objets différents, ce qui est l'observation fondatrice du
domaine. La recette stable du second degré est un savoir commun qui remonte aux
tout premiers jours du calcul automatique, et la leçon se généralise : réorganisez
l'algèbre pour ne jamais soustraire deux quantités presque égales que vous avez
calculées de façon inexacte.

\begin{refs}
\item Contexte : Annulation catastrophique --- Wikipédia (\url{https://fr.wikipedia.org/wiki/Annulation_catastrophique})
\item Contexte : Équation du second degré --- Wikipédia (\url{https://fr.wikipedia.org/wiki/%C3%89quation_du_second_degr%C3%A9})
\item Lecture : N. J. Higham, \textit{Accuracy and Stability of Numerical Algorithms}, ch. 1
\end{refs}

\bigskip

\begin{problem}[{La racine carrée babylonienne --- Lycée}]
\textbf{NUM-03.} Pour calculer $\sqrt{a}$ avec $a > 0$,
itérez la règle antique : faites la moyenne de votre estimation et de $a$ divisé par
cette estimation,
\[ x_{n+1} = \frac{1}{2}\Big( x_n + \frac{a}{x_n} \Big). \]
(a) Montrez que $\sqrt{a}$ est l'unique point fixe positif, et que
$x_n \ge \sqrt{a}$ pour tout $n \ge 1$, quelle que soit l'estimation initiale
$x_0 > 0$. (b) Démontrez que la suite décroît à partir de $n = 1$ et converge vers
$\sqrt{a}$. (c) Démontrez l'identité d'erreur
$x_{n+1} - \sqrt{a} = (x_n - \sqrt{a})^2/(2x_n)$, et concluez que le nombre de
chiffres corrects \textit{double} à peu près à chaque étape. (d) En partant de
$x_0 = 3/2$, calculez $\sqrt{2}$ à la main avec six décimales, en comptant le
nombre d'étapes nécessaires.
\end{problem}

\noindent La tablette d'argile babylonienne YBC 7289 (vers 1800--1600 av. J.-C.)
affiche $\sqrt{2}$ sous la forme $1;24,51,10$ en base 60 --- exact à environ six
décimales --- presque certainement obtenu par une règle de ce genre, que Héron
d'Alexandrie a consignée explicitement au premier siècle de notre ère. Le doublement
des chiffres de (c) est la plus ancienne apparition de la convergence quadratique, et
la méthode n'est autre que la méthode de Newton (NUM-04) deux mille ans avant Newton,
appliquée à $x^2 - a$.

\begin{refs}
\item Contexte : Extraction de racine carrée --- Wikipédia (\url{https://fr.wikipedia.org/wiki/Extraction_de_racine_carr%C3%A9e})
\item Contexte : YBC 7289 --- Wikipédia (\url{https://fr.wikipedia.org/wiki/YBC_7289})
\item Lecture : E. Süli et D. F. Mayers, \textit{An Introduction to Numerical Analysis}, ch. 1
\end{refs}

\bigskip

\begin{problem}[{La garantie de la dichotomie --- Lycée, short}]
\textbf{NUM-20.} Soit $f$ continue sur $[a,b]$ avec $f(a)f(b) < 0$ ; remplacez répétitivement l'intervalle par celle de ses deux moitiés qui présente encore un changement de signe aux extrémités. Démontrez qu'après $n$ étapes l'intervalle survivant a pour longueur $(b-a)/2^n$ et contient toujours une racine de $f$, et déterminez combien d'étapes garantissent six décimales correctes pour $f(x) = x^3 - x - 1$ sur $[1,2]$.
\end{problem}

\noindent Bernard Bolzano a employé exactement cet argument de bissection en 1817 pour donner la première démonstration du théorème des valeurs intermédiaires qui ne fasse pas appel à l'intuition géométrique --- la méthode est plus ancienne comme calcul que comme théorème. C'est le chercheur de racines le plus lent en usage courant, qui gagne un chiffre binaire par étape là où la méthode de Newton (NUM-04) double les chiffres, et c'est aussi le seul qui ne puisse pas échouer : il ne demande à $f$ que la continuité et un changement de signe. C'est pourquoi les routines sérieuses des bibliothèques, comme la méthode de Brent, gardent une étape de dichotomie en réserve derrière leurs itérations rapides.

\begin{refs}
\item Contexte : Méthode de dichotomie --- Wikipédia (\url{https://fr.wikipedia.org/wiki/M%C3%A9thode_de_dichotomie})
\item Contexte : Théorème des valeurs intermédiaires --- Wikipédia (\url{https://fr.wikipedia.org/wiki/Th%C3%A9or%C3%A8me_des_valeurs_interm%C3%A9diaires})
\item Lecture : E. Süli et D. F. Mayers, \textit{An Introduction to Numerical Analysis}, ch. 1
\end{refs}

\bigskip

\begin{problem}[{L'emboîtement de Horner --- Lycée, short}]
\textbf{NUM-21.} Le schéma de Horner évalue $p(x) = a_n x^n + \cdots + a_1 x + a_0$ par l'emboîtement
\[ \big( \cdots ((a_n x + a_{n-1})x + a_{n-2})x + \cdots \big) x + a_0 . \]
Démontrez par récurrence sur $n$ qu'il rend bien $p(x)$, et comptez ses multiplications et ses additions face à celles que demande la formation séparée de chaque puissance $x^k$.
\end{problem}

\noindent La règle porte le nom de William George Horner, qui l'a publiée en 1819, mais Paolo Ruffini la possédait en 1804, Newton l'utilisait vers 1670 et le mathématicien chinois Qin Jiushao l'a exposée en 1247 --- une chaîne de redécouvertes typique de ce qui est utile et élémentaire. Le décompte qu'on vous demande, $n$ multiplications au lieu d'environ $n^2/2$, n'est pas qu'une économie : V. Pan a démontré en 1966 qu'aucun schéma travaillant à partir des coefficients donnés ne peut évaluer un polynôme général de degré $n$ avec moins de $n$ multiplications ; c'est donc l'un des rares algorithmes dont on sait qu'ils sont optimaux.

\begin{refs}
\item Contexte : Méthode de Ruffini-Horner --- Wikipédia (\url{https://fr.wikipedia.org/wiki/M%C3%A9thode_de_Ruffini-Horner})
\item Lecture : Knuth, \textit{The Art of Computer Programming}, vol. 2, \S{}4.6.4
\end{refs}

\bigskip

\begin{problem}[{Quelle taille doit avoir $h$ ? --- Lycée}]
\textbf{NUM-22.} Un ordinateur stocke les nombres avec une erreur relative d'au plus $\varepsilon \approx 1{,}1 \times 10^{-16}$ (double précision), de sorte que les valeurs de $f$ qu'il renvoie ne sont correctes qu'à cette précision relative. Estimez $f'(x)$ par le quotient de différences
\[ D_h f(x) = \frac{f(x+h) - f(x)}{h}. \]
\begin{enumerate}
\item[(a)] À l'aide de la formule de Taylor, montrez que l'erreur en arithmétique exacte (erreur de troncature) vaut environ $\tfrac{1}{2} h\,|f''(x)|$.
\item[(b)] Montrez que l'erreur d'arrondi apportée par les deux valeurs inexactes de la fonction vaut environ $2\varepsilon |f(x)| / h$, et remarquez qu'elle \textit{croît} quand $h$ diminue.
\item[(c)] Minimisez la somme des deux en $h$. Concluez que le meilleur $h$ possible est de taille environ $\sqrt{\varepsilon}$ et que la meilleure précision atteignable n'est que d'environ $\sqrt{\varepsilon} \approx 10^{-8}$ --- la moitié de vos chiffres est perdue, quel que soit le soin de votre programme.
\item[(d)] Refaites l'analyse pour la différence centrée $\big(f(x+h) - f(x-h)\big)/(2h)$, dont l'erreur de troncature est en $O(h^2)$, et montrez que le $h$ optimal est maintenant de taille $\varepsilon^{1/3}$, pour une précision d'environ $\varepsilon^{2/3}$.
\item[(e)] Faites l'essai : avec $f = \sin$ en $x = 1$, dressez le tableau de l'erreur de $D_h f(1)$ pour $h = 10^{-1}, 10^{-2}, \dots, 10^{-16}$ et trouvez le fond du V.
\end{enumerate}

\end{problem}

\noindent Tout cours d'analyse enseigne que le quotient de différences tend vers la dérivée quand $h \to 0$ ; tout calcul numérique découvre que sur une machine réelle il fait le contraire dès que $h$ devient petit, parce que soustraire deux nombres presque égaux détruit les chiffres qui portaient l'information (l'annulation de NUM-02, cette fois avec de vrais dégâts). La courbe d'erreur en V de (e) est l'image la plus instructive de toute l'analyse numérique élémentaire : la dérivation est une opération mal conditionnée, qui amplifie le bruit qu'on lui donne, tandis que l'intégration le lisse. L'échappatoire moderne consiste à cesser complètement de faire des différences --- la dérivation automatique propage les règles exactes de dérivation à côté de l'arithmétique et n'a aucun $h$ à choisir, ce qui explique que ce soit elle, et non les différences finies, qui entraîne les réseaux de neurones (voir NUM-19).

\begin{refs}
\item Contexte : Dérivation numérique --- Wikipédia (\url{https://fr.wikipedia.org/wiki/D%C3%A9rivation_num%C3%A9rique})
\item Contexte : Dérivation automatique --- Wikipédia (\url{https://fr.wikipedia.org/wiki/D%C3%A9rivation_automatique})
\item Lecture : N. J. Higham, \textit{Accuracy and Stability of Numerical Algorithms}, ch. 1
\item Cours : MIT OCW 18.330 Introduction to Numerical Analysis (\url{https://ocw.mit.edu/courses/18-330-introduction-to-numerical-analysis-spring-2012/})
\end{refs}

\bigskip

\begin{problem}[{La variance que l'ordinateur calcule de travers --- Lycée, short}]
\textbf{NUM-32.} La formule en une passe des manuels pour la variance empirique de $x_1, \ldots, x_n$,
\[ s^2 = \frac{1}{n-1}\left( \sum_{i=1}^{n} x_i^2 \;-\; \frac{1}{n}\Big(\sum_{i=1}^{n} x_i\Big)^{\!2} \right), \]
est algébriquement correcte ; montrez qu'évaluée en double précision IEEE elle renvoie $0$ pour les quatre nombres $10^9 + 4,\ 10^9 + 7,\ 10^9 + 13,\ 10^9 + 16$, dont la variance empirique vaut exactement $30$, et identifiez quelle soustraction a détruit la réponse et quelle était la taille des quantités en présence comparée à celle de leur différence. Démontrez ensuite que les récurrences de Welford
\[ M_k = M_{k-1} + \frac{x_k - M_{k-1}}{k}, \qquad S_k = S_{k-1} + (x_k - M_{k-1})(x_k - M_k), \]
issues de $M_0 = S_0 = 0$, vérifient $M_k = \tfrac{1}{k}\sum_{i \le k} x_i$ et $S_k = \sum_{i \le k} (x_i - M_k)^2$ pour tout $k$, de sorte que $s^2 = S_n/(n-1)$ s'obtient en une seule passe sans jamais former $\sum x_i^2$.
\end{problem}

\noindent B. P. Welford a publié cette mise à jour dans une note d'une page de \textit{Technometrics} en 1962, et Knuth l'a mise dans \textit{The Art of Computer Programming}, d'où elle a gagné toutes les bibliothèques de statistique ; c'est la façon standard dont une base de données, un capteur ou un flux de traitement tient aujourd'hui une variance à jour. La panne qu'elle répare est l'annulation de NUM-02 à l'endroit où elle fait le plus de dégâts en pratique, car la formule naïve est celle qu'impriment la plupart des manuels et elle échoue en silence --- les premiers tableurs étaient tristement célèbres pour annoncer des variances négatives, et une variance négative se signale au moins d'elle-même, ce que ne fait pas une variance positive fausse. Notez ce que la forme de Welford fait différemment : elle travaille avec des écarts à la moyenne courante, qui sont petits, plutôt qu'avec des carrés bruts, qui sont énormes.

\begin{refs}
\item Contexte : Algorithme de calcul de la variance --- Wikipédia (\url{https://fr.wikipedia.org/wiki/Algorithme_de_calcul_de_la_variance})
\item Contexte : Annulation catastrophique --- Wikipédia (\url{https://fr.wikipedia.org/wiki/Annulation_catastrophique})
\item Article : B. P. Welford, \og{} Note on a method for calculating corrected sums of squares and products \fg{} (1962), Technometrics 4
\item Lecture : D. E. Knuth, \textit{The Art of Computer Programming}, vol. 2, \S{}4.2.2
\end{refs}

\bigskip

\begin{problem}[{La sommation compensée de Kahan --- Lycée}]
\textbf{NUM-33.} Additionner une liste de nombres est l'opération la plus courante du calcul scientifique, et la boucle naïve n'est pas la meilleure façon de le faire. Travaillez en virgule flottante binaire avec une unité d'arrondi $\varepsilon$ (pour les doubles IEEE, $\varepsilon = 2^{-53} \approx 1{,}1 \times 10^{-16}$), où chaque opération obéit à
\[ \mathrm{fl}(a \pm b) = (a \pm b)(1 + \delta), \qquad |\delta| \le \varepsilon . \]
\begin{enumerate}
\item[(a)] Montrez que l'addition en virgule flottante n'est pas associative : trouvez trois doubles tels que $(a + b) + c \ne a + (b + c)$ lorsque chaque addition est arrondie. Concluez que la somme d'une liste dépend de l'ordre dans lequel vous l'additionnez.
\item[(b)] Pour la somme naïve de gauche à droite $\hat{s}$ de $x_1, \ldots, x_n$, démontrez la majoration de l'erreur
\[ |\hat{s} - s| \;\le\; (n-1)\,\varepsilon \sum_{i=1}^{n} |x_i| \;+\; O(\varepsilon^2), \]
et exhibez des données sur lesquelles l'erreur croît réellement avec $n$.
\item[(c)] Démontrez le lemme de Dekker, le moteur de tout le sujet : si $|a| \ge |b|$ et $s = \mathrm{fl}(a + b)$, alors le nombre $(a - s) + b$ est \textit{exactement représentable} et se calcule exactement en deux opérations flottantes supplémentaires --- l'erreur d'arrondi commise par une addition est elle-même un nombre flottant, et vous pouvez la récupérer.
\item[(d)] L'algorithme de Kahan reporte cette erreur récupérée sur l'étape suivante : posez $s = 0$, $c = 0$ et, pour chaque $x$, calculez
\[ y = x - c, \qquad t = s + y, \qquad c = (t - s) - y, \qquad s = t, \]
chaque ligne étant arrondie. Déroulez-le à la main sur les trois nombres $1,\ \varepsilon,\ \varepsilon$ (avec $\varepsilon = 2^{-53}$ et l'arrondi au plus proche pair) et montrez que la boucle naïve renvoie exactement $1$ tandis que celle de Kahan renvoie exactement $1 + 2^{-52}$. Énoncez ensuite, et démontrez autant que vous le pouvez, le théorème de Kahan : la somme compensée vérifie $|\hat{s} - s| \le \big(2\varepsilon + O(n\varepsilon^2)\big)\sum_i |x_i|$ --- le terme dominant ne contient aucun $n$.
\item[(e)] Comparez une troisième stratégie : la sommation par paires additionne la première moitié, additionne la seconde moitié, puis additionne les deux résultats, récursivement. Démontrez que son erreur est majorée proportionnellement à $\varepsilon \log_2 n \sum_i |x_i|$, et expliquez pourquoi une bibliothèque pourrait en faire son choix par défaut tout en réservant la boucle de Kahan aux cas qui l'exigent.
\end{enumerate}

\end{problem}

\noindent William Kahan a publié cela sous forme d'une note d'une demi-page dans les \textit{Communications of the ACM} en 1965 --- une \og{} Pracnique \fg{}, le nom que la revue donnait aux techniques pratiques --- vingt ans avant de diriger la conception d'IEEE 754 (NUM-01) et d'en recevoir le prix Turing. Le résultat de (d) est le plus étonnant : une précision qui ne se dégrade pas quand la liste s'allonge, obtenue avec quatre opérations supplémentaires par élément et sans aucune précision supplémentaire, ce qui fait dire parfois que la sommation compensée achète la double précision au prix d'une erreur d'arrondi. Arnold Neumaier a amélioré l'algorithme en 1974 pour traiter le cas où le nouveau terme est bien plus grand que la somme courante, et les bibliothèques de tableaux modernes utilisent par défaut le schéma par paires de (e). Il y a un avertissement caché dans (a) pour qui écrit du code parallèle : une somme répartie sur un nombre différent de processeurs est un nombre différent, si bien qu'un programme peut ne pas reproduire sa propre réponse.

\begin{refs}
\item Contexte : Algorithme de sommation de Kahan --- Wikipédia (en anglais) (\url{https://en.wikipedia.org/wiki/Kahan_summation_algorithm})
\item Contexte : Epsilon d'une machine --- Wikipédia (\url{https://fr.wikipedia.org/wiki/Epsilon_d%27une_machine})
\item Article : W. Kahan, \og{} Pracniques: further remarks on reducing truncation errors \fg{} (1965), Communications of the ACM 8
\item Lecture : N. J. Higham, \textit{Accuracy and Stability of Numerical Algorithms}, ch. 4 (Summation)
\end{refs}

\bigskip


\section*{Optimisation et théorie des jeux}

\begin{problem}[{Le meilleur rectangle, sans le calcul différentiel --- Lycée}]
\textbf{OPT-01.} Aucune dérivée n'est autorisée dans ce problème.
\begin{enumerate}
\item[(a)] Démontrez que pour tous réels positifs ou nuls $a$ et $b$,
\[ \frac{a+b}{2} \ge \sqrt{ab}, \]
avec égalité exactement lorsque $a = b$.
\item[(b)] Déduisez-en que parmi tous les rectangles de périmètre donné, le carré a la plus grande aire.
\item[(c)] Un agriculteur dispose de $100$ m de clôture et veut enclore un enclos rectangulaire adossé à une rivière rectiligne, en ne clôturant que trois côtés. Déterminez, avec démonstration, les dimensions qui donnent l'aire maximale.
\end{enumerate}

\end{problem}

\noindent L'inégalité entre moyenne arithmétique et moyenne géométrique est le plus ancien outil d'optimisation qui soit : les \textit{Éléments} d'Euclide (livre VI) contiennent déjà le cas à deux variables sous un habit géométrique, et l'instinct isopérimétrique derrière (b) remonte à la légende de la reine Didon délimitant Carthage avec une peau de b\oe{}uf. Les maxima sous contrainte ont été démontrés rigoureusement deux mille ans avant l'existence du calcul différentiel ; retrouver cette capacité est tout l'objet de ce problème.

\begin{refs}
\item Contexte : Inégalité arithmético-géométrique --- Wikipédia (\url{https://fr.wikipedia.org/wiki/In%C3%A9galit%C3%A9_arithm%C3%A9tico-g%C3%A9om%C3%A9trique})
\item Contexte : Théorème isopérimétrique --- Wikipédia (\url{https://fr.wikipedia.org/wiki/Th%C3%A9or%C3%A8me_isop%C3%A9rim%C3%A9trique})
\item Lecture : Steele, \textit{The Cauchy--Schwarz Master Class}, ch. 2
\end{refs}

\bigskip

\begin{problem}[{Une boîte ouverte et l'inégalité arithmético-géométrique à trois variables --- Lycée}]
\textbf{OPT-02.} Une boîte à base carrée et sans couvercle doit contenir exactement $32\,000$ cm$^3$.
\begin{enumerate}
\item[(a)] Démontrez l'inégalité entre moyenne arithmétique et moyenne géométrique pour trois réels positifs ou nuls :
\[ \frac{x+y+z}{3} \ge \sqrt[3]{xyz}, \]
avec égalité exactement lorsque $x = y = z$. (La récurrence en avant puis en arrière de Cauchy est une voie possible ; trouvez-en une correcte, quelle qu'elle soit.)
\item[(b)] À l'aide de (a) --- et non du calcul différentiel --- déterminez, avec démonstration, le côté de la base et la hauteur qui minimisent l'aire totale de la base et des quatre côtés, et démontrez que votre configuration est l'unique minimiseur.
\end{enumerate}

\end{problem}

\noindent Cauchy a démontré l'inégalité arithmético-géométrique générale à n variables dans son \textit{Cours d'analyse} de 1821, par une récurrence qui bondit en avant vers les puissances de deux puis redescend pas à pas --- l'un des schémas de démonstration les plus imités de l'analyse. La partie (b) est le genre de problème de conception (le moins de matière pour une capacité fixée) qui remplissait les manuels des XVIIIe et XIXe siècles ; le résoudre par pure inégalité, en découpant l'aire de surface en morceaux de produit constant, donne un avant-goût de la façon dont pensent ceux qui optimisent, avant qu'ils ne saisissent un gradient.

\begin{refs}
\item Contexte : Inégalité arithmético-géométrique --- Wikipédia (\url{https://fr.wikipedia.org/wiki/In%C3%A9galit%C3%A9_arithm%C3%A9tico-g%C3%A9om%C3%A9trique})
\item Lecture : Niven, \textit{Maxima and Minima Without Calculus}
\item Lecture : Steele, \textit{The Cauchy--Schwarz Master Class}, ch. 2
\end{refs}

\bigskip

\begin{problem}[{Le problème du régime alimentaire, concrètement --- Lycée}]
\textbf{OPT-03.} L'aliment A coûte 50 centimes les 100 g et fournit 20 g de protéines et 200 kcal ; l'aliment B coûte 20 centimes les 100 g et fournit 2 g de protéines et 250 kcal. Le régime d'une journée doit apporter au moins 60 g de protéines et au moins 2000 kcal.
\begin{enumerate}
\item[(a)] Justifiez le modèle : minimiser $50a + 20b$ sous les contraintes $20a + 2b \ge 60$, $200a + 250b \ge 2000$, $a, b \ge 0$.
\item[(b)] Tracez la région admissible et trouvez le régime le moins cher.
\item[(c)] Certifiez l'optimalité : trouvez des multiplicateurs positifs ou nuls $y_1, y_2$ pour les deux contraintes nutritionnelles tels que $y_1 \cdot (\text{ligne des protéines}) + y_2 \cdot (\text{ligne des calories})$ démontre, en deux lignes d'algèbre, qu'aucun régime admissible ne peut coûter moins que le vôtre.
\end{enumerate}

\end{problem}

\noindent C'est le problème du régime alimentaire, posé par l'économiste George Stigler en 1945 : quelle est la façon la moins chère de rester en vie ? Stigler a résolu heuristiquement une version à 77 aliments et 9 nutriments (39,93 \$ par an aux prix de 1939) ; en 1947, l'équipe de Dantzig l'a résolue de nouveau avec la toute nouvelle méthode du simplexe --- neuf employés sur des calculatrices de bureau pendant environ 120 jours-personne --- et l'estimation de Stigler ne s'écartait que de 24 cents. La partie (c) est votre premier certificat dual : les multiplicateurs sont des prix pour les protéines et les calories, et ils sont le germe de toute la théorie d'OPT-06.

\begin{refs}
\item Contexte : Problème du régime de Stigler --- Wikipédia (en anglais) (\url{https://en.wikipedia.org/wiki/Stigler_diet})
\item Contexte : Optimisation linéaire --- Wikipédia (\url{https://fr.wikipedia.org/wiki/Optimisation_lin%C3%A9aire})
\item Article : Stigler, \og{} The Cost of Subsistence \fg{}, Journal of Farm Economics 27 (1945)
\item Cours : MIT OCW 15.053 Optimization Methods in Management Science (\url{https://ocw.mit.edu/courses/15-053-optimization-methods-in-management-science-spring-2013/})
\end{refs}

\bigskip

\begin{problem}[{Le problème des transports, concrètement --- Lycée}]
\textbf{OPT-04.} L'entrepôt 1 détient 30 caisses, l'entrepôt 2 en détient 25 ; les magasins X, Y, Z ont besoin de 20, 20 et 15 caisses. Expédier une caisse coûte, en dollars : de E1 vers (X, Y, Z) : 4, 6, 9 ; de E2 : 5, 3, 8.
\begin{enumerate}
\item[(a)] Formulez le plan d'expédition le moins cher comme un programme linéaire en les six inconnues $x_{ij} \ge 0$.
\item[(b)] Trouvez un plan optimal à la main.
\item[(c)] Démontrez que votre plan est optimal en attachant un nombre $u_i$ à chaque entrepôt et $v_j$ à chaque magasin avec $u_i + v_j \le c_{ij}$ pour chaque route, puis en comparant $\sum_i (\text{offre}_i) u_i + \sum_j (\text{demande}_j) v_j$ à votre coût.
\end{enumerate}

\end{problem}

\noindent Le problème des transports a été formulé par Frank Hitchcock en 1941 et, indépendamment et plus tôt, par Leonid Kantorovitch en Union soviétique, dont l'opuscule de 1939 sur la planification de la production lui valut le prix Nobel d'économie 1975 (partagé avec Tjalling Koopmans) ; son ancêtre continu est le mémoire de Gaspard Monge de 1781 sur les déblais et remblais des fortifications. Les potentiels $u, v$ de (c) sont des prix fictifs aux deux extrémités de chaque route --- encore la dualité, un an avant que quiconque l'eût nommée.

\begin{refs}
\item Contexte : Théorie du transport --- Wikipédia (\url{https://fr.wikipedia.org/wiki/Th%C3%A9orie_du_transport})
\item Lecture : Papadimitriou et Steiglitz, \textit{Combinatorial Optimization: Algorithms and Complexity}, ch. 7
\item Cours : MIT OCW 15.053 Optimization Methods in Management Science (\url{https://ocw.mit.edu/courses/15-053-optimization-methods-in-management-science-spring-2013/})
\end{refs}

\bigskip

\begin{problem}[{Le plus court chemin de Héron --- Lycée, short}]
\textbf{OPT-24.} Les points $A$ et $B$ sont du même côté d'une droite
$\ell$. Démontrez que parmi tous les chemins qui vont de $A$ à un point $P$ de $\ell$ puis
à $B$, le plus court est l'unique chemin pour lequel $AP$ et $PB$ font des angles égaux avec
$\ell$ --- et n'utilisez que la géométrie plane, pas de calcul différentiel.
\end{problem}

\noindent Héron d'Alexandrie a démontré cela dans sa \textit{Catoptrique} vers le premier
siècle, déduisant la loi de la réflexion de l'hypothèse que la lumière emprunte le plus court
chemin : le premier argument d'optimisation en physique, dix-sept siècles avant que Fermat ne le
généralise au moindre \textit{temps} et n'obtienne aussi la réfraction. La raison d'insister sur la
géométrie plane est que la démonstration géométrique montre \textit{pourquoi} le minimum est ce
qu'il est, là où une dérivée ne fait que le confirmer.

\begin{refs}
\item Contexte : Principe de Fermat --- Wikipédia (\url{https://fr.wikipedia.org/wiki/Principe_de_Fermat})
\item Lecture : Niven, \textit{Maxima and Minima Without Calculus}
\end{refs}

\bigskip

\begin{problem}[{Quand la gourmandise a raison, et quand elle a tort --- Lycée, short}]
\textbf{OPT-25.} Avec des pièces de 1, 5, 10 et 25 centimes, démontrez que la règle
gloutonne --- prendre répétitivement la plus grosse pièce ne dépassant pas ce qui reste --- règle tout
nombre entier de centimes avec le moins de pièces possible. Exhibez ensuite un système de pièces
où la gourmandise échoue, et démontrez qu'elle échoue au montant que vous désignez.
\end{problem}

\noindent Les algorithmes gloutons sont ceux que tout le monde invente en premier et que
presque personne ne démontre ; le travail honnête ici est l'analyse de cas montrant qu'une solution
optimale peut toujours être réarrangée en la solution gloutonne. Les systèmes de pièces où la
gourmandise fonctionne sont dits \textit{canoniques}, et savoir si un système donné est canonique se
décide en temps polynomial (Pearson, 2005) --- un résultat étonnamment récent pour une question si
ancienne. La question générale des structures qui rendent la gourmandise optimale a une belle
réponse : c'est OPT-30.

\begin{refs}
\item Contexte : Problème du rendu de monnaie --- Wikipédia (\url{https://fr.wikipedia.org/wiki/Probl%C3%A8me_du_rendu_de_monnaie})
\item Contexte : Algorithme glouton --- Wikipédia (\url{https://fr.wikipedia.org/wiki/Algorithme_glouton})
\end{refs}

\bigskip

\begin{problem}[{Le paradoxe de Braess et le prix de l'anarchie --- Lycée}]
\textbf{OPT-26.} Quatre mille automobilistes vont de $S$ à $T$. Il y a deux itinéraires,
$S \to A \to T$ et $S \to B \to T$. Les arêtes $S \to A$ et $B \to T$ prennent chacune $x/100$
minutes quand $x$ automobilistes les empruntent ; les arêtes $A \to T$ et $S \to B$ prennent chacune
45 minutes fixes.
\begin{enumerate}
\item[(a)] Trouvez la répartition du trafic pour laquelle aucun automobiliste seul ne peut raccourcir son
trajet en changeant d'itinéraire, et démontrez qu'alors chaque automobiliste a besoin de 65 minutes.
Démontrez que cette même répartition minimise aussi le temps de trajet total de tous les
automobilistes.
\item[(b)] Une nouvelle route $A \to B$ est ouverte, si rapide que son temps de parcours est nul.
Démontrez que le seul agencement où aucun automobiliste ne souhaite changer est que \textit{tout le
monde} emprunte $S \to A \to B \to T$, et que chaque automobiliste a désormais besoin de 80
minutes. Construire une route a rendu la situation pire pour les 4000 personnes.
\item[(c)] Montrez qu'un planificateur central, la nouvelle route étant donnée, peut malgré tout obtenir
une moyenne de 65 minutes. Le \textit{prix de l'anarchie} est le rapport du pire résultat égoïste au
meilleur résultat coordonné : calculez-le pour ce réseau.
\item[(d)] Expliquez la contradiction apparente avec vos propres mots : ajouter une option ne peut pas
nuire à un décideur isolé, alors pourquoi cela nuit-il à une foule de décideurs ?
\end{enumerate}

\end{problem}

\noindent Dietrich Braess a publié le paradoxe en 1968 dans une revue allemande de recherche
opérationnelle, et il n'a cessé de resurgir dans le monde réel : la circulation de Stuttgart s'est
améliorée après la \textit{fermeture} d'un nouveau tronçon de route, la 42e rue de New York
s'est mieux écoulée le jour où elle a été barrée en 1990, et la démolition de la voie express de
Cheonggyecheon à Séoul en 2005 est le plus grand exemple enregistré. Koutsoupias et Papadimitriou
ont baptisé le rapport de (c) prix de l'anarchie en 1999, et Roughgarden et Tardos ont démontré en
2002 que pour des latences de la forme $ax + b$ il ne dépasse jamais $4/3$ --- dans \textit{n'importe
quel} réseau, si embrouillé soit-il. L'égoïsme est borné ; il n'est simplement jamais
gratuit.

\begin{refs}
\item Contexte : Paradoxe de Braess --- Wikipédia (\url{https://fr.wikipedia.org/wiki/Paradoxe_de_Braess})
\item Contexte : Prix de l'anarchie --- Wikipédia (\url{https://fr.wikipedia.org/wiki/Prix_de_l%27anarchie})
\item Article : Roughgarden et Tardos, \og{} How bad is selfish routing? \fg{}, Journal of the ACM 49(2) (2002)
\item Lecture : Nisan, Roughgarden, Tardos et Vazirani (dir.), \textit{Algorithmic Game Theory}, ch. 18
\end{refs}

\bigskip

\begin{problem}[{Appariez le plus grand avec le plus grand --- Lycée, short}]
\textbf{OPT-36.} Soient $a_1 \le a_2 \le \cdots \le a_n$ et $b_1 \le b_2 \le \cdots \le b_n$ des nombres réels et soit $\sigma$ une permutation quelconque de $\{1, \ldots, n\}$. Démontrez l'inégalité de réarrangement
\[ \sum_{i=1}^{n} a_i b_{n+1-i} \;\le\; \sum_{i=1}^{n} a_i b_{\sigma(i)} \;\le\; \sum_{i=1}^{n} a_i b_i , \]
et déduisez-en que si $n$ tâches de durées connues doivent être exécutées, une par machine, sur $n$ machines de vitesses connues, le temps d'exécution total $\sum_i \ell_{\sigma(i)}/s_i$ est minimisé en confiant la tâche la plus longue à la machine la plus rapide.
\end{problem}

\noindent L'inégalité est le théorème 368 des \textit{Inequalities} de Hardy, Littlewood et Pólya (1934), même si le fait relevait du savoir commun bien avant ; l'inégalité de Tchebychev pour les sommes (1882) en est la version moyennée. Toute la démonstration tient en une observation : si deux paires sont appariées dans le mauvais ordre, les échanger ne peut pas diminuer la somme, et un nombre fini d'échanges met tout en ordre. Cet argument --- prendre une solution optimale et montrer qu'elle peut être réarrangée, pas à pas, en celle que produit votre règle --- est l'\textit{argument d'échange}, et c'est ainsi que presque tout algorithme glouton de cette page est démontré correct (OPT-25, OPT-30). C'est aussi la raison pour laquelle le problème d'affectation d'OPT-17, qui en général demande la méthode hongroise, se réduit à un tri lorsque le coût d'appariement de l'ouvrier $i$ avec la tâche $j$ se trouve être un produit $a_i b_j$.

\begin{refs}
\item Contexte : Inégalité de réarrangement --- Wikipédia (\url{https://fr.wikipedia.org/wiki/In%C3%A9galit%C3%A9_de_r%C3%A9arrangement})
\item Contexte : Inégalité de Tchebychev pour les sommes --- Wikipédia (\url{https://fr.wikipedia.org/wiki/In%C3%A9galit%C3%A9_de_Tchebychev_pour_les_sommes})
\item Lecture : Hardy, Littlewood et Pólya, \textit{Inequalities}, ch. 10
\end{refs}

\bigskip

\begin{problem}[{Le point de Fermat, et la forme d'une lame de savon --- Lycée}]
\textbf{OPT-37.} Étant donné trois points distincts $A, B, C$ du plan, trouvez le point $P$ qui minimise la distance totale
\[ d(P) \;=\; |PA| + |PB| + |PC| . \]
\begin{enumerate}
\item[(a)] Démontrez qu'un minimiseur existe et qu'il est unique. (Pour l'existence, remarquez que $d$ est continue et tend vers l'infini ; pour l'unicité, montrez que $d$ est convexe et examinez où elle cesse d'être strictement convexe.)
\item[(b)] Supposez qu'aucun angle du triangle $ABC$ ne vaut $120^\circ$ ou plus. Démontrez que le minimiseur est le point intérieur en lequel les trois segments $PA, PB, PC$ se rencontrent selon trois angles de $120^\circ$ exactement. Une voie propre : faites tourner le plan entier de $60^\circ$ autour de $A$, amenant $B$ en $B'$ et $P$ en $P'$ ; montrez que $|PA| + |PB| + |PC|$ égale la longueur du chemin brisé $C \to P \to P' \to B'$, et qu'un chemin brisé est le plus court lorsqu'il est droit.
\item[(c)] Démontrez que si un angle du triangle vaut $120^\circ$ ou plus, alors le minimiseur est ce sommet lui-même --- le point intérieur de (b) a disparu.
\item[(d)] Démontrez la construction de Torricelli : élevez un triangle équilatéral vers l'extérieur sur chaque côté de $ABC$ ; alors les trois cercles circonscrits à ces triangles passent par un point commun, qui est le minimiseur de (b), et les trois droites joignant chaque nouveau sommet extérieur au sommet opposé de $ABC$ passent également par lui et ont toutes la même longueur, égale à la valeur minimale $d(P)$.
\item[(e)] Quatre points, et une autre question. Pour les coins d'un carré unité, comparez la longueur totale des deux diagonales, $2\sqrt{2}$, à celle du réseau formé de deux points de jonction intérieurs reliés par un court segment, chaque jonction envoyant deux bras vers deux coins, tous les angles aux jonctions valant $120^\circ$. Calculez la longueur $1 + \sqrt{3}$ du second réseau, et concluez que la façon la moins chère de relier quatre points n'est pas de tout faire passer par un point unique. Expliquez pourquoi la règle des $120^\circ$ survit au changement de problème.
\end{enumerate}

\end{problem}

\noindent Fermat a posé le problème des trois points vers 1636 à la fin de son essai sur les maxima et les minima, comme un défi aux autres géomètres ; Evangelista Torricelli l'a résolu avant 1640 par la construction au cercle de la partie (d), et son élève Viviani a publié la solution en 1659. Thomas Simpson a ajouté les droites concourantes en 1750, et au XIXe siècle le nom de Jakob Steiner s'est attaché à la version en réseau de la partie (e). Le problème est l'ancêtre de la localisation d'installations --- où placer l'entrepôt, le transformateur, l'hôpital --- et la version à $n$ points, la médiane géométrique, est célèbre pour son absence de formule : Bajaj a démontré en 1988 que pour cinq points elle ne peut en général pas s'exprimer par radicaux, si bien qu'on la calcule plutôt par l'itération de Weiszfeld de 1937. Les lames de savon tendues entre des tiges plantées sur deux plaques de verre parallèles résolvent physiquement le problème de réseau en minimisant l'énergie de surface, et elles se rencontrent toujours par trois à $120^\circ$ : ce sont (b) et (e) confirmés par la nature.

\begin{refs}
\item Contexte : Point de Fermat --- Wikipédia (\url{https://fr.wikipedia.org/wiki/Point_de_Fermat})
\item Contexte : Médiane géométrique --- Wikipédia (\url{https://fr.wikipedia.org/wiki/M%C3%A9diane_g%C3%A9om%C3%A9trique})
\item Contexte : Problème de l'arbre de Steiner --- Wikipédia (\url{https://fr.wikipedia.org/wiki/Probl%C3%A8me_de_l%27arbre_de_Steiner})
\item Lecture : Courant et Robbins, \textit{What Is Mathematics?}, ch. VII
\end{refs}

\bigskip


\section*{Mathématiques discrètes et informatique}

\begin{problem}[{La parité : l'échiquier mutilé --- Lycée}]
\textbf{DCS-01.} \begin{enumerate}
\item[(a)] On retire d'un échiquier ordinaire $8 \times 8$ deux cases d'angle diagonalement opposées. Démontrez que les 62 cases restantes ne peuvent pas être pavées par 31 dominos, chacun recouvrant deux cases adjacentes.
\item[(b)] Démontrez que, dans toute assemblée de personnes, le nombre de celles qui ont serré un nombre impair de mains est pair.
\item[(c)] Énoncez, en une phrase chacun, l'unique invariant qui fait marcher chaque démonstration.
\end{enumerate}

\end{problem}

\noindent L'échiquier mutilé a été popularisé par Martin Gardner (1957) et, avant lui, par le manuel de logique de Max Black ; le fait sur les poignées de main est d'Euler, tiré de l'article de 1736 sur les ponts de Königsberg qui a fondé la théorie des graphes. La parité est l'invariant non trivial le plus simple des mathématiques, et ces deux problèmes en sont l'initiation canonique : une démonstration d'impossibilité qu'aucune quantité d'essais infructueux ne saurait fournir. John McCarthy proposa plus tard l'échiquier comme banc d'essai du raisonnement automatique, précisément parce que la courte démonstration humaine ne ressemble en rien à une recherche exhaustive.

\begin{refs}
\item Contexte : Problème de l'échiquier mutilé --- Wikipédia (\url{https://fr.wikipedia.org/wiki/Probl%C3%A8me_de_l%27%C3%A9chiquier_mutil%C3%A9})
\item Contexte : Lemme des poignées de main --- Wikipédia (\url{https://fr.wikipedia.org/wiki/Lemme_des_poign%C3%A9es_de_main})
\item Lecture : Lehman, Leighton \& Meyer, \textit{Mathematics for Computer Science}, chapitre sur les invariants --- gratuit sur MIT OCW 6.042J (\url{https://ocw.mit.edu/courses/6-042j-mathematics-for-computer-science-spring-2015/})
\end{refs}

\bigskip

\begin{problem}[{Pourquoi la recherche dichotomique est correcte --- Lycée}]
\textbf{DCS-02.} La recherche dichotomique cherche une valeur $t$ dans un tableau trié $a_1 \le a_2 \le \dots \le a_n$ en maintenant deux indices $\ell \le r$, en comparant $t$ à l'élément du milieu, puis en écartant la moitié qui ne peut pas contenir $t$.
\begin{enumerate}
\item[(a)] Écrivez l'algorithme avec précision, dans un pseudocode de votre cru.
\item[(b)] Énoncez un invariant de boucle --- une phrase vraie avant chaque itération --- et démontrez par récurrence qu'il se maintient et que l'algorithme répond toujours correctement.
\item[(c)] Démontrez que la boucle s'exécute au plus $\lfloor \log_2 n \rfloor + 1$ fois.
\item[(d)] Expliquez où votre démonstration s'effondre si le tableau n'est pas trié, et trouvez le bogue subtil de la formule \og{} évidente \fg{} du milieu $m = (\ell + r)/2$ lorsqu'on l'implémente avec des entiers de taille fixe.
\end{enumerate}

\end{problem}

\noindent Jon Bentley rapporte dans \textit{Programming Pearls} que, lorsqu'il demandait à des programmeurs professionnels d'écrire une recherche dichotomique, environ quatre-vingt-dix pour cent se trompaient ; le débordement du milieu de la question (d) a survécu des années durant dans la bibliothèque standard de Java. L'idée fut publiée pour la première fois par John Mauchly en 1946 et pourtant, d'après Knuth, la première version publiée correcte a attendu 1962. C'est le plus petit exemple complet, sur ce site, de la discipline qu'enseigne cette page : un algorithme est un théorème, et l'invariant de boucle en est l'hypothèse de récurrence.

\begin{refs}
\item Contexte : Recherche dichotomique --- Wikipédia (\url{https://fr.wikipedia.org/wiki/Recherche_dichotomique})
\item Contexte : Invariant de boucle --- Wikipédia (\url{https://fr.wikipedia.org/wiki/Invariant_de_boucle})
\item Lecture : Bentley, \textit{Programming Pearls}, ch. 4
\item Cours : MIT OCW 6.042J Mathematics for Computer Science (\url{https://ocw.mit.edu/courses/6-042j-mathematics-for-computer-science-spring-2015/})
\end{refs}

\bigskip

\begin{problem}[{Dénombrer les chaînes de bits --- Lycée}]
\textbf{DCS-03.} Une chaîne de bits est une suite finie de 0 et de 1.
\begin{enumerate}
\item[(a)] Démontrez qu'il y a exactement $2^n$ chaînes de bits de longueur $n$, et exactement $\binom{n}{k}$ d'entre elles comportant exactement $k$ uns.
\item[(b)] Soit $c_n$ le nombre de chaînes de bits de longueur $n$ sans deux 1 consécutifs. Calculez $c_1, \dots, c_5$, conjecturez la règle et démontrez que $c_n = c_{n-1} + c_{n-2}$ ; concluez que $c_n$ est un nombre de Fibonacci.
\item[(c)] Utilisez (a) pour démontrer l'identité $\binom{n}{0} + \binom{n}{1} + \dots + \binom{n}{n} = 2^n$ en dénombrant un même ensemble de deux façons.
\end{enumerate}

\end{problem}

\noindent Dénombrer des chaînes soumises à des contraintes est la plus ancienne des mathématiques de l'informatique --- de plusieurs siècles antérieure aux ordinateurs : le compte de la question (b) apparaît dans la prosodie de l'Inde ancienne (Pingala, puis plus tard Hemachandra, vers 1150), où des syllabes de longueur un et deux devaient paver un vers, anticipant Fibonacci de quelques décennies. Le double dénombrement, technique de la question (c), est l'une des méthodes de démonstration les plus discrètement puissantes de la combinatoire ; Aigner et Ziegler lui consacrent un chapitre dans \textit{Proofs from THE BOOK}.

\begin{refs}
\item Contexte : Suite de Fibonacci --- Wikipédia (\url{https://fr.wikipedia.org/wiki/Suite_de_Fibonacci})
\item Contexte : Preuve par double dénombrement --- Wikipédia (\url{https://fr.wikipedia.org/wiki/Preuve_par_double_d%C3%A9nombrement})
\item Lecture : Lehman, Leighton \& Meyer, \textit{Mathematics for Computer Science}, chapitres de dénombrement --- gratuit sur MIT OCW 6.042J (\url{https://ocw.mit.edu/courses/6-042j-mathematics-for-computer-science-spring-2015/})
\end{refs}

\bigskip

\begin{problem}[{Chevaliers et valets, formalisés --- Lycée}]
\textbf{DCS-04.} Sur une île, chaque habitant est un chevalier (qui dit toujours la vérité) ou un valet (qui ment toujours). Vous rencontrez A et B. A dit : \og{} B est un valet. \fg{} B dit : \og{} A et moi sommes de la même espèce. \fg{} (a) Introduisez des variables propositionnelles ($p$ = \og{} A est un chevalier \fg{}, $q$ = \og{} B est un chevalier \fg{}), traduisez chaque déclaration en une formule qui doit être vraie exactement lorsque celui qui la prononce est un chevalier, et résolvez le système obtenu par table de vérité. (b) Démontrez que votre réponse est la \textit{seule} qui soit cohérente. (c) Vous rencontrez maintenant C, qui dit \og{} Je suis un valet. \fg{} Démontrez qu'aucune affectation n'est cohérente, et expliquez soigneusement le rapport avec la phrase \og{} cette phrase est fausse \fg{}.
\end{problem}

\noindent Raymond Smullyan a bâti tout un genre littéraire sur ces énigmes dans \textit{What Is the Name of This Book?} (1978), et il s'en est servi comme d'un escalier doux menant jusqu'au théorème d'incomplétude de Gödel --- la question (c) en est la première marche. Le contenu sérieux, c'est la traduction elle-même : transformer une parole en formules dont la vérité se calcule est exactement ce que fait un compilateur, un vérificateur ou un solveur SAT, et c'est un savoir-faire qu'il vaut mieux acquérir sur des énigmes avant qu'il ne compte vraiment.

\begin{refs}
\item Contexte : Chevaliers et valets --- Wikipédia (en anglais) (\url{https://en.wikipedia.org/wiki/Knights_and_Knaves})
\item Contexte : Calcul des propositions --- Wikipédia (\url{https://fr.wikipedia.org/wiki/Calcul_des_propositions})
\item Lecture : Smullyan, \textit{What Is the Name of This Book?}
\item Cours : MIT OCW 6.042J Mathematics for Computer Science (\url{https://ocw.mit.edu/courses/6-042j-mathematics-for-computer-science-spring-2015/})
\end{refs}

\bigskip

\begin{problem}[{Combien de comparaisons pour trouver le plus grand ? --- Lycée, short}]
\textbf{DCS-29.} On vous donne $ n $ nombres distincts et vous ne pouvez
comparer que deux à la fois. Démontrez que $ n - 1 $ comparaisons suffisent pour trouver le
plus grand, et --- moitié plus difficile --- démontrez qu'aucune méthode ne peut toujours y parvenir en moins.
\end{problem}

\noindent La borne supérieure est une boucle que n'importe qui sait écrire ; la borne
inférieure est un véritable argument, et la voie habituelle consiste à voir chaque comparaison
comme éliminant exactement un candidat, de sorte qu'il faut $ n-1 $ éliminations pour n'en
laisser qu'un debout. C'est le plus petit exemple honnête de borne inférieure algorithmique --- un
énoncé qui ne porte pas sur une méthode mais sur toutes les méthodes possibles --- et cette
distinction est toute la théorie de la complexité en miniature. Suite naturelle : trouver le plus
grand \textit{et} le deuxième plus grand demande environ $ n + \log_2 n $ comparaisons, et non $ 2n $.

\begin{refs}
\item Contexte : Algorithme de sélection --- Wikipédia (\url{https://fr.wikipedia.org/wiki/Algorithme_de_s%C3%A9lection})
\item Lecture : Cormen, Leiserson, Rivest \& Stein, \textit{Introduction to Algorithms}, ch. 9
\item Voir aussi : Combinatoire et théorie des graphes (\url{pure-combinatorics.html}) pour le lemme des poignées de main
\end{refs}

\bigskip

\begin{problem}[{Un PGCD sans division --- Lycée, short}]
\textbf{DCS-30.} Pour des entiers strictement positifs $a, b$, démontrez les trois identités
\[ \gcd(2a, 2b) = 2\gcd(a,b), \qquad \gcd(2a, b) = \gcd(a,b) \ \text{ pour } b \text{ impair}, \qquad \gcd(a,b) = \gcd(a-b, b) \ \text{ pour } a > b. \]
Déduisez-en que ces seules règles calculent n'importe quel PGCD à l'aide de rien d'autre que des comparaisons, des soustractions et des divisions par deux, et démontrez que le procédé se termine toujours.
\end{problem}

\noindent Ces trois règles constituent l'algorithme binaire du PGCD, publié sous sa forme moderne par le physicien Josef Stein en 1967, mais déjà reconnaissable dans les \textit{Neuf Chapitres sur l'art mathématique} de la Chine des Han : \og{} si l'on peut, prendre la moitié ; sinon, poser le dénominateur et le numérateur, retrancher le plus petit du plus grand, et faire cela alternativement jusqu'à les rendre égaux. \fg{} La division coûte cher en matériel et la division par deux est gratuite en binaire : cette recette de simplification de fractions vieille de deux mille ans reste donc compétitive face à l'algorithme d'Euclide (DCS-11) au c\oe{}ur des bibliothèques arithmétiques réelles.

\begin{refs}
\item Contexte : Algorithme binaire de calcul du PGCD --- Wikipédia (\url{https://fr.wikipedia.org/wiki/Algorithme_binaire_de_calcul_du_PGCD})
\item Lecture : Knuth, \textit{The Art of Computer Programming}, vol. 2, \S{}4.5.2
\end{refs}

\bigskip

\begin{problem}[{Le vote majoritaire en un seul passage --- Lycée}]
\textbf{DCS-31.} Appelons \textit{élément majoritaire} de la liste $a_1, \dots, a_n$ une valeur qui y apparaît plus de $n/2$ fois. Parcourez la liste en ne conservant qu'un candidat $c$ et un compteur $k$, en partant de $k = 0$ : à la lecture de $a_i$, si $k = 0$ posez $c = a_i$ et $k = 1$ ; sinon, si $a_i = c$ augmentez $k$ de un ; sinon diminuez $k$ de un.
\begin{enumerate}
\item[(a)] Déroulez le parcours à la main sur $3,3,4,2,4,4,2,4,4$ puis sur $1,2,3,4,5$, en notant $(c,k)$ après chaque étape.
\item[(b)] Démontrez l'invariant qui porte l'algorithme : après lecture des $i$ premières entrées, celles-ci se décomposent en $k$ copies de $c$ et en paires d'entrées dont les deux membres diffèrent.
\item[(c)] Déduisez-en que si la liste possède un élément majoritaire, le $c$ final est cet élément.
\item[(d)] Montrez par un exemple que le $c$ final peut ne pas être majoritaire lorsqu'aucun élément ne l'est, de sorte qu'un second passage comptant les occurrences de $c$ est vraiment nécessaire. Concluez que la majorité se décide en temps $O(n)$ avec $O(1)$ mémoire au-delà de l'entrée.
\end{enumerate}

\end{problem}

\noindent Robert Boyer et J Strother Moore ont trouvé cela en 1980 au SRI International en réfléchissant aux systèmes tolérants aux pannes, où un \og{} vote \fg{} entre calculs redondants doit se tenir vite et presque sans mémoire ; ils l'ont fait circuler comme rapport technique et ne l'ont publié qu'en 1991, dans un volume en l'honneur de Woody Bledsoe. Le même tandem avait déjà donné au monde l'algorithme de recherche de chaîne de Boyer--Moore et le démonstrateur automatique de Boyer--Moore, et ils ont traité cet algorithme comme un cas d'essai pour la vérification mécanique --- destin bien choisi, car sa correction reste entièrement invisible tant qu'on n'a pas trouvé l'invariant de la question (b), et devient parfaitement évidente ensuite. Misra et Gries l'ont généralisé en 1982 pour trouver tous les éléments apparaissant plus de $n/k$ fois.

\begin{refs}
\item Contexte : Algorithme de vote majoritaire de Boyer--Moore --- Wikipédia (en anglais) (\url{https://en.wikipedia.org/wiki/Boyer%E2%80%93Moore_majority_vote_algorithm})
\item Article : Boyer \& Moore, \og{} MJRTY --- A Fast Majority Vote Algorithm \fg{}, dans \textit{Automated Reasoning: Essays in Honor of Woody Bledsoe} (1991)
\item Article : Misra \& Gries, \og{} Finding repeated elements \fg{}, Science of Computer Programming 2 (1982)
\end{refs}

\bigskip

\begin{problem}[{La location de skis : décider sans connaître l'avenir --- Lycée}]
\textbf{DCS-41.} Vous skiez pendant un nombre de jours que vous ne connaîtrez qu'à la fin de la saison, laquelle s'achèvera sans prévenir. Louer des skis coûte 1 par jour ; les acheter coûte $ b $ une fois pour toutes, après quoi skier est gratuit. Chaque matin vous devez décider, en sachant seulement combien de jours vous avez déjà skié. Appelons une stratégie $ c $-\textit{compétitive} si, sur toute saison possible, son coût total vaut au plus $ c $ fois le coût de la meilleure stratégie qui aurait connu d'avance la durée de la saison.
\begin{enumerate}
\item[(a)] Montrez que la meilleure stratégie sachant que la saison dure $ n $ jours paie $ \min(n, b) $.
\item[(b)] Démontrez que la stratégie \og{} louer les jours $ 1 $ à $ b-1 $, acheter le jour $ b $ \fg{} est $ \left(2 - \frac{1}{b}\right) $-compétitive.
\item[(c)] Démontrez qu'aucune stratégie déterministe n'atteint un rapport inférieur à $ 2 - \frac{1}{b} $ : quel que soit le jour où le skieur a décidé d'acheter, un adversaire qui met fin à la saison au bon moment lui impose ce rapport.
\item[(d)] Expliquez en un paragraphe pourquoi (c) est un énoncé sur toute stratégie concevable et non sur celle de (b), et décrivez ce qui change si la durée de la saison est au contraire tirée selon une loi de probabilité que vous connaissez.
\end{enumerate}

\end{problem}

\noindent Ce petit dilemme est entré en informatique par le matériel : l'article de 1988 de Karlin, Manasse, Rudolph et Sleator sur le \og{} competitive snoopy caching \fg{} traitait d'un cache multiprocesseur devant choisir entre payer indéfiniment un petit coût par accès distant ou payer une fois pour prendre possession d'un bloc mémoire ; et la notion de rapport de compétitivité --- qui mesure un algorithme en ligne contre un algorithme hors ligne omniscient, introduite par Sleator et Tarjan en 1985 pour la pagination --- est devenue la façon standard de donner un sens au pire cas des décisions prises sans l'avenir. La version aléatoire, où le skieur lance des pièces que l'adversaire ne voit pas, atteint le rapport strictement meilleur $ e/(e-1) \approx 1{,}58 $ ; cet écart est l'une des démonstrations les plus nettes qui soient de ce que le hasard aide véritablement contre un adversaire. Le même calcul est refait chaque jour par quiconque arbitre entre instances cloud à la demande et instances réservées, et par la pile TCP décidant combien de temps retarder un accusé de réception.

\begin{refs}
\item Contexte : Problème de la location de skis --- Wikipédia (\url{https://fr.wikipedia.org/wiki/Probl%C3%A8me_de_la_location_de_skis})
\item Contexte : Analyse compétitive (algorithme en ligne) --- Wikipédia (en anglais) (\url{https://en.wikipedia.org/wiki/Competitive_analysis_(online_algorithm)})
\item Article : Karlin, Manasse, Rudolph \& Sleator, \og{} Competitive snoopy caching \fg{}, Algorithmica 3 (1988)
\item Lecture : Borodin \& El-Yaniv, \textit{Online Computation and Competitive Analysis}, ch. 1
\end{refs}

\bigskip

\begin{problem}[{Les rotations se cachent dans une chaîne doublée --- Lycée, short}]
\textbf{DCS-42.} Soient $ x $ et $ y $ deux chaînes de même longueur $ n $ sur un alphabet quelconque. Démontrez que $ x $ est une rotation circulaire de $ y $ si et seulement si $ x $ apparaît comme facteur contigu de la concaténation $ yy $, et déduisez-en que la rotation se teste en temps $ O(n) $ dès qu'on dispose d'une recherche de facteur en temps linéaire.
\end{problem}

\noindent Un sens est un calcul d'indices qui tient en une ligne ; l'autre demande du soin, et le faire honnêtement, c'est dire exactement quelles positions d'occurrence dans $ yy $ sont possibles et pourquoi. Au-delà de l'astuce, la leçon est celle de la réduction : plutôt que d'inventer un algorithme pour un problème nouveau, on transforme l'entrée de sorte qu'un algorithme auquel on fait déjà confiance --- ici la recherche de Knuth--Morris--Pratt de DCS-44 --- y réponde à notre place. C'est le même geste, à une échelle infiniment plus petite, qui transforme un problème NP-complet en mille autres dans DCS-20.

\begin{refs}
\item Contexte : Algorithme de recherche de sous-chaîne --- Wikipédia (\url{https://fr.wikipedia.org/wiki/Algorithme_de_recherche_de_sous-cha%C3%AEne})
\item Contexte : Algorithme de Knuth--Morris--Pratt --- Wikipédia (\url{https://fr.wikipedia.org/wiki/Algorithme_de_Knuth-Morris-Pratt})
\end{refs}

\bigskip


\section*{Cryptographie et théorie de l'information}

\begin{problem}[{Cassez ce chiffre de César --- Lycée}]
\textbf{CRY-01.} Le message suivant a été chiffré par un chiffre
de César : chaque lettre du texte clair, rédigé en anglais, a été décalée vers l'avant dans
l'alphabet d'une même quantité secrète (Z revenant à A), les espaces restant intacts.
 \texttt{DRO ZBSCYXOBC WSCDKUO DRO CRKNYGC YX DRO GKVV PYB DRO DRSXQC
DROWCOVFOC}
\begin{enumerate}
\item[(a)] Retrouvez le texte clair et le décalage \textit{sans essayer aveuglément les 25 décalages} :
comptez les fréquences des lettres dans le texte chiffré, comparez-les aux fréquences
connues de l'anglais (E, T, A, O en sont les lettres les plus courantes) et justifiez votre
identification.
\item[(b)] Expliquez en général pourquoi tout chiffre qui remplace toujours une même lettre claire
par une même lettre chiffrée doit laisser fuir le profil de fréquences du texte clair, si
compliquée que soit la règle de remplacement.
\end{enumerate}

\end{problem}

\noindent Suétone rapporte que Jules César utilisait la version à décalage de trois
de ce chiffre pour sa correspondance militaire. L'attaque systématique --- l'analyse
fréquentielle --- a d'abord été décrite par le philosophe arabe al-Kind\={\i} dans le Bagdad du
neuvième siècle, dans un manuscrit sur le déchiffrement des messages cryptographiques ;
c'est sans doute la naissance de la cryptanalyse, et elle a tué la substitution simple pour
toujours. Toute définition moderne de sécurité est, au fond, la promesse que ce genre de
fuite statistique ne peut pas se produire.

\begin{refs}
\item Contexte : Chiffrement par décalage (chiffre de César) --- Wikipédia (\url{https://fr.wikipedia.org/wiki/Chiffrement_par_d%C3%A9calage})
\item Contexte : Analyse fréquentielle --- Wikipédia (\url{https://fr.wikipedia.org/wiki/Analyse_fr%C3%A9quentielle})
\item Lecture : Simon Singh, \textit{The Code Book}, ch. 1
\item Lecture : Katz \& Lindell, \textit{Introduction to Modern Cryptography}, ch. 1
\end{refs}

\bigskip

\begin{problem}[{Échauffements d'arithmétique modulaire --- Lycée}]
\textbf{CRY-02.} Écrivons $a \equiv b \pmod{n}$ pour signifier que
$n$ divise $a-b$.
\begin{enumerate}
\item[(a)] Démontrez que les congruences s'additionnent et se multiplient : si
$a \equiv b$ et $c \equiv d \pmod{n}$, alors $a+c \equiv b+d$ et
$ac \equiv bd \pmod{n}$.
\item[(b)] Démontrez qu'un carré parfait est toujours $\equiv 0$
ou $1 \pmod 4$, et déduisez-en qu'aucun entier de la forme $4k+3$ n'est somme de deux
carrés.
\item[(c)] Trouvez, avec démonstration, les deux derniers chiffres de $3^{2026}$. (Indice sur la
méthode, non sur la réponse : les puissances de 3 modulo 100 doivent finir par se répéter ---
pourquoi ?)
\end{enumerate}

\end{problem}

\noindent La notation des congruences est due à Gauss, qui ouvrit avec elle ses
\textit{Disquisitiones Arithmeticae} (1801) et transforma l'\og{} arithmétique de l'horloge \fg{} en
un calcul. Toute la cryptographie à clé publique --- Diffie--Hellman, RSA, courbes elliptiques ---
est du calcul dans exactement cette arithmétique : ces exercices de doigté sont donc le droit
d'entrée de tout ce qui suit. La théorie des nombres qui les sous-tend est développée plus
avant sur la page Théorie des nombres (\url{pure-number-theory.html}).

\begin{refs}
\item Contexte : Arithmétique modulaire --- Wikipédia (\url{https://fr.wikipedia.org/wiki/Arithm%C3%A9tique_modulaire})
\item Lecture : Hoffstein, Pipher \& Silverman, \textit{An Introduction to Mathematical Cryptography}, ch. 1
\end{refs}

\bigskip

\begin{problem}[{Quels chiffres affines sont vraiment des chiffres ? --- Lycée}]
\textbf{CRY-03.} Numérotons les lettres $A=0,\dots,Z=25$. Un
\textit{chiffre affine} de clé $(a,b)$ chiffre la lettre $x$ en
$ax+b \pmod{26}$.
\begin{enumerate}
\item[(a)] Montrez que cette règle peut être défaite (de sorte que le déchiffrement soit
possible) exactement lorsque $\gcd(a,26)=1$.
\item[(b)] Montrez concrètement ce qui se gâte pour
$a=13$ : trouvez deux lettres différentes qui se chiffrent en la même lettre chiffrée.
\item[(c)] Comptez le nombre de clés $(a,b)$ valides.
\item[(d)] Expliquez pourquoi, malgré ses centaines
de clés au lieu de 25, le chiffre affine tombe exactement sous la même attaque par analyse
fréquentielle que le chiffre de César.
\end{enumerate}

\end{problem}

\noindent Le chiffre affine est le plus petit exemple d'un thème qui traverse toute
la cryptographie : une application de chiffrement doit être une bijection, et savoir si une
formule donne une bijection est une question de théorie des nombres (ici, l'inversibilité
modulo 26). La même condition de PGCD, portée à des modules de 2048 bits, décide quels
exposants sont des clés RSA légales.

\begin{refs}
\item Contexte : Chiffre affine --- Wikipédia (\url{https://fr.wikipedia.org/wiki/Chiffre_affine})
\item Lecture : Hoffstein, Pipher \& Silverman, \textit{An Introduction to Mathematical Cryptography}, ch. 1
\end{refs}

\bigskip

\begin{problem}[{Un espace de clés immense n'est pas la sécurité --- Lycée}]
\textbf{CRY-04.} Un chiffre par substitution général remplace chaque
lettre par une autre selon une permutation arbitraire de l'alphabet de 26 lettres : il y a
donc $26!$ clés.
\begin{enumerate}
\item[(a)] Démontrez que $26! > 10^{26}$.
\item[(b)] Montrez qu'un ordinateur
essayant un milliard ($10^9$) de clés par seconde mettrait plus de $10^9$ années
à les essayer toutes.
\item[(c)] Et pourtant, les chiffres par substitution appliqués à un texte
anglais de quelques centaines de lettres se cassent couramment à la main. Réconciliez cela
avec (b) : dites précisément quelle propriété du chiffre l'attaque par analyse fréquentielle
exploite, et formulez une leçon de la forme \og{} un grand espace de clés est une condition
\_\_\_\_\_\_ de sécurité, mais non une condition \_\_\_\_\_\_ \fg{}.
\end{enumerate}

\end{problem}

\noindent C'est la plus ancienne leçon chèrement acquise du domaine : la recherche
exhaustive n'est pas la seule attaque. Edgar Allan Poe a mis en scène le cassage à la main
de chiffres par substitution dans \og{} Le Scarabée d'or \fg{} (1843), et al-Kind\={\i} possédait la
méthode un millénaire plus tôt. La cryptographie moderne répond par des définitions : un
schéma doit résister à \textit{toute} attaque efficace, et non seulement à l'attaque évidente ---
raison pour laquelle ce sont les démonstrations de sécurité, et non les décomptes de clés,
qui font la monnaie du domaine.

\begin{refs}
\item Contexte : Chiffrement par substitution --- Wikipédia (\url{https://fr.wikipedia.org/wiki/Chiffrement_par_substitution})
\item Contexte : Attaque par force brute --- Wikipédia (\url{https://fr.wikipedia.org/wiki/Attaque_par_force_brute})
\item Lecture : Katz \& Lindell, \textit{Introduction to Modern Cryptography}, ch. 1
\end{refs}

\bigskip

\begin{problem}[{Le chiffre de Vigenère et l'examen de Kasiski --- Lycée}]
\textbf{CRY-24.} Le chiffre de Vigenère fixe un mot-clé, disons \texttt{LEMON}, et
chiffre un message en décalant ses 1re, 2e, 3e, \dots{} lettres
par les 1re, 2e, 3e, \dots{} lettres du mot-clé, en répétant
celui-ci autant de fois qu'il le faut (avec $A = 0, \dots, Z = 25$ comme quantités de
décalage).
\begin{enumerate}
\item[(a)] Expliquez précisément pourquoi l'analyse fréquentielle à un seul comptage de CRY-01
échoue contre ce chiffre : une même lettre claire n'est plus toujours chiffrée par la même
lettre chiffrée.
\item[(b)] Montrez cependant que si le mot-clé a pour longueur $k$, alors les lettres du texte
chiffré aux positions $1, k+1, 2k+1, \dots$ sont toutes chiffrées par un seul et même
décalage de César --- et de même pour chacune des $k-1$ autres classes de positions.
\item[(c)] Démontrez l'observation qui fonde l'examen de Kasiski : si un bloc de lettres claires se
répète, et si les deux occurrences commencent à des positions dont la différence est un
multiple de $k$, alors le texte chiffré se répète aussi. Expliquez comment le recueil des
distances entre blocs chiffrés répétés et la prise de leurs diviseurs communs exposent la
longueur $k$ du mot-clé.
\item[(d)] Assemblez l'attaque complète : devinez $k$, scindez le texte chiffré en $k$ chiffres
de César, et cassez chacun par analyse fréquentielle. Quelle propriété de la langue du texte
clair l'attaque entière consomme-t-elle ?
\end{enumerate}

\end{problem}

\noindent L'idée polyalphabétique remonte à Alberti (1467) et à Bellaso
(1553) ; elle fut attribuée à tort à Blaise de Vigenère et régna trois siècles durant
comme \textit{le chiffre indéchiffrable}. Charles Babbage le cassa en privé vers 1854, et
Friedrich Kasiski publia la méthode en 1863, faisant de la cryptanalyse une science
systématique. La leçon a survécu à son époque : un chiffre à clé courte répétée n'est qu'un
entrelacement de nombreux chiffres faibles, et le masque jetable de CRY-05 est exactement la
limite dans laquelle le \og{} mot-clé \fg{} ne se répète jamais.

\begin{refs}
\item Contexte : Chiffre de Vigenère --- Wikipédia (\url{https://fr.wikipedia.org/wiki/Chiffre_de_Vigen%C3%A8re})
\item Contexte : Cryptanalyse du chiffre de Vigenère (examen de Kasiski) --- Wikipédia (\url{https://fr.wikipedia.org/wiki/Cryptanalyse_du_chiffre_de_Vigen%C3%A8re})
\item Lecture : Simon Singh, \textit{The Code Book}, ch. 2
\end{refs}

\bigskip

\begin{problem}[{Le double chiffrement n'apporte rien --- Lycée, short}]
\textbf{CRY-25.} Démontrez que chiffrer par un décalage de César de
$a$ puis chiffrer le résultat par un décalage de $b$ revient exactement au décalage de
César unique de $a + b \bmod 26$ --- de sorte que le double chiffrement de César (ou le
centuple) n'est pas plus fort que le simple. Concluez que les 26 décalages forment un groupe
pour la composition, et expliquez en une phrase pourquoi cette propriété de clôture est
précisément ce qui rend ici le chiffrement répété sans valeur.
\end{problem}

\noindent Une famille de chiffres close par composition s'appelle un groupe, et pour
de telles familles itérer le chiffrement n'achète rien. Savoir si DES avait ce défaut fut une
question ouverte sérieuse jusqu'à ce que Campbell et Wiener démontrent en 1992 que DES n'est
\textit{pas} un groupe --- raison pour laquelle Triple-DES l'a véritablement renforcé.

\begin{refs}
\item Contexte : Chiffrement par décalage (chiffre de César) --- Wikipédia (\url{https://fr.wikipedia.org/wiki/Chiffrement_par_d%C3%A9calage})
\item Contexte : Groupe (mathématiques) --- Wikipédia (\url{https://fr.wikipedia.org/wiki/Groupe_(math%C3%A9matiques)})
\end{refs}

\bigskip

\begin{problem}[{Quelle est la force d'une phrase secrète ? --- Lycée, short}]
\textbf{CRY-26.} On construit une phrase secrète en choisissant quatre
mots indépendamment et uniformément au hasard dans une liste publique de $7776 = 6^5$ mots.
Montrez que le nombre de phrases secrètes possibles vaut environ $3{,}7 \times 10^{15}$ ---
soit à peu près $4\log_2 7776 \approx 51{,}7$ bits d'entropie --- et comparez-le à un mot de
passe de huit lettres minuscules aléatoires ($26^8$ possibilités, environ
$37{,}6$ bits) : lequel est le plus difficile à deviner, et par quel facteur ?
\end{problem}

\noindent C'est le procédé Diceware (Arnold Reinhold, 1995) : la liste compte
$6^5$ mots afin que cinq jets de dés en sélectionnent un uniformément. L'idée, popularisée
par une bande dessinée d'xkcd, est celle de Shannon (CRY-11) déguisée --- ce qui compte est
l'entropie de la \textit{procédure de sélection}, non l'air compliqué du résultat, et le
caractère public de la procédure ne coûte rien.

\begin{refs}
\item Contexte : Diceware --- Wikipédia (\url{https://fr.wikipedia.org/wiki/Diceware})
\item Contexte : Robustesse d'un mot de passe --- Wikipédia (\url{https://fr.wikipedia.org/wiki/Robustesse_d%27un_mot_de_passe})
\end{refs}

\bigskip

\begin{problem}[{Clés de contrôle : attraper les erreurs que les gens font vraiment --- Lycée}]
\textbf{CRY-36.} Un ISBN-10 est une suite de chiffres $d_1 d_2 \dots d_{10}$, dont le
dernier peut aussi être le symbole X valant 10, vérifiant
\[ \sum_{i=1}^{10} i \, d_i \;\equiv\; 0 \pmod{11}. \]
Le dernier chiffre est choisi pour que cela soit vérifié.
\begin{enumerate}
\item[(a)] Démontrez que ce contrôle détecte toute erreur modifiant un seul chiffre, et toute erreur
échangeant deux chiffres distincts --- qu'ils soient adjacents ou non. Où exactement votre
démonstration utilise-t-elle que 11 est premier ?
\item[(b)] Le contrôle des cartes bancaires (Luhn) travaille modulo 10 : en lisant depuis la droite,
on double un chiffre sur deux, on retranche 9 de toute valeur doublée dépassant 9, et l'on
exige que le total soit $0 \bmod 10$. Montrez qu'il attrape toute erreur d'un chiffre et
toute transposition de chiffres \textit{adjacents} sauf une, et trouvez cette exception.
\item[(c)] Démontrez que l'exception est inévitable pour les schémas de cette forme : aucun contrôle
du type $\sum_i a_i d_i \equiv 0 \pmod{10}$, à poids entiers fixés $a_i$, ne peut
détecter à la fois toutes les erreurs d'un chiffre et toutes les transpositions adjacentes.
(Montrez que la détection des erreurs simples force tous les $a_i$ à être impairs, et voyez
ce que cela fait à $a_i - a_{i+1}$.)
\end{enumerate}

\end{problem}

\noindent En 1969 Jacobus Verhoeff fit une chose inhabituelle pour un mathématicien :
avant de concevoir un code, il collecta des données sur les fautes que commettent les êtres
humains en recopiant des nombres. Les chiffres isolément faux représentaient environ quatre
erreurs sur cinq, les transpositions adjacentes l'essentiel du reste --- ce sont donc les deux
classes qu'une clé de contrôle décimale doit attraper, et la question (c) montre que
l'arithmétique modulo 10 ne peut attraper les deux. La réponse de Verhoeff fut d'abandonner
tout à fait l'addition et de calculer dans le groupe diédral d'ordre 10, dont la
non-commutativité fournit exactement ce qui manque à l'addition modulaire. L'autre solution,
celle des libraires, fut de changer de module : l'ISBN-10 a pris 11, au prix d'un onzième
symbole X. Lorsque le commerce du livre a fusionné avec le système mondial de codes-barres en
2007, l'ISBN-13 a adopté les poids 1 et 3 modulo 10 et a discrètement abandonné la garantie
sur les transpositions.

\begin{refs}
\item Contexte : Clé de contrôle --- Wikipédia (\url{https://fr.wikipedia.org/wiki/Cl%C3%A9_de_contr%C3%B4le})
\item Contexte : Formule de Luhn --- Wikipédia (\url{https://fr.wikipedia.org/wiki/Formule_de_Luhn})
\item Contexte : International Standard Book Number --- Wikipédia (\url{https://fr.wikipedia.org/wiki/International_Standard_Book_Number})
\item Lecture : J. Verhoeff, \textit{Error Detecting Decimal Codes}, Mathematical Centre Tracts 29 (Amsterdam, 1969)
\end{refs}

\bigskip

\begin{problem}[{Arbres de Merkle : un reçu de la taille d'un logarithme --- Lycée}]
\textbf{CRY-41.} Un journal public contient $n = 2^k$ enregistrements $x_1, \dots, x_n$.
Hachez chaque enregistrement pour obtenir les feuilles d'un arbre binaire, puis hachez à
répétition les paires voisines --- en remplaçant deux valeurs $u, v$ par $h(u \,\|\, v)$,
où $\|$ désigne la concaténation --- jusqu'à ce qu'il ne reste qu'une valeur : la
\textit{racine}. Supposez seulement que la fonction de hachage $h$ est \textit{résistante aux
collisions}, au sens où personne ne peut exhiber deux entrées différentes de même
sortie.
\begin{enumerate}
\item[(a)] Démontrez qu'il y a exactement $k = \log_2 n$ niveaux au-dessus des feuilles, et
comptez le nombre total de calculs de hachage nécessaires pour construire l'arbre entier.
Remarquez que la racine publiée est une unique valeur de hachage, si grand que soit $n$.
\item[(b)] Une \textit{preuve d'inclusion} pour l'enregistrement $x_i$ est la liste des valeurs
s\oe{}urs rencontrées sur le chemin de la feuille $i$ jusqu'à la racine. Montrez que cette
liste compte exactement $k$ entrées, et décrivez précisément le calcul par lequel un
vérificateur qui ne connaît que la racine et $x_i$ la contrôle.
\item[(c)] Démontrez la propriété de sécurité : si un exploitant de journal malhonnête convainc le
vérificateur qu'un certain $y \ne x_i$ occupe la position $i$ d'un arbre de \textit{même}
racine, alors quelque part le long des deux chemins il existe deux entrées différentes de
$h$ ayant la même sortie, contredisant la résistance aux collisions. (Comparez les deux
chemins niveau par niveau et considérez le niveau le plus bas où les valeurs coïncident.)
\item[(d)] Des chiffres, puis une subtilité. Avec $n = 2^{30}$ enregistrements et des valeurs de
hachage de 32 octets, calculez la taille en octets d'une preuve d'inclusion et comparez-la au
coût de l'envoi du journal entier. Expliquez ensuite ce qui se gâte si les feuilles et les
n\oe{}uds internes sont hachés par exactement la même règle : montrez qu'un attaquant pourrait
présenter un n\oe{}ud interne comme s'il s'agissait d'un enregistrement, et dites ce qu'il faut
intégrer à l'entrée de $h$ pour l'empêcher.
\end{enumerate}

\end{problem}

\noindent Ralph Merkle a décrit les arbres de hachage dans sa thèse de doctorat à
Stanford en 1979 et les a brevetés la même année, comme moyen de signer de nombreux messages
avec une seule valeur publique ; l'idée a dormi jusqu'à ce que le monde des systèmes découvre
qu'elle résout le problème de prouver l'appartenance à une énorme collection à quelqu'un qui
n'est pas disposé à la télécharger. Git nomme chaque commit par un arbre de hachages, Bitcoin
résume les transactions de chaque bloc par une racine de Merkle, et Certificate Transparency ---
créé après que la compromission en 2011 de l'autorité de certification DigiNotar fut passée
inaperçue des semaines durant --- publie chaque certificat TLS émis dans un journal de Merkle en
ajout seul, de sorte qu'un navigateur peut exiger une preuve au lieu de se fier à une
promesse. La confusion décrite en (d) n'a rien d'hypothétique : la norme (RFC 6962) préfixe
les entrées de feuille par un octet nul et les entrées internes par un octet égal à un,
exactement pour cette raison.

\begin{refs}
\item Contexte : Arbre de Merkle --- Wikipédia (\url{https://fr.wikipedia.org/wiki/Arbre_de_Merkle})
\item Contexte : Certificate Transparency --- Wikipédia (en anglais) (\url{https://en.wikipedia.org/wiki/Certificate_Transparency})
\item Contexte : Fonction de hachage cryptographique --- Wikipédia (\url{https://fr.wikipedia.org/wiki/Fonction_de_hachage_cryptographique})
\item Lecture : Katz \& Lindell, \textit{Introduction to Modern Cryptography}, ch. 5
\end{refs}

\bigskip

\begin{problem}[{Pourquoi le mode ECB dessine le manchot --- Lycée, short}]
\textbf{CRY-42.} Un chiffrement par bloc chiffre des blocs de 16 octets
sous une clé $K$ par une bijection $E_K$ de l'ensemble de tous les blocs de 16 octets ; le
mode \textit{dictionnaire de codes électronique} (ECB) chiffre un long message en le découpant
en blocs $P_1, P_2, \dots$ et en envoyant $E_K(P_1), E_K(P_2), \dots$. Démontrez que
$P_i = P_j$ si et seulement si les $i$-ième et $j$-ième blocs chiffrés sont égaux --- pour
tout chiffrement par bloc, toute clé, et si robuste que soit $E_K$ --- et décrivez exactement
ce qu'une oreille indiscrète apprend, dès lors, du chiffrement ECB d'une image comportant de
larges plages d'une seule couleur.
\end{problem}

\noindent La démonstration standard est une image de Tux, le manchot de Linux,
chiffrée en mode ECB : les couleurs changent et le manchot reste parfaitement visible. ECB
fut l'un des quatre modes publiés avec DES en 1980 et il a survécu dans les logiciels depuis
lors, ce qui explique que cette image a probablement empêché plus de brèches réelles
qu'aucun théorème. L'acte d'accusation formel arrive plus loin sur cette page : un schéma ECB
est déterministe, et CRY-38(a) montre qu'aucun schéma déterministe ne peut satisfaire la
définition moderne de la sécurité, si incassable que soit le chiffrement sous-jacent.

\begin{refs}
\item Contexte : Mode d'opération (cryptographie) --- Wikipédia (\url{https://fr.wikipedia.org/wiki/Mode_d%27op%C3%A9ration_(cryptographie)})
\item Contexte : Chiffrement par bloc --- Wikipédia (\url{https://fr.wikipedia.org/wiki/Chiffrement_par_bloc})
\end{refs}

\bigskip


\section*{Physique mathématique}

\begin{problem}[{Les nombres impairs de Galilée --- Lycée}]
\textbf{PHY-30.} Un corps part du repos et tombe avec une accélération constante $ g $.
\begin{enumerate}
\item[(a)] Démontrez, à partir de la seule définition de l'accélération constante, que la distance
parcourue pendant le temps $ t $ vaut $ s = \tfrac{1}{2} g t^{2} $. Ne citez pas la formule ---
établissez-la, soit en moyennant la vitesse, soit en calculant l'aire sous un graphe
vitesse--temps, et dites quelle propriété de l'accélération constante votre argument utilise.
\item[(b)] Découpez la chute en intervalles de temps successifs égaux. Démontrez que les distances
parcourues pendant ces intervalles sont dans le rapport $ 1 : 3 : 5 : 7 : \cdots $, celui des
nombres impairs.
\item[(c)] Déduisez-en l'identité $ 1 + 3 + 5 + \cdots + (2n-1) = n^{2} $, et expliquez pourquoi
l'apparition d'un fait sur les carrés parfaits dans un énoncé sur la chute des corps n'est pas
une coïncidence.
\end{enumerate}

\end{problem}

\noindent Galilée ne pouvait pas mesurer le temps avec assez de précision pour tester
$ s \propto t^{2} $ directement ; il faisait donc rouler des billes sur des plans inclinés pour
ralentir le mouvement et les chronométrait en pesant l'eau s'écoulant d'un vase. Ce qu'il a
trouvé, c'est la règle des nombres impairs de la question (b) --- un motif dans des rapports, qui
ne demande aucun étalonnage d'horloge --- et c'est de là qu'il a inféré la loi du carré. Voilà une
leçon de méthode expérimentale qui vaut autant que la physique : il a remplacé une grandeur qu'il
ne savait pas mesurer par un rapport qu'il savait mesurer. La question (c) est la récompense : un
morceau d'arithmétique grecque antique qui tombe d'une expérience du dix-septième siècle.

\begin{refs}
\item Contexte : Équation du mouvement --- Wikipédia (\url{https://fr.wikipedia.org/wiki/%C3%89quation_du_mouvement})
\item Contexte : Galilée --- Wikipédia (\url{https://fr.wikipedia.org/wiki/Galil%C3%A9e_(savant)})
\item Lecture : Galileo Galilei, \textit{Discorsi e dimostrazioni matematiche intorno a due nuove scienze} (1638), Troisième Journée
\item Voir aussi : L'art de la démonstration (\url{proofs.html}) pour la somme des nombres impairs par récurrence
\end{refs}

\bigskip

\begin{problem}[{Le boulet de canon de Newton --- Lycée, short}]
\textbf{PHY-31.} Un boulet de canon est tiré horizontalement du sommet d'une
montagne, assez vite pour faire le tour de la Terre au ras du sol. En n'utilisant que
$ F = ma $, l'accélération centripète $ v^{2}/r $ et $ g \approx 9{,}8\ \mathrm{m\,s^{-2}} $,
établissez la vitesse nécessaire et la durée d'une orbite, et énoncez clairement quelle hypothèse
physique vous autorise à traiter la chute vers la Terre et l'orbite comme un seul et même
mouvement.
\end{problem}

\noindent Newton a dessiné ce canon dans \textit{A Treatise of the System of the World},
et l'image porte à elle seule toute l'idée de la gravitation universelle : un corps en orbite
n'est ni en apesanteur ni dispensé de la gravité, il tombe simplement et manque sa cible. Obtenir
environ $ 7{,}9\ \mathrm{km\,s^{-1}} $ et à peu près quatre-vingt-dix minutes à partir de deux
lignes d'algèbre est l'un des meilleurs rapports effort/résultat de toute la physique.

\begin{refs}
\item Contexte : Le canon de Newton --- Wikipédia (\url{https://fr.wikipedia.org/wiki/Canon_de_Newton})
\item Contexte : Orbite circulaire --- Wikipédia (\url{https://fr.wikipedia.org/wiki/Orbite_circulaire})
\end{refs}

\bigskip

\begin{problem}[{La portée d'un projectile, et la parabole de sûreté --- Lycée}]
\textbf{PHY-32.} Un projectile est lancé depuis un sol horizontal avec la vitesse $ v $ sous
un angle $ \theta $ au-dessus de l'horizontale et se déplace sous une accélération constante
dirigée vers le bas $ g $, sans résistance de l'air.
\begin{enumerate}
\item[(a)] En traitant séparément les mouvements horizontal et vertical et en éliminant le temps entre
les deux, démontrez que la trajectoire est la parabole
\[ y = x\tan\theta \;-\; \frac{g\,x^{2}}{2v^{2}\cos^{2}\theta}. \]
\item[(b)] Démontrez que le projectile revient au sol à la distance horizontale
\[ R(\theta) = \frac{v^{2}\sin 2\theta}{g}, \]
et donc que la portée est maximale pour $ \theta = 45^{\circ} $. Donnez la démonstration
deux fois --- une fois en dérivant $ R $, une fois en raisonnant directement sur la plus grande
valeur de $ \sin 2\theta $ --- puis dites pourquoi c'est le même argument sous un autre
habillage.
\item[(c)] Démontrez que $ \theta $ et $ 90^{\circ} - \theta $ donnent exactement la même portée,
de sorte que toute distance inférieure au maximum est atteignable de deux façons précisément.
Déterminez lequel des deux tirs reste le plus longtemps en l'air, et dans quel rapport.
\item[(d)] Maintenant, fixez $ v $ et laissez $ \theta $ parcourir tous les angles de tir.
Démontrez qu'un point $ (x, y) $ avec $ x > 0 $ peut être atteint si et seulement si
\[ y \;\le\; \frac{v^{2}}{2g} \;-\; \frac{g\,x^{2}}{2v^{2}}, \]
de sorte que toute la région accessible est bornée par une seconde parabole. (Posez
$ u = \tan\theta $ dans la question (a), regardez le résultat comme une équation du second
degré en $ u $, et demandez-vous quand elle admet une solution réelle.) Vérifiez que cette
frontière rencontre le sol à la portée maximale de (b) et l'axe vertical à la hauteur d'un tir
vertical.
\end{enumerate}

\end{problem}

\noindent Tartaglia fut le premier à braquer les mathématiques sur l'artillerie, dans
\textit{Nova Scientia} (1537) ; il affirmait que $ 45^{\circ} $ donne le tir le plus long, et il
avait raison, bien que sa trajectoire --- une course rectiligne, puis un arc, puis une chute
verticale --- fût fausse. Galilée a fourni la vraie courbe dans les \textit{Discorsi} (1638) en
composant un mouvement horizontal uniforme avec la chute uniformément accélérée de PHY-30, et il
en a imprimé des tables de tir. La question (d) est de Torricelli, dans \textit{De motu gravium}
(achevé en 1641, imprimé dans l'\textit{Opera geometrica} de 1644) : la \textit{parabola di sicurezza}
de l'artilleur, hors de laquelle vous êtes à l'abri de la pièce quelle que soit son élévation, et
que toute trajectoire possible touche exactement une fois. C'est la première enveloppe d'une
famille de courbes jamais calculée, trouvée des décennies avant que le calcul différentiel n'ait
un mot pour dire la chose.

\begin{refs}
\item Contexte : Trajectoire d'un projectile --- Wikipédia (\url{https://fr.wikipedia.org/wiki/Trajectoire_d%27un_projectile})
\item Contexte : Parabole de sûreté --- Wikipédia (\url{https://fr.wikipedia.org/wiki/Parabole_de_s%C3%BBret%C3%A9})
\item Contexte : Niccolò Fontana Tartaglia --- Wikipédia (\url{https://fr.wikipedia.org/wiki/Niccol%C3%B2_Fontana_Tartaglia})
\item Lecture : Galileo Galilei, \textit{Discorsi} (1638), Quatrième Journée ; Kleppner et Kolenkow, \textit{An Introduction to Mechanics}, ch. 1
\end{refs}

\bigskip

\begin{problem}[{Archimède à partir de la pression sur une face --- Lycée}]
\textbf{PHY-33.} Un liquide de masse volumique uniforme $ \rho $ repose sous la gravité
$ g $, avec la pression $ p_{0} $ à sa surface libre. Vous ne pouvez utiliser que deux
faits : un liquide au repos pousse perpendiculairement sur toute surface qu'il touche, avec une
force égale à (pression) $\times$ (aire) ; et toute portion du liquide est elle-même en équilibre. Rien
sur la poussée d'Archimède ne peut être supposé --- c'est précisément ce qu'il faut démontrer.
\begin{enumerate}
\item[(a)] En équilibrant les forces sur une mince colonne verticale de liquide, démontrez que la
pression à la profondeur $ h $ vaut
\[ p(h) = p_{0} + \rho g h, \]
et qu'elle ne dépend que de la profondeur --- ni de la forme du récipient, ni de la quantité de
liquide qu'il contient.
\item[(b)] Déduisez-en le paradoxe hydrostatique de Stevin : la force dirigée vers le bas sur un fond
plat d'aire $ A $ à la profondeur $ h $ vaut $ (p_{0} + \rho g h)A $, ce qui, pour un
récipient qui s'évase vers le haut, peut valoir plusieurs fois le poids de tout le liquide qu'il
contient. Expliquez sans esquive d'où vient la force manquante ; une réponse correcte nomme les
parois et dit dans quel sens elles poussent.
\item[(c)] Immergez un bloc rectangulaire de hauteur $ H $ et de section horizontale $ A $, faces
horizontales et verticales. Additionnez les forces de pression sur les six faces et démontrez
que leur total vaut $ \rho g A H $, dirigé vers le haut --- le poids du liquide déplacé par le
bloc.
\item[(d)] Étendez (c) à un corps de forme quelconque par l'argument suivant, que vous devez rendre
étanche plutôt que simplement répéter. Les forces de pression sur la surface du corps ne
dépendent que de la forme de cette surface et de sa profondeur, et nullement de ce qui se trouve
à l'intérieur ; imaginez donc le corps retiré et sa place remplie de liquide, et considérez
l'équilibre de cette portion. Énoncez et démontrez alors la loi de la flottaison --- un corps
flottant déplace son propre poids de liquide --- et calculez la fraction d'un iceberg située sous
la ligne de flottaison, en prenant $ \rho_{\text{glace}} = 917 $ et
$ \rho_{\text{eau de mer}} = 1025\ \mathrm{kg\,m^{-3}} $.
\end{enumerate}

\end{problem}

\noindent Archimède a démontré la loi de la flottaison dans \textit{Des corps flottants}
(vers 250 av. J.-C.), livre I, à partir d'un postulat sur l'équilibre des parties d'un fluide ;
l'argument de la question (d) est le sien, vieux de vingt-deux siècles et toujours le plus net
que quiconque ait trouvé. L'histoire du bain et de la couronne vient de Vitruve, qui écrit deux
cents ans plus tard, et ce n'est certainement pas ainsi que l'orfèvre a été confondu : mesurer le
débordement provoqué par une couronne avec la précision qu'exigerait la fraude est sans espoir.
Galilée l'a fait remarquer dans \textit{La Bilancetta} (1586), en soutenant qu'Archimède avait dû
peser la couronne immergée et se servir de la poussée elle-même comme instrument. Le
\textit{De Beghinselen des Waterwichts} de Stevin, de la même année, a fourni la question (b) --- la
première avancée réelle sur Archimède en dix-huit cents ans, et encore aujourd'hui le résultat
qui surprend le plus sûrement ceux qui croient savoir ce qu'est la pression.

\begin{refs}
\item Contexte : Poussée d'Archimède --- Wikipédia (\url{https://fr.wikipedia.org/wiki/Pouss%C3%A9e_d%27Archim%C3%A8de})
\item Contexte : Des corps flottants --- Wikipédia (\url{https://fr.wikipedia.org/wiki/Des_corps_flottants})
\item Contexte : Variation de pression verticale (paradoxe hydrostatique) --- Wikipédia (\url{https://fr.wikipedia.org/wiki/Variation_de_pression_verticale})
\item Lecture : Feynman, \textit{The Feynman Lectures on Physics (\url{https://www.feynmanlectures.caltech.edu/II_40.html}), vol. II, ch. 40 (hydrostatique, avant l'écoulement)}
\end{refs}

\bigskip

\begin{problem}[{Pourquoi le Soleil soulève la plus petite marée --- Lycée, short}]
\textbf{PHY-34.} Une lune de masse $ M $ se trouve à la distance $ d $ du
centre d'une planète de rayon $ R \ll d $ ; en travaillant dans le référentiel en chute libre,
non tournant, du centre de la planète, démontrez que l'accélération gravitationnelle résiduelle
au point le plus proche de la lune vaut $ 2GMR/d^{3} $ dirigée vers la lune, qu'au point
opposé elle a la même intensité mais pointe en sens inverse, et donc que l'effet de marée décroît
comme le \textit{cube} inverse de la distance plutôt que comme le carré inverse. En prenant
$ M_{\text{Soleil}} = 1{,}99\times 10^{30}\ \mathrm{kg} $ à
$ 1{,}50\times 10^{11}\ \mathrm{m} $ et $ M_{\text{Lune}} = 7{,}35\times 10^{22}\
\mathrm{kg} $ à $ 3{,}84\times 10^{8}\ \mathrm{m} $, calculez le rapport de l'effet de marée
du Soleil à celui de la Lune, et expliquez pourquoi le Soleil --- dont l'attraction sur la Terre
est quelque 180 fois plus forte que celle de la Lune --- n'arrive néanmoins qu'en second au bord de
la mer.
\end{problem}

\noindent Newton a expliqué les marées au livre III des \textit{Principia} (1687) comme la
différence entre l'attraction de la Lune sur l'océan proche, sur la Terre solide, et sur l'océan
lointain, et il la comptait parmi ses meilleures preuves qu'une seule loi gouverne la pomme, la
Lune et la mer. Le cube inverse, c'est là que loge la surprise : ce qui soulève une marée n'est
pas l'intensité d'une attraction mais la vitesse à laquelle cette attraction varie sur la largeur
d'une planète, et ce seul fait tranche la question Soleil contre Lune, explique les deux
bourrelets plutôt qu'un seul, et donne les vives-eaux quand les deux s'alignent et les mortes-eaux
quand ils ne s'alignent pas. Laplace a remplacé cette image statique par une théorie dynamique
dans les années 1770, parce qu'un océan réel est un système forcé et résonant qui ballotte dans
des bassins de forme malcommode ; aucune table des marées n'a jamais été calculée à partir de la
seule formule ci-dessus, et aucune ne l'a jamais été sans elle.

\begin{refs}
\item Contexte : Force de marée --- Wikipédia (\url{https://fr.wikipedia.org/wiki/Force_de_mar%C3%A9e})
\item Contexte : Théorie des marées --- Wikipédia (en anglais) (\url{https://en.wikipedia.org/wiki/Theory_of_tides})
\item Lecture : Taylor, \textit{Classical Mechanics}, \S{}9.2 (The Tides)
\end{refs}

\bigskip

\begin{problem}[{L'effet Doppler, démontré deux fois --- Lycée}]
\textbf{PHY-35.} Une source émet un son de fréquence $ f $ dans de l'air immobile, où le son
se propage à la vitesse $ c $. Traitez les deux cas suivants comme deux problèmes distincts, et
ne supposez pas d'avance qu'ils s'accordent.
\begin{enumerate}
\item[(a)] L'auditeur est immobile et la source s'approche à la vitesse $ v_{s} < c $. En suivant
les lieux d'émission des fronts d'onde successifs, et donc leur écartement dans l'air, démontrez
que la fréquence entendue est
\[ f_{\text{source mobile}} = f\,\frac{c}{c - v_{s}}. \]
\item[(b)] La source est immobile et l'auditeur s'approche à la vitesse $ v_{o} $. En comptant combien
de fronts d'onde l'auditeur dépasse par seconde, démontrez que la fréquence entendue est
\[ f_{\text{auditeur mobile}} = f\,\frac{c + v_{o}}{c}. \]
\item[(c)] Démontrez que ces deux expressions ne sont \textit{pas} égales lorsque $ v_{s} = v_{o} = v $,
qu'elles coïncident au premier ordre en $ v/c $, et qu'elles diffèrent d'abord à l'ordre
$ v^{2}/c^{2} $. Dites ensuite quel fait physique ce désaccord signale : la réponse tient en un
mot, c'est la raison pour laquelle les deux cas se distinguent expérimentalement, et c'est ce qui
rend le son différent de la lumière.
\item[(d)] Écrivez la formule lorsque la source et l'auditeur sont tous deux en mouvement. Montrez que
lorsque $ v_{s} \ge c $ les fronts d'onde s'accumulent le long d'un cône, et trouvez son
demi-angle en fonction de $ c $ et $ v_{s} $.
\item[(e)] La lumière n'a pas de milieu, si bien que pour elle les deux cas \textit{ne peuvent pas}
différer. Vérifiez que la formule relativiste pour un rapprochement,
\[ f' = f\,\sqrt{\frac{c+v}{c-v}}, \]
est exactement la moyenne géométrique de vos réponses à (a) et (b) --- puis expliquez pourquoi
\og{} couper la poire en deux entre les deux réponses classiques \fg{} est un slogan commode et non une
démonstration.
\end{enumerate}

\end{problem}

\noindent Doppler a proposé l'effet en 1842 dans un mémoire sur la lumière colorée des
étoiles doubles, en soutenant que la couleur d'une étoile trahit son mouvement. L'idée était
juste et cette application était fausse --- les vitesses stellaires décalent bien trop peu la
lumière visible pour en changer la couleur --- ce qui rappelle utilement qu'un principe correct ne
garantit rien quant au premier usage que son découvreur en fait. Buys Ballot a réglé
correctement le cas acoustique en 1845 en plaçant des trompettistes sur un wagon découvert de la
ligne au départ d'Utrecht, avec des musiciens à l'oreille exercée postés le long de la voie : la
hauteur montait à l'approche et descendait à l'éloignement, exactement comme prévu, et Buys
Ballot n'en resta pas moins sceptique quant à l'astronomie. Fizeau parvint indépendamment à la
même conclusion pour la lumière en 1848, ce qui explique qu'en France l'effet s'appelle
Doppler--Fizeau. L'asymétrie entre les questions (a) et (b) n'est pas une tache à faire
disparaître : c'est un milieu qui annonce sa propre existence, et l'échec expérimental à trouver
une telle asymétrie pour la lumière est l'un des faits qui ont fini par imposer la transformation
de Lorentz de PHY-09.

\begin{refs}
\item Contexte : Effet Doppler --- Wikipédia (\url{https://fr.wikipedia.org/wiki/Effet_Doppler})
\item Contexte : Christoph Buys Ballot --- Wikipédia (\url{https://fr.wikipedia.org/wiki/Christoph_Buys_Ballot})
\item Contexte : Effet Doppler relativiste --- Wikipédia (\url{https://fr.wikipedia.org/wiki/Effet_Doppler_relativiste})
\item Lecture : A. P. French, \textit{Vibrations and Waves} (MIT Introductory Physics Series)
\end{refs}

\bigskip

\begin{problem}[{Pourquoi les harmoniques sont des nombres entiers --- Lycée, short}]
\textbf{PHY-36.} Une corde de longueur $ L $, fixée à ses deux extrémités,
porte des ondes transversales de vitesse $ c = \sqrt{T/\mu} $, où $ T $ est la tension et
$ \mu $ la masse par unité de longueur. Démontrez qu'une onde stationnaire sinusoïdale ne peut
y subsister que lorsqu'un nombre entier de demi-longueurs d'onde tient entre les extrémités, de
sorte que $ \lambda_{n} = 2L/n $ et $ f_{n} = nc/(2L) = n f_{1} $ pour
$ n = 1, 2, 3, \dots $, et déduisez-en les lois de Mersenne : la hauteur d'une corde varie en
raison inverse de sa longueur, comme la racine carrée de sa tension, et en raison inverse de la
racine carrée de sa masse par unité de longueur.
\end{problem}

\noindent Les pythagoriciens savaient que diviser une corde par deux élève sa hauteur
d'une octave et que les consonances répondent à de petits rapports entiers ; ce qu'ils ne
pouvaient pas dire, c'est pourquoi les nombres entiers auraient quoi que ce soit à y voir.
Mersenne a publié les lois quantitatives dans l'\textit{Harmonie universelle} (1636), après avoir
employé des cordes assez longues et assez lourdes pour compter les vibrations à l'\oe{}il, et Brook
Taylor a le premier tiré la fréquence fondamentale de la mécanique en 1713. La réponse au
\og{} pourquoi des entiers \fg{} est celle qui revient sans cesse en physique : une onde enfermée entre
deux frontières ne peut prendre que les formes que les frontières permettent, et les formes
viennent en nombres entiers parce qu'on les compte. C'est exactement l'argument de PHY-08(b) avec
une corde de violon à la place d'une particule quantique, et PHY-22(d) arrive aux mêmes modes par
le chemin inverse, comme limite d'un chapelet de billes.

\begin{refs}
\item Contexte : Onde sur une corde vibrante --- Wikipédia (\url{https://fr.wikipedia.org/wiki/Onde_sur_une_corde_vibrante})
\item Contexte : Lois de Mersenne --- Wikipédia (\url{https://fr.wikipedia.org/wiki/Lois_de_Mersenne})
\item Contexte : Onde stationnaire --- Wikipédia (\url{https://fr.wikipedia.org/wiki/Onde_stationnaire})
\item Lecture : Feynman, \textit{The Feynman Lectures on Physics (\url{https://www.feynmanlectures.caltech.edu/I_49.html}), vol. I, ch. 49 (Modes)}
\end{refs}

\bigskip

\begin{problem}[{Archimède et la loi du levier --- Lycée}]
\textbf{PHY-56.} Une barre rigide de poids négligeable repose sur un point d'appui, avec les poids
$ w_{1} $ et $ w_{2} $ suspendus aux distances $ d_{1} $ et $ d_{2} $ de part et d'autre. Vous
ne pouvez supposer que ce qu'Archimède supposait : des poids égaux à des distances égales
s'équilibrent ; des poids égaux à des distances inégales ne s'équilibrent pas, le plus éloigné
descendant ; et tout poids peut être remplacé par deux demi-poids suspendus symétriquement de part et
d'autre de sa place initiale sans rien déranger.
\begin{enumerate}
\item[(a)] Démontrez la loi du levier
\[ w_{1} d_{1} \;=\; w_{2} d_{2} \]
dans le cas où $ w_{1}/w_{2} $ est un rapport d'entiers. (Une voie : découpez les deux poids en
petits morceaux égaux, utilisez le troisième postulat pour répartir chaque tas symétriquement le long
de la barre de sorte que les deux répartitions se fondent en une seule rangée régulièrement espacée,
puis appliquez le premier postulat à cette rangée.)
\item[(b)] Étendez la loi à des poids dont le rapport est irrationnel. Énoncez précisément l'hypothèse
supplémentaire dont votre extension a besoin --- aucune subdivision finie n'atteindra un rapport
incommensurable, il faut donc ajouter quelque chose sur la continuité --- et dites où vous l'avez
utilisée.
\item[(c)] Démontrez maintenant la même loi une seconde fois, par l'énergie. Faites tourner la barre d'un
petit angle $ \theta $ autour du point d'appui, montrez que les deux extrémités se déplacent des
hauteurs $ d_{1}\theta $ et $ d_{2}\theta $ au premier ordre, et imposez que l'énergie
potentielle totale ne varie pas au premier ordre. Dites ensuite laquelle de vos deux démonstrations
use de plus de géométrie et laquelle de plus de physique, et à laquelle vous confieriez un levier
légèrement tordu.
\item[(d)] Lisez la première démonstration d'un \oe{}il critique, comme le fit Ernst Mach en 1883. Son objection
était que l'argument de symétrie suppose en catimini ce qu'il prétend démontrer : que l'effet de
rotation d'un poids dépend de sa distance, et en dépend linéairement plutôt que, disons,
quadratiquement. Énoncez l'objection dans vos propres termes, décidez si vous la jugez fatale, et
identifiez exactement ce qu'un traitement purement logique devrait ajouter comme axiome.
\item[(e)] \og{} Donnez-moi un point d'appui et je soulèverai la Terre. \fg{} Supposez que l'homme puisse pousser
avec 700 N, et prenez pour poids de la Terre sa masse $ 5{,}97\times10^{24}\ \mathrm{kg} $
multipliée par $ 9{,}8\ \mathrm{m\,s^{-2}} $. Calculez le rapport des longueurs de bras qu'il lui
faut, et la distance que sa propre extrémité de la barre doit parcourir pour élever la charge d'un
centimètre ; comparez cette distance au diamètre de la Voie lactée, environ $ 10^{5} $
années-lumière. Énoncez ensuite en une phrase ce qu'un levier conserve réellement, et ce qu'il ne
fait qu'échanger.
\end{enumerate}

\end{problem}

\noindent Archimède a démontré la loi dans \textit{De l'équilibre des figures planes}, livre I,
propositions 6 et 7, vers 250 av. J.-C., à partir d'une courte liste de postulats sur l'équilibre et
les centres de gravité ; la vanterie de la question (e) est rapportée par Pappus d'Alexandrie quelque
cinq siècles plus tard, et Plutarque raconte l'histoire jumelle d'Archimède tirant seul un navire
chargé sur la plage pour le roi Hiéron. Le résultat était déjà ancien alors, sous une forme plus
fruste : la \textit{Mechanica} pseudo-aristotélicienne avait déjà expliqué le levier en notant que
l'extrémité éloignée de la barre balaie un cercle plus grand, ce qui est la question (c) en germe ---
l'argument par les vitesses virtuelles que Stevin et Galilée allaient affûter et que Lagrange
finirait par transformer en principe des travaux virtuels, l'ancêtre de la machinerie de moindre
action de PHY-04. La question (d) est une querelle philosophique véritable et non tranchée. Mach
admirait la démonstration et soutenait pourtant qu'elle était circulaire, puisque l'expérience dont
elle se nourrit --- que la distance compte, et en proportion --- est précisément le contenu du théorème ;
ses défenseurs répondent que les postulats joints à la symétrie imposent bel et bien la loi linéaire
et rien d'autre. En décider par vous-même vaut mieux que la formule, que vous connaissiez déjà.

\begin{refs}
\item Contexte : Levier --- Wikipédia (\url{https://fr.wikipedia.org/wiki/Levier_(m%C3%A9canique)})
\item Contexte : De l'équilibre des figures planes --- Wikipédia (\url{https://fr.wikipedia.org/wiki/De_l%27%C3%A9quilibre_des_figures_planes})
\item Contexte : Principe des puissances virtuelles --- Wikipédia (\url{https://fr.wikipedia.org/wiki/Principe_des_puissances_virtuelles})
\item Lecture : E. Mach, \textit{The Science of Mechanics} (1883), ch. 1 ; Feynman, \textit{The Feynman Lectures on Physics (\url{https://www.feynmanlectures.caltech.edu/I_04.html}), vol. I, ch. 4 (Conservation of Energy)}
\end{refs}

\bigskip

\begin{problem}[{Pourquoi le ciel est bleu --- Lycée, short}]
\textbf{PHY-57.} Une molécule bien plus petite que la longueur d'onde de la
lumière qui la traverse est polarisée par le champ électrique de cette lumière et rayonne de nouveau
comme un minuscule dipôle oscillant ; en admettant que le moment dipolaire induit est proportionnel
au champ, que la puissance rayonnée par un dipôle oscillant croît comme la puissance quatrième de sa
fréquence, et que la section efficace de diffusion $ \sigma $ ne peut donc dépendre que de la
polarisabilité $ \alpha $ (qui a les dimensions d'un volume) et de la longueur d'onde
$ \lambda $, démontrez par analyse dimensionnelle que
\[ \sigma \;\propto\; \frac{\alpha^{2}}{\lambda^{4}}, \]
et calculez le rapport de la diffusion à 450 nm à celle à 650 nm. Utilisez ensuite le résultat pour
rendre compte de trois observations --- le bleu du ciel de midi, le rouge du Soleil bas sur l'horizon,
et la forte polarisation de la lumière du ciel à quatre-vingt-dix degrés du Soleil --- et expliquez
pourquoi le ciel n'est néanmoins pas violet, bien que le violet soit diffusé plus fortement encore
que le bleu.
\end{problem}

\noindent John Tyndall a montré en 1869 qu'un faisceau traversant un tube de fines
particules en suspension diffuse latéralement la lumière bleue et laisse passer la rouge, et il a
reproduit un coucher de soleil sur une paillasse. Rayleigh a fourni la loi deux ans plus tard dans le
\textit{Philosophical Magazine}, essentiellement par l'argument ci-dessus --- l'un des tout premiers
endroits où le raisonnement dimensionnel sert à extraire une loi physique plutôt qu'à en vérifier
une, et un cousin de PHY-01 et PHY-21. Il a d'abord supposé que les diffuseurs étaient des
poussières ; en 1899 il avait conclu que les molécules de l'air suffisent à elles seules, et que le
ciel bleu est donc une preuve directe que la matière est granuleuse. Smoluchowski (1908) et Einstein
(1910) ont ensuite réécrit tout le calcul comme une diffusion par les fluctuations thermiques de
densité, ce qui explique que la même théorie rende aussi compte de la lueur laiteuse d'un fluide à
son point critique. La dernière clause du problème est celle qui piège les gens : la puissance
quatrième inverse n'est que la moitié de la réponse, et l'autre moitié met en jeu le spectre propre
du Soleil et la réponse de l'\oe{}il humain.

\begin{refs}
\item Contexte : Diffusion Rayleigh --- Wikipédia (\url{https://fr.wikipedia.org/wiki/Diffusion_Rayleigh})
\item Contexte : Couleur du ciel --- Wikipédia (\url{https://fr.wikipedia.org/wiki/Couleur_du_ciel})
\item Contexte : John Tyndall --- Wikipédia (\url{https://fr.wikipedia.org/wiki/John_Tyndall_(physicien)})
\item Lecture : Feynman, \textit{The Feynman Lectures on Physics (\url{https://www.feynmanlectures.caltech.edu/I_32.html}), vol. I, ch. 32 (Radiation Damping. Light Scattering)}
\end{refs}

\bigskip


\section*{Casse-tête}

\begin{problem}[{Neuf boules et une balance --- Lycée}]
\textbf{PUZ-01.} Vous disposez de neuf boules d'apparence identique. Huit pèsent
exactement le même poids ; une est légèrement plus lourde. Votre seul instrument est une
balance à plateaux, qui compare le contenu de ses deux plateaux et vous indique quel côté
est le plus lourd, ou qu'ils s'équilibrent. Elle ne donne aucune valeur numérique.
\begin{enumerate}
\item[(a)] Trouvez une procédure qui identifie la boule lourde en utilisant la balance exactement
deux fois, et démontrez qu'elle fonctionne toujours --- quels que soient les résultats.
\item[(b)] Démontrez qu'une seule pesée ne suffit pas. Démontrez ensuite la borne générale : avec
$ k $ pesées vous ne pouvez pas distinguer plus de $ 3^{k} $ possibilités, parce que
chaque pesée n'a que trois issues. Déduisez-en le plus grand nombre de boules qu'une
procédure à $ k $ pesées peut traiter dans cette version \og{} une boule lourde \fg{}.
\item[(c)] Changez l'énoncé : on sait que la boule intruse a un poids différent, mais on ne vous
dit \textit{pas} si elle est plus lourde ou plus légère. Montrez que deux pesées ne
suffisent alors que pour trois boules, et expliquez exactement où l'argument de comptage
de la question (b) doit être modifié.
\end{enumerate}

\end{problem}

\noindent C'est l'ancêtre de tous les casse-tête de pesée, et sa maison naturelle
est la théorie de l'information. Chaque usage de la balance produit l'un de trois
résultats, donc $ k $ pesées peuvent porter au plus $ \log_2 3^k $ bits d'information ;
une énigme à $ N $ réponses équiprobables demande $ 3^k \geq N $. Ce qui rend la
question (a) satisfaisante, c'est que la borne n'est pas seulement respectée mais
exactement atteinte --- neuf vaut précisément $ 3^2 $, de sorte que la procédure doit tirer
toute l'information de chaque pesée et ne peut se permettre aucune comparaison gaspillée.
La question (c) est là où l'intuition de la plupart des gens casse : ignorer le sens coûte
plus cher qu'il n'y paraît, parce que les issues doivent désormais identifier à la fois
\textit{quelle} boule et \textit{dans quel sens}, et la comptabilité change. La version en
taille industrielle --- douze pièces, sens inconnu, trois pesées --- a tellement circulé
pendant la Seconde Guerre mondiale qu'on plaisantait en y voyant une arme secrète destinée
à faire perdre son temps aux mathématiciens alliés.

\begin{refs}
\item Contexte : Problème de pesée --- Wikipédia (\url{https://fr.wikipedia.org/wiki/Probl%C3%A8me_de_pes%C3%A9e})
\item Contexte : Théorie de l'information --- Wikipédia (\url{https://fr.wikipedia.org/wiki/Th%C3%A9orie_de_l%27information})
\item Lecture : Martin Gardner, \textit{The Scientific American Book of Mathematical Puzzles and Diversions} (1959)
\item Voir aussi : Cryptographie et théorie de l'information (\url{applied-crypto.html}), pour l'entropie comme borne de comptage
\end{refs}

\bigskip

\begin{problem}[{Douze pièces, sens inconnu --- Lycée}]
\textbf{PUZ-02.} Douze pièces sont identiques, à ceci près qu'exactement l'une d'elles
est fausse et diffère des autres par le poids. Vous ignorez si elle est plus lourde ou
plus légère. Vous disposez d'une balance à plateaux et de trois pesées.
\begin{enumerate}
\item[(a)] Trouvez la pièce fausse \textit{et} déterminez si elle est lourde ou légère, en trois
pesées. Démontrez que votre procédure traite tous les cas.
\item[(b)] Il y a $ 24 $ réponses possibles (douze pièces, deux sens) et trois pesées donnent
$ 27 $ issues. Expliquez pourquoi la marge est si mince, et démontrez que treize pièces
est le vrai maximum pour trois pesées si vous devez aussi annoncer le sens --- et qu'avec
treize vous pouvez trouver la pièce sans toujours pouvoir nommer le sens.
\item[(c)] Concevez une procédure où les trois pesées sont décidées \textit{à l'avance}, avant de
voir la moindre issue. Une telle solution non adaptative existe ; en trouver une est plus
difficile que la version adaptative, et c'est la forme qui se généralise aux codes
correcteurs d'erreurs.
\end{enumerate}

\end{problem}

\noindent C'est le casse-tête le plus difficile que la plupart des gens résolvent
sans aide, et le moment d'illumination --- comprendre qu'une pièce peut passer d'un plateau
à l'autre entre deux pesées, ou être mise de côté entièrement --- reste véritablement en
mémoire. La question (c) est la profonde : un schéma de pesée non adaptatif attribue à
chaque pièce une suite fixe de positions, c'est-à-dire précisément un mot de code, et
l'exigence que toutes les issues soient distinguables est une condition de distance
minimale. Casse-tête de pesée et codes correcteurs d'erreurs sont le même sujet vu de deux
côtés, lien rendu explicite par les travaux de Richard Hamming à la fin des années 1940.
Freeman Dyson en a publié une analyse élégante en 1946, ayant reçu l'énigme alors qu'il
travaillait à la recherche opérationnelle du temps de guerre.

\begin{refs}
\item Contexte : Problème de pesée --- Wikipédia (\url{https://fr.wikipedia.org/wiki/Probl%C3%A8me_de_pes%C3%A9e})
\item Contexte : Code de Hamming --- Wikipédia (\url{https://fr.wikipedia.org/wiki/Code_de_Hamming})
\item Article : Freeman J. Dyson, \og{} The problem of the pennies \fg{}, \textit{The Mathematical Gazette} 30 (1946)
\item Voir aussi : Cryptographie et théorie de l'information (\url{applied-crypto.html})
\end{refs}

\bigskip

\begin{problem}[{Deux mèches, quarante-cinq minutes --- Lycée, short}]
\textbf{PUZ-03.} Vous avez deux mèches et un briquet. Chaque mèche
brûle exactement une heure d'un bout à l'autre, mais aucune ne brûle à vitesse uniforme,
de sorte qu'une demi-mèche ne fait pas une demi-heure. Mesurez quarante-cinq minutes, et
démontrez que votre méthode est exacte quelle que soit l'irrégularité de la
combustion.
\end{problem}

\noindent L'astuce, une fois vue, est inoubliable, et la raison pour laquelle elle
marche mérite d'être énoncée précisément : brûler une mèche par les deux bouts la consume
en la moitié de sa \textit{durée totale de combustion}, quelle que soit la répartition de
cette durée le long de sa longueur. C'est un énoncé sur une intégrale que l'on divise par
deux, non sur le milieu de la mèche --- et garder ces deux choses distinctes, c'est tout le
casse-tête.

\begin{refs}
\item Contexte : Mèche (pyrotechnie) --- Wikipédia (\url{https://fr.wikipedia.org/wiki/M%C3%A8che_%28pyrotechnie%29})
\item Lecture : Peter Winkler, \textit{Mathematical Puzzles: A Connoisseur's Collection}
\end{refs}

\bigskip

\begin{problem}[{Les cent casiers --- Lycée, short}]
\textbf{PUZ-04.} Un couloir compte $ 100 $ casiers fermés. L'élève
$ 1 $ bascule l'état de tous les casiers, l'élève $ 2 $ celui d'un casier sur deux,
l'élève $ 3 $ celui d'un casier sur trois, et ainsi de suite jusqu'à l'élève
$ 100 $. Déterminez, avec démonstration, exactement quels casiers sont ouverts à la
fin.
\end{problem}

\noindent Un casier finit ouvert précisément lorsque son numéro a un nombre impair
de diviseurs, et la réponse transforme un casse-tête en petit théorème : les diviseurs
s'apparient selon $ d \leftrightarrow n/d $, donc le compte est impair exactement quand
un diviseur est apparié avec lui-même. Le casse-tête vous demande en réalité de découvrir
pourquoi les carrés parfaits sont particuliers, et c'est l'un des exemples les plus purs
d'argument de comptage qu'un enfant de douze ans peut mener à bout et qu'un théoricien des
nombres respecte encore.

\begin{refs}
\item Contexte : Fonction somme des diviseurs --- Wikipédia (\url{https://fr.wikipedia.org/wiki/Fonction_somme_des_puissances_k-i%C3%A8mes_des_diviseurs})
\item Voir aussi : Théorie des nombres (\url{pure-number-theory.html})
\end{refs}

\bigskip

\begin{problem}[{Le pont et la lampe --- Lycée}]
\textbf{PUZ-05.} Quatre personnes doivent traverser de nuit un pont branlant. Le pont
supporte au plus deux personnes à la fois, et quiconque traverse doit porter l'unique
lampe. Seules, les quatre mettent $ 1 $, $ 2 $, $ 5 $ et $ 10 $ minutes ; une
paire traverse à l'allure du plus lent. Quelqu'un doit toujours rapporter la lampe.
\begin{enumerate}
\item[(a)] Faites passer les quatre en $ 17 $ minutes.
\item[(b)] Démontrez que $ 17 $ est optimal --- aucun ordonnancement ne fait mieux.
\item[(c)] La stratégie gloutonne évidente, où la personne la plus rapide fait à chaque fois la
navette avec la lampe, prend $ 19 $ minutes. Expliquez précisément pourquoi la gloutonnerie
échoue ici, et dégagez le principe général : les deux plus lents doivent traverser ensemble,
de sorte que le plus grand de leurs temps ne soit payé qu'une fois.
\end{enumerate}

\end{problem}

\noindent La question (b) est la partie mathématique. Produire un ordonnancement en
$ 17 $ minutes est affaire d'essais ; démontrer qu'aucun ordonnancement ne fait mieux
demande de raisonner sur tous les ordonnancements à la fois, et la voie habituelle consiste
à observer que la structure de toute solution est forcée --- les traversées alternent avec
les retours, de sorte qu'avec quatre personnes il y a exactement trois traversées aller et
deux retours. Ce casse-tête est un favori en entretien précisément parce que l'écart entre
(a) et (b) sépare ceux qui ont trouvé une réponse de ceux qui la comprennent. C'est aussi
une véritable instance de problème d'ordonnancement où l'heuristique gloutonne naturelle
est démontrablement sous-optimale, raison pour laquelle il figure dans les cours
d'algorithmique.

\begin{refs}
\item Contexte : Problème du pont et de la lampe --- Wikipédia (en anglais) (\url{https://en.wikipedia.org/wiki/Bridge_and_torch_problem})
\item Contexte : Algorithme glouton --- Wikipédia (\url{https://fr.wikipedia.org/wiki/Algorithme_glouton})
\item Voir aussi : Optimisation et théorie des jeux (\url{applied-optimization.html})
\end{refs}

\bigskip

\begin{problem}[{Le loup, la chèvre et le chou --- Lycée}]
\textbf{PUZ-06.} Un fermier doit faire traverser une rivière à un loup, une chèvre et
un chou dans une barque qui ne porte que le fermier et au plus un des trois. Laissés seuls
ensemble, le loup mange la chèvre et la chèvre mange le chou.
\begin{enumerate}
\item[(a)] Trouvez une suite de traversées qui fonctionne.
\item[(b)] Modélisez le casse-tête par un graphe : chaque sommet est une configuration licite ---
qui se trouve sur quelle rive, et où est la barque --- et l'on relie deux sommets lorsqu'une
seule traversée transforme l'un en l'autre. Dessinez ce graphe et identifiez chaque
solution comme un chemin dedans.
\item[(c)] Démontrez que votre solution est la plus courte, et déterminez combien de solutions
distinctes de longueur minimale existent.
\end{enumerate}

\end{problem}

\noindent Ce casse-tête a plus de mille ans --- il apparaît dans les
\textit{Propositiones ad Acuendos Juvenes}, un recueil attribué à Alcuin d'York vers l'an
800, ce qui en fait l'un des plus anciens problèmes récréatifs consignés en Europe. La
question (b) est la raison de sa présence sur ce site : dès que vous dessinez le graphe des
configurations, une histoire de bétail devient un objet fini sur lequel vous pouvez
raisonner exhaustivement, et des questions comme \og{} existe-t-il une solution plus
courte ? \fg{} deviennent des questions de longueur de chemin. Cette traduction --- du récit à
l'espace d'états --- est exactement ce que fait un informaticien quand il planifie une
recherche, et elle vaut d'être effectuée à la main au moins une fois.

\begin{refs}
\item Contexte : Problème du loup, de la chèvre et du chou --- Wikipédia (\url{https://fr.wikipedia.org/wiki/Probl%C3%A8me_du_loup%2C_de_la_ch%C3%A8vre_et_du_chou})
\item Contexte : Espace d'états --- Wikipédia (\url{https://fr.wikipedia.org/wiki/Espace_d%27%C3%A9tats})
\item Voir aussi : Combinatoire et théorie des graphes (\url{pure-combinatorics.html}), Mathématiques discrètes et informatique (\url{applied-discrete-cs.html})
\end{refs}

\bigskip

\begin{problem}[{Chevaliers, valets et une seule question --- Lycée}]
\textbf{PUZ-07.} Sur une île, chaque habitant est soit un chevalier, qui dit toujours
la vérité, soit un valet, qui ment toujours.
\begin{enumerate}
\item[(a)] Vous rencontrez deux insulaires, A et B. A déclare : \og{} Nous sommes tous les deux des
valets. \fg{} Déterminez ce qu'est chacun d'eux, et démontrez que votre déduction est la seule
possibilité.
\item[(b)] Vous arrivez à une bifurcation dont une branche mène au salut. Un seul insulaire s'y
tient, et vous pouvez poser une seule question à réponse oui-ou-non. Trouvez une question
qui révèle la route sûre, et démontrez qu'elle fonctionne que l'insulaire soit chevalier
ou valet.
\item[(c)] Généralisez : montrez qu'une question de la forme \og{} si je vous demandais $ Q $,
répondriez-vous oui ? \fg{} fait toujours émerger la vérité au sujet de $ Q $, et expliquez
pourquoi composer un menteur avec lui-même produit de l'honnêteté.
\end{enumerate}

\end{problem}

\noindent Raymond Smullyan a bâti une carrière sur ces insulaires, et la question
(c) met à nu le mécanisme qui les sous-tend tous : la réponse d'un valet est la négation de
la vérité, de sorte qu'enfouir la question d'un niveau supplémentaire applique la négation
deux fois, et le mensonge s'annule lui-même. C'est une observation de théorie des groupes
déguisée --- les comportements de véracité et de mensonge se composent comme les éléments
d'un groupe d'ordre deux. La même astuce autoréférentielle anime l'énigme logique la plus
difficile jamais posée, le problème des trois dieux de George Boolos, où les dieux
répondent dans une langue inconnue.

\begin{refs}
\item Contexte : Chevaliers et valets --- Wikipédia (en anglais) (\url{https://en.wikipedia.org/wiki/Knights_and_knaves})
\item Contexte : L'Énigme la plus difficile du monde --- Wikipédia (\url{https://fr.wikipedia.org/wiki/L%27%C3%89nigme_la_plus_difficile_du_monde})
\item Lecture : Raymond Smullyan, \textit{What Is the Name of This Book?} (1978) --- publié en français sous le titre \textit{Quel est le titre de ce livre ?}
\item Voir aussi : Logique, théorie des ensembles et fondements (\url{pure-foundations.html})
\end{refs}

\bigskip

\begin{problem}[{Paver un damier avec des pièces 1$\times$4 --- Lycée, short}]
\textbf{PUZ-08.} On veut paver un plateau $ 10 \times 10 $ par des
pièces $ 1 \times 4 $ placées horizontalement ou verticalement. Démontrez que c'est
impossible.
\end{problem}

\noindent Contrairement à l'échiquier mutilé, un argument à deux couleurs ne suffit
pas ici ; il vous faut un coloriage plus astucieux, et le trouver est tout le casse-tête.
Les arguments de coloriage sont la forme la moins coûteuse de démonstration d'impossibilité
en mathématiques --- vous inventez une quantité que tout coup licite préserve, puis vous
montrez que la cible la viole --- et apprendre à chercher le bon invariant plutôt que la
bonne construction est un véritable changement dans la manière d'attaquer les
problèmes.

\begin{refs}
\item Contexte : Polyomino --- Wikipédia (\url{https://fr.wikipedia.org/wiki/Polyomino})
\item Voir aussi : L'art de la démonstration (\url{proofs.html}), pour l'échiquier mutilé
\end{refs}

\bigskip

\begin{problem}[{Deux brocs et aucune graduation --- Lycée}]
\textbf{PUZ-17.} Vous disposez de deux brocs sans graduation, de capacités entières $ a $ et $ b $ litres, d'un robinet fournissant de l'eau sans limite, et d'un évier. Les seules opérations autorisées sont : remplir un broc à ras bord, vider complètement un broc, et verser d'un broc dans l'autre jusqu'à ce que la source soit vide ou que la destination soit pleine.
\begin{enumerate}
\item[(a)] Avec $ a = 3 $ et $ b = 5 $, laissez exactement $ 4 $ litres dans l'un des brocs. Écrivez la suite complète des opérations et justifiez chaque étape.
\item[(b)] Déterminez, avec démonstration, exactement quelles quantités $ d $ sont mesurables --- c'est-à-dire pour quels $ d $ une suite finie d'opérations autorisées laisse exactement $ d $ litres dans un broc. Démontrez les deux sens : que tout $ d $ de ce type est atteignable, et qu'aucun autre $ d $ ne l'est.
\item[(c)] La version de Tartaglia n'a ni robinet ni évier, de sorte que le total est conservé : un broc de $ 8 $ litres est plein de vin, et des brocs de $ 5 $ et $ 3 $ litres sont vides. Déterminez, avec démonstration, toutes les répartitions de vin atteignables et, en particulier, décidez si le vin peut être partagé en deux parts égales.
\end{enumerate}

\end{problem}

\noindent Les casse-tête de transvasement sont médiévaux, et le problème du vin $ 8 $--$ 5 $--$ 3 $ de la question (c) est assez ancien pour figurer dans les \textit{Problèmes plaisants et délectables} de Claude Gaspard Bachet de Méziriac, en 1612 ; la version $ 3 $--$ 5 $ de la question (a) a touché un très large public grâce à \textit{Une journée en enfer} en 1995. Les mathématiques derrière la question (b) sont la théorie des équations diophantiennes linéaires : tout état atteignable est une combinaison entière des deux capacités, si bien que la question est de savoir quels entiers de telles combinaisons produisent --- c'est le contenu du théorème de Bachet-Bézout, que Bachet lui-même connaissait dans le cas entier bien avant que Bézout n'attache son nom à la version polynomiale. La question (c) est un problème réellement différent, car la conservation du total vous confine à un ensemble d'états de dimension deux ; M. C. K. Tweedie a montré en 1939 que toute la famille des casse-tête à total conservé cède à une seule image géométrique, bel exemple de changement de représentation rendant inutile une recherche exhaustive.

\begin{refs}
\item Contexte : Casse-tête de transvasement --- Wikipédia (en anglais) (\url{https://en.wikipedia.org/wiki/Water_pouring_puzzle})
\item Contexte : Théorème de Bachet-Bézout --- Wikipédia (\url{https://fr.wikipedia.org/wiki/Th%C3%A9or%C3%A8me_de_Bachet-B%C3%A9zout})
\item Article : M. C. K. Tweedie, \og{} A graphical method of solving Tartaglian measuring puzzles \fg{}, \textit{The Mathematical Gazette} 23 (1939), 278--282
\item Voir aussi : Théorie des nombres (\url{pure-number-theory.html})
\end{refs}

\bigskip

\begin{problem}[{Les missionnaires et les cannibales --- Lycée}]
\textbf{PUZ-18.} Trois missionnaires et trois cannibales se tiennent sur une rive de rivière. Une barque porte au plus deux personnes et ne peut pas traverser à vide : quelqu'un doit ramer dans chaque sens. Si, à un instant quelconque, les cannibales sont plus nombreux que les missionnaires sur l'une des rives où se trouve au moins un missionnaire, les missionnaires sont mangés.
\begin{enumerate}
\item[(a)] Trouvez une suite de traversées qui amène les six sains et saufs sur l'autre rive.
\item[(b)] Formez le graphe dont les sommets sont les états licites --- le nombre de missionnaires et le nombre de cannibales sur la rive de départ, ainsi que la position de la barque --- et dont les arêtes relient les états qui diffèrent d'une traversée. Démontrez que votre solution est un plus court chemin dans ce graphe, et comptez combien de solutions de longueur minimale il existe.
\item[(c)] Déterminez, avec démonstration, si quatre missionnaires et quatre cannibales peuvent traverser dans une barque de deux places. Déterminez ensuite le plus grand $ n $ pour lequel $ n $ missionnaires et $ n $ cannibales peuvent traverser quand la barque en porte trois.
\end{enumerate}

\end{problem}

\noindent Ce casse-tête descend d'un problème de passage de rivière des \textit{Propositiones ad Acuendos Juvenes}, le recueil attribué à Alcuin d'York vers l'an 800, où il apparaît sous la forme de trois frères voyageant chacun avec une s\oe{}ur, avec une barque de deux places. Il a circulé dans le nord de l'Europe à partir du treizième siècle comme problème des \og{} maris jaloux \fg{}, et n'a acquis ses missionnaires et ses cannibales qu'à la fin du dix-neuvième. Son importance moderne est autre : Saul Amarel l'a utilisé en 1968 comme exemple central d'un article influent sur la façon dont le choix de la représentation détermine si un problème de recherche est traitable, et il illustre depuis, dans tous les manuels, la recherche dans un espace d'états en intelligence artificielle. La question (c) est là où le casse-tête cesse d'être une histoire pour devenir un théorème --- la capacité de la barque et le nombre de paires interagissent d'une manière qu'Ian Pressman et David Singmaster ont analysée avec soin en 1989.

\begin{refs}
\item Contexte : Problèmes de passage de rivière --- Wikipédia (\url{https://fr.wikipedia.org/wiki/Probl%C3%A8mes_de_passage_de_rivi%C3%A8re})
\item Contexte : Propositiones ad acuendos juvenes --- Wikipédia (\url{https://fr.wikipedia.org/wiki/Propositiones_ad_acuendos_juvenes})
\item Article : Ian Pressman \& David Singmaster, \og{} `The Jealous Husbands' and `The Missionaries and Cannibals' \fg{}, \textit{The Mathematical Gazette} 73 (1989)
\item Voir aussi : Combinatoire et théorie des graphes (\url{pure-combinatorics.html}), Mathématiques discrètes et informatique (\url{applied-discrete-cs.html})
\end{refs}

\bigskip

\begin{problem}[{Une seule lecture sur une balance numérique --- Lycée}]
\textbf{PUZ-19.} Dix sacs contiennent chacun un grand nombre de pièces. Les pièces authentiques pèsent exactement $ 10 $ grammes ; toutes les pièces d'exactement un sac sont fausses et pèsent $ 9 $ grammes. Votre instrument est une balance numérique qui affiche le poids total exact de ce que vous y posez, et vous ne pouvez l'utiliser qu'une fois.
\begin{enumerate}
\item[(a)] Identifiez le sac de fausses pièces en une seule pesée, et démontrez que votre procédure distingue les dix possibilités.
\item[(b)] Supposons maintenant qu'un nombre inconnu des dix sacs --- peut-être aucun, peut-être tous --- contienne des fausses pièces. Déterminez, avec démonstration, si une pesée suffit encore, et donnez en conséquence soit une procédure, soit un argument d'impossibilité.
\item[(c)] Modifiez l'énoncé pour que ce soient des pièces individuelles, et non des sacs entiers, qui puissent être fausses : parmi $ n $ pièces, un sous-ensemble inconnu est léger, et chaque pesée affiche le poids total exact de tout sous-ensemble que vous choisissez. Supposez que tous les sous-ensembles à peser doivent être fixés à l'avance. Démontrez les meilleures bornes supérieure et inférieure que vous pouvez sur le nombre de pesées nécessaires, en fonction de $ n $.
\end{enumerate}

\end{problem}

\noindent Une balance à deux plateaux renvoie l'un de trois symboles : elle porte donc au plus $ \log_2 3 $ bits par usage ; une balance qui affiche un nombre peut en principe en porter un nombre non borné, et le casse-tête consiste à convertir cette capacité brute en un code utilisable. Les questions (a) et (b) relèvent d'un folklore d'origine incertaine, circulant depuis des décennies comme question d'entretien, mais la question (c) est un vrai problème de recherche : c'est le problème de pesée de pièces posé par Paul Erd\H{o}s et Alfréd Rényi en 1963, dont la réponse asymptotique a demandé plusieurs années et plusieurs auteurs. La même idée --- regrouper les échantillons pour qu'une mesure résolve d'un coup de nombreux objets --- a été inventée indépendamment par Robert Dorfman en 1943 pour dépister la syphilis chez les soldats américains, et fait aujourd'hui l'objet de la théorie des tests groupés, revenue sur le devant de la scène lors des campagnes de dépistage à grande échelle.

\begin{refs}
\item Contexte : Tests groupés --- Wikipédia (\url{https://fr.wikipedia.org/wiki/Tests_group%C3%A9s})
\item Contexte : Problème de pesée --- Wikipédia (\url{https://fr.wikipedia.org/wiki/Probl%C3%A8me_de_pes%C3%A9e})
\item Article : P. Erd\H{o}s \& A. Rényi, \og{} On two problems of information theory \fg{}, \textit{Publications of the Mathematical Institute of the Hungarian Academy of Sciences} 8 (1963), 229--243
\item Voir aussi : Cryptographie et théorie de l'information (\url{applied-crypto.html})
\end{refs}

\bigskip

\begin{problem}[{Cent pièces dans le noir --- Lycée, short}]
\textbf{PUZ-20.} Cent pièces sont posées sur une table dans une pièce entièrement obscure ; exactement dix d'entre elles sont côté pile visible, et vous ne pouvez déterminer par le toucher ni par aucun autre moyen dans quel sens une pièce est tournée, bien que vous puissiez déplacer les pièces et les retourner. Séparez-les en deux groupes, pas nécessairement de même taille, de sorte que les deux groupes contiennent le même nombre de piles, et démontrez que votre procédure fonctionne pour toute disposition initiale possible.
\end{problem}

\noindent Le difficile n'est pas de trouver une man\oe{}uvre astucieuse, mais de remarquer quel genre d'objet on vous demande : puisque vous n'apprenez jamais rien sur la disposition, votre procédure est une recette fixe qui doit réussir simultanément pour les $ \binom{100}{10} $ dispositions. Cela vous force à chercher une quantité dont le comportement sous vos coups est le même quoi que vous ne puissiez voir --- un argument d'invariance mené à l'envers, et une bonne première rencontre avec l'idée qu'une démonstration peut maîtriser tous les cas à la fois sans en examiner aucun.

\begin{refs}
\item Contexte : Invariant (mathématiques) --- Wikipédia (\url{https://fr.wikipedia.org/wiki/Invariant})
\item Voir aussi : L'art de la démonstration (\url{proofs.html})
\end{refs}

\bigskip

\begin{problem}[{Le traiteur paresseux et le pâtissier négligent --- Lycée}]
\textbf{PUZ-21.} Une crêpe est un disque. Une coupe est une droite entière qui la traverse, et les morceaux sont les régions connexes qui restent ; les morceaux ne sont jamais réarrangés entre deux coupes.
\begin{enumerate}
\item[(a)] Déterminez, avec démonstration, le plus grand nombre de morceaux que l'on peut obtenir avec $ n $ coupes, et démontrez votre formule par récurrence sur $ n $.
\item[(b)] Une autre question : marquez $ n $ points sur la circonférence d'un disque, en position générale de sorte que trois des cordes qu'ils déterminent ne passent jamais par un même point intérieur, et tracez les $ \binom{n}{2} $ cordes. Déterminez, avec démonstration, le nombre de régions en lesquelles le disque est découpé, et comparez votre réponse à $ 2^{n-1} $.
\item[(c)] En dimension trois : un gâteau cubique est coupé par $ n $ plans, toujours sans réarrangement. Déterminez, avec démonstration, le plus grand nombre de morceaux, et expliquez le lien entre cette réponse et celle trouvée en (a).
\end{enumerate}

\end{problem}

\noindent La question (a) est la plus ancienne des trois et parle en réalité des arrangements de droites du plan, sujet ouvert par Jakob Steiner en 1826 et aujourd'hui standard en géométrie algorithmique, où la complexité d'un arrangement détermine le temps d'exécution de bien des algorithmes. La question (b) est la configuration dont Leo Moser s'est servi en 1949 pour mettre en garde contre la généralisation à partir d'une poignée de cas calculés ; c'est l'exemple le plus cité en mathématiques d'une récurrence qui n'existe pas, et l'intérêt de l'exercice est d'obtenir le compte par une méthode qui ne repose pas sur le repérage d'un motif. La question (c) appartient à la même famille, et les trois réponses forment ensemble un petit motif de coefficients binomiaux qui dit quelque chose de structurel sur la dimension.

\begin{refs}
\item Contexte : Suite du traiteur paresseux --- Wikipédia (\url{https://fr.wikipedia.org/wiki/Suite_du_traiteur_paresseux})
\item Contexte : Découpage d'un disque en régions --- Wikipédia (en anglais) (\url{https://en.wikipedia.org/wiki/Dividing_a_circle_into_areas})
\item Contexte : Suite des nombres gâteaux --- Wikipédia (\url{https://fr.wikipedia.org/wiki/Suite_des_nombres_g%C3%A2teaux})
\item Voir aussi : Combinatoire et théorie des graphes (\url{pure-combinatorics.html})
\end{refs}

\bigskip

\begin{problem}[{Des tétrominos en T sur un plateau --- Lycée}]
\textbf{PUZ-22.} Un tétromino en T est le polyomino à quatre cases formé d'une rangée de trois cases avec une quatrième case accolée au milieu de l'un des longs côtés. Les tuiles peuvent être tournées, ne peuvent pas se chevaucher, et doivent recouvrir le plateau exactement, sans laisser de case découverte.
\begin{enumerate}
\item[(a)] Pavez un plateau $ 4 \times 4 $ avec quatre tétrominos en T.
\item[(b)] Démontrez qu'un plateau $ 10 \times 10 $ n'admet aucun pavage par des tétrominos en T, bien que ses $ 100 $ cases soient divisibles par $ 4 $.
\item[(c)] Déterminez, avec démonstration, exactement quels rectangles $ m \times n $ admettent un pavage par des tétrominos en T.
\end{enumerate}

\end{problem}

\noindent Solomon Golomb a forgé le mot \og{} polyomino \fg{} lors d'un exposé de 1953 au Harvard Mathematics Club et a passé les quarante années suivantes à faire de ces formes une discipline ; la question (c) a été réglée par D. W. Walkup en 1965, dans une note de trois pages de l'\textit{American Mathematical Monthly}. Ce qui rend le tétromino en T instructif, c'est que l'obstruction naïve échoue : le compte des cases est correct, et l'argument ordinaire du damier noir et blanc qui tue le plateau mutilé ne tue pas immédiatement celui-ci, de sorte qu'il faut réfléchir plus sérieusement à ce à quoi sert un coloriage. Une démonstration d'impossibilité de ce genre est une affirmation sur une infinité de pavages candidats, établie sans en examiner aucun, et apprendre à fabriquer la bonne pondération est un ajout définitif à la boîte à outils de qui résout des problèmes.

\begin{refs}
\item Contexte : Tétromino --- Wikipédia (\url{https://fr.wikipedia.org/wiki/T%C3%A9tromino})
\item Contexte : Polyomino --- Wikipédia (\url{https://fr.wikipedia.org/wiki/Polyomino})
\item Article : D. W. Walkup, \og{} Covering a rectangle with T-tetrominoes \fg{}, \textit{American Mathematical Monthly} 72 (1965), 986--988
\item Lecture : Solomon W. Golomb, \textit{Polyominoes: Puzzles, Patterns, Problems, and Packings} (2e éd., 1994)
\end{refs}

\bigskip

\begin{problem}[{L'anniversaire de Cheryl --- Lycée}]
\textbf{PUZ-31.} Albert et Bernard veulent connaître la date d'anniversaire de Cheryl.
Elle leur dit que c'est l'une de dix dates : 15 mai, 16 mai, 19 mai ; 17 juin, 18 juin ;
14 juillet, 16 juillet ; 14 août, 15 août, 17 août. Elle dit ensuite à Albert le mois, en
privé, et à Bernard le quantième, en privé. Chacun sait que l'autre a été informé, et
chacun sait que l'autre raisonne parfaitement. Ce qui suit est dit à voix haute, dans
l'ordre.

\noindent Albert : \og{} Je ne sais pas quand Cheryl est née, et je sais que Bernard ne le sait pas
non plus. \fg{} Bernard : \og{} Au début je ne savais pas, mais maintenant je sais. \fg{} Albert :
\og{} Alors je le sais aussi. \fg{}
\begin{enumerate}
\item[(a)] Déterminez la date d'anniversaire de Cheryl, et justifiez chaque élimination : pour
chacune des dix dates, dites laquelle des trois déclarations l'écarte, et pourquoi.
\item[(b)] Écrivez chaque déclaration comme une condition exacte. La première phrase d'Albert,
par exemple, est un énoncé qui ne porte que sur le mois ; rendez cela précis, et énoncez
ce que la phrase de Bernard affirme au sujet du quantième, compte tenu de ce qu'Albert a
déjà dit.
\item[(c)] Le casse-tête est fragile. Construisez une liste de dates de votre cru sur laquelle
les mêmes trois déclarations ne pourraient pas toutes être faites de façon véridique, puis
une seconde liste sur laquelle elles peuvent toutes être faites tout en laissant deux dates
possibles. Concluez que l'unicité de la réponse est une propriété de la liste particulière
de Cheryl, non du format.
\end{enumerate}

\end{problem}

\noindent Ce problème est paru comme question 24 de l'olympiade mathématique des
écoles de Singapour et d'Asie, le 8 avril 2015, rédigé par le Dr Joseph Yeo Boon Wooi, du
National Institute of Education de Singapour ; un présentateur de télévision, Kenneth Kong,
l'a mis en ligne deux jours plus tard et il a circulé dans le monde entier en une semaine,
en grande partie parce que quantité d'adultes étaient convaincus qu'il ne pouvait pas être
résolu. Il peut l'être, et la raison en est le moteur de toute cette famille de problèmes :
une déclaration d'\textit{ignorance} est informative. Chaque locuteur, en avouant ne pas
savoir, apprend à l'autre que certaines possibilités sont absentes, et les possibilités
sont élaguées tour après tour jusqu'à ce qu'une seule date survive. C'est exactement le
mécanisme des insulaires aux yeux bleus de PUZ-11, mais assez petit pour que l'argument
tout entier tienne sur une page --- ce qui en fait le meilleur premier exercice de
raisonnement sur ce que les autres savent.

\begin{refs}
\item Contexte : Anniversaire de Cheryl --- Wikipédia (\url{https://fr.wikipedia.org/wiki/Anniversaire_de_Cheryl})
\item Contexte : Logique épistémique dynamique --- Wikipédia (en anglais) (\url{https://en.wikipedia.org/wiki/Dynamic_epistemic_logic})
\item Voir aussi : PUZ-11 ci-dessus, et Logique, théorie des ensembles et fondements (\url{pure-foundations.html})
\end{refs}

\bigskip

\begin{problem}[{L'énigme du zèbre --- Lycée}]
\textbf{PUZ-32.} Cinq maisons s'alignent. Chacune est d'une couleur différente,
habitée par un homme d'une nationalité différente, et chaque homme boit une boisson
différente, fume une marque de cigarettes différente et garde un animal différent. On vous
donne quinze faits.

\noindent (1) L'Anglais habite la maison rouge. (2) L'Espagnol possède le chien. (3) On boit du
café dans la maison verte. (4) L'Ukrainien boit du thé. (5) La maison verte est
immédiatement à droite de la maison ivoire. (6) Le fumeur d'Old Gold possède les
escargots. (7) On fume des Kools dans la maison jaune. (8) On boit du lait dans la maison
du milieu. (9) Le Norvégien habite la première maison. (10) L'homme qui fume des
Chesterfield habite la maison voisine de celle de l'homme au renard. (11) On fume des
Kools dans la maison voisine de celle où l'on garde le cheval. (12) Le fumeur de Lucky
Strike boit du jus d'orange. (13) Le Japonais fume des Parliament. (14) Le Norvégien
habite à côté de la maison bleue. (15) Il y a cinq maisons.
\begin{enumerate}
\item[(a)] Déterminez qui boit de l'eau et qui possède le zèbre, et démontrez que les quinze
faits n'admettent qu'une seule affectation complète.
\item[(b)] Reformulez le casse-tête comme un problème de satisfaction de contraintes : prenez
cinq variables par catégorie (couleur, nationalité, boisson, marque, animal), chacune à
valeurs dans les positions de maison $ \{1,2,3,4,5\} $, imposez une contrainte de
différence deux à deux à l'intérieur de chaque catégorie, et traduisez chaque fait en
contrainte. Montrez que le nombre d'affectations à parcourir aveuglément est
$ (5!)^{5} = 24\,883\,200\,000 $, et décrivez ce que fait gagner l'application locale de
chaque contrainte.
\item[(c)] Ni l'eau ni le zèbre ne sont mentionnés dans aucun des quinze faits, et pourtant le
casse-tête interroge sur les deux. Expliquez pourquoi c'est légitime. Construisez ensuite
un petit casse-tête de même forme --- trois catégories sur trois positions, contraintes de
différence deux à deux, et une poignée de faits \og{} à côté de \fg{} et \og{} même maison \fg{} --- qui
possède exactement deux solutions, puis un autre qui n'en possède aucune.
\end{enumerate}

\end{problem}

\noindent L'énigme a été imprimée dans \textit{Life International} le 17 décembre
1962, la réponse paraissant en mars suivant ; on l'attribue souvent à Einstein ou à Lewis
Carroll, sans le moindre élément de preuve dans l'un ou l'autre cas. Son importance réelle
est d'être devenue l'exemple traité standard du problème de satisfaction de contraintes, le
modèle dans lequel un ensemble fini de variables à domaines finis doit satisfaire des
relations données. Ce cadre couvre le coloriage de graphes, la construction d'emplois du
temps, le sudoku et la vérification de circuits, et il possède une théorie profonde : Feder
et Vardi ont conjecturé dans les années 1990 que tout langage de contraintes fixé est soit
résoluble en temps polynomial, soit NP-complet, sans rien entre les deux, et la conjecture
a été démontrée indépendamment par Andrei Bulatov et Dmitriy Zhuk en 2017. La question (b)
est le moment où une histoire de cigarettes devient cette discipline : une fois les
contraintes écrites, la question \og{} comment explorer cela efficacement ? \fg{} ne parle plus de
zèbres.

\begin{refs}
\item Contexte : Énigme d'Einstein (énigme du zèbre) --- Wikipédia (\url{https://fr.wikipedia.org/wiki/%C3%89nigme_d%27Einstein})
\item Contexte : Problème de satisfaction de contraintes --- Wikipédia (\url{https://fr.wikipedia.org/wiki/Probl%C3%A8me_de_satisfaction_de_contraintes})
\item Lecture : Stuart Russell \& Peter Norvig, \textit{Artificial Intelligence: A Modern Approach}, le chapitre sur les problèmes de satisfaction de contraintes
\item Voir aussi : Mathématiques discrètes et informatique (\url{applied-discrete-cs.html})
\end{refs}

\bigskip

\begin{problem}[{Les vingt questions --- Lycée, short}]
\textbf{PUZ-33.} Un joueur choisit secrètement un élément d'un
ensemble fini $ S $ connu des deux ; l'autre doit l'identifier en posant des questions de
la forme \og{} votre élément appartient-il à $ A $ ? \fg{} pour toute partie
$ A \subseteq S $ de son choix, chaque réponse étant véridique. Déterminez, avec
démonstration, le plus petit nombre de questions qui suffise toujours, en fonction de
$ |S| $, les questions ultérieures pouvant dépendre des réponses antérieures.
\end{problem}

\noindent Le jeu de société est ancien --- Hannah More raconte y avoir joué lors
d'un dîner londonien dans les années 1780, et son cousin hongrois, le Barkochba, se joue à
Budapest depuis 1911 --- mais les mathématiques sont celles de Shannon. Une réponse par oui
ou par non porte au plus un bit, de sorte que $ q $ questions peuvent séparer au plus
$ 2^{q} $ possibilités ; c'est la même comptabilité que celle qui régit la balance de
PUZ-01, où chaque pesée porte $ \log_2 3 $ bits au lieu d'un. La moitié facile consiste à
exhiber une stratégie. La moitié qui mérite d'être traitée avec soin est la borne
inférieure, car elle doit mettre en échec toutes les stratégies adaptatives à la fois, et
pas seulement les évidentes.

\begin{refs}
\item Contexte : Le jeu des vingt questions --- Wikipédia (en anglais) (\url{https://en.wikipedia.org/wiki/Twenty_questions})
\item Contexte : Entropie de Shannon --- Wikipédia (\url{https://fr.wikipedia.org/wiki/Entropie_de_Shannon})
\item Voir aussi : Cryptographie et théorie de l'information (\url{applied-crypto.html})
\end{refs}

\bigskip

\begin{problem}[{Le problème de Josèphe --- Lycée}]
\textbf{PUZ-34.} Numérotez $ n $ personnes $ 1, 2, \ldots, n $ dans le sens des
aiguilles d'une montre autour d'un cercle. En commençant à compter à la personne $ 1 $,
une personne sur deux encore dans le cercle est retirée --- la personne $ 2 $ d'abord,
puis $ 4 $, et ainsi de suite tour après tour --- jusqu'à ce qu'il n'en reste qu'une.
Notons $ J(n) $ le numéro de ce survivant.
\begin{enumerate}
\item[(a)] Calculez $ J(n) $ pour $ n = 1, 2, \ldots, 16 $ et décrivez le motif que vous
observez.
\item[(b)] Démontrez une récurrence exprimant $ J(2m) $ et $ J(2m+1) $ en fonction de
$ J(m) $, et servez-vous-en pour obtenir une forme close de $ J(n) $ valable pour tout
$ n $. Démontrez cette forme close par récurrence.
\item[(c)] Retirez maintenant une personne sur $ k $ au lieu d'une sur deux. Démontrez que
\[ J_k(n) = \big( J_k(n-1) + k - 1 \big) \bmod n \; + \; 1, \qquad J_k(1) = 1, \]
et expliquez pourquoi cela donne un algorithme plutôt qu'une formule. Pour quels $ k $
pouvez-vous trouver quoi que ce soit d'aussi net que la réponse à la question (b) ?
\end{enumerate}

\end{problem}

\noindent Flavius Josèphe raconte dans \textit{La Guerre des Juifs} (livre III,
chapitre 8) comment, après la chute de Yodfat en 67, lui et quarante compagnons se sont
cachés dans une grotte et ont convenu de mourir par tirage au sort plutôt que de se rendre
aux Romains --- et comment lui et un autre survécurent, \og{} par chance ou peut-être par la main
de Dieu \fg{}. La formalisation par comptage est postérieure : Bachet de Méziriac en a donné
une version en 1612, et une refonte médiévale place quinze chrétiens et quinze Turcs sur un
navire en train de sombrer, un homme sur neuf étant jeté par-dessus bord. Ce qui fait du cas
$ k=2 $ un petit classique, c'est que la réponse, une fois trouvée, est étonnamment simple
--- assez simple pour que Graham, Knuth et Patashnik ouvrent \textit{Concrete Mathematics} avec
elle, s'en servant pour enseigner comment deviner un motif à partir de données, le démontrer
par récurrence, puis voir pourquoi la démonstration fonctionne. La question (c) est le
contrepoids honnête : le cas général $ k $ est calculable mais obstinément informe.

\begin{refs}
\item Contexte : Problème de Josèphe --- Wikipédia (\url{https://fr.wikipedia.org/wiki/Probl%C3%A8me_de_Jos%C3%A8phe})
\item Lecture : Ronald Graham, Donald Knuth \& Oren Patashnik, \textit{Concrete Mathematics}, \S{}1.3
\item Voir aussi : Combinatoire et théorie des graphes (\url{pure-combinatorics.html}), L'art de la démonstration (\url{proofs.html})
\end{refs}

\bigskip

\begin{problem}[{Les deux enveloppes, scellées --- Lycée}]
\textbf{PUZ-45.} Deux enveloppes scellées contiennent de l'argent ; l'une en contient
exactement deux fois plus que l'autre. Vous en prenez une au hasard et, sans l'ouvrir, on
vous propose d'échanger.
\begin{enumerate}
\item[(a)] Voici un argument en faveur de l'échange. Soit $ A $ la somme contenue dans votre
enveloppe. L'autre enveloppe contient $ 2A $ ou $ A/2 $, chacun avec probabilité
$ 1/2 $, de sorte que son contenu espéré vaut
\[ \tfrac{1}{2}(2A) + \tfrac{1}{2}\!\left(\tfrac{A}{2}\right) = \tfrac{5}{4}A > A . \]
Le même argument s'applique à l'autre enveloppe, si bien que chacune vaut plus que l'autre.
Localisez l'étape exacte à partir de laquelle ceci cesse d'être une probabilité valide, et
justifiez votre verdict en nommant un espace probabilisé, les variables aléatoires qui y
sont définies, et la quantité sur laquelle on conditionne réellement.
\item[(b)] Construisez un modèle honnête : soit $ X $ la variable aléatoire donnant la plus
petite somme, et supposons que les deux enveloppes contiennent $ X $ et $ 2X $.
Démontrez que si $ E[X] $ est finie, le gain espéré de l'échange est nul.
\item[(c)] Montrez que la question (b) ne liquide pas entièrement le casse-tête. Construisez une
loi pour $ X $ telle que, après avoir ouvert votre enveloppe et vu la somme $ a $,
l'espérance conditionnelle de l'autre enveloppe dépasse strictement $ a $ --- et cela pour
toute valeur de $ a $ que vous pourriez voir. Démontrez ensuite que toute loi ayant cette
propriété vérifie $ E[X] = \infty $, et comparez la situation au paradoxe de
Saint-Pétersbourg.
\end{enumerate}

\end{problem}

\noindent La forme la plus ancienne en est le paradoxe des cravates de Maurice
Kraitchik, en 1943, où deux hommes soutiennent chacun que la cravate de l'autre a plus de
chances d'être la plus chère ; la version aux cartes d'Erwin Schrödinger est parvenue à
l'imprimé via le \textit{Miscellany} de Littlewood en 1953, Martin Gardner a popularisé un jeu
de portefeuilles en 1958 en admettant qu'il ne savait pas dire proprement ce qui clochait,
et Barry Nalebuff a donné la formulation moderne aux enveloppes en 1989. Ce qui vaut à ce
casse-tête sa longue vie, c'est la question (c) : l'argument naïf n'est pas seulement
bâclé, il est \textit{presque} réalisable, et les lois qui le réalisent sont exactement celles
d'espérance infinie, là où la valeur espérée cesse d'être un guide pour l'action. La
littérature n'est toujours pas unanime sur la phrase du sophisme qui est la coupable, bien
que les mathématiques soient entièrement réglées dès qu'un modèle est fixé --- bonne
illustration de la différence entre une question de probabilité et une question de
mots.

\begin{refs}
\item Contexte : Paradoxe des deux enveloppes --- Wikipédia (\url{https://fr.wikipedia.org/wiki/Paradoxe_des_deux_enveloppes})
\item Contexte : Paradoxe de Saint-Pétersbourg --- Wikipédia (\url{https://fr.wikipedia.org/wiki/Paradoxe_de_Saint-P%C3%A9tersbourg})
\item Article : Barry Nalebuff, \og{} Puzzles: The other person's envelope is always greener \fg{}, \textit{Journal of Economic Perspectives} 3 (1989), 171--181
\item Voir aussi : Probabilités (\url{applied-probability.html}), pour le paradoxe de Saint-Pétersbourg traité en entier
\item Voir aussi : PUZ-41 (\url{puzzles.html#PUZ-41}), la version dans laquelle vous ouvrez d'abord l'enveloppe --- là où le paradoxe devient réellement subtil.
\end{refs}

\bigskip

\begin{problem}[{Le jeu de Penney --- Lycée}]
\textbf{PUZ-46.} Le joueur A annonce une suite de trois résultats de pile ou face --- disons
PFP --- et la montre au joueur B, qui annonce alors une autre suite de trois. Une pièce
équilibrée est lancée de façon répétée jusqu'à ce que l'une des deux suites apparaisse en
trois lancers consécutifs ; celui qui l'a annoncée gagne.
\begin{enumerate}
\item[(a)] Démontrez que B, qui joue en second, peut toujours répondre par une suite qui apparaît
la première avec une probabilité strictement supérieure à $ 1/2 $, quel qu'ait été le
choix de A. Déterminez la meilleure réponse à chacune des huit suites de longueur trois, et
démontrez chacune de vos huit affirmations.
\item[(b)] Déduisez-en que \og{} bat \fg{} n'est pas une relation transitive sur les suites : exhibez
$ S_1, S_2, S_3, S_4 $ de longueur trois, chacune battant la suivante et $ S_4 $ battant
$ S_1 $. Démontrez ensuite qu'il n'existe aucun cycle de longueur trois parmi les huit
suites, de sorte que quatre est le plus court cycle disponible.
\item[(c)] Calculez deux fois l'une des probabilités de victoire, par des voies véritablement
différentes, et vérifiez que les réponses concordent : une fois par analyse au premier pas
sur l'automate fini dont les états enregistrent quelle portion de chaque suite est
actuellement appariée, et une fois par un argument de martingale du type employé pour les
temps d'attente de motifs.
\item[(d)] Montrez que le phénomène ne porte pas seulement sur les temps d'attente. Démontrez que,
parmi les suites de longueur trois, aucun renversement ne se produit : si $ S $ a un temps
d'attente espéré strictement plus court que $ T $ considéré isolément, alors $ T $ ne
gagne pas la course en face-à-face. Trouvez ensuite des suites $ S $ et $ T $ de
longueur \textit{quatre} pour lesquelles $ S $ a le temps d'attente espéré le plus court et
pourtant $ T $ apparaît la première avec une probabilité strictement supérieure à
$ 1/2 $, et expliquez quelle quantité, autre que les deux temps d'attente, décide de la
course.
\end{enumerate}

\end{problem}

\noindent Walter Penney a publié ce jeu dans le \textit{Journal of Recreational
Mathematics} en octobre 1969, et Martin Gardner l'a répandu dans \textit{Scientific
American} comme fleuron de ses paradoxes non transitifs. Son charme est d'être un jeu à
pièce équilibrée dans lequel le second joueur possède un avantage décisif, ce qui heurte
l'instinct même que heurtent les dés intransitifs : nous attendons de \og{} a plus de chances de
gagner contre \fg{} que ce soit un ordre, et ce n'en est pas un. Les mathématiques
sous-jacentes tiennent à la corrélation entre motifs qui se chevauchent --- PPP porte PP à la
fois comme préfixe et comme suffixe, et se fait donc concurrence à elle-même, alors que PPF
non --- et cette structure de chevauchement est exactement ce que l'algorithme du nombre
directeur de Conway condense en une formule. La question (d) est la forme la plus nette de
la surprise : temps d'attente espéré et victoire en face-à-face sont des comparaisons
différentes, et une suite peut gagner l'une et perdre l'autre.

\begin{refs}
\item Contexte : Jeu de Penney --- Wikipédia (en anglais) (\url{https://en.wikipedia.org/wiki/Penney%27s_game})
\item Contexte : Jeu intransitif --- Wikipédia (en anglais) (\url{https://en.wikipedia.org/wiki/Intransitive_game})
\item Lecture : Martin Gardner, \textit{The Colossal Book of Mathematics} (chapitre sur les paradoxes non transitifs)
\item Voir aussi : Probabilités (\url{applied-probability.html}), pour les temps d'attente de motifs par arrêt optionnel, et pour les dés intransitifs
\end{refs}

\bigskip

\begin{problem}[{Quatorze, quinze, et le coup qui ne vient jamais --- Lycée, short}]
\textbf{PUZ-47.} Le taquin contient des tuiles numérotées de $ 1 $ à
$ 15 $ dans un cadre $ 4 \times 4 $ comportant une case vide, et un coup fait glisser
orthogonalement une tuile dans la case vide. Démontrez qu'en partant des tuiles rangées dans
l'ordre, la case vide en dernier, aucune suite de coups n'atteint jamais la position où
$ 14 $ et $ 15 $ sont échangés et où tout le reste est à sa place initiale.
\end{problem}

\noindent Noyes Palmer Chapman, receveur des postes à Canastota, dans l'État de New
York, semble en avoir fabriqué la première version vers 1874 ; Sam Loyd s'en est attribué
l'invention à partir de 1891 et a offert un prix de mille dollars pour une solution à la
position où précisément $ 14 $ et $ 15 $ sont permutés, laquelle est impossible ---
l'argent n'a jamais couru le moindre risque. Woolsey Johnson et William Story l'avaient déjà
démontré en 1879, dans l'une des toutes premières applications de la signature d'une
permutation à une question extérieure à l'algèbre : exactement la moitié des $ 16! $
dispositions sont atteignables, et les deux moitiés ne se rencontrent jamais. L'outil à
saisir est l'homomorphisme $ \operatorname{sgn} : S_{16} \to \{\pm 1\} $, assorti de
l'observation qu'un simple glissement fait quelque chose de prévisible à la position de la
case vide autant qu'à la permutation.

\begin{refs}
\item Contexte : Taquin --- Wikipédia (\url{https://fr.wikipedia.org/wiki/Taquin})
\item Contexte : Signature d'une permutation --- Wikipédia (\url{https://fr.wikipedia.org/wiki/Signature_d%27une_permutation})
\item Article : Wm. W. Johnson \& W. E. Story, \og{} Notes on the ``15'' puzzle \fg{}, \textit{American Journal of Mathematics} 2 (1879), 397--404
\item Article : Aaron F. Archer, \og{} A modern treatment of the 15 puzzle \fg{}, \textit{American Mathematical Monthly} 106 (1999), 793--799
\end{refs}

\bigskip

\begin{problem}[{Le jeu de Nim, premier jeu jamais résolu --- Lycée}]
\textbf{PUZ-48.} Le Nim se joue avec plusieurs tas de jetons. Un coup retire un nombre
strictement positif de jetons d'un seul tas, et le joueur qui prend le tout dernier jeton
gagne. Pour des tas de tailles $ a_1, \dots, a_k $, définissez la \textit{somme de Nim}
$ a_1 \oplus \cdots \oplus a_k $ comme le résultat de l'écriture des tailles en binaire et
de l'addition des colonnes modulo $ 2 $, sans retenue.
\begin{enumerate}
\item[(a)] Démontrez le théorème de Bouton : le joueur qui doit jouer perd sous jeu parfait si et
seulement si $ a_1 \oplus \cdots \oplus a_k = 0 $. Votre démonstration doit également
produire le coup gagnant.
\item[(b)] Déduisez-en le cas à deux tas, et énoncez la stratégie obtenue en une phrase qu'un
enfant pourrait suivre et utiliser.
\item[(c)] Résolvez le jeu de soustraction à un tas où un coup retire $ 1 $, $ 2 $ ou $ 3 $
jetons : déterminez quelles tailles de tas sont perdantes pour le joueur qui doit jouer,
démontrez-le, et généralisez au retrait d'entre $ 1 $ et $ m $ jetons.
\item[(d)] Nim misère : c'est maintenant le joueur qui prend le dernier jeton qui \textit{perd}.
Déterminez les positions perdantes, démontrez votre réponse, et identifiez précisément
quelle étape de votre argument de la question (a) s'effondre et où les deux jeux commencent
à diverger.
\end{enumerate}

\end{problem}

\noindent Charles Bouton, à Harvard, a publié la théorie complète dans les
\textit{Annals of Mathematics} en 1901-1902, et a donné son nom au jeu ; des variantes se
jouaient depuis très longtemps, et le jeu ressemble de près au \textit{ji\v{a}n shíz\v{\i}} chinois,
\og{} ramasser des pierres \fg{}, bien que l'origine ne soit pas établie. L'article de Bouton
compte bien au-delà du jeu : c'est la première fois que quelqu'un prend un jeu à deux
joueurs de pure adresse et en produit une stratégie complète, démontrable et calculable, et
il l'a fait en trouvant un invariant arithmétique --- la somme de Nim --- qu'aucun coup ne peut
préserver à moins que celui qui joue ne le veuille. L'idée s'est révélée si féconde que tout
jeu impartial se révèle être un tas de Nim déguisé, ce qui fait le contenu de PUZ-52.
Westinghouse a construit le Nimatron, une machine à relais qui jouait au Nim contre les
visiteurs de l'Exposition universelle de New York, l'une des toutes premières machines à
jouer jamais construites.

\begin{refs}
\item Contexte : Jeux de Nim --- Wikipédia (\url{https://fr.wikipedia.org/wiki/Jeux_de_Nim})
\item Contexte : Nimber --- Wikipédia (\url{https://fr.wikipedia.org/wiki/Nimber})
\item Article : C. L. Bouton, \og{} Nim, a game with a complete mathematical theory \fg{}, \textit{Annals of Mathematics} (2) 3 (1901--1902), 35--39
\item Lecture : Elwyn Berlekamp, John Conway \& Richard Guy, \textit{Winning Ways for Your Mathematical Plays}, vol. 1
\end{refs}

\bigskip

\begin{problem}[{La barque, l'ancre et le niveau du bassin --- Lycée}]
\textbf{PUZ-59.} Une petite barque flotte dans un bassin fermé de section horizontale
$ A $ ; aucune eau n'entre ni ne sort. Dans la barque repose une ancre en fer de masse
$ m $ et de masse volumique $ \rho_a $, supérieure à la masse volumique $ \rho_w $ de
l'eau. Le batelier soulève l'ancre par-dessus bord ; elle coule et vient reposer au fond du
bassin.
\begin{enumerate}
\item[(a)] Déterminez, avec démonstration, si le niveau de l'eau dans le bassin monte, descend ou
reste le même, et donnez exactement la variation de niveau en fonction de $ m $,
$ \rho_a $, $ \rho_w $ et $ A $.
\item[(b)] Refaites le calcul lorsque l'objet jeté par-dessus bord est un bloc de liège de masse
$ m $ et de masse volumique inférieure à $ \rho_w $, qui s'éloigne en flottant à la
surface. Démontrez votre réponse.
\item[(c)] Refaites-le encore lorsque l'ancre n'est pas lâchée mais descendue au bout d'une corde
et laissée suspendue à la barque, sans toucher le fond. Déterminez le niveau dans ce cas et
démontrez-le.
\item[(d)] Enfin, la barque prend l'eau et coule au fond avec l'ancre encore à bord. Déterminez la
variation de niveau, puis énoncez un principe unique dont découlent les quatre réponses,
avec une démonstration que c'est bien le cas.
\end{enumerate}

\end{problem}

\noindent Archimède a écrit \textit{Des corps flottants} vers 250 av. J.-C., et c'est
le plus ancien ouvrage conservé qui mérite le nom de physique mathématique : il part d'un
postulat sur la pression dans un fluide et en déduit, comme théorèmes, les conditions dans
lesquelles les corps flottent, coulent et reposent stablement. La proposition 5 du livre I ---
qu'un corps flottant déplace son propre poids de fluide --- est tout ce casse-tête, et la
difficulté n'est pas la physique mais la comptabilité, car \og{} déplace son propre poids \fg{} et
\og{} déplace son propre volume \fg{} sont des énoncés différents et le casse-tête passe de l'un à
l'autre sans prévenir. Les architectes navals appellent encore \textit{déplacement} la masse
d'un navire, précisément pour cette raison. Un lecteur qui a terminé les quatre questions
devrait pouvoir dire aussitôt, et démontrer, ce qu'il advient du niveau de la mer quand la
glace de mer flottante fond, et pourquoi une calotte reposant sur la terre ferme est une
autre question --- ce qui n'est pas une curiosité mais l'un des calculs porteurs de la science
du climat.

\begin{refs}
\item Contexte : Poussée d'Archimède --- Wikipédia (\url{https://fr.wikipedia.org/wiki/Pouss%C3%A9e_d%27Archim%C3%A8de})
\item Contexte : Flottabilité --- Wikipédia (\url{https://fr.wikipedia.org/wiki/Flottabilit%C3%A9})
\item Lecture : Archimède, \textit{Des corps flottants}, livre I --- dans T. L. Heath (dir.), \textit{The Works of Archimedes} (1897)
\item Lecture : Jearl Walker, \textit{The Flying Circus of Physics} --- un livre entier de questions de cette forme
\end{refs}

\bigskip

\begin{problem}[{Le singe et le chasseur --- Lycée}]
\textbf{PUZ-60.} Une sarbacane est pointée \textit{directement} sur un singe suspendu à une
branche, à une distance horizontale $ d $ et à une hauteur $ h $ au-dessus de la bouche
du canon. À l'instant où la fléchette quitte le canon à la vitesse $ v_0 $, le singe lâche
prise et tombe. La pesanteur est uniforme, d'intensité $ g $, et le sol se trouve à une
hauteur $ H $ sous la branche.
\begin{enumerate}
\item[(a)] En travaillant dans le référentiel du sol, résolvez les deux trajectoires et déterminez,
avec démonstration, l'ensemble exact des vitesses initiales $ v_0 $ pour lesquelles la
fléchette atteint le singe avant que le singe n'atteigne le sol. Exprimez la condition sous
forme d'inégalité en $ d $, $ h $, $ H $ et $ g $.
\item[(b)] Passez maintenant à un référentiel lui-même en chute libre, et obtenez le même résultat
une seconde fois. Identifiez précisément l'étape où votre argument utilise l'égalité des
masses inerte et gravitationnelle, et dites ce qui changerait si la fléchette et le singe
avaient des rapports différents entre les deux.
\item[(c)] Rétablissez la résistance de l'air sous la forme d'un frottement linéaire, de sorte que
chaque corps subisse une accélération $ -b\vec{v} $ avec la même constante $ b $.
Déterminez si la conclusion de la question (a) survit, et démontrez votre réponse.
\item[(d)] Remplacez le frottement linéaire par un frottement quadratique, une accélération
$ -c\,|\vec{v}|\,\vec{v} $ avec la même constante $ c $ pour les deux corps. Montrez par
un argument explicite que le raisonnement de la question (b) échoue désormais, et calculez
la distance de manqué au premier ordre en $ c $.
\end{enumerate}

\end{problem}

\noindent C'est la plus ancienne démonstration de cours de physique encore
pratiquée : un électroaimant retient un singe en peluche, la sarbacane coupe le circuit en
tirant, et toute la salle regarde. Son contenu est la décomposition galiléenne du mouvement
des projectiles dans les \textit{Discorsi} de 1638 en un mouvement horizontal uniforme composé
avec un mouvement vertical uniformément accéléré --- la première fois que quelqu'un traitait
une trajectoire courbe comme la somme de deux trajectoires simples. La question (b) est une
miniature de ce qu'Einstein appelait en 1907 la pensée la plus heureuse de sa vie : un
observateur en chute libre ne sent pas son propre poids, et la gravitation peut donc être
localement éliminée par changement de référentiel. Les questions (c) et (d) sont là parce
que l'astuce de la chute libre est souvent enseignée comme si elle était inconditionnelle,
et elle ne l'est pas : elle survit exactement aux forces qui dépendent des deux corps de la
même manière, et il revient au lecteur de découvrir par lui-même quelle loi de frottement
qualifie. Ce site ne vous dira pas ce que montre la démonstration ; demandez à un professeur
de physique de la faire, et prédisez le résultat par écrit d'abord.

\begin{refs}
\item Contexte : Le singe et le chasseur --- Wikipédia (en anglais) (\url{https://en.wikipedia.org/wiki/Monkey_and_hunter})
\item Contexte : Principe d'équivalence --- Wikipédia (\url{https://fr.wikipedia.org/wiki/Principe_d%27%C3%A9quivalence})
\item Cours : MIT OCW 8.01SC --- Classical Mechanics (\url{https://ocw.mit.edu/courses/8-01sc-classical-mechanics-fall-2016/})
\item Voir aussi : Physique mathématique (\url{applied-physics.html}), pour les nombres impairs de Galilée et la naissance de la cinématique
\end{refs}

\bigskip

\begin{problem}[{La bicyclette tirée en arrière par sa pédale basse --- Lycée, short}]
\textbf{PUZ-61.} Une bicyclette est debout, maintenue verticale mais
libre de rouler, une pédale au point le plus bas de son cercle ; une corde attachée à cette
pédale est tirée horizontalement vers l'arrière, assez doucement pour que les deux roues
roulent sans glisser et que la machine ne bascule jamais. Déterminez, avec démonstration,
dans quel sens la bicyclette se déplace, et trouvez la condition exacte, portant sur la
longueur de manivelle $ r $, le rayon $ R $ de la roue motrice et le rapport de
transmission $ g $ --- tours de pédale par tour de roue ---, qui décide de la réponse.
\end{problem}

\noindent D. E. Daykin, à Reading, a mis la question par écrit sous le titre \og{} The
bicycle problem \fg{} dans \textit{Mathematics Magazine} en janvier 1972 ; Martin Gardner l'a
reprise dans \textit{Mathematical Carnival} (1975) et a rapporté que ses lecteurs se
divisaient nettement, plusieurs d'entre eux ne changeant d'avis qu'après être allés à
l'atelier pour essayer. L'outil qui tranche est l'axe instantané de rotation d'une roue qui
roule --- le point du solide qui est momentanément au repos --- d'où la vitesse de tout autre
point découle par un simple produit vectoriel ; de manière équivalente, on peut se demander
quelle courbe la pédale décrit par rapport au sol pendant que la bicyclette roule, à savoir
une trochoïde, et si cette courbe recule jamais. La condition demandée dans la seconde
moitié de la question est le contenu véritable : le casse-tête a un seuil, et une bicyclette
dont le braquet le dépasse ne se comporte pas comme une bicyclette dont le braquet lui est
inférieur, raison pour laquelle discuter de \og{} la \fg{} réponse sans nommer un braquet, c'est
discuter de rien.

\begin{refs}
\item Contexte : Axe instantané de rotation --- Wikipédia (\url{https://fr.wikipedia.org/wiki/Axe_instantan%C3%A9_de_rotation})
\item Contexte : Trochoïde --- Wikipédia (\url{https://fr.wikipedia.org/wiki/Trocho%C3%AFde})
\item Article : D. E. Daykin, \og{} The bicycle problem \fg{}, \textit{Mathematics Magazine} 45 (1972), no 1, p. 1
\item Lecture : Martin Gardner, \textit{Mathematical Carnival} (1975), p. 182
\end{refs}

\bigskip

\begin{problem}[{Le slinky qui tombe --- Lycée}]
\textbf{PUZ-62.} Un slinky de masse totale $ M $ et de constante de raideur $ k $ pend
au repos par sa spire supérieure, étiré par son propre poids, et est lâché sans vitesse
initiale à $ t = 0 $. Modélisez-le par un ressort sans masse de constante $ k $ portant
une masse répartie uniformément selon une coordonnée $ s \in [0,1] $ qui étiquette les
spires de haut en bas.
\begin{enumerate}
\item[(a)] Trouvez la configuration d'équilibre avant le lâcher : la position de la spire $ s $
en fonction de $ s $, et l'allongement statique total en fonction de $ M $, $ g $ et
$ k $.
\item[(b)] Démontrez qu'à partir de l'instant du lâcher le centre de masse obéit exactement à
\[ x_{\mathrm{cm}}(t) \;=\; x_{\mathrm{cm}}(0) + \tfrac{1}{2} g t^{2} \]
et déduisez-en une contrainte rigoureuse reliant le mouvement de la spire supérieure à celui
de la spire inférieure.
\item[(c)] Écrivez l'équation du mouvement des perturbations longitudinales de ce modèle,
identifiez la vitesse à laquelle elles se propagent, et calculez le temps qu'une
perturbation met pour aller de la spire supérieure à la spire inférieure. Comparez ce temps
à celui que le slinky met pour tomber de sa propre longueur étirée initiale.
\item[(d)] Déterminez, avec démonstration, ce que fait la spire inférieure entre l'instant du
lâcher et l'instant où la spire supérieure la rejoint. Énoncez exactement sur quelle
idéalisation du modèle repose votre conclusion, et décrivez en quoi un slinky réel --- dont
les spires se heurtent et ne peuvent pas se traverser --- doit s'écarter du modèle.
\end{enumerate}

\end{problem}

\noindent Richard James, ingénieur naval travaillant sur des suspensions à ressort
pour instruments de bord, a fait tomber un ressort de torsion d'une étagère en 1943 et l'a
regardé marcher ; lui et sa femme Betty en ont vendu quatre cents en quatre-vingt-dix
minutes chez Gimbels, à Philadelphie, en 1945. La version en chute a été analysée par M. G.
Calkin dans l'\textit{American Journal of Physics} en 1993, puis de nouveau, avec vidéo à
haute cadence et un bien meilleur traitement des spires, par R. C. Cross et M. S. Wheatland
dans la même revue en 2012 --- leur apport est que les spires situées derrière la perturbation
mettent un temps fini à se refermer, là où le modèle antérieur rendait cela instantané. Ce
qui vaut à cette question une page ici, c'est que deux arguments complètement différents s'y
appliquent et doivent être conciliés : le théorème du centre de masse, qui est exact et dit
quelque chose de l'objet entier, et l'équation des ondes, qui est approchée et dit quelque
chose des parties. La question (d) est là où ils se rencontrent. Prédisez la réponse, puis
filmez un slinky avec un téléphone en cadence rapide ; l'équipement coûte quelques euros et
l'expérience prend un après-midi.

\begin{refs}
\item Contexte : Slinky --- Wikipédia (\url{https://fr.wikipedia.org/wiki/Slinky})
\item Contexte : Équation des ondes --- Wikipédia (\url{https://fr.wikipedia.org/wiki/%C3%89quation_des_ondes})
\item Article : M. G. Calkin, \og{} Motion of a falling spring \fg{}, \textit{American Journal of Physics} 61 (1993), 261--264
\item Article : R. C. Cross \& M. S. Wheatland, \og{} Modeling a falling slinky \fg{}, \textit{American Journal of Physics} 80 (2012), 1051--1060 --- arXiv:1208.4629 (\url{https://arxiv.org/abs/1208.4629}) (attention : le résumé donne la réponse à la question (d))
\end{refs}

\bigskip

\begin{problem}[{L'oiseau buveur, ou comment encadrer une estimation --- Lycée}]
\textbf{PUZ-63.} L'oiseau buveur est un jouet de verre scellé : une tête recouverte de
feutre, un bulbe formant le corps, et un tube qui les relie, contenant du dichlorométhane et
sa propre vapeur, et rien d'autre. Mouillez la tête de feutre, posez l'oiseau à côté d'un
verre d'eau, et il bascule en avant, trempe son bec, se redresse, et recommence --- pendant
des jours.
\begin{enumerate}
\item[(a)] Identifiez la source chaude et la source froide du cycle et exprimez chaque température
en fonction de grandeurs que vous pourriez mesurer dans la pièce. Établissez la borne de
Carnot sur le rendement de toute machine fonctionnant entre elles, et évaluez-la pour une
humidité intérieure typique.
\item[(b)] Estimez la puissance mécanique que délivre le jouet. Choisissez vos données d'entrée ---
masse, géométrie du pivot, amplitude, période --- et, pour chacune, donnez un intervalle que
vous êtes prêt à défendre plutôt qu'un nombre unique ; propagez les intervalles dans le
calcul et donnez la puissance en watts avec des bornes explicites.
\item[(c)] Estimez le débit d'évaporation de l'eau à la tête, convertissez-le en puissance
thermique, et bornez ainsi le rendement réel du jouet. Comparez avec la question (a),
donnez le rapport avec son incertitude, et identifiez la donnée d'entrée unique qui domine
la largeur de votre intervalle.
\item[(d)] Prédisez, avec justification, ce que fait l'oiseau dans chacune des situations
suivantes, et concevez une expérience qui distinguerait votre prédiction de sa négation :
enfermé dans un bocal d'air déjà saturé de vapeur d'eau ; dans un fort courant d'air ; avec
le verre d'eau retiré ; dans une enceinte à vide.
\end{enumerate}

\end{problem}

\noindent Miles V. Sullivan, chimiste aux Bell Telephone Laboratories, a breveté la
forme moderne en 1946 (brevet américain 2 402 463), et Arthur M. Hillery en avait breveté
une version l'année précédente ; mais le jouet est bien plus ancien que l'un et l'autre, et
l'on a montré à Albert et Elsa Einstein un \og{} petit oiseau insatiable \fg{} chinois à leur
arrivée à Shanghai en 1922. L'oiseau est le meilleur argument de bureau qu'une machine
thermique a besoin d'une \textit{différence} de température, et non simplement de chaleur, et
trouver d'où vient cette différence --- elle n'est fournie par rien que l'on branche --- est
tout le casse-tête. Les questions (b) et (c) forment un problème de Fermi au sens strict où
Enrico Fermi l'entendait. On se souvient de lui pour avoir demandé combien d'accordeurs de
piano travaillent à Chicago, mais ce qu'il enseignait réellement était la discipline du
report des bornes : donnez un intervalle pour chaque entrée, propagez-les, et sachez ensuite
dire quelle entrée vous devriez aller mesurer sérieusement. À Trinity, en juillet 1945, il a
laissé tomber des bouts de papier au passage de l'onde de choc et a estimé sur-le-champ la
puissance de la première bombe atomique, à un facteur près que les mesures instrumentées ont
ensuite respecté.

\begin{refs}
\item Contexte : Oiseau buveur --- Wikipédia (\url{https://fr.wikipedia.org/wiki/Oiseau_buveur})
\item Contexte : Estimation de Fermi --- Wikipédia (\url{https://fr.wikipedia.org/wiki/Estimation_de_Fermi})
\item Contexte : Cycle de Carnot --- Wikipédia (\url{https://fr.wikipedia.org/wiki/Cycle_de_Carnot})
\item Voir aussi : Physique mathématique (\url{applied-physics.html}), pour la thermodynamique derrière la borne de Carnot
\end{refs}

\bigskip

\begin{problem}[{Le cube peint --- Lycée}]
\textbf{PUZ-73.} Un cube plein de côté $ n \ge 2 $ est construit à partir de
$ n^{3} $ cubes unités. Toute sa surface extérieure est peinte, puis il est de nouveau
démonté en cubes unités.
\begin{enumerate}
\item[(a)] Déterminez, avec démonstration, combien de petits cubes portent de la peinture sur
exactement trois faces, sur exactement deux, sur exactement une, et sur aucune, chacun en
fonction de $ n $. Vérifiez que vos quatre décomptes totalisent $ n^{3} $.
\item[(b)] Démontrez algébriquement l'identité
\[ n^{3} \;=\; (n-2)^{3} + 6(n-2)^{2} + 12(n-2) + 8 \]
puis dites exactement quel trait du cube compte chacun des quatre termes. Généralisez à
$ d $ dimensions : démontrez que le nombre de cellules unités du $ d $-cube de côté
$ n $ qui touchent le bord exactement selon une face de dimension $ k $ vaut
$ 2^{\,d-k}\binom{d}{k}(n-2)^{k} $, et déduisez-en l'identité
\[ n^{d} \;=\; \sum_{k=0}^{d} \binom{d}{k} 2^{\,d-k} (n-2)^{k} . \]
\item[(c)] Une autre question sur le même objet. Une coupe est une unique tranche plane à travers
tout ce qui se trouve alors devant la lame ; entre deux coupes, les morceaux peuvent être
réempilés à votre guise. Déterminez, avec démonstration, le plus petit nombre de coupes qui
divise un cube $ 3 \times 3 \times 3 $ en ses $ 27 $ cubes unités, puis déterminez ce
plus petit nombre pour un cube $ n \times n \times n $ général.
\end{enumerate}

\end{problem}

\noindent Les questions (a) et (b) sont le même énoncé raconté deux fois, et c'est
tout l'intérêt de les mettre côte à côte : un exercice de comptage qu'un enfant peut faire
en regardant un cube se révèle être la décomposition en couches d'un hypercube, où le bord
du $ d $-cube se scinde en $ 2^{\,d-k}\binom{d}{k} $ faces de chaque dimension $ k $.
Ludwig Schläfli a dégagé cette structure combinatoire pour les polytopes en toutes
dimensions entre 1850 et 1852, dans un manuscrit si en avance sur son public qu'il n'a été
publié intégralement qu'en 1901, neuf ans après sa mort ; l'identité binomiale de la
question (b) est l'ombre arithmétique de sa classification. La question (c) est d'une tout
autre nature et forme la moitié la plus difficile. Produire une suite courte de coupes est
affaire d'essais, mais montrer qu'aucune suite n'est plus courte oblige à maîtriser d'un
coup tout réempilement possible, et ce genre de minoration se démontre presque toujours en
trouvant une caractéristique du travail achevé dont une seule coupe ne peut fournir qu'une
quantité bornée. Notez ce qui change d'une question à l'autre : (a) et (b) portent sur un
objet fixé, tandis que (c) porte sur une recherche non bornée parmi des procédures, et c'est
la seconde qui a fait de l'optimisation combinatoire une discipline.

\begin{refs}
\item Contexte : Hypercube --- Wikipédia (\url{https://fr.wikipedia.org/wiki/Hypercube})
\item Contexte : Formule du binôme de Newton --- Wikipédia (\url{https://fr.wikipedia.org/wiki/Formule_du_bin%C3%B4me_de_Newton})
\item Lecture : H. S. M. Coxeter, \textit{Regular Polytopes} (3e éd., Dover, 1973), chap. VII sur Schläfli
\item Voir aussi : Combinatoire et théorie des graphes (\url{pure-combinatorics.html}), Géométrie (\url{pure-geometry.html})
\end{refs}

\bigskip

\begin{problem}[{Les poignées de main, et ce qu'une soirée ne peut pas faire --- Lycée}]
\textbf{PUZ-74.} Lors d'une réunion, certaines paires de personnes se serrent la main.
Personne ne se serre la main à soi-même, et aucune paire de personnes ne se serre la main
plus d'une fois.
\begin{enumerate}
\item[(a)] Démontrez que le nombre de personnes qui serrent un nombre impair de mains est pair.
Énoncez le résultat comme un théorème sur les graphes et démontrez-le dans ce langage.
\item[(b)] Démontrez que si au moins deux personnes sont présentes, alors deux d'entre elles
serrent le même nombre de mains.
\item[(c)] Voici maintenant $ n $ couples mariés qui se rencontrent, l'hôte et l'hôtesse compris,
de sorte que $ 2n $ personnes sont présentes. Personne ne serre la main de son propre
conjoint. L'hôte demande à chacune des $ 2n-1 $ autres personnes combien de mains elle a
serrées, et reçoit $ 2n-1 $ réponses différentes. Déterminez, avec démonstration, combien
de mains l'hôtesse a serrées, en fonction de $ n $, puis déterminez combien l'hôte
lui-même en a serrées.
\item[(d)] Décidez, avec démonstration, si la réponse à la question (c) reste forcée lorsque la
règle sur les conjoints est supprimée, et si elle reste forcée lorsque l'hôte est autorisé à
figurer parmi les personnes qu'il interroge.
\end{enumerate}

\end{problem}

\noindent La question (a) est le lemme des poignées de main, et c'est le plus
ancien théorème de la théorie des graphes : Euler a démontré la formule de la somme des
degrés dans l'article de 1736 sur les ponts de Königsberg qui a fondé la discipline, et tout
argument de parité ultérieur sur les graphes en descend. La question (b) est le principe des
tiroirs appliqué à un ensemble de degrés possibles trop petit d'un élément, et si elle vaut
d'être traitée avec soin, c'est que le principe des tiroirs évident échoue d'exactement un ---
il vous faut remarquer pourquoi deux degrés particuliers ne peuvent pas coexister. La
question (c) circule partout dans les recueils de problèmes sous le nom de \og{} problème de la
soirée \fg{}, sans attribution établie, et elle mérite sa réputation : les données semblent bien
trop maigres pour déterminer quoi que ce soit, la réponse se révèle indépendante de savoir
quel couple est lequel, et la démonstration naturelle est une récurrence sur $ n $ où l'on
épluche les réponses extrêmes deux par deux. La question (d) est là parce qu'un casse-tête
dont vous n'avez pas testé les hypothèses est un casse-tête que vous n'avez pas compris ;
découvrez laquelle des conditions est porteuse.

\begin{refs}
\item Contexte : Lemme des poignées de main --- Wikipédia (\url{https://fr.wikipedia.org/wiki/Lemme_des_poign%C3%A9es_de_main})
\item Contexte : Principe des tiroirs --- Wikipédia (\url{https://fr.wikipedia.org/wiki/Principe_des_tiroirs})
\item Lecture : Miklós Bóna, \textit{A Walk Through Combinatorics}, chapitres sur la récurrence et sur les graphes
\item Voir aussi : Combinatoire et théorie des graphes (\url{pure-combinatorics.html}), L'art de la démonstration (\url{proofs.html})
\end{refs}

\bigskip

\begin{problem}[{Les zéros à la fin d'une factorielle --- Lycée, short}]
\textbf{PUZ-75.} Déterminez, avec démonstration, combien de zéros
figurent à la fin du développement décimal de $ 100! $, et donnez une formule pour le
nombre de zéros terminaux de $ n! $ valable pour tout entier strictement positif $ n $.
Démontrez ensuite que tout entier positif ou nul n'apparaît pas comme un tel décompte, et
caractérisez exactement quels entiers apparaissent.
\end{problem}

\noindent La formule qu'on vous demande est celle de Legendre : l'exposant d'un
nombre premier $ p $ dans $ n! $ vaut
\[ \nu_p(n!) \;=\; \sum_{i \ge 1} \left\lfloor \frac{n}{p^{i}} \right\rfloor
\;=\; \frac{n - s_p(n)}{p - 1}, \]
où $ s_p(n) $ est la somme des chiffres de $ n $ en base $ p $ ; Adrien-Marie Legendre
l'a consignée dans sa \textit{Théorie des nombres}, et on l'appelle aussi formule de Polignac.
La seconde forme est la plus intéressante, car elle convertit une question de divisibilité
en une question de chiffres, et la dernière partie du problème --- quels décomptes sont sautés
--- se résout entièrement en regardant les développements en base cinq. La même idée de
comptage de chiffres est le théorème de Kummer de 1852, selon lequel l'exposant de $ p $
dans $ \binom{m+n}{n} $ est le nombre de retenues lorsqu'on additionne $ m $ et $ n $
en base $ p $, l'un des faits les plus élégants de la théorie élémentaire des nombres, et
invisible tant qu'on n'écrit pas la factorielle de la seconde façon.

\begin{refs}
\item Contexte : Formule de Legendre --- Wikipédia (\url{https://fr.wikipedia.org/wiki/Formule_de_Legendre})
\item Contexte : Zéro terminal --- Wikipédia (en anglais) (\url{https://en.wikipedia.org/wiki/Trailing_zero})
\item Voir aussi : Théorie des nombres (\url{pure-number-theory.html})
\end{refs}

\bigskip

\begin{problem}[{L'araignée et la mouche --- Lycée}]
\textbf{PUZ-76.} Une pièce rectangulaire mesure $ 30 $ pieds de long, $ 12 $ pieds de
large et $ 12 $ pieds de haut. Une araignée se tient sur l'un des murs d'extrémité
$ 12 \times 12 $, au milieu de ce mur et à un pied sous le plafond. Une mouche se tient
sur le mur d'extrémité opposé, au milieu et à un pied au-dessus du sol, et ne bouge pas.
L'araignée peut ramper sur les murs, le sol et le plafond, mais ne peut pas se laisser
tomber dans les airs.
\begin{enumerate}
\item[(a)] Déterminez, avec démonstration, la longueur du plus court chemin de l'araignée à la
mouche.
\item[(b)] Démontrez le principe sur lequel repose tout le problème : un plus court chemin sur la
surface d'une boîte est un segment de droite dans tout dépliage plan de la suite de faces
qu'il traverse. Énumérez ensuite toutes les suites de faces qu'un candidat au plus court
chemin pourrait emprunter, dépliez chacune, et justifiez que votre liste est complète ---
c'est là, et non dans l'arithmétique, que sont les mathématiques.
\item[(c)] Une fourmi marche sur la surface d'un cube unité, d'un sommet au sommet diagonalement
opposé. Déterminez la longueur d'un plus court chemin, et déterminez combien de plus courts
chemins distincts il existe.
\item[(d)] Montrez que le dépliage évident n'est pas toujours le bon : trouvez une boîte de côtés
$ a, b, c $ et deux points de sa surface pour lesquels le plus court chemin de surface ne
se trouve \textit{pas} dans le dépliage qu'une première tentative choisirait. Expliquez
ensuite précisément pourquoi une géodésique sur une boîte ne peut jamais passer par un
sommet, et pourquoi ce fait rend le problème combinatoire plutôt qu'affaire de calcul
différentiel.
\end{enumerate}

\end{problem}

\noindent Henry Ernest Dudeney a posé ce problème dans un journal anglais en 1903,
et c'est le plus connu de ses casse-tête géométriques pour la bonne raison que la première
réponse de presque tout le monde est fausse, et fausse par un chemin qui paraît manifestement
optimal. La technique du dépliage est bien plus ancienne que le casse-tête :
l'\textit{Underweysung der Messung} d'Albrecht Dürer, en 1525, dessine les solides platoniciens
et plusieurs solides archimédiens comme des patrons plans à découper et à plier,
apparemment la première fois que quelqu'un mettait des polyèdres sur le papier de cette
manière. Ce que Dudeney ajoute, c'est l'observation qu'un patron n'est pas seulement une aide
à la fabrication mais un changement de coordonnées dans lequel \og{} le plus court \fg{} devient
\og{} tout droit \fg{}. La question (d) ouvre sur un sujet vivant --- savoir si tout polyèdre convexe
possède seulement un patron qui ne se recouvre pas lui-même est un problème non résolu,
repris sur cette page en PUZ-86 (\url{puzzles.html#PUZ-86}) --- et il vaut la peine de savoir que
l'innocente affaire d'aplatir une boîte résiste aux mathématiciens depuis cinq siècles.

\begin{refs}
\item Contexte : Patron (géométrie) --- Wikipédia (\url{https://fr.wikipedia.org/wiki/Patron_%28g%C3%A9om%C3%A9trie%29}), y compris sa section sur les plus courts chemins
\item Contexte : Spider and Fly Problem --- Wolfram MathWorld (en anglais) (\url{https://mathworld.wolfram.com/SpiderandFlyProblem.html}) (attention : donne la réponse)
\item Lecture : Henry E. Dudeney, \textit{Amusements in Mathematics} (1917) --- son propre recueil des casse-tête de journaux
\item Lecture : Erik D. Demaine \& Joseph O'Rourke, \textit{Geometric Folding Algorithms: Linkages, Origami, Polyhedra} (Cambridge, 2007), partie III
\item Voir aussi : PUZ-86 (\url{puzzles.html#PUZ-86}) ci-dessous, et Géométrie (\url{pure-geometry.html})
\end{refs}

\bigskip

\begin{problem}[{Quand les aiguilles d'une horloge se superposent --- Lycée, short}]
\textbf{PUZ-77.} Sur une horloge analogique ordinaire, les deux
aiguilles tournent régulièrement, la grande aiguille exactement douze fois plus vite que la
petite. Déterminez, avec démonstration, tous les instants d'une période de douze heures où
les deux aiguilles pointent exactement dans la même direction, démontrez que votre liste est
complète, puis décidez --- avec démonstration --- si une trotteuse tournant régulièrement les
rejoint jamais à un instant autre que midi.
\end{problem}

\noindent Les casse-tête sur les aiguilles d'une horloge comptent parmi les plus
anciens de la littérature populaire, et Dudeney comme Sam Loyd en ont imprimé plusieurs.
Dépouillée du cadran, la première moitié est un énoncé sur des angles modulo un tour complet :
les deux aiguilles se superposent exactement lorsque la différence de leurs angles est un
multiple entier de $ 2\pi $, de sorte qu'une question géométrique devient une congruence
linéaire et que les coïncidences tombent à des instants régulièrement espacés dont
l'espacement n'est pas un nombre entier de minutes. Ce dernier fait est ce qui surprend, et
ce qui fait d'un cadran d'horloge une correcte première rencontre avec l'arithmétique modulo
un. La question des trois aiguilles est d'une tout autre nature : elle demande si trois
fonctions affines du temps peuvent partager une solution commune, et une fois les deux
conditions écrites vous constaterez que la réponse tient à une question de rationalité
plutôt qu'à quoi que ce soit d'horloger.

\begin{refs}
\item Contexte : Problème de l'angle des aiguilles --- Wikipédia (en anglais) (\url{https://en.wikipedia.org/wiki/Clock_angle_problem})
\item Contexte : Arithmétique modulaire --- Wikipédia (\url{https://fr.wikipedia.org/wiki/Arithm%C3%A9tique_modulaire})
\item Voir aussi : Théorie des nombres (\url{pure-number-theory.html})
\end{refs}

\bigskip

\begin{problem}[{Le plateau amputé et le triomino en L --- Lycée}]
\textbf{PUZ-78.} Un \textit{triomino en L} est la forme de trois carrés unités joints en L ;
un \textit{triomino en I} en compte trois alignés. Un plateau est \textit{amputé} lorsqu'on lui a
retiré exactement une case. Les pièces peuvent être tournées, ne peuvent pas se chevaucher,
et doivent recouvrir chaque case restante exactement une fois.
\begin{enumerate}
\item[(a)] Démontrez que, pour tout $ n \ge 1 $ et pour tout choix de la case retirée, le plateau
amputé $ 2^{n} \times 2^{n} $ peut être pavé par des triominos en L. Votre démonstration
doit être une récurrence dont vous puissiez décrire le pas inductif en une phrase.
\item[(b)] Déterminez, avec démonstration, exactement quels rectangles $ m \times n $ --- sans case
retirée --- peuvent être pavés par des triominos en L. Les deux sens sont exigés : une
construction pour les rectangles qui marchent et un argument d'impossibilité pour les
autres.
\item[(c)] Déterminez, avec démonstration, exactement quels plateaux amputés $ m \times n $
peuvent être pavés par des triominos en L. Là où la réponse dépend de la case
\textit{précisément} retirée, dites exactement comment, et démontrez-le.
\item[(d)] Remplacez maintenant le L par le triomino droit en I. Déterminez quels rectangles
$ m \times n $ il pave, et démontrez que le théorème de la question (a) échoue pour lui :
trouvez, pour chaque $ n \ge 1 $, les positions de la case retirée pour lesquelles le
plateau amputé $ 2^{n} \times 2^{n} $ n'admet aucun pavage par des triominos en I. Le
coloriage qui tranche n'est pas le noir et blanc.
\end{enumerate}

\end{problem}

\noindent Solomon Golomb a démontré l'énoncé de la question (a) en 1954, à l'époque
même où il forgeait le mot \og{} polyomino \fg{} ; Martin Gardner l'a porté à un très large lectorat,
et il est depuis la réclame de référence, en classe, pour la récurrence, parce que le pas
inductif contient exactement une idée et que cette idée est visible dès qu'on l'a. Ce qui
rend le reste de ce problème digne d'être posé, c'est que (a) n'est pas représentative. Le
cas $ 2^{n} \times 2^{n} $ est le cas facile ; les rectangles généraux demandent une
construction plus un véritable argument d'impossibilité, et la version amputée de la
question (c) demande les deux ensemble et comporte des exceptions. La question (d) est la
leçon la plus tranchante qui soit sur les arguments de coloriage : deux formes ayant le même
nombre de cases peuvent se comporter tout autrement, et l'invariant qui tue l'une est
invisible au coloriage ordinaire du damier. Comparez PUZ-08 (\url{#PUZ-08}) et
PUZ-22 (\url{#PUZ-22}) sur cette page, où la même difficulté apparaît pour la pièce
$ 1 \times 4 $ et pour le tétromino en T.

\begin{refs}
\item Contexte : Triomino --- Wikipédia (\url{https://fr.wikipedia.org/wiki/Triomino})
\item Contexte : Polyomino --- Wikipédia (\url{https://fr.wikipedia.org/wiki/Polyomino})
\item Lecture : Solomon W. Golomb, \textit{Polyominoes: Puzzles, Patterns, Problems, and Packings} (2e éd., Princeton, 1994)
\item Voir aussi : Combinatoire et théorie des graphes (\url{pure-combinatorics.html}) pour le cas $ 2^{n} \times 2^{n} $ à lui seul, et PUZ-22 (\url{#PUZ-22}) ci-dessus
\end{refs}

\bigskip

\begin{problem}[{Trois interrupteurs, une seule visite --- Lycée, short}]
\textbf{PUZ-87.} Trois interrupteurs situés au rez-de-chaussée sont
reliés un à un à trois lampes à incandescence identiques dans un grenier sans fenêtre ; vous
pouvez man\oe{}uvrer les interrupteurs aussi longtemps que vous le voulez, selon le motif que
vous voulez, mais vous ne pouvez monter au grenier qu'une seule fois et ne jamais y
retourner. Trouvez une procédure qui identifie toujours quel interrupteur commande quelle
lampe, démontrez qu'elle est correcte, puis déterminez le plus grand nombre $ n $ de
lampes qu'une visite unique pourrait résoudre si chaque lampe présente $ m $ états
distinguables à un observateur qui entre.
\end{problem}

\noindent Ce casse-tête relève du folklore --- il circule comme question d'entretien
technique depuis au moins les années 1990 et figure dans la plupart des recueils modernes
sans auteur attribué --- mais la comptabilité qui le sous-tend est le registre même qui régit
PUZ-01 et PUZ-33. Vous n'avez droit qu'à un regard, de sorte que tout ce que vous faites
auparavant doit organiser les choses pour que l'unique observation que vous finirez par faire
prenne une valeur différente pour chacun des $ n! $ câblages possibles ; la seconde moitié
de la question vous demande d'écrire cette condition sous forme d'inégalité et de la
résoudre. Si ce casse-tête est plus qu'un exercice de comptage, c'est qu'il vous force à dire
ce qu'\textit{est} une observation. Un canal porte l'information dont ses sorties peuvent
réellement être distinguées, non l'information que son concepteur voulait lui faire porter,
et cet écart n'est pas une plaisanterie : c'est tout le sujet des attaques par canal
auxiliaire en cryptographie, où l'on a lu des secrets dans la durée, la consommation
électrique, le son et la température de machines qui ne transmettaient rien du tout.

\begin{refs}
\item Contexte : Théorie de l'information --- Wikipédia (\url{https://fr.wikipedia.org/wiki/Th%C3%A9orie_de_l%27information})
\item Contexte : Attaque par canal auxiliaire --- Wikipédia (\url{https://fr.wikipedia.org/wiki/Attaque_par_canal_auxiliaire})
\item Voir aussi : PUZ-01 (\url{#PUZ-01}) et PUZ-33 (\url{#PUZ-33}) pour la même borne de comptage, et Cryptographie et théorie de l'information (\url{applied-crypto.html})
\end{refs}

\bigskip

\begin{problem}[{Les pirates se partagent le butin --- Lycée}]
\textbf{PUZ-88.} Cinq pirates, hiérarchisés par ancienneté et se connaissant les uns les
autres, doivent partager $ 100 $ pièces d'or indivisibles. Le pirate survivant le plus
ancien propose un partage ; tous les pirates survivants, le proposant compris, votent
ensuite. Si au moins la moitié vote pour, le partage est adopté et la partie s'arrête ; sinon
le proposant est jeté par-dessus bord et le suivant dans l'ordre d'ancienneté fait une
proposition, sous la même règle. Chaque pirate est parfaitement rationnel, sait que tous les
autres le sont, et préfère --- dans cet ordre de priorité --- survivre, puis recevoir plus d'or,
puis, toutes choses égales par ailleurs, voir un compagnon jeté par-dessus bord.
\begin{enumerate}
\item[(a)] Déterminez, avec démonstration, ce que propose le pirate le plus ancien et montrez que
c'est accepté. Construisez l'argument de bas en haut à partir du cas d'un seul pirate.
\item[(b)] Démontrez que votre réponse est l'unique équilibre parfait en sous-jeux du jeu sous forme
extensive à information parfaite correspondant. Montrez ensuite que la préférence de
troisième rang fait un vrai travail : remplacez \og{} préfère voir un pirate jeté par-dessus
bord \fg{} par \og{} préfère ne pas le voir \fg{}, et déterminez comment la réponse change.
\item[(c)] Généralisez à $ n $ pirates et $ g $ pièces. Déterminez, avec démonstration, pour
quels $ n $ exactement le proposant survit, et décrivez le motif de survie une fois que
$ n $ dépasse $ 2g $.
\item[(d)] Changez la règle de vote de sorte qu'une proposition ne passe qu'à la majorité stricte
des survivants, la voix du proposant ne départageant plus les égalités. Refaites (a) et (c)
sous cette règle et dites précisément quelle étape de votre argument précédent échoue.
\end{enumerate}

\end{problem}

\noindent Ian Stewart a présenté ce casse-tête à un large public dans sa chronique
de \textit{Scientific American} de mai 1999, \og{} A puzzle for pirates \fg{}, où il rapportait aussi
l'extension de Steve Omohundro à un nombre quelconque de pirates et de pièces ; une version
était parue l'année précédente dans la contribution de Bruce Talbot Coram à \textit{The Theory of
Institutional Design}. Les mathématiques en jeu sont le raisonnement rétrograde sur un jeu
fini à information parfaite --- la méthode qu'Ernst Zermelo a introduite en 1913 dans le
premier théorème jamais démontré sur les échecs --- et le casse-tête est la démonstration la
plus nette de ce que le raisonnement rétrograde donne vu de l'intérieur : le pouvoir du
pirate le plus ancien sur les autres est entièrement emprunté à un sous-jeu qui ne sera
jamais joué, et le vote de chaque pirate est émis contre une hypothèse. Le résultat est
célèbre pour être à la fois inattaquable et incroyable, ce qui lui vaut de figurer aux côtés
du jeu du mille-pattes dans la littérature sur la question de savoir si les gens raisonnent
réellement à rebours. La question (c) est là où cela cesse d'être une anecdote : la réponse
pour $ n $ général a une véritable structure, et la trouver exige d'organiser la récurrence
plutôt que d'allonger le tableau.

\begin{refs}
\item Contexte : Jeu du pirate --- Wikipédia (\url{https://fr.wikipedia.org/wiki/Jeu_du_pirate})
\item Contexte : Raisonnement rétrograde --- Wikipédia (\url{https://fr.wikipedia.org/wiki/Raisonnement_r%C3%A9trograde})
\item Lecture : Ian Stewart, \og{} A puzzle for pirates \fg{}, \textit{Scientific American}, mai 1999 (attention : la chronique traite intégralement le cas des cinq pirates)
\item Voir aussi : Optimisation et théorie des jeux (\url{applied-optimization.html})
\end{refs}

\bigskip

\begin{problem}[{Les deux tiers de la moyenne --- Lycée}]
\textbf{PUZ-89.} Chacun de $ n \ge 3 $ joueurs écrit simultanément et indépendamment un
nombre réel de $ [0, 100] $. Le prix va à celui qui est le plus proche des
$ \tfrac{2}{3} $ de la moyenne des $ n $ nombres, les égalités étant partagées à parts
égales. Chaque joueur veut le prix et sait que tous les autres le veulent aussi.
\begin{enumerate}
\item[(a)] Déterminez, avec démonstration, tous les équilibres de Nash en stratégies pures, puis
tous les équilibres de Nash en stratégies mixtes.
\item[(b)] Démontrez que le jeu se résout par élimination itérée des stratégies strictement
dominées. Déterminez exactement quelle part de l'intervalle survit après $ r $ tours
d'élimination, et combien de tours sont nécessaires pour atteindre l'ensemble
d'équilibre.
\item[(c)] Modélisez des raisonneurs imparfaits. Dites qu'un joueur est de \textit{niveau $ 0 $}
s'il tire son nombre uniformément dans $ [0,100] $, et de \textit{niveau $ k $} s'il joue
une meilleure réponse à une population entièrement composée de joueurs de niveau
$ (k-1) $. Calculez explicitement le choix de niveau $ k $ en fonction de $ k $ et de
$ n $, et déterminez ce que la moyenne d'une population observée de choix vous apprend sur
la distribution des niveaux qui la composent.
\item[(d)] Remplacez $ \tfrac{2}{3} $ par une constante $ p > 0 $. Déterminez, avec
démonstration, l'ensemble des équilibres de Nash pour chaque tel $ p $, et décidez pour
quels $ p $ le jeu se résout encore par élimination itérée des stratégies strictement
dominées. Traitez $ p = 1 $ et $ p > 1 $ séparément et expliquez pourquoi ces cas sont
réellement différents.
\end{enumerate}

\end{problem}

\noindent Alain Ledoux a inventé ce jeu en 1981 comme épreuve de départage pour le
magazine français \textit{Jeux et Stratégie}, en demandant à quelque quatre mille lecteurs de
deviner un entier entre un et un milliard ; l'article de Rosemarie Nagel paru en 1995 dans
l'\textit{American Economic Review} en a fait l'instrument de laboratoire standard pour mesurer
ce que les économistes appellent la profondeur de raisonnement, et le quotidien danois
\textit{Politiken} l'a lancé en 2005 avec 19 196 participations. Son ancêtre est le chapitre 12
de la \textit{Théorie générale} de Keynes (1936), où l'investissement professionnel est comparé
à un concours de journal dans lequel il faut désigner non pas les plus jolis visages, mais
ceux que l'on croit que les autres concurrents désigneront --- \og{} et il en est, je crois, qui
pratiquent le quatrième degré, le cinquième et au-delà \fg{}. L'analyse d'équilibre des questions
(a) et (b) est un exercice court. Ce qui vaut à ce jeu sa célébrité, c'est que presque
personne ne joue l'équilibre, et que la distance entre les deux n'est pas du bruit mais une
quantité mesurable : la question (c) vous demande de construire l'instrument qui la
mesure.

\begin{refs}
\item Contexte : p-Concours de beauté (deviner les 2/3 de la moyenne) --- Wikipédia (\url{https://fr.wikipedia.org/wiki/P-Concours_de_beaut%C3%A9})
\item Contexte : Concours de beauté de Keynes --- Wikipédia (\url{https://fr.wikipedia.org/wiki/Concours_de_beaut%C3%A9_de_Keynes})
\item Article : Rosemarie Nagel, \og{} Unraveling in guessing games: an experimental study \fg{}, \textit{American Economic Review} 85 (1995), 1313--1326
\item Voir aussi : Optimisation et théorie des jeux (\url{applied-optimization.html})
\end{refs}

\bigskip

\begin{problem}[{Le jeu de Sim --- Lycée}]
\textbf{PUZ-90.} Six points sont tracés en position générale, et les $ 15 $ segments
qui les joignent deux à deux sont laissés incolores. Deux joueurs alternent, Rouge
commençant ; un coup colorie un segment incolore dans la couleur du joueur qui joue. Un
joueur qui complète un triangle dont les trois côtés sont de sa propre couleur perd
immédiatement.
\begin{enumerate}
\item[(a)] Démontrez que la partie ne peut jamais être nulle. C'est-à-dire : démontrez que tout
coloriage des arêtes du graphe complet $ K_6 $ en deux couleurs contient un triangle dont
les trois arêtes sont de la même couleur ; puis exhibez un coloriage de $ K_5 $ en deux
couleurs sans un tel triangle, de sorte que le nombre de Ramsey $ R(3,3) $ vaut exactement
$ 6 $.
\item[(b)] Déduisez de la question (a) et du théorème de Zermelo que l'un des deux joueurs possède
une stratégie qui gagne contre toute défense. Déterminez lequel, et décrivez une analyse
exhaustive qui trancherait la question ; le point a été réglé pour la première fois par
ordinateur en 1974, et une stratégie assez simple pour être menée par une personne n'a été
publiée qu'en 2020.
\item[(c)] Jouez la même partie sur cinq points. Démontrez qu'une partie nulle est désormais
possible, et déterminez, avec démonstration, l'issue de la partie à cinq points sous jeu
parfait.
\item[(d)] Généralisez. Énoncez la version jouée sur $ K_n $ dans laquelle un joueur perd en
complétant un $ K_r $ monochrome, déterminez pour quels couples $ (n, r) $ la partie ne
peut pas être nulle, puis expliquez soigneusement pourquoi savoir qu'une partie ne peut pas
être nulle ne vous apprend absolument rien sur celui qui la gagne.
\end{enumerate}

\end{problem}

\noindent Gustavus Simmons --- cryptographe aux Sandia National Laboratories, plus
connu pour avoir fondé la théorie des codes d'authentification --- a publié \og{} The game of SIM \fg{}
dans le \textit{Journal of Recreational Mathematics} en 1969. Ce qu'il avait bâti est la plus
petite conséquence jouable du théorème de Ramsey : la propriété d'absence de nulle de la
question (a) est précisément le fait mondain que, parmi six personnes quelconques, trois se
connaissent mutuellement ou trois sont mutuellement étrangères. Frank Ramsey a démontré le
théorème général dans un article de 1930 sur un problème de décision en logique formelle, où
il n'était qu'un lemme ; il est mort la même année, à vingt-six ans, et la discipline
combinatoire qui porte aujourd'hui son nom est née tout entière de ce lemme. La question (d)
est la morale honnête. La théorie de Ramsey garantit que de la structure doit apparaître, et
le garantit sans la produire --- Erd\H{o}s avait coutume de dire que si des extraterrestres
exigeaient $ R(5,5) $ nous devrions mobiliser les ordinateurs du monde et l'obtenir, et
que s'ils exigeaient $ R(6,6) $ nous devrions attaquer les extraterrestres --- et le même
divorce entre existence et construction traverse PUZ-56 et PUZ-92 sur cette page.

\begin{refs}
\item Contexte : Sim (jeu de papier-crayon) --- Wikipédia (\url{https://fr.wikipedia.org/wiki/Sim_%28jeu%29})
\item Contexte : Théorème de Ramsey --- Wikipédia (\url{https://fr.wikipedia.org/wiki/Th%C3%A9or%C3%A8me_de_Ramsey})
\item Article : G. J. Simmons, \og{} The game of SIM \fg{}, \textit{Journal of Recreational Mathematics} 2 (1969), no 2
\item Voir aussi : Combinatoire et théorie des graphes (\url{pure-combinatorics.html})
\end{refs}

\bigskip

\begin{problem}[{La majorité qui tourne en rond --- Lycée, short}]
\textbf{PUZ-91.} Trois électeurs classent trois candidats : le premier
préfère $ A \succ B \succ C $, le deuxième $ B \succ C \succ A $, le troisième
$ C \succ A \succ B $. Vérifiez qu'une majorité stricte préfère $ A $ à $ B $, qu'une
majorité stricte préfère $ B $ à $ C $, et qu'une majorité stricte préfère $ C $ à
$ A $ ; déterminez ensuite, avec démonstration, le plus petit nombre de candidats et le
plus petit nombre d'électeurs pour lesquels un tel cycle de majorités deux à deux peut se
produire.
\end{problem}

\noindent Le marquis de Condorcet a publié ceci dans son \textit{Essai sur l'application
de l'analyse à la probabilité des décisions rendues à la pluralité des voix} en 1785, un
livre écrit pour donner aux élections un fondement mathématique ; il était mathématicien,
abolitionniste de la première heure et partisan du vote des femmes, et il est mort en prison
en 1794, sous la Terreur. Ramon Llull avait trouvé le même cycle au treizième siècle en
concevant des procédures d'élection des dignitaires de l'Église, dans des manuscrits restés
perdus jusqu'au vingt et unième. Le contenu, c'est que \og{} la volonté de la majorité \fg{} peut ne
définir aucun ordre : la préférence majoritaire peut ne pas être transitive, de sorte qu'il
peut n'exister aucun candidat que l'électorat préfère collectivement à tous les autres. Ce
n'est pas non plus un accident rare --- sous le modèle standard où le classement de chaque
électeur est tiré uniformément et indépendamment, la probabilité d'un cycle à trois candidats
tend vers environ neuf pour cent à mesure que l'électorat grandit, et croît avec le nombre de
candidats. Tout ce qui est dans PUZ-98 naît de cet unique exemple.

\begin{refs}
\item Contexte : Paradoxe de Condorcet --- Wikipédia (\url{https://fr.wikipedia.org/wiki/Paradoxe_de_Condorcet})
\item Contexte : Théorie du choix social --- Wikipédia (\url{https://fr.wikipedia.org/wiki/Th%C3%A9orie_du_choix_social})
\item Voir aussi : PUZ-98 (\url{puzzles.html#PUZ-98}) ci-dessous, et Optimisation et théorie des jeux (\url{applied-optimization.html})
\end{refs}

\bigskip


\section*{Problèmes ouverts}

\begin{problem}[{Les trois premiers zéros, à partir d'une table --- Lycée, short}]
\textbf{OPEN-43.} Dans une table publiée des zéros de la fonction zêta
de Riemann (Andrew Odlyzko tient en ligne les tables de référence), relevez les parties
imaginaires des trois premiers zéros non triviaux --- vous devriez trouver environ
$ 14{,}1347 $, $ 21{,}0220 $ et $ 25{,}0109 $ --- et vérifiez que la table donne
chaque zéro avec une partie réelle exactement égale à $ \tfrac{1}{2} $. Expliquez
ensuite en une phrase pourquoi $ 10^{13} $ confirmations de ce genre ne démontreraient
toujours pas l'hypothèse de Riemann.
\end{problem}

\noindent Savoir où vivent les données est une vraie compétence. L'étude
numérique des zéros de zêta va des calculs à la main de Riemann lui-même à Alan Turing,
qui a calculé des zéros sur le Manchester Mark 1 en 1950, jusqu'aux immenses tables
modernes d'Odlyzko. Manipuler trois nombres réels rend OPEN-01 concret comme aucun résumé
ne saurait le faire.

\begin{refs}
\item Contexte : Hypothèse de Riemann --- Wikipédia (\url{https://fr.wikipedia.org/wiki/Hypoth%C3%A8se_de_Riemann})
\item Tables : Andrew Odlyzko --- Tables of zeros of the Riemann zeta function (\url{https://www-users.cse.umn.edu/~odlyzko/zeta_tables/index.html})
\end{refs}

\bigskip

\begin{problem}[{L'application $ 3n+1 $ : calculez, et voyez le problème --- Lycée}]
\textbf{OPEN-17.} Définissez l'application de Collatz sur les entiers strictement
positifs par
\[ T(n) = \begin{cases} n/2, & n \text{ pair},\\ 3n+1, & n \text{ impair}.\end{cases} \]
\begin{enumerate}
\item[(a)] Calculez la trajectoire complète de $ n = 27 $ jusqu'à ce qu'elle atteigne $ 1 $ ;
notez le nombre de pas qu'elle demande et la plus grande valeur atteinte. (Vous devriez
trouver que le voyage est bien plus long et bien plus haut que la taille de $ 27 $ ne le
laisse supposer --- c'est tout l'intérêt.)
\item[(b)] Définissez le \textit{temps d'arrêt total} de $ n $ comme le nombre de pas nécessaires
pour atteindre $ 1 $. Tabulez-le pour $ n = 1, \dots, 30 $ et trouvez quelle valeur
initiale inférieure à $ 30 $ met le plus longtemps.
\item[(c)] Expliquez avec vos propres mots pourquoi un calcul, aussi loin qu'il soit poussé, ne
pourra jamais démontrer la conjecture.
\end{enumerate}

\end{problem}

\noindent Lothar Collatz a consigné le problème dans ses carnets en 1937 ; il a
depuis été rattaché aux noms d'Ulam, Kakutani, Thwaites, Hasse et Syracuse, signe du
nombre de personnes qui l'ont redécouvert indépendamment. Kakutani plaisantait en disant
que c'était un complot pour ralentir la recherche mathématique américaine. La trajectoire
de $ 27 $ est la démonstration classique : elle grimpe jusqu'à $ 9232 $ et prend
$ 111 $ pas, sans raison visible. Toute valeur jusqu'à environ $ 2^{68} $, soit à peu
près $ 2{,}95 \times 10^{20} $, a été vérifiée comme atteignant $ 1 $ --- et cette
vérification ne nous dit rien sur le nombre suivant.

\begin{refs}
\item Contexte : Conjecture de Syracuse (Collatz) --- Wikipédia (\url{https://fr.wikipedia.org/wiki/Conjecture_de_Syracuse})
\item Lecture : Jeffrey C. Lagarias (dir.), \textit{The Ultimate Challenge: The 3x+1 Problem}, AMS (2010)
\item Lecture : Jeffrey C. Lagarias, \og{} The 3x+1 problem: An annotated bibliography \fg{}, arXiv:math/0309224 (\url{https://arxiv.org/abs/math/0309224})
\end{refs}

\bigskip

\begin{problem}[{L'application $ 3n-1 $ a des cycles --- trouvez-les --- Lycée}]
\textbf{OPEN-18.} Changez un signe. Définissez
\[ S(n) = \begin{cases} n/2, & n \text{ pair},\\ 3n-1, & n \text{ impair}.\end{cases} \]
\begin{enumerate}
\item[(a)] Montrez que $ S $ possède un cycle contenant $ 5 $, en calculant la trajectoire de
$ 5 $ jusqu'à son retour.
\item[(b)] Trouvez un second cycle, plus long, en partant de $ 17 $.
\item[(c)] Expliquez maintenant la portée de tout cela : la conjecture $ 3n+1 $ affirme
qu'aucun cycle autre que $ 1 \to 4 \to 2 \to 1 $ n'existe, et voici une application
presque identique où cette affirmation est \textbf{fausse}. Que vous apprend cela sur toute
démonstration proposée de Collatz qui fonctionnerait aussi bien pour $ 3n-1 $ ?
\end{enumerate}

\end{problem}

\noindent C'est l'exercice le plus utile de la page, et il ne coûte que de
l'arithmétique. Tout argument en faveur de la conjecture $ 3n+1 $ doit, quelque part,
utiliser une propriété qui distingue $ +1 $ de $ -1 $ ; s'il ne le fait pas, le même
argument démontrerait quelque chose de faux. Des générations de \og{} démonstrations \fg{}
amateurs de Collatz meurent exactement ici, et confronter votre propre raisonnement à
l'application $ 3n-1 $ est le contrôle d'honnêteté le moins coûteux des mathématiques.
Notez aussi que $ 3n-1 $ sur les entiers positifs revient au même que $ 3n+1 $ sur les
négatifs, raison pour laquelle le comportement de l'application de Collatz sur les entiers
négatifs fait partie du tableau standard.

\begin{refs}
\item Contexte : Conjecture de Syracuse (Collatz) --- Wikipédia (\url{https://fr.wikipedia.org/wiki/Conjecture_de_Syracuse})
\item Lecture : Lagarias, \textit{The Ultimate Challenge}, panorama introductif
\end{refs}

\bigskip

\begin{problem}[{Le temps d'arrêt total de 97 --- Lycée, short}]
\textbf{OPEN-47.} À l'aide de l'application de Collatz $ T $
d'OPEN-17, calculez le temps d'arrêt total de $ n = 97 $ --- le nombre de pas nécessaires
pour atteindre $ 1 $ --- et notez la plus grande valeur que sa trajectoire atteint en
chemin. Comparez ces deux nombres à ceux qui correspondent à $ n = 27 $, et dites ce que
la comparaison suggère quant à la prédiction d'une trajectoire à partir de la taille de sa
valeur initiale.
\end{problem}

\noindent Le faire une fois, à la main ou en cinq lignes de code, vaut mieux que
lire n'importe quelle description du problème. Deux choses deviennent visibles : à quelle
hauteur une trajectoire peut errer au-dessus de son point de départ, et comment des
valeurs initiales différentes sont canalisées vers des chemins communs, de sorte que des
nombres tout à fait étrangers l'un à l'autre finissent par descendre la même échelle
jusqu'à $ 1 $. La suite des temps d'arrêt totaux est cataloguée sous OEIS A006577, et
son graphe --- hérissé, sans structure, n'obéissant à aucune formule que quiconque ait
trouvée --- est la raison pour laquelle Erd\H{o}s disait que les mathématiques ne sont
peut-être pas prêtes pour ce problème.

\begin{refs}
\item Contexte : Conjecture de Syracuse (Collatz) --- Wikipédia (\url{https://fr.wikipedia.org/wiki/Conjecture_de_Syracuse})
\item Données : OEIS A006577 (\url{https://oeis.org/A006577}) --- nombre de pas pour atteindre 1 dans le problème $ 3x+1 $
\end{refs}

\bigskip

\begin{problem}[{Goldbach à la main, de 4 à 60 --- Lycée, short}]
\textbf{OPEN-48.} Pour tout nombre pair $ n $ avec
$ 4 \leq n \leq 60 $, écrivez $ n $ comme somme de deux nombres premiers et notez
combien de représentations non ordonnées $ n = p + q $ avec $ p \leq q $ vous trouvez.
Décrivez ensuite ce qu'il advient de ce décompte quand $ n $ grandit, et énoncez
clairement pourquoi aucune quantité de vérifications de ce genre ne règle OPEN-23.
\end{problem}

\noindent C'est le calcul que faisaient Goldbach et Euler en 1742, et il demande
une vingtaine de minutes avec une liste des nombres premiers jusqu'à $ 60 $. Le nombre
de représentations ne se contente pas de rester positif, il croît --- son graphique est
connu sous le nom de comète de Goldbach, et sa forme est prédite assez précisément par
l'heuristique de la méthode du cercle de Hardy--Littlewood. Cette croissance explique
pourquoi presque tout le monde croit à la conjecture, et aussi pourquoi personne ne sait
la démontrer : un décompte asymptotique grand en moyenne ne dit rien sur la possibilité
qu'un nombre pair isolé soit malchanceux.

\begin{refs}
\item Contexte : Conjecture de Goldbach --- Wikipédia (\url{https://fr.wikipedia.org/wiki/Conjecture_de_Goldbach})
\item Contexte : Comète de Goldbach --- Wikipédia (\url{https://fr.wikipedia.org/wiki/Com%C3%A8te_de_Goldbach})
\end{refs}

\bigskip


\section*{L'art de la démonstration}

\begin{problem}[{Démonstration directe ou contraposée ? --- Lycée}]
\textbf{PRF-01.} Une implication \og{} si P alors Q \fg{} peut se démontrer directement
(on suppose P, on en déduit Q) ou par contraposée (on suppose Q fausse, on en déduit que
P est fausse). Pour chacun des énoncés ci-dessous, décidez laquelle des deux formes donne
la démonstration la plus nette, et rédigez-la.
\begin{enumerate}
\item[(a)] Si n est un entier impair, alors $ n^2 $ est impair.
\item[(b)] Si $ n^2 $ est pair pour un entier n, alors n est pair.
\item[(c)] Si x et y sont des réels strictement positifs avec $ xy >25 $, alors $ x >5 $ ou
$ y >5 $.
\item[(d)] Si la somme de deux entiers vaut au moins 100, alors l'un au moins des deux vaut au
moins 50.
\item[(e)] Formulez ensuite, en une ou deux phrases, une règle empirique pour reconnaître quand
la contraposée sera la voie la plus facile.
\end{enumerate}

\end{problem}

\noindent Qu'une implication et sa contraposée soient vraies ou fausses ensemble
était déjà clair pour Aristote, et ce fut un lieu commun de la logique médiévale
(\textit{modus tollens}). Choisir la bonne forme est la première décision véritable que prend
celui qui rédige une démonstration : un énoncé dont la conclusion contient un \og{} ou \fg{}, ou
dont l'hypothèse est difficile à saisir, devient souvent docile dès qu'on le retourne.
Toute introduction à la démonstration commence ici.

\begin{refs}
\item Contexte : Proposition contraposée --- Wikipédia (\url{https://fr.wikipedia.org/wiki/Proposition_contrapos%C3%A9e})
\item Lecture : D. J. Velleman, \textit{How to Prove It}, ch. 3
\item Lecture : R. Hammack, \textit{Book of Proof}, ch. 4--5 (gratuit en ligne (\url{https://richardhammack.github.io/BookOfProof/}))
\end{refs}

\bigskip

\begin{problem}[{La racine carrée de 2 est irrationnelle --- Lycée}]
\textbf{PRF-02.} Démontrez que $ \sqrt{2} $ est irrationnel : il
n'existe pas d'entiers strictement positifs p et q tels que $ \sqrt{2} = p/q $. (Forme
suggérée : le raisonnement par l'absurde --- supposez qu'une telle fraction existe et partez
en chasse d'une impossibilité.)
\end{problem}

\noindent Le plus vieux coup de tonnerre des mathématiques. Les pythagoriciens, pour
qui \og{} tout est nombre \fg{} voulait dire les entiers et leurs rapports, ont découvert que la
diagonale d'un carré est incommensurable avec son côté --- la légende veut que la découverte
ait été imputée à Hippase de Métaponte, noyé en mer pour l'avoir révélée. Aristote cite
l'argument pair-impair comme exemple type du raisonnement par l'impossible, et une version
en clôt le livre X des \textit{Éléments} d'Euclide. C'est le modèle du raisonnement par
l'absurde, et le problème suivant demande de quoi ce modèle est réellement fait.

\begin{refs}
\item Contexte : Racine carrée de deux --- Wikipédia (\url{https://fr.wikipedia.org/wiki/Racine_carr%C3%A9e_de_deux})
\item Lecture : G. H. Hardy \& E. M. Wright, \textit{An Introduction to the Theory of Numbers}, ch. 4
\item Lecture : R. Hammack, \textit{Book of Proof}, ch. 6
\end{refs}

\bigskip

\begin{problem}[{Deux autres irrationnels, et la forme du canevas --- Lycée}]
\textbf{PRF-03.} PRF-02 a démontré que $ \sqrt{2} $ est irrationnel ; l'argument
est un canevas réutilisable, et ce problème vous fait l'en extraire.
\begin{enumerate}
\item[(a)] Démontrez que $ \sqrt{3} $ est irrationnel.
\item[(b)] Démontrez que $ \log_2 3 $ est irrationnel.
\item[(c)] Prenez maintenant du recul : énoncez précisément le canevas commun aux trois
démonstrations (PRF-02 comprise) --- ce que l'on suppose, quel genre d'impossibilité on
produit, et quel fait arithmétique sur les entiers fait tourner le moteur.
\item[(d)] Faites tourner le canevas sur $ \sqrt{4} $ et désignez la ligne exacte où il casse.
S'il ne cassait pas, quelque chose irait très mal --- pourquoi ?
\end{enumerate}

\end{problem}

\noindent Une démonstration qu'on ne sait exécuter qu'une fois est un tour de main ;
une démonstration qu'on sait exécuter sur toute une famille est une méthode. Extraire le
canevas réutilisable d'un argument déjà rédigé --- et le mettre à l'épreuve sur un cas où la
conclusion est fausse --- voilà comment les mathématiciens apprennent réellement les
techniques. La question (b) est une favorite parce que le canevas se transplante sur un
énoncé de logarithmes où nulle racine carrée n'est en vue : l'impossibilité y devient un
conflit de parités entre puissances de 2 et puissances de 3.

\begin{refs}
\item Contexte : Raisonnement par l'absurde --- Wikipédia (\url{https://fr.wikipedia.org/wiki/Raisonnement_par_l%27absurde})
\item Lecture : R. Hammack, \textit{Book of Proof}, ch. 6
\item Cours : MIT OCW 6.042J Mathematics for Computer Science (en anglais) (\url{https://ocw.mit.edu/courses/6-042j-mathematics-for-computer-science-fall-2010/})
\end{refs}

\bigskip

\begin{problem}[{Récurrence : les n premiers nombres impairs --- Lycée}]
\textbf{PRF-04.} Pour tout entier strictement positif n,
\[ 1 + 3 + 5 + \cdots + (2n-1) = n^2 . \]
\begin{enumerate}
\item[(a)] Démontrez l'identité par récurrence, en énonçant intégralement l'initialisation et
l'hérédité.
\item[(b)] Trouvez ensuite une seconde démonstration qui n'utilise aucune récurrence --- une façon
de \textit{voir} l'identité --- et comparez ce que chacune des deux exige du lecteur.
\end{enumerate}

\end{problem}

\noindent L'identité elle-même est très ancienne : les pythagoriciens disposaient des
cailloux en équerres (gnomons) autour d'un carré et le regardaient croître. La récurrence
comme principe de démonstration explicitement énoncé remonte d'ordinaire à Francesco
Maurolico (1575) --- qui a démontré cette identité même --- et à Pascal ; le nom \og{} induction
mathématique \fg{} a été fixé par Augustus De Morgan en 1838. La récurrence est la manière
standard de gravir une échelle infinie barreau par barreau, et voici son premier exercice
traditionnel.

\begin{refs}
\item Contexte : Raisonnement par récurrence --- Wikipédia (\url{https://fr.wikipedia.org/wiki/Raisonnement_par_r%C3%A9currence})
\item Lecture : R. Hammack, \textit{Book of Proof}, ch. 10
\item Cours : MIT OCW 6.042J Mathematics for Computer Science (en anglais) (\url{https://ocw.mit.edu/courses/6-042j-mathematics-for-computer-science-fall-2010/})
\end{refs}

\bigskip

\begin{problem}[{Trouvez la faille : tous les chevaux sont de la même couleur --- Lycée}]
\textbf{PRF-05.} Voici une \og{} démonstration \fg{} par récurrence du fait que tous les
chevaux sont de la même couleur. \textit{Initialisation : tout ensemble de 1 cheval ne contient
que des chevaux d'une seule couleur. Hérédité : supposons que tout ensemble de n chevaux
soit monochrome. Étant donné n+1 chevaux, mettez-en un de côté ; les n restants partagent
une couleur. Remettez-le et mettez de côté un cheval différent ; de nouveau les n restants
partagent une couleur. Les deux groupes se recouvrent, donc les n+1 chevaux partagent une
seule et même couleur. Donc, par récurrence, tous les chevaux sont de la même couleur.}
La conclusion est absurde, donc la démonstration est fausse. Votre tâche n'est pas de le
constater, mais de localiser la défaillance.
\begin{enumerate}
\item[(a)] Identifiez la valeur exacte de n à laquelle l'argument d'hérédité casse pour la
première fois, et dites précisément quelle phrase y échoue, et pourquoi.
\item[(b)] Expliquez pourquoi \og{} mais les chevaux ont des couleurs variées \fg{} n'est pas une réponse
recevable.
\end{enumerate}

\end{problem}

\noindent Cette parabole a été popularisée par George Pólya dans les années 1950 (sa
version d'origine démontrait que toutes les jeunes filles ont les yeux de la même couleur).
C'est le vaccin classique contre une maladie précise : une hérédité valable pour les grands
n mais qui suppose en silence quelque chose de faux pour un unique petit n. Déboguer une
démonstration est un savoir-faire distinct de celui de l'écrire, et il s'apprend le mieux sur
des arguments conçus pour échouer.

\begin{refs}
\item Contexte : Tous les chevaux sont de la même couleur --- Wikipédia (en anglais) (\url{https://en.wikipedia.org/wiki/All_horses_are_the_same_color})
\item Lecture : G. Pólya, \textit{Induction and Analogy in Mathematics} (1954)
\end{refs}

\bigskip

\begin{problem}[{Trouvez la faille : une démonstration que 1 = 2 --- Lycée}]
\textbf{PRF-06.} Considérez la \og{} démonstration \fg{} classique suivante. \textit{Soient
$ a = b $ non nuls. Alors $ a^2 = ab $, donc $ a^2 - b^2 = ab - b^2 $, donc
$ (a-b)(a+b) = b(a-b) $, donc $ a + b = b $, donc $ 2b = b $, et par conséquent
$ 2 = 1 $.}
\begin{enumerate}
\item[(a)] Identifiez l'étape exacte qui est illicite et énoncez la règle générale qu'elle
viole.
\item[(b)] Énoncez la leçon sous forme d'un principe précis sur la manipulation des équations :
quelles opérations préservent la vérité sans condition, et lesquelles portent des
conditions annexes ?
\item[(c)] Construisez votre propre démonstration fallacieuse --- d'un énoncé faux de votre choix ---
qui dissimule le même coup illicite au moins deux étapes plus profond que celle-ci.
\end{enumerate}

\end{problem}

\noindent Les sophismes algébriques de ce genre circulent comme folklore scolaire
depuis le dix-neuvième siècle ; E. A. Maxwell en a rassemblé les classiques dans
\textit{Fallacies in Mathematics} (1959). Ce sont plus que des plaisanteries : chaque ligne
d'une vraie démonstration est une affirmation assortie d'une justification, et le plus
rapide moyen d'intérioriser cette discipline est de regarder une chaîne d'équations à l'air
assuré y faire entrer un mensonge en fraude.

\begin{refs}
\item Contexte : Pseudo-démonstration d'égalité entre nombres --- Wikipédia (\url{https://fr.wikipedia.org/wiki/Pseudo-d%C3%A9monstration_d%27%C3%A9galit%C3%A9_entre_nombres})
\item Lecture : E. A. Maxwell, \textit{Fallacies in Mathematics} (Cambridge, 1959)
\end{refs}

\bigskip

\begin{problem}[{Principe des tiroirs : cinq points entiers --- Lycée}]
\textbf{PRF-07.} Un point entier du plan est un point dont les coordonnées sont
entières.
\begin{enumerate}
\item[(a)] Démontrez que parmi 5 points entiers quelconques, il en existe deux dont le milieu est
lui-même un point entier.
\item[(b)] Montrez que 5 ne peut pas être remplacé par 4.
\item[(c)] Déterminez ensuite, avec démonstration, le plus petit nombre de points entiers de
l'espace de dimension trois qui force l'existence d'une telle paire.
\end{enumerate}

\end{problem}

\noindent Le principe des tiroirs --- si l'on range plus de k objets dans k tiroirs,
l'un des tiroirs en contient deux --- a l'air trop évident pour être une arme. Dirichlet a
montré le contraire, en se servant de son \textit{Schubfachprinzip} (\og{} principe des tiroirs \fg{},
1834) pour démontrer des faits profonds sur l'approximation des irrationnels par des
fractions. Tout l'art réside dans le choix des tiroirs ; ce classique du point milieu,
standard des olympiades depuis des décennies, en est l'illustration la plus limpide
possible.

\begin{refs}
\item Contexte : Principe des tiroirs --- Wikipédia (\url{https://fr.wikipedia.org/wiki/Principe_des_tiroirs})
\item Lecture : A. Engel, \textit{Problem-Solving Strategies}, ch. 4 (le principe des tiroirs)
\end{refs}

\bigskip

\begin{problem}[{Double dénombrement : le lemme des poignées de main --- Lycée}]
\textbf{PRF-08.} Démontrez que, dans toute assemblée, le nombre de
personnes ayant serré un nombre impair de mains est pair. (De façon équivalente : dans tout
graphe fini, la somme des degrés des sommets vaut deux fois le nombre d'arêtes, de sorte que
le nombre de sommets de degré impair est pair.) Arme suggérée : comptez une même quantité ---
les participations à une poignée de main --- de deux manières différentes.
\end{problem}

\noindent C'est le premier théorème de la théorie des graphes, et il figure dans le
premier article de la théorie des graphes : la résolution par Euler, en 1736, du problème des
ponts de Königsberg. La technique, le double dénombrement, est l'un des gestes les plus
élégants de la combinatoire --- une égalité démontrée non par manipulation, mais en décrivant
un même ensemble depuis deux points de vue. Si simple soit-il, le lemme des poignées de main
tue quantité de constructions à l'air plausible (essayez de dessiner un graphe ayant
exactement trois sommets de degré 3 et aucun autre).

\begin{refs}
\item Contexte : Lemme des poignées de main --- Wikipédia (\url{https://fr.wikipedia.org/wiki/Lemme_des_poign%C3%A9es_de_main})
\item Contexte : Preuve par double dénombrement --- Wikipédia (\url{https://fr.wikipedia.org/wiki/Preuve_par_double_d%C3%A9nombrement})
\item Cours : MIT OCW 6.042J Mathematics for Computer Science (en anglais) (\url{https://ocw.mit.edu/courses/6-042j-mathematics-for-computer-science-fall-2010/})
\end{refs}

\bigskip

\begin{problem}[{Invariants : l'échiquier mutilé --- Lycée}]
\textbf{PRF-09.} Retirez deux cases d'angle diagonalement opposées d'un échiquier
ordinaire 8 $\times$ 8, ce qui laisse 62 cases. Un domino recouvre exactement deux cases partageant
un côté.
\begin{enumerate}
\item[(a)] Démontrez que l'échiquier restant ne peut pas être pavé par 31 dominos.
\item[(b)] Tranchez ensuite la question réciproque : si l'on retire à la place une case noire et
une case blanche --- n'importe laquelle de ces paires --- les 62 cases restantes peuvent-elles
toujours être pavées ? Déterminez la réponse, avec démonstration.
\end{enumerate}

\end{problem}

\noindent Posé par le philosophe Max Black dans \textit{Critical Thinking} (1946) et
chéri depuis lors ; Gomory a répondu à la seconde question par un argument resté célèbre. La
méthode --- trouver une quantité ou une configuration que tout coup licite préserve, puis
montrer que le départ et l'arrivée n'en donnent pas la même valeur --- c'est le principe
d'invariance, une bête de somme qui va des mathématiques récréatives jusqu'aux lois de
conservation. John McCarthy a proposé l'échiquier mutilé à la jeune communauté de
l'intelligence artificielle en 1964 comme cas d'école, précisément parce que la recherche
exhaustive y est immense alors qu'une seule idée clôt l'affaire à l'instant.

\begin{refs}
\item Contexte : Problème de l'échiquier mutilé --- Wikipédia (\url{https://fr.wikipedia.org/wiki/Probl%C3%A8me_de_l%27%C3%A9chiquier_mutil%C3%A9})
\item Lecture : A. Engel, \textit{Problem-Solving Strategies}, ch. 1--2
\end{refs}

\bigskip

\begin{problem}[{Ce que la démonstration par l'exemple ne peut pas faire --- Lycée}]
\textbf{PRF-10.} Trois exercices sur ce que la vérification d'exemples peut et ne peut
pas établir.
\begin{enumerate}
\item[(a)] Énoncez précisément ce qui cloche dans la \og{} démonstration par l'exemple \fg{} : pourquoi
aucune liste finie de cas vérifiés ne peut-elle établir un énoncé de la forme \og{} pour tout
n \fg{} ? Dites aussi ce que la vérification d'exemples \textit{peut} légitimement démontrer.
\item[(b)] Le polynôme $ n^2 + n + 41 $ produit des nombres premiers pour
$ n = 0, 1, 2, \dots, 39 $ --- quarante d'affilée. Trouvez, avec démonstration, le plus
petit n pour lequel il échoue, et expliquez pourquoi il était condamné à échouer, sans rien
calculer.
\item[(c)] Fermat croyait que tout nombre $ F_k = 2^{2^k} + 1 $ est premier, et les cinq
premiers le sont. Montrez que $ 641 \mid F_5 $ en utilisant les identités
$ 641 = 5^4 + 2^4 = 5 \cdot 2^7 + 1 $ plutôt qu'une division posée.
\end{enumerate}

\end{problem}

\noindent Les exemples sont ce qui fait naître les conjectures, et les exemples sont
ce qui a tué celle-ci : en 1732 le jeune Euler a factorisé $ F_5 = 4294967297 $, démolissant
la conjecture à laquelle Fermat s'était engagé avec le plus d'assurance. (Le polynôme de la
question (b) est lui aussi d'Euler, et date de 1772.) \textit{Mathematics and Plausible
Reasoning} de Pólya est la méditation classique sur le juste rapport entre indices et
démonstration : les exemples suggèrent ; seul l'argument établit.

\begin{refs}
\item Contexte : Nombre de Fermat --- Wikipédia (\url{https://fr.wikipedia.org/wiki/Nombre_de_Fermat})
\item Lecture : G. Pólya, \textit{Mathematics and Plausible Reasoning} (1954)
\end{refs}

\bigskip

\begin{problem}[{Quand une figure est-elle une démonstration ? --- Lycée}]
\textbf{PRF-11.} La figure classique de réarrangement pour le théorème de Pythagore
place quatre copies d'un triangle rectangle de côtés de l'angle droit a et b et d'hypoténuse
c à l'intérieur d'un carré de côté $ a + b $, de deux manières différentes. L'une laisse à
découvert un unique carré incliné d'aire $ c^2 $ ; l'autre laisse deux carrés d'aires
$ a^2 $ et $ b^2 $.
\begin{enumerate}
\item[(a)] Dessinez les deux dispositions et rédigez, en phrases complètes, l'argument que les
figures encodent --- y compris chaque fait que la figure vous demande d'admettre sur parole
(que les pièces ne se chevauchent pas, que les régions découvertes sont bien des carrés,
que l'aire survit au réarrangement).
\item[(b)] Proposez des critères : quand un dessin doit-il compter comme une démonstration, et
quand n'est-il qu'un indice ?
\item[(c)] Le célèbre paradoxe de découpage dans lequel un carré 8 $\times$ 8 est coupé en quatre pièces
puis réassemblé en un rectangle 5 $\times$ 13 \og{} démontre \fg{} que 64 = 65. Expliquez exactement où
cette figure ment, et ce que vos critères de la question (b) en disent.
\end{enumerate}

\end{problem}

\noindent Les démonstrations du théorème de Pythagore par réarrangement sont très
anciennes --- l'une accompagne le \textit{Zhoubi Suanjing} en Chine, et l'on rapporte que
Bh\=askara II, dans l'Inde du douzième siècle, dessina la figure et écrivit un seul mot :
\og{} Vois ! \fg{}. Les mathématiciens débattent depuis lors de la confiance que l'on peut accorder à
une figure ; la réponse honnête est qu'un bon dessin est un argument comprimé dont la
décompression incombe au lecteur. \textit{Proofs Without Words} de Nelsen en rassemble une
centaine de beaux spécimens sur lesquels s'exercer.

\begin{refs}
\item Contexte : Preuve sans mots --- Wikipédia (\url{https://fr.wikipedia.org/wiki/Preuve_sans_mots})
\item Contexte : Théorème de Pythagore --- Wikipédia (\url{https://fr.wikipedia.org/wiki/Th%C3%A9or%C3%A8me_de_Pythagore})
\item Lecture : R. B. Nelsen, \textit{Proofs Without Words: Exercises in Visual Thinking} (MAA, 1993)
\end{refs}

\bigskip

\begin{problem}[{La symétrie comme stratégie : des pièces sur une table ronde --- Lycée}]
\textbf{PRF-37.} Deux joueurs posent tour à tour des pièces circulaires identiques sur
une table ronde. Chaque pièce doit reposer entièrement sur la table sans recouvrir aucune
pièce déjà posée (le contact est autorisé), et le joueur qui n'a plus de coup licite
perd.
\begin{enumerate}
\item[(a)] Démontrez que le premier joueur peut toujours gagner, quelles que soient les tailles
de la table et des pièces : décrivez une stratégie et démontrez qu'elle ne manque jamais
de coups avant l'adversaire.
\item[(b)] Désignez les faits géométriques exacts qu'utilise votre démonstration --- où,
précisément, l'argument a-t-il besoin de la symétrie de la table, et pourquoi le premier
coup doit-il être le centre ?
\item[(c)] Pour une table carrée, une table rectangulaire non carrée et une table en L, décidez
si votre démonstration survit ; partout où elle échoue, identifiez l'étape exacte qui
meurt.
\end{enumerate}

\end{problem}

\noindent Un classique des mathématiques récréatives qui distille un principe
sérieux : une symétrie du plateau peut être enrôlée comme stratégie, car l'image miroir de
tout coup licite de l'adversaire est licite à coup sûr --- pourvu que l'on se soit d'abord
emparé de l'unique point que la symétrie laisse fixe. Les stratégies d'appariement et de
miroir de ce genre décident quantité de jeux à deux joueurs sans le moindre calcul, et cette
même habitude de pensée --- argumenter une fois, laisser la symétrie répéter l'argument à
votre place --- raccourcit les démonstrations partout en mathématiques, depuis chaque \og{} sans
perte de généralité \fg{} en géométrie jusqu'au dénombrement par involution de PRF-47.

\begin{refs}
\item Contexte : Argument de vol de stratégie --- Wikipédia (en anglais) (\url{https://en.wikipedia.org/wiki/Strategy-stealing_argument})
\item Contexte : La symétrie en mathématiques --- Wikipédia (en anglais) (\url{https://en.wikipedia.org/wiki/Symmetry_in_mathematics})
\item Lecture : E. R. Berlekamp, J. H. Conway \& R. K. Guy, \textit{Winning Ways for Your Mathematical Plays} (1982)
\item Lecture : A. Engel, \textit{Problem-Solving Strategies}, ch. 13 (les jeux)
\end{refs}

\bigskip

\begin{problem}[{Démontrez ou réfutez : irrationnel + irrationnel --- Lycée, short}]
\textbf{PRF-38.} Démontrez ou réfutez : la somme de deux nombres
irrationnels est irrationnelle. Faites ensuite de même pour le produit de deux nombres
irrationnels.
\end{problem}

\noindent \og{} Démontrez ou réfutez \fg{} est un genre à part entière : la moitié de la
tâche consiste à décider de quel côté on se range, et une réfutation est une démonstration
elle aussi --- un contre-exemple bien choisi, présenté avec le soin dû à un théorème (PRF-43
en est l'atelier complet). Remarquez à quel point les deux charges sont inégales : une
affirmation universelle meurt d'un seul exemple.

\begin{refs}
\item Contexte : Nombre irrationnel --- Wikipédia (\url{https://fr.wikipedia.org/wiki/Nombre_irrationnel})
\item Lecture : R. Hammack, \textit{Book of Proof}, ch. 9 (la réfutation)
\end{refs}

\bigskip

\begin{problem}[{Trois démonstrations que 6 divise n$^3$ $-$ n --- Lycée, short}]
\textbf{PRF-39.} Démontrez que $ n^3 - n $ est divisible par 6 pour
tout entier n --- trois fois, par trois arguments véritablement différents (candidats : la
récurrence ; la factorisation en entiers consécutifs ; l'arithmétique modulo 2 et modulo
3).
\end{problem}

\noindent Une miniature de la discipline récurrente de cette page (PRF-15, PRF-18,
PRF-24) : une vérité, plusieurs démonstrations, et une comparaison de ce que chacune voit.
L'argument par entiers consécutifs se généralise aux coefficients binomiaux ; l'argument
modulaire est le germe du petit théorème de Fermat, qui fait pour $ n^p - n $ ce que ce
problème fait pour $ n^3 - n $ (PRF-18).

\begin{refs}
\item Contexte : Arithmétique modulaire --- Wikipédia (\url{https://fr.wikipedia.org/wiki/Arithm%C3%A9tique_modulaire})
\item Lecture : R. Hammack, \textit{Book of Proof}, ch. 10
\end{refs}

\bigskip

\begin{problem}[{Six personnes à une soirée, en dix phrases --- Lycée, short}]
\textbf{PRF-40.} Démontrez que parmi six personnes quelconques, il en
existe trois qui se connaissent deux à deux, ou trois qui sont deux à deux étrangères l'une
à l'autre --- en dix phrases au plus. Montrez ensuite que six est optimal : disposez cinq
personnes de sorte qu'aucun des deux trios n'existe.
\end{problem}

\noindent Le théorème des amis et des étrangers --- R(3,3) = 6 dans la notation de
Ramsey --- figurait au concours Putnam de 1953 et ouvre toute visite guidée de la théorie de
Ramsey, l'étude de l'ordre qu'aucune quantité de désordre ne peut éviter. Le budget de
phrases est ici le véritable exercice : un argument correct s'étale facilement, et le
comprimer sans rien perdre, c'est l'écriture de démonstration sous sa forme la plus pure. La
frontière est étonnamment proche : la valeur de R(5,5) est toujours inconnue.

\begin{refs}
\item Contexte : Théorème des amis et des étrangers --- Wikipédia (\url{https://fr.wikipedia.org/wiki/Th%C3%A9or%C3%A8me_des_amis_et_des_%C3%A9trangers})
\item Contexte : Théorème de Ramsey --- Wikipédia (\url{https://fr.wikipedia.org/wiki/Th%C3%A9or%C3%A8me_de_Ramsey})
\end{refs}

\bigskip

\begin{problem}[{La réciproque n'est pas gratuite --- Lycée}]
\textbf{PRF-46.} Une affirmation de la forme \og{} P si et seulement si Q \fg{} énonce deux
implications distinctes --- $ P \Rightarrow Q $ et sa réciproque $ Q \Rightarrow P $ --- et
aucune des deux ne découle de l'autre.
\begin{enumerate}
\item[(a)] Démontrez qu'un entier n est divisible par 3 si et seulement si la somme de ses
chiffres décimaux est divisible par 3. Rédigez les deux sens comme deux arguments
clairement étiquetés, et marquez à quel sens appartient chaque étape de votre travail.
\item[(b)] Pour chacun des énoncés ci-dessous, écrivez sa réciproque et décidez, par une
démonstration ou par un contre-exemple, si cette réciproque est vraie : (i) si un
quadrilatère est un carré, alors ses diagonales sont perpendiculaires ; (ii) si n est un
nombre premier supérieur à 2, alors n est impair ; (iii) si un entier est divisible par 24,
alors il est divisible à la fois par 4 et par 6.
\item[(c)] Certaines équivalences se démontrent d'un seul trait, par une chaîne d'étapes toutes
réversibles. Déterminez toutes les solutions réelles de
\[ \sqrt{x+7} \;=\; x+1, \]
puis repérez l'étape exacte de la méthode habituelle \og{} on élève les deux membres au
carré \fg{} qui n'est \textit{pas} réversible, et expliquez pourquoi cette seule étape oblige la
méthode à s'achever par une vérification.
\end{enumerate}

\end{problem}

\noindent Euclide traite déjà un théorème et sa réciproque comme deux travaux
distincts : la proposition I.5 des \textit{Éléments} démontre qu'un triangle isocèle a des
angles à la base égaux, et I.6 --- proposition entièrement séparée, munie de sa propre
démonstration --- en établit la réciproque. La plupart des erreurs d'étudiants du genre \og{} je
l'ai démontré à l'envers \fg{} sont des erreurs de réciproque, et le phénomène des racines
parasites de la question (c) est la même faute en habits algébriques : élever au carré est
une voie à sens unique, de sorte que la chaîne d'équations démontre seulement que \textit{si}
une solution existe, elle figure parmi les candidates trouvées. L'abréviation \og{} iff \fg{} est de
Paul Halmos, et la convention selon laquelle le mot \og{} si \fg{} dans une \textit{définition}
signifie tacitement \og{} si et seulement si \fg{} est de celles que tout lecteur doit assimiler
tôt.

\begin{refs}
\item Contexte : Condition nécessaire et suffisante --- Wikipédia (\url{https://fr.wikipedia.org/wiki/Condition_n%C3%A9cessaire_et_suffisante})
\item Contexte : Implication réciproque --- Wikipédia (\url{https://fr.wikipedia.org/wiki/Implication_r%C3%A9ciproque})
\item Lecture : R. Hammack, \textit{Book of Proof}, le chapitre sur la démonstration des énoncés non conditionnels (gratuit en ligne (\url{https://richardhammack.github.io/BookOfProof/}))
\item Lecture : D. J. Velleman, \textit{How to Prove It}, ch. 3
\end{refs}

\bigskip

\begin{problem}[{Deux quantificateurs, deux sens --- Lycée, short}]
\textbf{PRF-54.} Décidez, avec démonstration, si les deux phrases de
chacune des paires ci-dessous disent la même chose : (i) \og{} pour tout entier x il existe un
entier y tel que $ y > x $ \fg{} contre \og{} il existe un entier y tel que $ y > x $ pour tout
entier x \fg{} ; (ii) \og{} tout réel strictement positif est plus petit qu'un certain entier \fg{}
contre \og{} un certain entier est plus grand que tout réel strictement positif \fg{}. Énoncez
ensuite, en une phrase, la règle générale sur l'échange de \og{} pour tout \fg{} et \og{} il existe \fg{}, et
identifiez le seul sens de cet échange qui soit toujours valide.
\end{problem}

\noindent L'ordre des quantificateurs est de loin la première source de
démonstrations fausses qui ont l'air justes, et l'ennui est que la langue courante le
masque : \og{} tout le monde aime quelqu'un \fg{} et \og{} quelqu'un est aimé de tout le monde \fg{} ne
diffèrent que par l'accent lorsqu'on les dit, et totalement par le sens lorsqu'on les écrit.
Le \textit{Begriffsschrift} de Frege (1879) fut la première notation capable d'exprimer sans
ambiguïté une quantification imbriquée, et ce n'est pas un hasard si l'analyse rigoureuse ---
où continuité, continuité uniforme et convergence ne se distinguent les unes des autres que
par l'ordre des quantificateurs --- est devenue possible au même siècle. PRF-44 met le même
savoir-faire au travail sur la négation d'un énoncé en $\varepsilon$--$\delta$.

\begin{refs}
\item Contexte : Calcul des prédicats --- Wikipédia (\url{https://fr.wikipedia.org/wiki/Calcul_des_pr%C3%A9dicats})
\item Lecture : D. J. Velleman, \textit{How to Prove It}, ch. 2
\item Lecture : R. Hammack, \textit{Book of Proof}, ch. 2 (gratuit en ligne (\url{https://richardhammack.github.io/BookOfProof/}))
\end{refs}

\bigskip

\begin{problem}[{Le principe des tiroirs déguisé : inventer les tiroirs --- Lycée}]
\textbf{PRF-55.} Le principe des tiroirs est trivial à énoncer et difficile à employer,
car toute la difficulté est d'inventer les tiroirs. Dans PRF-07 les tiroirs étaient presque
visibles ; dans chacune des questions ci-dessous, ils ne vous sont pas donnés du tout.
\begin{enumerate}
\item[(a)] Démontrez que parmi $ n+1 $ nombres quelconques choisis dans
$ \{1, 2, \dots, 2n\} $, l'un d'eux en divise un autre.
\item[(b)] Démontrez que toute liste de n entiers (pas nécessairement distincts) contient un bloc
non vide de termes consécutifs dont la somme est divisible par n.
\item[(c)] Démontrez le théorème d'Erd\H{o}s--Szekeres : toute suite de $ n^2 + 1 $ réels distincts
contient une sous-suite croissante de longueur $ n+1 $ ou une sous-suite décroissante de
longueur $ n+1 $.
\item[(d)] Démontrez le théorème d'approximation de Dirichlet : pour tout réel $\alpha$ et tout entier
strictement positif N, il existe des entiers p et q avec $ 1 \le q \le N $ et
\[ \left| \alpha - \frac{p}{q} \right| < \frac{1}{qN}. \]
\item[(e)] Voici maintenant le but de l'exercice. Pour chacune des questions (a)--(d), écrivez
explicitement quels étaient les objets à ranger, quels étaient les tiroirs, et combien il y
en avait de chaque --- puis dites, en une phrase par question, ce qui rendait ces tiroirs
difficiles à voir.
\end{enumerate}

\end{problem}

\noindent Dirichlet a introduit le \textit{Schubfachprinzip} en 1834 non comme une
curiosité mais comme un outil, et la question (d) est ce pour quoi il l'a bâti : un énoncé
absolument non évident sur l'approximation des irrationnels, démontré en laissant tomber
$ N+1 $ parties fractionnaires dans N tiroirs. La question (c) est le théorème d'Erd\H{o}s et
Szekeres de 1935, dont la démonstration la plus habile --- d'ordinaire attribuée à Seidenberg
(1959) --- étiquette chaque terme par un couple de nombres et laisse le principe des tiroirs
conclure en une ligne. Le schéma est le même dans tous les cas, et il mérite d'être nommé :
les tiroirs ne sont pas les objets qu'on vous a remis, mais un invariant calculé à partir
d'eux --- une partie impaire, un reste, un couple de longueurs, une partie fractionnaire.

\begin{refs}
\item Contexte : Théorème d'Erd\H{o}s-Szekeres --- Wikipédia (\url{https://fr.wikipedia.org/wiki/Th%C3%A9or%C3%A8me_d%27Erd%C5%91s-Szekeres})
\item Contexte : Approximation diophantienne --- Wikipédia (\url{https://fr.wikipedia.org/wiki/Approximation_diophantienne})
\item Lecture : A. Engel, \textit{Problem-Solving Strategies}, ch. 4 (le principe des tiroirs)
\item Lecture : M. Aigner \& G. M. Ziegler, \textit{Proofs from THE BOOK}, le chapitre sur le principe des tiroirs et le double dénombrement
\end{refs}

\bigskip


\vfill
\noindent\emph{Hors de la Caverne} --- des probl\`emes, pas de r\'eponses.
\end{document}
